\documentclass[11pt]{article}
\usepackage{geometry}
\geometry{letterpaper,margin=1in}

\usepackage{graphicx}
\usepackage{caption}
\usepackage{subcaption}
\usepackage{amssymb}
\usepackage{csquotes}
\usepackage{amsmath}
\usepackage{amsthm}
\usepackage{mathtools}
\usepackage{float}
\usepackage{longtable}
\usepackage{physics}
\usepackage{enumitem}
\usepackage{silence}
\WarningFilter{hyperref}{Suppressing link with empty target}
\usepackage{xr-hyper}
\usepackage[hidelinks]{hyperref}
\makeatletter\def\XR@addURL#1{\XR@noanchor#1{}{}{}{}\\}
\def\XR@noanchor#1#2#3#4#5\\{{#1}{#2}{#3}{}{}}\makeatother
\externaldocument[I-]{../part1/main}

\MakeOuterQuote{"}
\setlength{\emergencystretch}{5em}

\newtheorem{theorem}{Theorem}[section]
\newtheorem{definition}[theorem]{Definition}
\newtheorem{proposition}[theorem]{Proposition}
\newtheorem{lemma}[theorem]{Lemma}
\newtheorem{corollary}[theorem]{Corollary}
\newtheorem{hypothesis}[theorem]{Hypothesis}


\newcommand{\R}{\mathbb{R}}
\newcommand{\C}{\mathbb{C}}
\newcommand{\Z}{\mathbb{Z}}
\newcommand{\N}{\mathbb{N}}

\newcommand{\KBS}{K_{BS}}
\newcommand{\Knear}{K_{BS}^{\mathrm{nr}}}
\newcommand{\Kfar}{K_{BS}^{\mathrm{fr}}}
\newcommand{\er}{\textbf{e}_r}
\newcommand{\etheta}{\textbf{e}_\theta}
\newcommand{\DT}{\mathcal{D}_T}
\newcommand{\LipO}{\mathrm{Lip}_0}
\newcommand{\Trunc}[1]{\mathrm{Tr}_{#1}}
\newcommand{\Ren}{\Upsilon}
\newcommand{\Circ}{\Gamma}
\newcommand{\jap}[1]{\langle #1\rangle}
\newcommand{\Ip}[2]{I_{#1}(#2)}
\newcommand{\pout}{p^{\mathrm{out}}}
\newcommand{\pin}{p^{\mathrm{in}}}
\newcommand{\Qm}{Q_m}
\newcommand{\flow}{\mathcal{X}}
\newcommand{\Orot}{\mathsf{O}}
\DeclareMathOperator{\sgn}{sgn}
\DeclareMathOperator{\divop}{div}
\DeclareMathOperator{\supp}{supp}
\DeclareMathOperator{\spec}{spec}
\DeclareMathOperator{\wind}{wind}
\DeclareMathOperator{\RN}{RN}

\newcommand{\xiL}{\xi_L}
\newcommand{\xia}{\xi_a}
\newcommand{\xib}{\xi_b}
\newcommand{\xiR}{\xi_R}
\newcommand{\prof}{\varpi}
\newcommand{\omc}{\overline{\omega}_{\mathrm{c}}}%
\newcommand{\Rc}{R_{\mathrm{c}}}
\newcommand{\Gone}{\overline{G}}
\newcommand{\Gmax}{G_{\max}}
\newcommand{\bomega}{\overline{\omega}}
\newcommand{\NN}{\mathsf{N}}
\newcommand{\NNm}{\mathsf{N}^{-}}
\newcommand{\Nflat}{N^{\flat}}

\newcommand{\Lin}{L}
\newcommand{\Lsec}[1]{L_{#1}}
\newcommand{\Green}{\mathcal{G}}
\newcommand{\kcoef}{k}
\newcommand{\Vric}{V}
\newcommand{\Evans}{\mathcal{E}}
\newcommand{\Wr}{\mathcal{W}}
\newcommand{\Mres}{M}
\newcommand{\lamst}{\lambda_*}
\newcommand{\phist}{\phi_*}
\newcommand{\vphist}{\varphi_*}
\newcommand{\ellst}{\ell_*}
\newcommand{\Riesz}{\Pi}
\newcommand{\Jsimp}{J}
\newcommand{\Thl}[1]{\Theta_{#1}}
\newcommand{\epsst}{\epsilon_*}
\newcommand{\Zedge}{Z}
\newcommand{\Rect}{\mathcal{R}}
\newcommand{\Cell}{\mathcal{C}}
\newcommand{\sigb}{\sigma_\flat}
\newcommand{\Lam}{\Lambda_{\mathrm{en}}}

\newcommand{\sbar}{\bar s}
\newcommand{\Mvort}{M_{\mathrm{vort}}}
\newcommand{\Cmom}{C_{\mathrm{mom}}}
\newcommand{\Cstar}{C_{*}}
\newcommand{\maj}{\mathfrak{a}}
\newcommand{\majst}{\mathfrak{a}_*}
\newcommand{\wlin}{w^{\mathrm{lin}}}
\newcommand{\hin}{h^{\mathrm{in}}}
\newcommand{\hout}{h^{\mathrm{out}}}
\newcommand{\etin}{\eta^{\mathrm{in}}}
\newcommand{\etout}{\eta^{\mathrm{out}}}
\newcommand{\etmax}{\eta_{\max}}
\newcommand{\KO}{K_{O}}
\newcommand{\Pprod}{\mathsf{P}}
\newcommand{\Ptil}{\tilde{\mathsf{P}}}
\newcommand{\CDuh}{C^{\mathrm{Duh}}_\mu}
\newcommand{\Cls}{C_{\ell_*}}
\newcommand{\Cin}{C_{\mathrm{in}}}
\newcommand{\Cp}{C_{\mathrm{p}}}
\newcommand{\Cnab}{C_{\nabla}}
\newcommand{\cout}{c_{\mathrm{out}}}
\newcommand{\Abar}{\bar A}
\newcommand{\Klo}{K_{\mathrm{lo}}}
\newcommand{\Kup}{K_{\mathrm{up}}}
\newcommand{\Chand}{C_{\mathrm{hand}}}
\newcommand{\Cseed}{C_{\mathrm{seed}}}
\newcommand{\Cdeep}{C_{\mathrm{deep}}}
\newcommand{\Cpre}{C_{\mathrm{pre}}}
\newcommand{\Cop}{C_{O'}}
\newcommand{\Tpin}{T_{\mathrm{pin}}}
\newcommand{\rT}{\mathsf{r}}
\newcommand{\Tz}{\mathcal{T}_0}
\newcommand{\Tstar}{T_\star}
\newcommand{\Tend}{T_{\mathrm{end}}}
\newcommand{\tsh}{t^\sharp}
\newcommand{\tash}{\tau^\sharp}
\newcommand{\ksh}{\kappa^\sharp}
\newcommand{\kseed}{\kappa_{\mathrm{seed}}}
\newcommand{\kL}{\kappa_{\mathrm{L}}}
\newcommand{\tauK}{\tau_{\mathrm{K}}}
\newcommand{\Ain}{A^{\mathrm{in}}}
\newcommand{\cin}{c^{\mathrm{in}}}
\newcommand{\qs}{q_{\mathrm{s}}}
\newcommand{\Errg}{\mathrm{Err}_g}
\newcommand{\Errpre}{\mathrm{Err}_{\mathrm{pre}}}
\newcommand{\err}{\mathrm{err}}
\newcommand{\Qbar}{\bar Q_K}
\newcommand{\pK}{p_K}
\newcommand{\dH}{\delta_0^{\mathrm{H}}}
\newcommand{\dkick}{\delta_0^{\mathrm{kick}}}
\newcommand{\drd}{\delta_0^{\mathrm{rd}}}
\newcommand{\rhohat}{\hat\rho}
\newcommand{\Lconf}{\Lambda^{\mathrm{conf}}}
\newcommand{\Phase}{\mathcal{A}}
\newcommand{\Pin}{\mathcal{P}}
\newcommand{\Iint}{\mathcal{I}}
\newcommand{\mpole}{\mathfrak{p}}
\newcommand{\Psione}{\Psi_1}
\newcommand{\chiseed}{\chi_{\mathrm{s}}}
\newcommand{\chiout}{\chi_{\mathrm{o}}}

\newcommand{\kst}{k_*}
\newcommand{\NT}{N_{\mathrm{T}}}
\newcommand{\epsend}{\varepsilon_{\mathrm{end}}}



\title{Non-Uniqueness of Weak Solutions to the Unforced\\
Two-Dimensional Euler Equations\\ with Vorticity in $L^1\cap L^p$,\\
II: The Nonlinear Construction}
\hypersetup{pdftitle={Non-Uniqueness of Weak Solutions to the Unforced
Two-Dimensional Euler Equations with Vorticity in L1 and Lp,
II: The Nonlinear Construction},
pdfauthor={Matei P. Coiculescu}}
\author{Matei P. Coiculescu\thanks{
Department of Mathematics, New York University,
Email: \texttt{m.coiculescu@nyu.edu}
}}
\date{Work in progress: draft of \today}

\begin{document}
\maketitle

\begin{abstract}
We prove that the incompressible Euler equations in the plane, without external
force, admit two distinct renormalized weak solutions with the same initial
vorticity $f$, where $f$ is radial and nonnegative and lies in
$(L^1\cap L^q)(\mathbb{R}^2)$ for every $q<\infty$ but is not bounded. One
solution is radially symmetric and stationary. The other is $3$-fold rotationally symmetric,
is equimeasurable with $f$ at every time, and departs from $f$ in finite time;
the two solutions are distinct in $L^\infty_t(L^1\cap L^q)_x$ on $[0,T]\times\mathbb{R}^2$ for
every $q\in(4/3,\infty)$ and every sufficiently large $T$. The initial
vorticity is an explicit superposition of rescaled copies of one linearly
unstable vortex ring, with radii that decrease to zero and amplitudes that
increase geometrically, and the second solution is a limit of solutions in
which an instability placed in a small ring is relayed outward through all the
larger rings in a bounded time. The proof, using the compactness criterion and the linear
analysis of Part I, reduces the theorem to three explicit
statements, two about the zeros of an explicit holomorphic function and one
inequality between explicit numbers; we verify the first two with rigorous
interval arithmetic and prove the third without machine computation.
\end{abstract}

\section{Introduction}

\subsection{The Main Results}

We consider the incompressible Euler equations in the plane, in vorticity form:
\begin{equation}
\label{eq:euler}
\partial_t \omega + v\cdot\nabla \omega = 0,\qquad
v = \KBS\ast \omega,\qquad
\omega(t=0) = \omega_0 .
\end{equation}
Here $\KBS(x)=x^\perp/(2\pi|x|^2)$, $x^\perp=(-x_2,x_1)$, denotes the
Biot-Savart kernel, so that $v$ is divergence free and
$\partial_1v_2-\partial_2v_1=\omega$. There is no external force, and we call
$\omega_0$ the datum.

This paper is the second of two. Together with Part~I~\cite{PartI}, it proves
that, without force in the plane, uniqueness fails among renormalized weak
solutions with vorticity in $L^\infty_t(L^1\cap L^q)_x$, for every
$q\in(4/3,\infty)$ at once and for an explicit radial datum. Numbers with the
prefix I refer to statements, sections and equations of Part~I. Section~1 of
Part~I states the question and discusses the literature, including the
non-uniqueness with a radial force of Vishik \cite{Vishik1}, \cite{Vishik2}, in
the account of \cite{ABCDGJK}, and the results obtained without force by convex
integration. It also describes the construction heuristically. We do not
repeat this discussion.

For $T>0$ and $1\le q\le\infty$ we write $X_q(T)$ for the space
$L^\infty_t(L^1\cap L^q)_x$ of functions on $(0,T)\times\R^2$. Renormalized weak solutions in $X_q(T)$, $4/3<q<\infty$ (after
DiPerna and Lions \cite{DiPernaLions}), and equimeasurability are defined in
Definitions~I.\ref{I-def:weak} and~I.\ref{I-def:equimeasurable}. The ring $\Phi$ of
Definition~I.\ref{I-def:ring} is a closed formula in $\log|y|$ times a fixed smooth
cutoff; it is radial, $\Phi\in C_c^\infty(\R^2)$, $0\le\Phi<1$,
$\Phi\not\equiv0$, and $\supp\Phi\subset\{a\le|y|\le1\}$ for an explicit
$a\in(0,1)$ (Lemma~I.\ref{I-lem:ring-facts}). For $A_1>0$, $\mu\in(1,2]$ and
$\sbar\in(0,a/16]$ we put
\begin{equation}
\label{eq:datum}
\begin{gathered}
f=\sum_{k\ge1}A_k\,\Phi(\cdot/r_k),\qquad A_k=A_1\mu^{k-1},\\
r_1=1,\qquad r_{k+1}=s_kr_k,\qquad s_k=\sbar(2\mu)^{-k}\qquad(k\ge1).
\end{gathered}
\end{equation}
This is Definition~\ref{def:datum}. We call $A_k\Phi(\cdot/r_k)$ ring $k$ of
$f$. It is supported in the annulus $\{ar_k\le|x|\le r_k\}$, and these annuli
are pairwise disjoint (Proposition~\ref{prop:integrability}). The main theorem
of the two papers is the following; it is Theorem~I.\ref{I-thm:main}, restated.

\begin{theorem}
\label{thm:main}
Let $\Phi$ and $a$ be as in Definition~I.\ref{I-def:ring}. There are
$\mu\in(1,\frac{17}{16}]$ and $\sbar\in(0,a/16]$ such that, for every $A_1>0$,
the function $f$ of Equation~\eqref{eq:datum} has the following properties. It
is radial and nonnegative, $f\in (L^1\cap L^q)(\R^2)$ for every $q\in[1,\infty)$,
and $f\notin L^\infty(\R^2)$. There is $\Tstar\in(0,\infty)$ such that, for
every $T>\Tstar$ and every $q\in(4/3,\infty)$, Equation~\eqref{eq:euler} with
$\omega_0=f$ has two distinct renormalized weak solutions in $X_q(T)$, in the
sense of Definition~I.\ref{I-def:weak}:
\begin{enumerate}[label=\textup{(\roman*)}]
\item the stationary solution $\omega(t)=f$, $t\in[0,T]$;
\item a solution $\omega_\infty$ such that
  $\omega_\infty\in C([0,\infty);L^p(\R^2))$ for every $p\in[1,\infty)$ and,
  for every $t\ge0$, $\omega_\infty(t)$ is equimeasurable with $f$ and
  $3$-fold rotationally symmetric.
\end{enumerate}
More precisely, there is $t_0\in[0,\Tstar)$ such that $\omega_\infty$ and $f$
are distinct elements of $X_q(T')$ for every $T'>t_0$ and every
$q\in(4/3,\infty)$.
\end{theorem}

We do not treat $q\le4/3$. No symmetry is
imposed on the class; the $3$-fold rotational symmetry of $\omega_\infty$ comes from the
construction, where it removes the angular modes $\ell\notin3\Z$, among them the
modes $\ell=\pm1$ that contain the translations of a ring, and leaves the
modes $\ell\in3\Z$, whose spectra Part~I controls under
Hypothesis~I.\ref{I-hyp:N}.

The proof of Theorem~\ref{thm:main} is divided between Part~I and this paper.
Part~I proves a non-uniqueness criterion (Section~I.\ref{I-sec:compact}) and the
linear analysis of the ring (Section~I.\ref{I-sec:ring}), and it uses no machine
computation. The main conclusions of its linear analysis assume three explicit
statements, collected in Section~I.\ref{I-sec:ring} as
Hypothesis~I.\ref{I-hyp:N}. The statements \textup{(N1)} and \textup{(N2)} locate
the zeros in the mode $\ell=3$ of the Evans function of the ring, an explicit
holomorphic function, for every $G\in[0,\frac1{50}]$, where $2\pi G$ is the
circulation of a point vortex at the origin (in the reference units of
Definition~I.\ref{I-def:ring}; this is $G\in[0,\Gone]$ in ring units). The
statement \textup{(N3)} is an inequality between explicit numbers controlling
the modes $\ell=\pm6$ and $|\ell|\ge9$. This paper proves that
\textup{(N1)}--\textup{(N3)} imply Theorem~\ref{thm:main}, by a nonlinear
construction followed by the criterion of Part~I, and it proves
\textup{(N1)}--\textup{(N3)}. The implication is the following reduction,
which we prove in Section~\ref{sec:bootstrap}.

\begin{proposition}
\label{prop:reduction}
Suppose that \textup{(N1)}--\textup{(N3)} of Hypothesis~I.\ref{I-hyp:N} hold. Then
Theorem~\ref{thm:main} holds.
\end{proposition}

The proof of Proposition~\ref{prop:reduction} uses no machine computation.
Appendix~\ref{app:verify} proves \textup{(N3)} by hand and verifies
\textup{(N1)} and \textup{(N2)} with rigorous interval arithmetic, the only
machine-computed part of the proof (code: the package
\texttt{euler-ring-certificate}, version~1.0.0, in the public repository
\url{https://github.com/mcoiculescu/NonUniqueEuler};
Section~\ref{app:reproduce}). It thus proves the following theorem.

\begin{theorem}
\label{thm:N-holds}
Hypothesis~I.\ref{I-hyp:N} holds: the statements \textup{(N1)}, \textup{(N2)} and
\textup{(N3)} are true.
\end{theorem}

\begin{proof}[Proof of Theorem~\ref{thm:main}]
By Theorem~\ref{thm:N-holds}, the hypotheses \textup{(N1)}--\textup{(N3)}
hold, so Theorem~\ref{thm:main} follows from Proposition~\ref{prop:reduction}.
\end{proof}

\subsection{Results of Part I}
\label{sec:II-partI}

Part~I's introduction summarizes what this paper uses in two theorems: the
conditional non-uniqueness theorem (Theorem~I.\ref{I-thm:conditional}, the form
of Corollary~I.\ref{I-cor:defect} given by Remark~I.\ref{I-rem:conditional}) and
the linear instability of the ring under \textup{(N1)}--\textup{(N3)}
(Theorem~I.\ref{I-thm:linear}).

In detail, we use the following definitions and results of Part~I, which we
cite by their numbers there and do not restate. From
Section~I.\ref{I-sec:prelim}, we use the properties of the Biot-Savart kernel
(Lemma~I.\ref{I-lem:BS-basic}) and the multipoles, equimeasurability and
renormalized weak solutions (Definitions~I.\ref{I-def:multipole},
I.\ref{I-def:equimeasurable} and~I.\ref{I-def:weak}, with
Equations~(I.\ref{I-eq:weak}) and~(I.\ref{I-eq:renorm})). From
Section~I.\ref{I-sec:compact}, we use the fact that a radial function in
$(L^1\cap L^q)(\R^2)$, $4/3<q<\infty$, is a stationary renormalized weak
solution (Lemma~I.\ref{I-lem:radial-stationary}), the $2$-admissible sequences
(Definition~I.\ref{I-def:adm}), and the non-uniqueness criterion of
Corollary~I.\ref{I-cor:defect}, which rests on Theorem~I.\ref{I-thm:criterion}
and which we apply in the proof of Proposition~\ref{prop:reduction}. From
Section~I.\ref{I-sec:ring}, we use the ring and the annuli $N$, $N'$
(Definition~I.\ref{I-def:ring}, Lemma~I.\ref{I-lem:ring-facts}); the mode
settings, the decomposition into angular modes, the Evans function and its
simple zeros, and the semigroup (Definition~I.\ref{I-def:sector-setting},
Lemmas~I.\ref{I-lem:sector}, I.\ref{I-lem:evans-basic}
and~I.\ref{I-lem:semigroup-construct}, Theorem~I.\ref{I-thm:evans-simple}); and
Hypothesis~I.\ref{I-hyp:N}. Under \textup{(N1)}--\textup{(N3)} we use, from the
same section, the spectrum of the linearization (Proposition~I.\ref{I-prop:SU}),
the eigendata (Definition~I.\ref{I-def:eigendata}), the continuity in $G$ of
$\lamst(G)$ and of the sizes $|\pout_*(G)|$, $|\beta(G)|$,
$\|\ellst(G)\|_{(X^0)^*}$ and $\|\phist(G)\|_{H^k}$ of the eigendata, the
semigroup bounds and the uniformity of the constants in $G$
(Propositions~I.\ref{I-prop:U} and~I.\ref{I-prop:semigroup},
Lemma~I.\ref{I-lem:uniform}).

\subsection{Strategy}

We describe the construction with references to the statements of this paper
that carry it out. The background is the datum $f$, which is radial, hence
stationary (Lemma~I.\ref{I-lem:radial-stationary}). Its inner rings have larger
amplitudes and rotate faster, and the task is to
carry an instability from an inner ring outward through every ring in a time
independent of the depth, uniformly in the number of interacting rings.

The linear input is the instability of one ring. The linearization of
Equation~\eqref{eq:euler} at $\Phi$ plus a point vortex of small circulation
$2\pi G$ at the origin, which stands for everything inside the ring, is the
operator $\Lin(G)$ of Theorem~I.\ref{I-thm:linear}. On $3$-fold rotationally
symmetric perturbations it has an eigenvalue $\lamst(G)$ of the mode $\ell=3$,
algebraically simple there, with real part $\sigma(G)>0$, and its conjugate
$\overline{\lamst(G)}$ in the mode $\ell=-3$; the rest of the spectrum has
smaller real part (Proposition~I.\ref{I-prop:SU}). Ring $k$ of $f$ evolves on
the time scale $1/A_k$ and feels the inner rings mostly through their
circulation $2\pi G_kA_kr_k^2$, with $G_k\ge0$ small when $\sbar$ is small
(Lemma~\ref{lem:inner-circ}). In ring-$k$ units (Section~\ref{sec:datum}), the linearization at
ring $k$ plus this point vortex is therefore $\Lin(G_k)$.

The approximate datum $f^{(n)}$ keeps the first $n$ rings of $f$ and replaces
ring $n$, the innermost of them, by a rearrangement of it whose coefficient
along the unstable eigenfunction has a small size of order $\kseed$, fixed
independently of $n$ (Definition~\ref{def:approx}, Lemma~\ref{lem:seed}). As
$r_n$ decreases faster than $A_n$ increases, $f^{(n)}\to f$ in
$(L^1\cap L^q)(\R^2)$ for every $q<\infty$ (Lemma~\ref{lem:adm}). Each
$f^{(n)}$ is smooth with compact support, so it has a unique global smooth
solution $\omega^{(n)}$ (Theorem~\ref{thm:smooth}), which is $3$-fold
rotationally symmetric (Lemma~\ref{lem:adm}).

In $\omega^{(n)}$ the instability passes outward like a relay. The unstable
coefficient of ring $k$ grows at the rate $A_k\sigma(G_k)$ to a fixed size
$\delta_0$, and we then say that ring $k$ saturates. Its mode-$3$ moment drives
ring $k-1$ through a field smaller by a factor of order $s_{k-1}^5$
(Lemma~\ref{lem:coupling}). This drive outgrows the free growth of ring $k-1$
by a margin of order $\sigma(0)(\mu-1)/\mu$. The forced part of the
coefficient of ring $k-1$ is a truncated Laplace transform of the drive
(Lemma~\ref{lem:coupling}, Equation~\eqref{eq:kick-A}). Because of this
margin, it is dominated by the end of the growth and pinned after saturation. Its modulus lies between two multiples of the size of
the drive at saturation divided by $\mu-1$, with constants independent of $k$
and $n$ (Proposition~\ref{prop:kick} and its proof).

As $\mu$ is fixed, ring $k-1$ therefore starts with a coefficient of order
$s_{k-1}^5$ (Step~3 in the proof of Proposition~\ref{prop:bootstrap}), and it
saturates after a physical time of about
$5\log(1/s_{k-1})/(\sigma(G_{k-1})A_{k-1})$. Since $\log(1/s_k)$ grows
linearly in $k$ and $A_k$ geometrically, the total time is bounded by a
$\Tstar$ independent of $n$ (Lemma~\ref{lem:time}).

The phase clock of Definition~\ref{def:clock} fixes, from $\omega^{(n)}$
itself, the time at which each ring becomes active and the time at which it is
expected to saturate. Until it becomes active, a ring stays close to $\Phi$,
although forced by the fields inside and outside it (Theorem~\ref{thm:quiet}).
The active ring follows its unstable eigenfunction (Theorem~\ref{thm:active}).
Once the ring just inside it has saturated, the material inside each ring
stays in a small disk (Proposition~\ref{prop:confine}). The hypotheses of these three statements are
verified jointly, for all rings and all times until the outermost ring
saturates, by one bootstrap with constants independent of $n$ (proof of
Proposition~\ref{prop:bootstrap}).

When the outermost ring saturates, at $T^{(n)}\le\Tstar$, its moment
$\int\omega^{(n)}(T^{(n)})\Psione$, where $\Psione(x)=\chiout(|x|)\,\bar x^{\,3}$
with a fixed real cutoff $\chiout\in C_c^\infty((a/4,3))$, has modulus at least
$c_0>0$ independent of $n$ (proof of Proposition~\ref{prop:reduction}), while
it vanishes for $f$. This lower bound is the moment hypothesis of
Theorem~I.\ref{I-thm:conditional}. The $\omega^{(n)}$ are bounded in
$(L^1\cap L^2)(\R^2)$ uniformly in time, and compactness of renormalized
solutions in $L^2$ (Section~I.\ref{I-sec:compact}) gives a subsequence
converging to a renormalized weak solution $\omega_\infty$ with datum $f$. The
moment bound passes to the limit and separates $\omega_\infty$ from $f$
(Corollary~I.\ref{I-cor:defect}).

Section~\ref{sec:II-prelim} collects notation and preliminaries.
Section~\ref{sec:single} studies, in ring units, one ring of the approximate
solutions, driven by the field of the vorticity inside and outside it: the quiet
ring, the active ring and the seeded ring.
Section~\ref{sec:datum} defines the datum, the approximants $f^{(n)}$ and the
phase clock. Section~\ref{sec:several} bounds the fields that the rings induce
on one another, and it confines the material inside a ring after the ring just
inside it has saturated. Section~\ref{sec:kick} computes the kick that a ring
receives from the rings inside it. Section~\ref{sec:relay} proves the relay
lemmas, which are the estimates on the rings that the bootstrap combines, each
under explicit hypotheses; it also proves the time bound of Lemma~\ref{lem:time}.
Section~\ref{sec:param} chooses the parameters so that the conditions of the
relay lemmas on them hold; the choices are made in the order of
Table~\ref{tab:param}, and none of them depends on $n$, $A_1$ or
$\omega^{(n)}$ (Lemma~\ref{lem:nonempty}). Section~\ref{sec:bootstrap} proves
the joint bootstrap (Proposition~\ref{prop:bootstrap}) and
Proposition~\ref{prop:reduction}. Appendices~\ref{app:verify}
and~\ref{app:smooth} prove Theorems~\ref{thm:N-holds} and~\ref{thm:smooth}, respectively.
Appendix~\ref{app:deps} lists, for each statement on the path to
Theorem~\ref{thm:main}, the statements of this paper and of Part~I that its
proof uses and the parts of \textup{(N1)}--\textup{(N3)} on which it depends.

The main idea of a linear instability in physical coordinates driving the
non-uniqueness mechanism, the use of non-monotone unstable radial vortices and of
a multiscale datum, and the overall strategy of the proof are due to the author.
Within this strategy, the detailed proofs and estimates, the code for the
computer-assisted verification and the text were largely produced with
Claude Opus 5.5 (Anthropic), under the direction of the author. The author has
checked all proofs by hand and takes full responsibility for the paper.

\paragraph{Acknowledgments.} I am supported at New York University by a Junior
Fellowship of the Simons Society of Fellows, funded by Simons Foundation
International and administered by the Simons Foundation. I would also like to thank
Tristan Buckmaster for generously offering his time and advice on this and many
other projects.

\subsection{Preliminaries}
\label{sec:II-prelim}

This subsection collects the notation that we take from Part~I and the
preliminaries used only in this paper. We use all notation and conventions of
Section~I.\ref{I-sec:prelim}, among them the convention on constants and
$O(\cdot)$, which we follow. From
Section~I.\ref{I-sec:ring} we use the convention $m=3$, the ring units and the
number $\Gone$ of Definition~I.\ref{I-def:ring}, the annuli $N$ and $N'$ of
Section~I.\ref{I-sec:ring-def}, the spaces $X^s$ and the operators $\Lin(G)$ on
them, fixed before Proposition~I.\ref{I-prop:SU}, and the constants $c_N$, $C_N$,
$\cout$, $\|\ellst\|$, $\Cls$, $\Cp$ and $\Cnab(k)$ of
Lemma~I.\ref{I-lem:uniform}. The results and definitions of
Part~I that we use are listed in Section~\ref{sec:II-partI}.

In ring units we also use the closed annuli
\[
\NNm=\{\tfrac78a\le|y|\le\tfrac54\},\qquad \NN=\{\tfrac34a\le|y|\le\tfrac32\},\qquad
\Nflat=\{\tfrac58a\le|y|\le\tfrac74\}.
\]
They mark where the support of a ring lies at the stages of the construction:
in $\NNm$ for a quiet ring and at the start of the active phase
(Theorems~\ref{thm:quiet} and~\ref{thm:active}), in $\NN$ during the active
phase, and in $\Nflat$ under the bootstrap assumption of
Section~\ref{sec:bootstrap}, which the bootstrap improves to $\NN$. All of them
lie in $N$, the common support of the spaces $X^s$ of Part~I, and the strict
inclusion $\NN\subset\operatorname{int}\Nflat$ leaves the room that the bootstrap
needs.

\begin{lemma}
\label{lem:annuli}
Let $a\in(0,1)$ and define the annuli above, and $N$, $N'$, with this $a$. Then
$\{a\le|y|\le1\}\subset\NNm\subset\NN\subset\operatorname{int}\Nflat\subset\Nflat\subset N\subset N'$.
Let $S\subset\{a\le|y|\le1\}$ \textup{(}respectively $S\subset\NNm$\textup{)},
and let $S'$ be a set such that every $y'\in S'$ satisfies
$\big||y'|-|y|\big|\le a/16$ for some $y\in S$. Then $S'\subset\NNm$
\textup{(}respectively $S'\subset\NN$\textup{)}.
\end{lemma}
\begin{proof}
The inner radii $a$, $\frac78a$, $\frac34a$, $\frac58a$, $\frac12a$, $\frac14a$
strictly decrease and the outer radii $1$, $\frac54$, $\frac32$, $\frac74$, $2$, $3$
strictly increase, and the last two radii of each list are those of $N$ and $N'$;
this gives the chain. If $a\le|y|\le1$ and
$||y'|-|y||\le a/16$, then $\frac{15}{16}a\le|y'|\le1+\frac a{16}$, and
$\frac{15}{16}a\ge\frac78a$ and $1+\frac a{16}\le\frac54$ because $a\le4$. If
instead $y\in\NNm$, then $\frac{13}{16}a\le|y'|\le\frac54+\frac a{16}$, and
$\frac{13}{16}a\ge\frac34a$ and $\frac54+\frac a{16}\le\frac32$, for the same
reason.
\end{proof}

Away from the support of a vorticity, complex variables turn the Biot-Savart
law into a Cauchy integral, which can be expanded in powers of the complex
variable.
We call $V_1-iV_2$ the \textbf{complex form} of a vector field $V$, and wherever
we use complex forms we write $z=x_1+ix_2$ (this $z$ is not the number
$\lamst(0)$ of Definition~I.\ref{I-def:eigendata}). Since $(\KBS)_1(x)-i(\KBS)_2(x)=(-x_2-ix_1)/(2\pi|x|^2)=1/(2\pi iz)$,
the velocity of every $g\in L^1(\R^2)$ has the complex form
\begin{equation}
\label{eq:BS-complex}
(\KBS\ast g)_1(z)-i(\KBS\ast g)_2(z)=\frac1{2\pi i}\int\frac{g(\zeta)}{z-\zeta}\dd A(\zeta)
\qquad(z\in\C\setminus\supp g),
\end{equation}
where $\dd A$ is the area element, the integral converges absolutely and, by differentiation under the
integral sign, defines a holomorphic function of $z\in\C\setminus\supp g$. Let
$V$ be a real vector field on an open set $U$ whose complex form $F$ is
holomorphic on $U$. Then $\partial_{x_1}F=F'$ and $\partial_{x_2}F=iF'$, so every
partial derivative of order $q$ of $F$ is $i^{\gamma_2}F^{(q)}$ for the
corresponding multi-index $\gamma$. Fix one of the $2^q$ ordered $q$-tuples of
directions. Since $V_1=\Re F$ and $V_2=-\Im F$, the corresponding derivatives of
$V_1$ and $V_2$ at $x$ are the real part and minus the imaginary part of the
corresponding derivative of $F$ at $z$, a number of modulus $|F^{(q)}(z)|$; so this
pair of entries of $\nabla^qV(x)$ has Euclidean norm $|F^{(q)}(z)|$. Summing the
squares over the $2^q$ tuples gives
\begin{equation}
\label{eq:BS-complex-deriv}
|\nabla^qV(x)|=2^{q/2}\,|F^{(q)}(z)|\qquad(z\in U,\ q\ge0).
\end{equation}

The next two lemmas bound the field of a rotationally symmetric vorticity away
from its support: in a disk inside a hole of the support, and outside a disk that
contains the support, after the monopole is removed. Section~\ref{sec:several}
applies them to the rings outside and inside a given ring, respectively.

\begin{lemma}
\label{lem:hollow}
Let $m_0\ge2$ be an integer and $R>0$. Let $\nu\in (L^1\cap L^\infty)(\R^2)$ be
real and $m_0$-fold rotationally symmetric with $\supp\nu\subset\{|y|\ge R\}$, and let
$\bar\nu(y)=\frac1{2\pi}\int_0^{2\pi}\nu(\Orot_\phi y)\dd\phi$ be its angular
average. Then $\KBS\ast\bar\nu=0$ on $B_R$, and for $|x|\le R/2$,
\[
|\KBS\ast\nu(x)|\le\tfrac4{3\pi}\|\nu\|_{L^1}|x|^{m_0-1}R^{-m_0}.
\]
\end{lemma}

\begin{proof}
Let $|z|<R$, so that $|z-\zeta|\ge R-|z|>0$ for $\zeta\in\supp\nu$, and let $u=\KBS\ast\nu$. For $|\zeta|\ge R$ the series
$(z-\zeta)^{-1}=-\sum_{l\ge1}z^{l-1}\zeta^{-l}$ converges uniformly in $\zeta$, with ratio
at most $|z|/R$, so Equation~\eqref{eq:BS-complex} gives, term by term,
$u_1-iu_2=\frac{i}{2\pi}\sum_{l\ge1}c_l(\nu)z^{l-1}$ with
$c_l(\nu)=\int\nu(\zeta)\zeta^{-l}\dd A(\zeta)$. The substitution $\zeta\mapsto e^{2\pi i/m_0}\zeta$ and the $m_0$-fold rotational
symmetry of $\nu$ show $c_l(\nu)=e^{-2\pi
il/m_0}c_l(\nu)$, so $c_l(\nu)=0$ unless $m_0\mid l$. The angular average
$\bar\nu$ is radial and, like $\nu$, supported in $\{|y|\ge R\}$, so the same
expansion applies to it, and integration in the angle
shows $c_l(\bar\nu)=0$ for $l\ge1$; this proves $\KBS\ast\bar\nu=0$ on $B_R$. It
also gives $c_l(\nu)=c_l(\nu-\bar\nu)$, so $|c_l(\nu)|\le\|\nu-\bar\nu\|_{L^1}R^{-l}$,
as $|\zeta|\ge R$ on the support of $\nu-\bar\nu$.
Since $\nu$ is real, $|u|=|u_1-iu_2|$. Only the terms with $l=jm_0$, $j\ge1$,
remain, so for $|x|\le R/2$ these bounds and the sum of a geometric series
give
\[
|u(x)|\le\frac{\|\nu-\bar\nu\|_{L^1}}{2\pi}\sum_{j\ge1}|x|^{jm_0-1}R^{-jm_0}
\le\frac{\|\nu-\bar\nu\|_{L^1}}{2\pi(1-2^{-m_0})}|x|^{m_0-1}R^{-m_0},
\]
and $1/(1-2^{-m_0})\le\frac43$ because $m_0\ge2$. Finally
$\|\nu-\bar\nu\|_{L^1}\le2\|\nu\|_{L^1}$, because averaging over rotations does
not increase the $L^1$ norm.
\end{proof}

For an integer $m\ge1$ and $g$ as in Definition~I.\ref{I-def:multipole}, let
$\pout_m[g]$ and $\Qm[g]$ be the multipoles defined there.
If $\supp g\subset\overline B_{\varrho}$, the expansion
\[
\log|x-y|=\log|x|-\Re\sum_{k\ge1}\frac1k\Big(\frac yx\Big)^k\quad(|y|<|x|)
\]
shows that the angular mode $m$ of the stream function
$(2\pi)^{-1}\log|\cdot|\ast g$ of $g$ equals $-(2m)^{-1}\pout_m[g]\,r^{-m}$
for $r>\varrho$; the series may be integrated term by term against $g$ because,
for fixed $|x|>\varrho$, it converges uniformly in $y\in\overline B_\varrho$. For real $g$ this is the angular mode $m$ of the exterior
expansion after Definition~I.\ref{I-def:multipole}.
Section~\ref{sec:kick} uses it for the field that the inner rings induce on a
ring.

\begin{lemma}
\label{lem:ext-multipole}
Let $m\ge1$ and $q\ge0$ be integers, and let $\varrho>0$. Let $g\in L^1(\R^2)$ be
real and $m$-fold rotationally symmetric with $\supp g\subset\overline{B}_\varrho$, and let $E$ be the
velocity field $\KBS\ast g$ minus its monopole
$\frac{1}{2\pi}\big(\int g\big)\frac{x^\perp}{|x|^2}$. Then for $|x|\ge2\varrho$,
\[
|\nabla^qE(x)|\le C_{m,q}\,\|g\|_{L^1}\,\varrho^m|x|^{-m-1-q},\qquad
C_{m,q}=\frac{2^q}{2\pi}\sup_{0<\rho\le1/2}\rho^{-m}\sum_{l\in m\N}(l+q)^q\rho^{l}<\infty .
\]
\end{lemma}

\begin{proof}
Let $|z|>\varrho$. For $|\zeta|\le\varrho$ the series
$(z-\zeta)^{-1}=\sum_{l\ge0}\zeta^lz^{-l-1}$ converges uniformly in $\zeta$, so
Equation~\eqref{eq:BS-complex} gives
$(\KBS\ast g)_1-i(\KBS\ast g)_2=\frac1{2\pi i}\sum_{l\ge0}Q_lz^{-l-1}$ with
$Q_l=\int g(\zeta)\zeta^l\dd A(\zeta)$. The term $l=0$, which is $(\int g)/(2\pi iz)$, is the
monopole, and $Q_l=0$ unless $m\mid l$, because the substitution $\zeta\mapsto
e^{2\pi i/m}\zeta$ and the $m$-fold rotational symmetry of $g$ give $Q_l=e^{2\pi il/m}Q_l$. Hence the complex form of $E$ is
$F(z)=\frac1{2\pi i}\sum_{l\in m\N}Q_lz^{-l-1}$, where
$|Q_l|\le\|g\|_{L^1}\varrho^l$ because $|\zeta|\le\varrho$ on $\supp g$. The series converges locally uniformly on
$\{|z|>\varrho\}$, so $F$ is holomorphic there, may be differentiated term by
term, and, as $E$ is real,
Equation~\eqref{eq:BS-complex-deriv} gives
$|\nabla^qE(x)|=2^{q/2}|F^{(q)}(z)|\le2^q|F^{(q)}(z)|$. With
$|\partial_z^qz^{-l-1}|=(l+1)\cdots(l+q)|z|^{-l-1-q}\le(l+q)^q|z|^{-l-1-q}$ and
$\rho=\varrho/|z|\le\frac12$,
\[
|F^{(q)}(z)|\le\frac{\|g\|_{L^1}}{2\pi}|z|^{-1-q}\sum_{l\in m\N}(l+q)^q\rho^l
\le\frac{C_{m,q}}{2^q}\|g\|_{L^1}\rho^m|z|^{-1-q},
\]
and $\rho^m|z|^{-1-q}=\varrho^m|x|^{-m-1-q}$. The constant is finite, because
$\rho^{-m}\sum_{l\in m\N}(l+q)^q\rho^l=\sum_{j\ge1}(mj+q)^q\rho^{m(j-1)}$ is
nondecreasing in $\rho$ and finite at $\rho=\frac12$.
\end{proof}

The approximate solutions $\omega^{(n)}$ of Definition~\ref{def:approx} are the
global smooth solutions with the smooth, compactly supported data $f^{(n)}$,
given by the following classical theorem, which is proved in
Appendix~\ref{app:smooth}. Besides the properties (1)--(3) of the solution, the
proof of Theorem~\ref{thm:main} uses the flow of item~(4), the uniqueness of
item~(5), which identifies a rescaled or time-translated solution with the
solution of its datum, the symmetry of item~(6), and item~(7), by which these solutions are renormalized weak solutions
in the sense of Part~I.

\begin{theorem}[cf. \cite{Wolibner}, \cite{Holder}, \cite{Kato1967}]
\label{thm:smooth}
Let $\omega_0\in C_c^\infty(\R^2)$ be real with $\supp\omega_0\subset B_{R_0}$.
There is a unique $\omega\in C^\infty(\R\times\R^2)$ such that:
\begin{enumerate}[label=\textup{(\arabic*)}]
\item $\supp\omega(t)\subset B_{R(t)}$ for every $t\in\R$, with
  $R(t)=R_0+(2\|\omega_0\|_{L^\infty}\|\omega_0\|_{L^1}/\pi)^{1/2}|t|$;
\item $v:=\KBS\ast\omega\in C^\infty(\R\times\R^2)$, $v$ and every derivative of
  $v$ are bounded on $[-T,T]\times\R^2$ for every $T>0$, and $\divop v=0$;
\item $\partial_t\omega+v\cdot\nabla\omega=0$ on $\R\times\R^2$ and $\omega(0)=\omega_0$.
\end{enumerate}
Moreover:
\begin{enumerate}[label=\textup{(\arabic*)},resume]
\item the flow $\flow$ of $v$, $\partial_t\flow(t,y)=v(t,\flow(t,y))$,
  $\flow(0,y)=y$, is defined for all $t\in\R$; each $\flow_t=\flow(t,\cdot)$ is
  a $C^\infty$ diffeomorphism of $\R^2$ with $\det D\flow_t\equiv1$; the maps
  $\flow$ and $(t,x)\mapsto\flow_t^{-1}(x)$ lie in $C^\infty(\R\times\R^2)$; and
  $\omega(t)=\omega_0\circ\flow_t^{-1}$; consequently
  $\int F(\omega(t))\dd x=\int F(\omega_0)\dd x$ for every continuous $F$ with
  $F(0)=0$;
\item if $I\ni0$ is an interval and $\tilde\omega\in C^1(I\times\R^2)$ satisfies
  $\tilde\omega(0)=\omega_0$, $\partial_t\tilde\omega+(\KBS\ast\tilde\omega)\cdot\nabla\tilde\omega=0$
  on $I\times\R^2$, and, for each compact subinterval $I'\subset I$,
  $\supp\tilde\omega(t)\subset B_{\tilde R}$ for $t\in I'$ with $\tilde R$
  depending on $I'$, then $\tilde\omega=\omega$ on $I\times\R^2$;
\item if $\omega_0\circ\Orot_\theta=\omega_0$ for a rotation $\Orot_\theta\in SO(2)$,
  then $\omega(t)\circ\Orot_\theta=\omega(t)$ for every $t\in\R$;
\item for every $T>0$ and every $q\in(4/3,\infty)$, the restriction of $\omega$
  to $[0,T)$ lies in $L^\infty_t(L^1\cap L^\infty)_x$ and is a renormalized
  weak solution in $X_q(T)$ with datum $\omega_0$, in the sense of
  Definition~I.\ref{I-def:weak}.
\end{enumerate}
\end{theorem}

The energy estimates of Section~\ref{sec:single}, the construction of
Appendix~\ref{app:smooth} and the proof of Lemma~\ref{lem:Hk-calculus} below use
the following Sobolev inequalities.

\begin{lemma}
\label{lem:embed}
The following embeddings hold on $\R^2$, with explicit constants.
\begin{enumerate}[label=\textup{(E\arabic*)}]
\item For $g\in H^1$, $\|g\|_{L^4}\le\|g\|_{H^1}$.
\item For $g\in H^2$, $g$ has a bounded continuous
  representative, and $\|g\|_{L^\infty}\le\|g\|_{H^2}$.
\item For an integer $k\ge0$ and $g\in H^{2k}$,
  $\|g\|_{H^k}^2\le3^k\|g\|_{L^2}\|g\|_{H^{2k}}$.
\end{enumerate}
\end{lemma}

\begin{proof}
(E1) Let first $g\in C_c^\infty(\R^2)$. Since
$|g(x)|^2=\int_{-\infty}^{x_1}\partial_1|g|^2(t,x_2)\dd t$, we have
\[
|g(x)|^2\le\int_\R|\partial_1|g|^2|(t,x_2)\dd t=:a(x_2),
\]
and likewise
$|g(x)|^2\le\int_\R|\partial_2|g|^2|(x_1,t)\dd t=:b(x_1)$. The product of
these two bounds, integrated over $x$, gives
$\int|g|^4\le\int a\int b$. Since $|\partial_j|g|^2|\le2|g||\partial_jg|$,
the Cauchy-Schwarz inequality, and then the inequality between the arithmetic
and geometric means, give
$\int|g|^4\le4\|g\|_{L^2}^2\allowbreak\|\partial_1g\|_{L^2}\|\partial_2g\|_{L^2}\le2\|g\|_{L^2}^2\|\nabla g\|_{L^2}^2$.
By density of $C_c^\infty$ in $H^1$ the inequality holds on $H^1$. Then, by the
same inequality between the means,
$\|g\|_{L^4}\le2^{1/4}\|g\|_{L^2}^{1/2}\|\nabla g\|_{L^2}^{1/2}\le2^{-1/4}\|g\|_{H^1}$.

(E2) Since $\widehat{\Delta g}=-|\zeta|^2\hat g$, Plancherel's theorem gives
$\|(1+|\zeta|^2)\hat g\|_{L^2}=2\pi\|g-\Delta g\|_{L^2}\le2\sqrt3\,\pi\|g\|_{H^2}$,
where the inequality follows from the triangle inequality for the three terms of
$g-\Delta g$ and the Cauchy-Schwarz inequality.
As $\int_{\R^2}(1+|\zeta|^2)^{-2}\dd\zeta=\pi$, Cauchy-Schwarz gives
$\|\hat g\|_{L^1}\le\sqrt\pi\,\|(1+|\zeta|^2)\hat g\|_{L^2}$. Thus $\hat g\in L^1$, the
Fourier inversion formula gives a bounded continuous representative, and
$\|g\|_{L^\infty}\le(2\pi)^{-2}\|\hat g\|_{L^1}\le\frac{\sqrt3}{2\sqrt\pi}\|g\|_{H^2}\le\|g\|_{H^2}$.

(E3) Put $P_k(\zeta)=\sum_{|\gamma|\le k}\zeta^{2\gamma}$. Every monomial
$\zeta^{2\gamma}$ with $|\gamma|\le k$ occurs in the expansion of
$(1+\zeta_1^2+\zeta_2^2)^k$ exactly once, with a multinomial coefficient between
$1$ and $3^k$; hence $P_k\le(1+|\zeta|^2)^k\le3^kP_k$. By Plancherel's theorem and
Cauchy-Schwarz,
\[
\|g\|_{H^k}^2\le\frac{\int(1+|\zeta|^2)^k|\hat g|^2\dd\zeta}{(2\pi)^2}
\le\frac{\|\hat g\|_{L^2}\|(1+|\zeta|^2)^{k}\hat g\|_{L^2}}{(2\pi)^2}
\le3^k\|g\|_{L^2}\|g\|_{H^{2k}} .\qedhere
\]
\end{proof}

In Section~\ref{sec:single} the perturbation of a ring is supported in a fixed
annulus, and the field that transports it is controlled only on an open
neighborhood of that annulus. The following calculus in $H^k$ is stated for
such data: $w$ is supported in a compact set $K$, and the field $U$ need only be
given on an open set $\mathcal O\supset K$.

\begin{lemma}
\label{lem:Hk-calculus}
Let $k\ge3$ be an integer, $K\subset\R^2$ compact, $\mathcal O\supset K$ open and
bounded, and $\chi_K\in C_c^\infty(\mathcal O)$ with $\chi_K=1$ on a neighborhood of
$K$. There is $C=C(k,\chi_K)$ such that:
\begin{enumerate}[label=\textup{(\arabic*)}]
\item for $w\in H^{k+1}(\R^2)$ with $\supp w\subset K$, a vector field
  $U\in W^{1,\infty}(\mathcal O)\cap H^k(\mathcal O)$, and every multi-index $\alpha$ with
  $|\alpha|\le k$,
  \[
  \|[\partial^\alpha,U\cdot\nabla]w\|_{L^2}\le C\big(\|U\|_{W^{1,\infty}(\mathcal O)}+\|U\|_{H^k(\mathcal O)}\big)\|w\|_{H^k};
  \]
  if moreover $U$ is real and $\divop U=0$ on $\mathcal O$, then
  $|\Re\langle\partial^\alpha(U\cdot\nabla w),\partial^\alpha w\rangle_{L^2}|\le
  C\big(\|U\|_{W^{1,\infty}(\mathcal O)}+\|U\|_{H^k(\mathcal O)}\big)\|w\|_{H^k}^2$, and for
  real $w$ the inner product is real;
\item $\|gh\|_{H^j}\le C\|g\|_{H^j}\|h\|_{H^j}$ for every integer $j$ with
  $2\le j\le k$ and all $g,h\in H^j(\R^2)$;
\item for every integer $j\ge0$, every $\rho>0$ and every $\chi\in C_c^\infty(B_\rho)$
  there is $C'=C'(j,\chi,\rho)$ such that
  $\|\chi\,\KBS\ast w\|_{H^{j+1}}\le C'\|w\|_{H^j}$ for every $w\in H^j(\R^2)$ with
  $\supp w\subset B_\rho$.
\end{enumerate}
Here $[\partial^\alpha,U\cdot\nabla]w=\partial^\alpha(U\cdot\nabla w)-U\cdot\nabla\partial^\alpha w$,
with both terms extended by $0$ outside $\mathcal O$.
\end{lemma}

\begin{proof}
We prove (2) first. Let $g,h\in C_c^\infty(\R^2)$, $|\gamma|\le j$ and
$\beta\le\gamma$. If $\beta=0$ or $\beta=\gamma$, then
$\|\partial^\beta g\,\partial^{\gamma-\beta}h\|_{L^2}\le\|g\|_{H^j}\|h\|_{H^j}$ by
Lemma~\ref{lem:embed}\textup{(E2)}, since $j\ge2$. If
$1\le|\beta|\le|\gamma|-1$, then by H\"older's inequality and
Lemma~\ref{lem:embed}\textup{(E1)},
$\|\partial^\beta g\,\partial^{\gamma-\beta}h\|_{L^2}\le\|\partial^\beta g\|_{H^1}\|\partial^{\gamma-\beta}h\|_{H^1}
\le\|g\|_{H^{|\beta|+1}}\|h\|_{H^{|\gamma|-|\beta|+1}}\le\|g\|_{H^j}\|h\|_{H^j}$.
The Leibniz formula
$\partial^\gamma(gh)=\sum_{\beta\le\gamma}\binom\gamma\beta\partial^\beta g\,\partial^{\gamma-\beta}h$
and $\sum_{\beta\le\gamma}\binom\gamma\beta=2^{|\gamma|}$ give (2) for smooth $g,h$,
with $C=C(k)$. For $g,h\in H^j$ take $g_n,h_n\in C_c^\infty$ converging to
$g,h$ in $H^j$; by the bound just proved, $g_nh_n$ is Cauchy in $H^j$, and by
(E2) it converges to $gh$ in $L^2$, so $gh\in H^j$ and (2) holds.

We now prove (1). Put $\tilde U=\chi_KU$, extended by $0$ outside $\mathcal O$; by the
Leibniz rule, $\|\tilde U\|_{W^{1,\infty}(\R^2)}\le C\|U\|_{W^{1,\infty}(\mathcal O)}$ and
$\|\tilde U\|_{H^k(\R^2)}\le C\|U\|_{H^k(\mathcal O)}$, with $C=C(k,\chi_K)$. Since
$\chi_K=1$ on a neighborhood $V$ of $K$ and $w=0$ outside $K$,
$[\partial^\alpha,U\cdot\nabla]w=[\partial^\alpha,\tilde U\cdot\nabla]w$, and
\[
[\partial^\alpha,\tilde U\cdot\nabla]w=\sum_{0<\beta\le\alpha}\binom\alpha\beta\,
\partial^\beta\tilde U\cdot\nabla\partial^{\alpha-\beta}w :
\]
with $\tilde U_n\in C_c^\infty$ in place of
$\tilde U$, this identity holds in $\mathcal D'$ by the Leibniz formula. Since
$\tilde U\in H^k(\R^2)$ by the bound above, there are such $\tilde U_n$ with
$\tilde U_n\to\tilde U$ in $H^k$. As $w\in H^{k+1}$, we have
$\nabla\partial^{\alpha-\beta}w\in L^2$ for $\beta\le\alpha$; so each product in
the identity, as
well as $\tilde U_n\cdot\nabla w$ and $\tilde U_n\cdot\nabla\partial^\alpha w$,
converges in $L^1$, and the identity passes to the limit. We bound the terms on
the right according to $|\beta|$. Write each of them as
$\sum_l\partial^\beta\tilde U_l\,\partial_l\partial^{\alpha-\beta}w$. If $|\beta|=1$,
its $L^2$ norm is at most $2\|\nabla\tilde U\|_{L^\infty}\|w\|_{H^k}$. If
$|\beta|=|\alpha|\ge2$, it is at most
$\sum_l\|\partial^\alpha\tilde U_l\|_{L^2}\|\partial_lw\|_{L^\infty}\le2\|\tilde U\|_{H^k}\|w\|_{H^3}
\le2\|\tilde U\|_{H^k}\|w\|_{H^k}$, by (E2) and $k\ge3$. If
$2\le|\beta|\le|\alpha|-1$, then by H\"older's inequality and (E1) it is at most
$\sum_l\|\partial^\beta\tilde U_l\|_{H^1}\|\partial_l\partial^{\alpha-\beta}w\|_{H^1}
\le2\|\tilde U\|_{H^{|\beta|+1}}\|w\|_{H^{|\alpha|-|\beta|+2}}\le2\|\tilde U\|_{H^k}\|w\|_{H^k}$,
because $|\beta|+1\le k$ and $|\alpha|-|\beta|+2\le k$. Summing over $\beta$ gives
the commutator bound. For the second claim, put $g=\partial^\alpha w\in H^1$,
which vanishes outside $K$. Then
$\partial^\alpha(U\cdot\nabla w)=[\partial^\alpha,\tilde U\cdot\nabla]w+\tilde U\cdot\nabla g$
and
$\Re\langle\tilde U\cdot\nabla g,g\rangle_{L^2}=\frac12\int\tilde U\cdot\nabla|g|^2\dd x
=-\frac12\int(\divop\tilde U)|g|^2\dd x=0$:
the integration by parts holds because $|g|^2\in W^{1,1}$ and $\tilde U$ is
Lipschitz with compact support, and the last integral vanishes because
$\divop\tilde U=\chi_K\divop U+U\cdot\nabla\chi_K=0$ on $V$ and $g=0$ outside $K$.
Hence the real part is at most
$\|[\partial^\alpha,\tilde U\cdot\nabla]w\|_{L^2}\|g\|_{L^2}$, by the
Cauchy-Schwarz inequality, and the commutator bound gives the claim, because
the $L^2$ norm of $g$ is at most the $H^k$ norm of $w$.

We now prove (3). Let first $w\in C_c^\infty(B_\rho)$ and $u=\KBS\ast w$. By
Lemma~I.\ref{I-lem:BS-basic}(d), and (c) with $R=\rho$,
$\|\nabla u\|_{H^j}=\|w\|_{H^j}$ and $\|u\|_{L^2(B_\rho)}\le2\rho\|w\|_{L^2}$.
For $|\gamma|\le j+1$, in the Leibniz formula
$\partial^\gamma(\chi u)=\sum_{\beta\le\gamma}\binom\gamma\beta\partial^{\gamma-\beta}\chi\,\partial^\beta u$
the term $\beta=0$ has $L^2$ norm at most
$\|\partial^\gamma\chi\|_{L^\infty}\|u\|_{L^2(B_\rho)}$, because $\chi$ vanishes
outside $B_\rho$, and every other term at most
$\|\partial^{\gamma-\beta}\chi\|_{L^\infty}\|\nabla u\|_{H^j}$. This gives (3) for
smooth $w$. For general $w$, $\supp w$ is a compact subset of $B_\rho$, so
mollification gives $w_n\in C_c^\infty(B_\rho)$ with $w_n\to w$ in $H^j$; then
$\chi\KBS\ast w_n$ is Cauchy in $H^{j+1}$ and, by Lemma~I.\ref{I-lem:BS-basic}(c),
converges to $\chi\KBS\ast w$ in $L^2$, so (3) holds for $w$.
\end{proof}

\section{Nonlinear Estimates for One Ring}
\label{sec:single}

In this section we study, in ring units, one ring of the approximate solutions: a
ring close to $\Phi$, whose deviation from $\Phi$ is driven by its own
nonlinearity and by the field $h$ of the vorticity inside and outside it, apart
from the point-vortex part of the inner vorticity, which enters the linearization
$\Lin(G)$ (Lemma~\ref{lem:reduction}). We prove three results. A ring that the
instability has not yet reached stays close to $\Phi$, and its deviation equals
the linear response to $h$ up to an error of second order
(Theorem~\ref{thm:quiet}). A ring whose component along
the unstable eigenfunction grows follows that eigenfunction until the envelope of
this component reaches a small size (Theorem~\ref{thm:active}). And the innermost ring of
the approximate datum can be given a small component along the unstable
eigenfunction by a rearrangement of $\Phi$ (Lemma~\ref{lem:seed}).
Sections~\ref{sec:relay} and~\ref{sec:bootstrap} apply Theorems~\ref{thm:quiet}
and~\ref{thm:active} to every ring, in its own units, and Section~\ref{sec:kick} also uses
Theorem~\ref{thm:quiet}. Lemma~\ref{lem:seed} gives the innermost ring of the
datum of Definition~\ref{def:approx}.

Throughout this section we assume \textup{(N1)}--\textup{(N3)}, we work in ring
units, and we use the objects of Definition~I.\ref{I-def:eigendata}: the unstable
eigenvalue $\lamst(G)$ of $\Lin(G)$ in the mode $3$ and its eigenfunction
$\phist(G)$, normalized in $L^2(\R^2)$; the functional $\ellst(G)$ that gives the
coefficient of a function along $\phist(G)$; the Riesz projection $\Riesz(G)$
onto the span of $\phist(G)$ and $\overline{\phist(G)}$; and the growth rates
$\sigma(G)=\Re\lamst(G)$ and $\sigb=\min_{G\in[0,\Gone]}\sigma(G)$. When $G$ is fixed
we write $\lamst=\lamst(G)$, $\phist=\phist(G)$, $\ellst=\ellst(G)$,
$\Riesz=\Riesz(G)$. The proofs of
Theorems~\ref{thm:quiet} and~\ref{thm:active} follow Grenier \cite{Grenier}: an
approximate solution of high order is built from the linear semigroup, and the order is so large that a crude energy estimate for the
remainder suffices.

We use the norms of $H^k$, $W^{1,\infty}$ and $C^k(N')$ fixed in
Section~I.\ref{I-sec:prelim}, and the embedding $\|f\|_{L^\infty}\le\|f\|_{H^2}$ of
Lemma~\ref{lem:embed}\textup{(E2)}. A vector field $h$ on a neighborhood of $N'$
is \textit{$3$-fold rotationally symmetric} if $h(\Orot x)=\Orot h(x)$.

The equation for the deviation of a ring from $\Phi$
(Lemma~\ref{lem:reduction}\textup{(3)} below) assumes that the ring stays in
the annulus $N$, and the continuity arguments of this section keep it there by
bounding how far its material moves radially. Let $\flow$ be the flow
of a solution of Theorem~\ref{thm:smooth}, let $\omega_{\mathrm r}(0)$ have compact
support, let $\omega_{\mathrm r}(t)=\omega_{\mathrm r}(0)\circ\flow_t^{-1}$, and let
$\delta_{\mathrm r}(t)$ be the supremum of $\big||\flow_\tau(y)|-|y|\big|$ over
$\tau\in[0,t]$ and $y\in\supp\omega_{\mathrm r}(0)$ ($0$ if the support is empty);
it is a maximum if the support is nonempty, it is continuous and
nondecreasing, and $\delta_{\mathrm r}(0)=0$, because $(\tau,y)\mapsto\flow_\tau(y)$
is continuous. We say that \textit{every trajectory of $\omega_{\mathrm r}$ moves
radially by at most $\delta$ on $[0,t]$} if $\delta_{\mathrm r}(t)\le\delta$. Since
$\supp\omega_{\mathrm r}(\tau)=\flow_\tau(\supp\omega_{\mathrm r}(0))$, the support
of $\omega_{\mathrm r}(\tau)$, $\tau\le t$, then lies within radial distance $\delta$
of $\supp\omega_{\mathrm r}(0)$.

For $G\in[0,\Gone]$ and an integer $k\ge0$, $e^{t\Lin(G)}$ denotes the semigroup
of Lemma~I.\ref{I-lem:semigroup-construct} for the ring in ring units. It preserves
each mode, hence the subspace $X^k$ of the space $Y^k$ of that lemma
consisting of the functions whose mode components vanish outside $3\Z$. It maps
real-valued functions to real-valued functions, since it is the sum of a
Dyson-Phillips series built from the two operators of that lemma, which both do:
the transport group $e^{t\Lin^{\rm tr}}$, which rotates each circle $\{|y|=r\}$ by
the angle $\Omega(r;G)t$, and the coupling operator $\mathcal K$, here
$\mathcal Kw=-(\KBS\ast w)\cdot\nabla\Phi$. The approximate solutions of this
section are Duhamel integrals of this semigroup. The next lemma shows that they are
solutions of the forced linear equation, differentiable in time as functions with
values in $X^{k-1}$ when the data lie in $X^k$.

\begin{lemma}
\label{lem:duhamel}
Let $G\in[0,\Gone]$, $T>0$, and let $k\ge1$ be an integer. Let $f_0\in X^k$ and
$g\in C([0,T];X^k)$, and put
$u(t)=e^{t\Lin(G)}f_0+\int_0^te^{(t-\tau)\Lin(G)}g(\tau)\dd\tau$. Then
$u\in C([0,T];X^k)\cap C^1([0,T];X^{k-1})$, and $\partial_tu=\Lin(G)u+g$ on
$[0,T]$. If $f_0$ and $g$ are real-valued, so is $u$.
\end{lemma}
\begin{proof}
We first prove the continuity of $u$ in $X^k$, and then its differentiability in
$X^{k-1}$, where the generator of the semigroup acts on $X^k$ as $\Lin(G)$.
The semigroup $S(t)=e^{t\Lin(G)}$ is strongly continuous on $X^k$ and, by the
uniform boundedness principle, bounded on $[0,T]$. So $\tau\mapsto S(t-\tau)g(\tau)$
is continuous into $X^k$: the argument given below for the continuity of
$(s,\tau)\mapsto S(s-\tau)H(\tau)$ in $X^{k-1}$ uses only the strong continuity of
$S$ and a bound on $\|S(t)\|$ for $t$ in a compact interval, and both hold in
$X^k$. Hence
$\int_0^tS(t-\tau)g(\tau)\dd\tau=\int_0^tS(\tau)g(t-\tau)\dd\tau$
is continuous in $t$ by dominated convergence, and $u\in C([0,T];X^k)$. Moreover
$X^k\hookrightarrow X^{k-1}$, and $S$ is a $C_0$-semigroup on $X^{k-1}$, since it
is one on $Y^{k-1}$ and $X^{k-1}=Y^{k-1}_{3\Z}$ is a closed invariant subspace,
both by Lemma~I.\ref{I-lem:semigroup-construct}\textup{(2)}. By
Lemma~I.\ref{I-lem:semigroup-construct}\textup{(2)}, the generator $A$ of $S$ on
$X^{k-1}$ (the part in $X^{k-1}$ of the generator on $Y^{k-1}$ of that lemma)
satisfies $X^k\subset D(A)$, $Ay=\Lin(G)y$ and $\|Ay\|_{X^{k-1}}\le C\|y\|_{X^k}$ for
$y\in X^k$; here $\Lin(G)$ maps $X^k$ into $X^{k-1}$ because it preserves modes
(Lemma~I.\ref{I-lem:sector}\textup{(2)}). Below, $\|\cdot\|$ is the norm of $X^{k-1}$ and
the operator norm on $X^{k-1}$.

All integrals below are Riemann integrals of continuous functions with values
in a Banach space; Fubini's theorem holds for them, because it holds after
applying any continuous linear functional. We extend $g$ by $g(t)=g(T)$ for
$T<t\le2T$. Since each $t\mapsto S(t)x$, $x\in X^{k-1}$, is continuous, hence bounded on $[0,2T]$, the
uniform boundedness principle gives $M_T:=\sup_{0\le t\le2T}\|S(t)\|<\infty$.

Let $x\in D(A)$. For $t\ge0$,
$\epsilon^{-1}(S(\epsilon)-I)S(t)x=S(t)\epsilon^{-1}(S(\epsilon)x-x)\to S(t)Ax$, so $S(t)x\in D(A)$,
$AS(t)x=S(t)Ax$, and $S(t)Ax$ is the right derivative of $t\mapsto S(t)x$. For
$0<\epsilon\le t\le2T$,
$\epsilon^{-1}(S(t)x-S(t-\epsilon)x)-S(t)Ax=S(t-\epsilon)\big(\epsilon^{-1}(S(\epsilon)x-x)-Ax\big)+\big(S(t-\epsilon)-S(t)\big)Ax\to0$,
because $\|S(t-\epsilon)\|\le M_T$. Thus $t\mapsto S(t)x$ lies in $C^1([0,2T];X^{k-1})$ with
derivative $S(t)Ax$, and
\begin{equation}
\label{eq:semigroup-FTC}
S(r)x-x=\int_0^rS(s)Ax\dd s\qquad(x\in D(A),\ 0\le r\le2T).
\end{equation}

Let $H\in C([0,2T];X^{k-1})$. The identity
$S(s'-\tau')H(\tau')-S(s-\tau)H(\tau)=S(s'-\tau')(H(\tau')-H(\tau))
+(S(s'-\tau')-S(s-\tau))H(\tau)$, whose first term has norm at most
$M_T\|H(\tau')-H(\tau)\|$ and whose second tends to $0$ by strong continuity,
shows that $(s,\tau)\mapsto S(s-\tau)H(\tau)$ is continuous on
$\{0\le\tau\le s\le2T\}$. Moreover
$s\mapsto\int_0^sS(s-\tau)H(\tau)\dd\tau$ is continuous on $[0,2T]$: for
$0\le s<s'\le2T$ the difference of its values is
$\int_s^{s'}S(s'-\tau)H(\tau)\dd\tau+\int_0^sS(s-\tau)(S(s'-s)-I)H(\tau)\dd\tau$,
of norm at most $M_T\sup\|H\|(s'-s)+2TM_T\sup_\tau\|(S(s'-s)-I)H(\tau)\|$; and the
operators $S(\epsilon)-I$, $0\le\epsilon\le2T$, are bounded by $M_T+1$ and tend strongly to
$0$ as $\epsilon\to0$, hence uniformly on the compact set $H([0,2T])$.

Now put $V(t)=\int_0^tS(t-\tau)g(\tau)\dd\tau$ and
$W(t)=\int_0^tS(t-\tau)Ag(\tau)\dd\tau$. Since
$\|A(g(\tau)-g(\tau'))\|\le C\|g(\tau)-g(\tau')\|_{X^k}$, $Ag$ is continuous
into $X^{k-1}$, so $W$ is continuous by the previous paragraph with $H=Ag$. By
Equation~\eqref{eq:semigroup-FTC} with $x=g(\tau)\in X^k\subset D(A)$ and Fubini's theorem,
\[
V(t)=\int_0^t\Big(g(\tau)+\int_\tau^tS(s-\tau)Ag(\tau)\dd s\Big)\dd\tau
=\int_0^t\big(g(s)+W(s)\big)\dd s ,
\]
so $V\in C^1([0,2T];X^{k-1})$ and $V'=g+W$. For $t\in[0,T]$ and $0<\epsilon\le T$ we have
$V(t+\epsilon)=S(\epsilon)V(t)+\int_t^{t+\epsilon}S(t+\epsilon-\tau)g(\tau)\dd\tau$, and $\epsilon^{-1}$ times the
last integral tends to $g(t)$ as $\epsilon\to0$, by the continuity of
$(s,\tau)\mapsto S(s-\tau)g(\tau)$ at $(t,t)$. Hence
$\epsilon^{-1}(S(\epsilon)V(t)-V(t))\to V'(t)-g(t)=W(t)$; that is, $V(t)\in D(A)$ and
$AV(t)=W(t)$. With the case $x=f_0\in X^k\subset D(A)$ above, $u=S(\cdot)f_0+V$ lies in
$C^1([0,T];X^{k-1})$, $u(t)\in D(A)$ with $Au(t)=S(t)Af_0+W(t)$, and
$\partial_tu=S(\cdot)Af_0+g+W=Au+g=\Lin(G)u+g$, since $u(t)\in X^k$. Finally, $S(t)$
preserves real-valued functions.
\end{proof}

\subsection{The Perturbation Equation}

We split a smooth solution into a ring $\omega_{\mathrm r}$, the vorticity
$\omega_{\mathrm i}$ inside it and the vorticity $\omega_{\mathrm o}$ outside it,
each transported by the full flow, and write the Euler equations as an equation
for the deviation $w=\omega_{\mathrm r}-\Phi$ of the ring from $\Phi$. The
vorticity inside the ring enters the linearization $\Lin(G)$ as a point vortex at
the origin with its circulation $2\pi G$; the rest of its field, and the field of
the vorticity outside the ring, form the field $h$, which forces the equation.

\begin{lemma}
\label{lem:reduction}
Let $\omega$ be a solution given by Theorem~\ref{thm:smooth} whose datum
$\omega(0)$ is $3$-fold rotationally symmetric, and let $\flow_t$ be its flow. Let
$\omega(0)=\omega_{\mathrm r}(0)+\omega_{\mathrm i}(0)+\omega_{\mathrm o}(0)$ with
pairwise disjoint supports, $\omega_{\mathrm r}(0)$, $\omega_{\mathrm i}(0)$,
$\omega_{\mathrm o}(0)$ each $3$-fold rotationally symmetric, and put
$\omega_\bullet(t)=\omega_\bullet(0)\circ\flow_t^{-1}$ for
$\bullet\in\{\mathrm r,\mathrm i,\mathrm o\}$. Put
$G=\frac1{2\pi}\int\omega_{\mathrm i}(0)$. Let $T>0$ and assume that for every
$t\in[0,T]$, $\supp\omega_{\mathrm i}(t)\subset\overline B_{a/8}$ and
$\supp\omega_{\mathrm o}(t)\cap B_6=\emptyset$. Put
\[
h(t)=\KBS\ast\big(\omega_{\mathrm i}(t)+\omega_{\mathrm o}(t)\big)-G\,\frac{x^\perp}{|x|^2}
\qquad\text{on a neighborhood of }N'.
\]
Then:
\begin{enumerate}[label=\textup{(\arabic*)}]
\item for every $t\in[0,T]$, $h(t)$ is smooth, divergence-free and
  $3$-fold rotationally symmetric on a neighborhood of $N'$; and $t\mapsto h(t)\in C^k(N')$ is
  continuous on $[0,T]$ for every $k$;
\item $\omega_{\mathrm r}\in C^\infty(\R\times\R^2)$, and on $[0,T]\times N'$ the
  velocity $v=\KBS\ast\omega$ is
  \[
  v=\Omega(|x|;G)\,x^\perp+h+\KBS\ast(\omega_{\mathrm r}-\Phi),
  \]
  with $\Omega$ of Definition~I.\ref{I-def:sector-setting} for $(\Phi,G)$; consequently,
  if $t\in[0,T]$ and $\flow_\tau(y)\in N'$ for every $\tau\in[0,t]$, then
  $|\flow_t(y)|-|y|=\int_0^t\big((h+\KBS\ast(\omega_{\mathrm r}-\Phi))\cdot\er\big)(\tau,\flow_\tau(y))\dd\tau$;
\item if $\supp\omega_{\mathrm r}(t)\subset N$ for $t\in[0,T]$, then the function
  $w=\omega_{\mathrm r}-\Phi$ is real-valued, it lies in $C^1([0,T];X^k)$ for
  every $k$, and
  \begin{equation}
  \label{eq:ring}
  \partial_tw=\Lin(G)w+B(w,w)-h\cdot\nabla w-h\cdot\nabla\Phi,\qquad B(v,w)=-(\KBS\ast v)\cdot\nabla w,
  \end{equation}
  on $[0,T]\times\R^2$, where $h\cdot\nabla w$ and $h\cdot\nabla\Phi$ denote the
  functions equal to these products on $N'$ and to $0$ elsewhere.
\end{enumerate}
\end{lemma}
\begin{proof}
Put $K_\bullet=\supp\omega_\bullet(0)$. These compact sets are pairwise
disjoint, so each $\omega_\bullet(0)$ agrees with $\omega(0)$ near $K_\bullet$ and
vanishes near the other two; hence $\omega_\bullet(0)\in C_c^\infty(\R^2)$. The
map $(t,x)\mapsto\flow_t^{-1}(x)$ is the value at time $0$ of the solution of
$\dot z=v(\tau,z)$ with $z(t)=x$. Since $v$ is smooth and bounded on
$[-T',T']\times\R^2$ for every $T'>0$ (Theorem~\ref{thm:smooth}\textup{(2)}),
Lemma~\ref{lem:flow}\textup{(a)}, \textup{(b)} (Appendix~\ref{app:smooth}) shows that this map
is defined and
$C^\infty$ on $\R\times\R^2$. Hence $\omega_\bullet\in C^\infty(\R\times\R^2)$,
$\supp\omega_\bullet(t)=\flow_t(K_\bullet)$, and $\int\omega_{\mathrm i}(t)=2\pi G$
because $\flow_t$ preserves area. By Theorem~\ref{thm:smooth}\textup{(6)},
$\omega(t)$ is $3$-fold rotationally symmetric. Since $\KBS(\Orot x)=\Orot\KBS(x)$, we get
$(\KBS\ast g)(\Orot x)=\Orot(\KBS\ast g)(x)$ for every $3$-fold rotationally symmetric
$g\in (L^1\cap L^\infty)(\R^2)$; so $v(t,\Orot x)=\Orot v(t,x)$, uniqueness for the flow
gives $\flow_t(\Orot y)=\Orot\flow_t(y)$, and each $\omega_\bullet(t)$ is $3$-fold
rotationally symmetric.

\emph{Proof of (1).} By continuity of the flow on compact sets, and since
$\flow_t(K_{\mathrm i})=\supp\omega_{\mathrm i}(t)\subset\overline B_{a/8}$,
$\rho_{\mathrm i}:=\max\{|\flow_t(y)|:t\in[0,T],\,y\in K_{\mathrm i}\}\le a/8<a/4$
($\rho_{\mathrm i}:=0$ if $K_{\mathrm i}=\emptyset$), and $|\flow_t(y)|\ge6$ for
$t\in[0,T]$ and $y\in K_{\mathrm o}$. So the open annulus
$U=\{\frac12(\rho_{\mathrm i}+\frac a4)<|x|<4\}\supset N'$ has distance at least
$\min(\frac12(\frac a4-\rho_{\mathrm i}),2)>0$ from
$\supp(\omega_{\mathrm i}(t)+\omega_{\mathrm o}(t))$ for $t\in[0,T]$. Since
$|\partial^\beta\KBS(z)|\le C_\beta|z|^{-1-|\beta|}$, and $\omega_{\mathrm i}$,
$\omega_{\mathrm o}$ are continuous with supports in a fixed compact set for
$t\in[0,T]$ (Theorem~\ref{thm:smooth}\textup{(1)}), we may differentiate under the
integral sign: $h(t)\in C^\infty(U)$, every $\partial^\beta_xh$ is continuous on
$[0,T]\times U$, and in particular $t\mapsto h(t)\in C^k(N')$ is continuous. The
fields $x\mapsto\KBS(x-y)=\nabla^\perp_x\big(\frac1{2\pi}\log|x-y|\big)$ and
$x^\perp/|x|^2$ are divergence-free away from $y$ and from $0$; hence
so is $h(t)$ on $U$, and it is $3$-fold rotationally symmetric by the equivariance above and
$(\Orot x)^\perp=\Orot x^\perp$.

\emph{Proof of (2).} For a radial $g\in C_c^\infty(\R^2)$, $\KBS\ast g$ is
equivariant under every rotation, so $(\KBS\ast g)\cdot\er$ and
$(\KBS\ast g)\cdot\etheta$ are constant on each $C_r=\{|x|=r\}$, $r>0$. Fix $r>0$,
orient $C_r$ counterclockwise, and let $|y|\ne r$. Along $C_r$,
$\KBS(x-y)\cdot\er\,\dd s=-\frac1{2\pi}\dd\log|x-y|$ and
$\KBS(x-y)\cdot\etheta\,\dd s=\frac1{2\pi}\dd\arg(x-y)$; the first integrates to
$0$, the second to the winding number of $C_r$ around $y$, which is $1$ if
$|y|<r$ and $0$ if $|y|>r$. Integrating against $g(y)\dd y$ gives the values $0$
and $\Circ_g(r)/(2\pi r)$; Fubini's theorem applies because
$\oint_{C_r}|\KBS(x-y)|\dd s(x)$ is bounded by a constant, depending on $r$, times
$1+|\log\operatorname{dist}(y,C_r)|$, which is integrable in $y$ over compact sets,
and $\{|y|=r\}$ has measure $0$. Hence $\KBS\ast g=\Circ_g(|x|)\,x^\perp/(2\pi|x|^2)$.
With $g=\Phi$ and the definition of $h$, on $[0,T]\times U$,
\[
v=\KBS\ast\Phi+\KBS\ast(\omega_{\mathrm r}-\Phi)+\KBS\ast(\omega_{\mathrm i}+\omega_{\mathrm o})
=\Big(\frac{\Circ_\Phi(|x|)}{2\pi|x|^2}+\frac G{|x|^2}\Big)x^\perp+\KBS\ast(\omega_{\mathrm r}-\Phi)+h,
\]
and the bracket is $\Omega(|x|;G)$. Since $\partial_\tau\flow_\tau(y)=v(\tau,\flow_\tau(y))$
and $x^\perp\cdot\er=0$, we get $\frac{\dd}{\dd\tau}|\flow_\tau(y)|=
\big((h+\KBS\ast(\omega_{\mathrm r}-\Phi))\cdot\er\big)(\tau,\flow_\tau(y))$ while
$\flow_\tau(y)\in N'$, which gives the integral formula.

\emph{Proof of (3).} Since $\supp\Phi\subset\{a\le|y|\le1\}\subset N$
(Lemma~I.\ref{I-lem:ring-facts}), $w(t)$ is a smooth, real, $3$-fold rotationally symmetric function
supported in $N$. The derivatives $\partial^\alpha_xw$ and
$\partial_t\partial^\alpha_xw$ are uniformly continuous on the compact set
$[0,T]\times N$ and vanish outside it, so the difference quotients of
$t\mapsto w(t)$ converge in every $H^k$, and $w\in C^1([0,T];X^k)$.
Differentiating $\omega_{\mathrm r}(t,\flow_t(y))=\omega_{\mathrm r}(0,y)$ in $t$ gives
$\partial_t\omega_{\mathrm r}+v\cdot\nabla\omega_{\mathrm r}=0$. On $N$ we insert (2),
$x^\perp\cdot\nabla\Phi=\partial_\theta\Phi=0$ and $x^\perp\cdot\nabla w=\partial_\theta w$:
\[
\partial_tw=-\Omega(\cdot;G)\,\partial_\theta w-(\KBS\ast w)\cdot\nabla\Phi-(\KBS\ast w)\cdot\nabla w-h\cdot\nabla w-h\cdot\nabla\Phi .
\]
The first two terms are $\Lin(G)w$ by Definition~I.\ref{I-def:sector-setting}, and the
third is $B(w,w)$. Outside $N$ both sides of Equation~\eqref{eq:ring} vanish.
\end{proof}

In the setting of Lemma~\ref{lem:reduction}, put $w=\omega_{\mathrm r}-\Phi$. Then
$t\mapsto\|w(t)\|_{H^k}$ is continuous
on $[0,T]$ for every $k$, since $\omega_{\mathrm r}$ is smooth and its support stays
in a fixed compact set (Theorem~\ref{thm:smooth}\textup{(1)}). We fix once and for
all the auxiliary objects of the estimates below. The functions we estimate are
supported in $N$, and the velocities that act on them are needed only on a
neighborhood of $N$. We therefore multiply the Biot-Savart velocity by a compactly
supported cutoff $\chi_B$ that does not change it on $N'$, so that
Lemma~\ref{lem:Hk-calculus}\textup{(3)} applies, and we use a cutoff $\chi_N$ around $N$ in Lemma~\ref{lem:Hk-calculus}\textup{(1)}
and \textup{(2)}. Let $N_1=\{\frac a4<|y|<3\}$ be
the interior of $N'$. Fix $\chi_B\in C_c^\infty(B_4)$ with $\chi_B=1$ on $B_{7/2}$,
and $\chi_N\in C_c^\infty(N_1)$ with $\chi_N=1$ on a neighborhood of $N$. For each
integer $j\ge0$ fix a constant $C^{\mathrm{el}}_j$ of
Lemma~\ref{lem:Hk-calculus}\textup{(3)} for $j$, $\chi_B$ and $\rho=4$; it depends
only on $j$. For each integer $k\ge2$ fix a constant $C_P(k)$ of
Lemma~\ref{lem:Hk-calculus} for $\max(k,3)$, the compact set $N$, the open set
$N_1$ and $\chi_N$; it depends only on $k$ and $a$.

\begin{lemma}
\label{lem:bilinear}
Let $k\ge2$ be an integer. Put $C_B=C_B(k)=2C_P(k)C^{\mathrm{el}}_{k-1}$ and
$C_{\mathrm{BS}}=2C^{\mathrm{el}}_1$, and let $C_h=C_h(k)$ be the supremum of
$\|h\cdot\nabla w\|_{H^k}$ over all $w\in X^{k+1}$ and
$h\in C^k(N';\R^2)$ with $\|w\|_{H^{k+1}}\le1$ and $\|h\|_{C^k(N')}\le1$. Then
$C_h$ is finite and depends only on $k$ and $a$, and for $v,w\in X^{k+1}$ and
every vector field $h\in C^k(N';\R^2)$:
\begin{enumerate}[label=\textup{(\arabic*)}]
\item $\|B(v,w)\|_{H^k}\le C_B\|v\|_{H^k}\|w\|_{H^{k+1}}$;
\item $\|h\cdot\nabla w\|_{H^k}\le C_h\|h\|_{C^k(N')}\|w\|_{H^{k+1}}$;
\item $\|h\cdot\nabla\Phi\|_{H^k}\le\Cnab\|h\|_{C^k(N')}$, with $\Cnab=\Cnab(k)$ of
  Lemma~I.\ref{I-lem:uniform};
\item $\sup_{|x|\le3}|(\KBS\ast v)(x)|\le C_{\mathrm{BS}}\|v\|_{H^1}$.
\end{enumerate}
The functions in \textup{(1)}--\textup{(3)} are supported in $N$, and they lie in
$X^k$ if $h$ is $3$-fold rotationally symmetric.
\end{lemma}
\begin{proof}
Since $\nabla w$ is supported in $N$,
$B(v,w)=-\sum_{i=1,2}(\chi_B\KBS\ast v)_i\,\partial_iw$. By
Lemma~\ref{lem:Hk-calculus}\textup{(2)} (with $\max(k,3)$ in place of its $k$, and
$j=k$) and \textup{(3)},
\[
\|B(v,w)\|_{H^k}\le2C_P(k)\|\chi_B\KBS\ast v\|_{H^k}\max_i\|\partial_iw\|_{H^k}\le2C_P(k)C^{\mathrm{el}}_{k-1}\|v\|_{H^{k-1}}\|w\|_{H^{k+1}},
\]
which is (1). Next, $h\cdot\nabla w=(\chi_Nh)\cdot\nabla w$, and the Leibniz rule
$\partial^\alpha\big((\chi_Nh)\cdot\nabla w\big)=\sum_{\beta\le\alpha}\binom\alpha\beta
\partial^\beta(\chi_Nh)\cdot\partial^{\alpha-\beta}\nabla w$ with
$|\partial^\beta(\chi_Nh)|\le2^{|\beta|}\|\chi_N\|_{C^k}\|h\|_{C^k(N')}$ bounds
$\|h\cdot\nabla w\|_{H^k}$ by a constant depending only on $k$ and $\chi_N$, hence
on $k$ and $a$, times $\|h\|_{C^k(N')}\|w\|_{H^{k+1}}$. So $C_h$ is finite, depends
only on $k$ and $a$, and (2) holds by homogeneity in $h$ and in $w$. Item (3) is the
last inequality of Lemma~I.\ref{I-lem:uniform}. For (4), $\KBS\ast v=\chi_B\KBS\ast v$ on
$B_{7/2}$, and Lemma~\ref{lem:embed}\textup{(E2)}, applied to each component, and
Lemma~\ref{lem:Hk-calculus}\textup{(3)} give
$\|\chi_B\KBS\ast v\|_{L^\infty}\le2C^{\mathrm{el}}_1\|v\|_{H^1}=C_{\mathrm{BS}}\|v\|_{H^1}$,
which gives (4), since $\{|x|\le3\}\subset B_{7/2}$. The functions in
(1)--(3) vanish outside $\supp\nabla w\cup\supp\nabla\Phi\subset N$. If $v$, $w$ and
$h$ are $3$-fold rotationally symmetric, so are the vector fields $\KBS\ast v$ (by the
equivariance in the proof of Lemma~\ref{lem:reduction}), $\nabla w$ and
$\nabla\Phi$, hence the three scalar products.
\end{proof}

Lemma~\ref{lem:energy} below compares the deviation $w$ of a ring with an
approximate solution $w_{\mathrm{app}}$ by an energy estimate in $H^s$. Its rate
$\Lam(s)$, defined next, bounds the transport of $w-w_{\mathrm{app}}$ by the
velocity and the coupling through $\nabla(\Phi+w_{\mathrm{app}})$; the entry
$2\max_{G\in[0,\Gone]}\sigma(G)$ in its definition makes $\Lam(s)\ge2\sigma(G)$ for every $G$,
which the proof of Lemma~\ref{lem:window} uses.
This rate involves only norms on $N'$. Let $s\ge3$ be
an integer, let $n_s$ be the number of multi-indices $\alpha\in\N_0^2$ with
$|\alpha|\le s$, and put $C_{\mathrm{HK}}=C_P(s)$. The function
$\Omega(r;G)=\Omega(r;0)+Gr^{-2}$ is smooth on $[\frac a4,3]$ and affine in $G$, so
\[
C_\Omega:=\sup_{G\in[0,\Gone]}\Big(\|\Omega(|x|;G)x^\perp\|_{W^{1,\infty}(N_1)}+\|\Omega(|x|;G)x^\perp\|_{H^s(N_1)}\Big)<\infty .
\]
Let $C_{N'}$ be the least constant such that
$\|h\|_{W^{1,\infty}(N_1)}+\|h\|_{H^s(N_1)}\le C_{N'}\|h\|_{C^s(N')}$ for every
$h\in C^s(N';\R^2)$; it is finite and depends only on $s$ and $a$. Put
$A_s=C_\Omega+C_{N'}+C^{\mathrm{el}}_s+6C^{\mathrm{el}}_2$ and
\begin{equation}
\label{eq:Lambda}
\Lam(s)=\max\Big(2\max_{G\in[0,\Gone]}\sigma(G),\ n_sC_{\mathrm{HK}}A_s+2C_{\mathrm{HK}}C^{\mathrm{el}}_{s-1}\big(\|\Phi\|_{H^{s+1}}+1\big)\Big).
\end{equation}

\begin{lemma}
\label{lem:energy}
Let $s\ge3$ be an integer, and let $\Lam(s)$ be the number of
Equation~\eqref{eq:Lambda}; it satisfies $\Lam(s)\ge2\max_{G\in[0,\Gone]}\sigma(G)$,
where the maximum exists by Proposition~I.\ref{I-prop:U}\textup{(a)}, and it depends
only on $\Phi$, $\Gone$ and $s$. Let $G\in[0,\Gone]$
and $t>0$, and let the assumptions of Lemma~\ref{lem:reduction}\textup{(3)} hold on
$[0,t]$, with $w=\omega_{\mathrm r}-\Phi$. Let
$w_{\mathrm{app}}\in C([0,t];X^{s+1})\cap C^1([0,t];X^s)$ be real-valued, and put
$R_{\mathrm{app}}=\partial_tw_{\mathrm{app}}-\big(\Lin(G)w_{\mathrm{app}}+B(w_{\mathrm{app}},w_{\mathrm{app}})-h\cdot\nabla w_{\mathrm{app}}-h\cdot\nabla\Phi\big)$
and $v=w-w_{\mathrm{app}}$. If on $[0,t]$ we have $\|w\|_{H^s}\le1$,
$\|w_{\mathrm{app}}\|_{H^{s+1}}\le1$ and $\|h\|_{C^{s+1}(N')}\le1$, then for
$\tau\in[0,t]$,
\[
\|v(\tau)\|_{H^s}\le e^{\Lam(s)\tau}\|v(0)\|_{H^s}+\int_0^\tau e^{\Lam(s)(\tau-\tau')}\|R_{\mathrm{app}}(\tau')\|_{H^s}\dd\tau' .
\]
\end{lemma}
\begin{proof}
On $N_1$ put $U=\Omega(|x|;G)x^\perp+h+\KBS\ast w$, and for $f$ supported in $N$ let
$U\cdot\nabla f$ mean $(\chi_NU)\cdot\nabla f$. By
Lemma~\ref{lem:reduction}\textup{(3)},
$\partial_tw+U\cdot\nabla w=-(\KBS\ast w)\cdot\nabla\Phi-h\cdot\nabla\Phi$, and by
the definition of $R_{\mathrm{app}}$ the same identity holds for $w_{\mathrm{app}}$,
with $\KBS\ast w_{\mathrm{app}}$ in place of $\KBS\ast w$ in $U$ and with
$R_{\mathrm{app}}$ added on the right. Since
$(\KBS\ast w)\cdot\nabla w-(\KBS\ast w_{\mathrm{app}})\cdot\nabla w_{\mathrm{app}}=(\KBS\ast w)\cdot\nabla v+(\KBS\ast v)\cdot\nabla w_{\mathrm{app}}$,
subtraction gives
\[
\partial_tv+U\cdot\nabla v=-(\KBS\ast v)\cdot\nabla(\Phi+w_{\mathrm{app}})-R_{\mathrm{app}},
\]
where every term lies in $C([0,t];H^s)$, and
$v\in C^1([0,t];H^s)\cap C([0,t];H^{s+1})$ is real and supported in $N$.

The field $U$ is divergence-free on $N_1$: $\divop h=0$ by
Lemma~\ref{lem:reduction}\textup{(1)}, $\divop(\KBS\ast w)=0$, and
$\divop(f(|x|)x^\perp)=f'(|x|)\frac{x}{|x|}\cdot x^\perp=0$. By
Lemma~\ref{lem:Hk-calculus}\textup{(3)} and Lemma~\ref{lem:embed}\textup{(E2)},
applied to each entry, and $\|w\|_{H^s}\le1$,
$\|\KBS\ast w\|_{H^s(N_1)}\le\|\chi_B\KBS\ast w\|_{H^{s+1}}\le C^{\mathrm{el}}_s$ and
$\|\KBS\ast w\|_{W^{1,\infty}(N_1)}\le6\|\chi_B\KBS\ast w\|_{H^3}\le6C^{\mathrm{el}}_2$.
With the definitions of $C_\Omega$ and $C_{N'}$ and $\|h\|_{C^s(N')}\le1$, this gives
$\|U\|_{W^{1,\infty}(N_1)}+\|U\|_{H^s(N_1)}\le A_s$. Since $U$ and $v$ are real and
$\divop U=0$ on $N_1$, Lemma~\ref{lem:Hk-calculus}\textup{(1)} gives
$|\langle\partial^\alpha(U\cdot\nabla v),\partial^\alpha v\rangle|\le C_{\mathrm{HK}}A_s\|v\|_{H^s}^2$
for $|\alpha|\le s$. Since $\nabla(\Phi+w_{\mathrm{app}})$ is supported in $N$,
Lemma~\ref{lem:Hk-calculus}\textup{(2)} and \textup{(3)} give
\[
\begin{aligned}
\|(\KBS\ast v)\cdot\nabla(\Phi+w_{\mathrm{app}})\|_{H^s}
&\le2C_{\mathrm{HK}}\|\chi_B\KBS\ast v\|_{H^s}\big(\|\Phi\|_{H^{s+1}}+\|w_{\mathrm{app}}\|_{H^{s+1}}\big)\\
&\le2C_{\mathrm{HK}}C^{\mathrm{el}}_{s-1}\big(\|\Phi\|_{H^{s+1}}+1\big)\|v\|_{H^s}.
\end{aligned}
\]
We apply $\partial^\alpha$ to the equation of $v$, take the $L^2$ inner product with
$\partial^\alpha v$, and sum over $|\alpha|\le s$. The Cauchy-Schwarz inequality and
Equation~\eqref{eq:Lambda} give
$\frac12\frac{\dd}{\dd\tau}\|v\|_{H^s}^2\le\Lam(s)\|v\|_{H^s}^2+\|R_{\mathrm{app}}\|_{H^s}\|v\|_{H^s}$.
To divide by $\|v\|_{H^s}$, which may vanish, we regularize it. For $\epsilon>0$
the function $y_\epsilon=(\|v\|_{H^s}^2+\epsilon^2)^{1/2}$ is $C^1$, as
$v\in C^1([0,t];H^s)$ and $y_\epsilon\ge\epsilon$. Since
$y_\epsilon y_\epsilon'=\frac12\frac{\dd}{\dd\tau}\|v\|_{H^s}^2$ and
$\|v\|_{H^s}\le y_\epsilon$, it satisfies
$y_\epsilon'\le\Lam(s)y_\epsilon+\|R_{\mathrm{app}}\|_{H^s}$, so
$(e^{-\Lam(s)\tau}y_\epsilon)'\le e^{-\Lam(s)\tau}\|R_{\mathrm{app}}\|_{H^s}$.
Integrating from $0$ to $\tau$, and letting $\epsilon\to0$ with
$\|v\|_{H^s}\le y_\epsilon$ and $y_\epsilon(0)\to\|v(0)\|_{H^s}$, proves the lemma.
\end{proof}

\emph{Approximate solutions.} Following Grenier, we approximate the deviation $w$
by the first $J$ terms $w^{[1]},\dots,w^{[J]}$ of its formal expansion in powers of
the size of the data, in which $h$ and the initial value count as terms of first
order: $w^{[1]}$ solves the linear equation forced by $-h\cdot\nabla\Phi$, with an
initial value $w_0$ that is $0$ or $w(0)$ in our applications, and $w^{[j]}$,
$j\ge2$, solves the linear equation with initial value $0$, forced by the terms of
order $j$ of $B(w,w)-h\cdot\nabla w$. The
residual of the sum then consists of terms of order larger than $J$
(Equation~\eqref{eq:hier-res}), and Lemma~\ref{lem:energy} controls the difference
between $w$ and the sum. Let the assumptions of Lemma~\ref{lem:reduction} hold
on $[0,T]$, assume $G\in[0,\Gone]$, and put $\Lin=\Lin(G)$ and $F=-h\cdot\nabla\Phi$.
By Lemma~\ref{lem:reduction}\textup{(1)} and Lemma~\ref{lem:bilinear}, $F$ is real
and lies in $C([0,T];X^k)$ for every $k$; so, by Lemma~\ref{lem:duhamel}, for every
real $w_0\in\bigcap_kX^k$ the function
$e^{t\Lin}w_0+\int_0^te^{(t-\tau)\Lin}F(\tau)\dd\tau$ is real, lies in
$C([0,T];X^k)\cap C^1([0,T];X^{k-1})$ for every $k\ge1$, and solves
$\partial_tu=\Lin u+F$. If $\supp\omega_{\mathrm r}(0)\subset N$, this applies to
$w_0=w(0)=\omega_{\mathrm r}(0)-\Phi$, which is then smooth
(Lemma~\ref{lem:reduction}\textup{(2)}), real, $3$-fold rotationally symmetric and supported in $N$,
hence lies in every $X^k$. Let $w^{[1]}$ be such a function, let $J\ge1$ be an integer,
and for $2\le j\le J$ define recursively
\begin{equation}
\label{eq:hier}
w^{[j]}(t)=\int_0^te^{(t-\tau)\Lin}g^{[j]}(\tau)\dd\tau,\qquad
g^{[j]}=\sum_{i+l=j}B(w^{[i]},w^{[l]})-h\cdot\nabla w^{[j-1]},
\end{equation}
where here and below $i,l\ge1$ in the sums. By induction on $j$,
Lemma~\ref{lem:bilinear} (whose bounds, with the continuity of $t\mapsto h(t)$ in
every $C^k(N')$, give continuity in time) and Lemma~\ref{lem:duhamel} show that
$g^{[j]}$ and $w^{[j]}$ are real, $g^{[j]}\in C([0,T];X^k)$ and
$w^{[j]}\in C([0,T];X^k)\cap C^1([0,T];X^{k-1})$ for every $k\ge1$, and
$\partial_tw^{[j]}=\Lin w^{[j]}+g^{[j]}$. Put $w_{\mathrm{app}}=\sum_{j=1}^Jw^{[j]}$.
Summing the equations of the $w^{[j]}$, and using $F+h\cdot\nabla\Phi=0$,
$\sum_{j=2}^{J}\sum_{i+l=j}B(w^{[i]},w^{[l]})=B(w_{\mathrm{app}},w_{\mathrm{app}})-\sum_{i,l\le J,\,i+l>J}B(w^{[i]},w^{[l]})$
and $\sum_{j=2}^{J}h\cdot\nabla w^{[j-1]}=h\cdot\nabla w_{\mathrm{app}}-h\cdot\nabla w^{[J]}$,
we find that the residual of $w_{\mathrm{app}}$ in Lemma~\ref{lem:energy} is
\begin{equation}
\label{eq:hier-res}
R_{\mathrm{app}}=-\sum_{i,l\le J,\,i+l>J}B(w^{[i]},w^{[l]})+h\cdot\nabla w^{[J]} .
\end{equation}

\subsection{A Quiet Ring}

A ring is quiet while the instability has not reached it: it starts exactly at
$\Phi$, so its deviation is created by the field $h$ alone, and
Theorem~\ref{thm:quiet} shows that it stays small. We bound a norm of the field by a
function $\eta$ of time, and its effect by the Duhamel majorant of $\eta$, defined
next. Since the semigroup bound of Proposition~I.\ref{I-prop:semigroup} grows like
$e^{\sigma(G)t}$, the majorant bounds, up to a constant factor, the linear
response to a forcing of size $\eta$.

\begin{definition}
\label{def:majorant}
Let $G\in[0,\Gone]$, $T>0$, and let $\eta\colon[0,T]\to[0,\infty)$ be integrable.
The \textbf{Duhamel majorant} of $\eta$ is
$\maj(t)=\int_0^te^{\sigma(G)(t-\tau)}\eta(\tau)\dd\tau$, $t\in[0,T]$.
\end{definition}

\begin{lemma}
\label{lem:majorant}
Let $G$, $T$, $\eta$ and $\maj$ be as in Definition~\ref{def:majorant}. Then for
$0\le\tau\le t\le T$:
\textup{(M1)} $e^{\sigma(G)(t-\tau)}\maj(\tau)\le\maj(t)$;
\textup{(M2)} $\int_\tau^t\eta\le\maj(t)$;
\textup{(M3)} $\int_0^t\maj\le\maj(t)/\sigma(G)$.
\end{lemma}
\begin{proof}
Write $\sigma=\sigma(G)$, which is positive by Proposition~I.\ref{I-prop:SU}. We have
$\maj(t)=e^{\sigma(t-\tau)}\maj(\tau)+\int_\tau^te^{\sigma(t-v)}\eta(v)\dd v$, and the
integral is nonnegative, which gives (M1), and at least $\int_\tau^t\eta$, which
gives (M2). By (M1), $\maj(\tau)\le e^{-\sigma(t-\tau)}\maj(t)$; integrating over
$\tau\in[0,t]$ gives (M3).
\end{proof}

\begin{theorem}
\label{thm:quiet}
Assume \textup{(N1)}--\textup{(N3)}. Let $s_h\ge3$ and $N_Q\ge2$ be integers with
$N_Q\sigb\ge\Lam(s_h)$, and put $M_h=s_h+N_Q+1$. There are $\majst\in(0,1]$ and
$C_Q<\infty$, depending only on $(\Phi,\Gone,s_h,N_Q)$, such that the following
holds. Let $\omega$,
$\omega_{\mathrm r},\omega_{\mathrm i},\omega_{\mathrm o}$, $G$ and $h$ be as in
Lemma~\ref{lem:reduction} on $[0,T]$, with $\omega_{\mathrm r}(0)=\Phi$ and
$G\in[0,\Gone]$, and put $w=\omega_{\mathrm r}-\Phi$. Let $\eta$ be measurable with
$\|h(t)\|_{C^{M_h}(N')}\le\eta(t)\le1$ on $[0,T]$, let $\maj$ be its Duhamel
majorant, and assume $\maj\le\majst$ on $[0,T]$. Then for every $t\in[0,T]$:
\begin{enumerate}[label=\textup{(\alph*)}]
\item $\supp\omega_{\mathrm r}(t)\subset\NNm$;
\item $\|w(t)-\wlin(t)\|_{H^{s_h}}\le C_Q\maj(t)^2$, where
  $\wlin(t)=\int_0^te^{(t-\tau)\Lin(G)}\big(-h(\tau)\cdot\nabla\Phi\big)\dd\tau$
  is the \textbf{linear response};
\item $\|w(t)\|_{H^{s_h}}\le C_Q\maj(t)$.
\end{enumerate}
\end{theorem}
\begin{proof}
Throughout, $\sigma=\sigma(G)\ge\sigb$ and $\Lin=\Lin(G)$. We compare $w$ with the
approximate solution of Equation~\eqref{eq:hier} with $J=N_Q$ and $w^{[1]}=\wlin$,
bound its terms by powers of $\maj$, and close with Lemma~\ref{lem:energy} in a
continuity argument on the radial displacement and on $\|w\|_{H^{s_h}}$. Both
terms of the bound on the residual contain the factor $\maj^{N_Q}$
(Equation~\eqref{eq:Q-res} below), and the condition $N_Q\sigb\ge\Lam(s_h)$ lets
the growth of this factor absorb the growth $e^{\Lam(s_h)t}$ allowed by the
energy estimate.

\emph{Constants.} Put $\sigma'_0=\frac12(\frac9{200}+\sigb)\in(\frac9{200},\sigb)$ and
\[
C_L=\max_{0\le k\le M_h}C_S(k,\sigma'_0),\qquad
C_{\mathrm{bil}}=\max_{2\le k\le M_h}\max\big(C_B(k),C_h(k),\Cnab(k)\big),
\]
with $C_S$ of Proposition~I.\ref{I-prop:semigroup} and the constants of
Lemma~\ref{lem:bilinear}; so $\|e^{t\Lin}\|_{X^k\to X^k}\le C_Le^{\sigma t}$ for
$k\le M_h$ and $t\ge0$. For $1\le j\le N_Q$ put $k_j=M_h-j$, so that
$k_{N_Q}=s_h+1$ and $k_j+1=k_{j-1}$. Put $C^{(1)}=C_LC_{\mathrm{bil}}$ and, for
$2\le j\le N_Q$,
\[
\begin{gathered}
C^{(j)}=C_LC_{\mathrm{bil}}\Big(\frac1{(j-1)\sigb}\sum_{i+l=j}C^{(i)}C^{(l)}+C^{(j-1)}\Big),\\
C_{\mathrm{res}}=C_{\mathrm{bil}}\Big(\sum_{i,l\le N_Q,\,i+l>N_Q}C^{(i)}C^{(l)}+C^{(N_Q)}\Big),\\
C_Q=\sum_{j=1}^{N_Q}C^{(j)}+C_{\mathrm{res}}\big(1+\sigb^{-1}\big),\qquad
C_\Delta=1+\frac{C_{\mathrm{BS}}C_Q}{\sigb},\\
\majst=\min\Big(1,\frac1{2C_Q},\frac a{16C_\Delta}\Big).
\end{gathered}
\]
The names $C_{\mathrm{bil}}$, $C^{(j)}$ and $C_{\mathrm{res}}$ are local to this proof;
later in this section they denote other numbers.
These numbers depend only on $(\Phi,\Gone,s_h,N_Q)$, and $C_\Delta\majst\le a/16$.

\emph{Approximate solution.} Let $w^{[j]}$, $1\le j\le N_Q$, be the functions of
Equation~\eqref{eq:hier} with $J=N_Q$ and $w^{[1]}=\wlin$ (the case $w_0=0$), so that
they are regular by Lemmas~\ref{lem:duhamel} and \ref{lem:bilinear}. By
Lemma~\ref{lem:bilinear}\textup{(3)},
$\|F(\tau)\|_{H^{M_h}}\le C_{\mathrm{bil}}\eta(\tau)$. We claim that
\begin{equation}
\label{eq:Q-hier}
\|w^{[j]}(t)\|_{H^{k_j}}\le C^{(j)}\maj(t)^j\qquad(t\in[0,T],\ 1\le j\le N_Q).
\end{equation}
For $j=1$,
$\|\wlin(t)\|_{H^{k_1}}\le C_L\int_0^te^{\sigma(t-\tau)}C_{\mathrm{bil}}\eta(\tau)\dd\tau=C^{(1)}\maj(t)$.
Let $j\ge2$ and assume Equation~\eqref{eq:Q-hier} for all indices less than $j$.
Since $k_j+1=k_{j-1}\le k_l$ for $l\le j-1$, Lemma~\ref{lem:bilinear} gives
$\|g^{[j]}(\tau)\|_{H^{k_j}}\le C_{\mathrm{bil}}\big(\sum_{i+l=j}C^{(i)}C^{(l)}\maj(\tau)^j+C^{(j-1)}\eta(\tau)\maj(\tau)^{j-1}\big)$.
By (M1), $e^{\sigma(t-\tau)}\maj(\tau)^j\le e^{-(j-1)\sigma(t-\tau)}\maj(t)^j$, whose
integral over $[0,t]$ is at most $\maj(t)^j/((j-1)\sigb)$; and
$e^{\sigma(t-\tau)}\eta(\tau)\maj(\tau)^{j-1}\le\eta(\tau)\maj(t)^{j-1}$, whose integral
is at most $\maj(t)^j$ by (M2). With $\|e^{(t-\tau)\Lin}\|\le C_Le^{\sigma(t-\tau)}$ this
gives Equation~\eqref{eq:Q-hier} for $j$. Since $k_j\ge s_h+1$ and $\maj\le1$,
Equations~\eqref{eq:hier-res} and \eqref{eq:Q-hier} and Lemma~\ref{lem:bilinear} give
\begin{equation}
\label{eq:Q-res}
\|R_{\mathrm{app}}(\tau)\|_{H^{s_h}}\le C_{\mathrm{res}}\big(\maj(\tau)^{N_Q+1}+\eta(\tau)\maj(\tau)^{N_Q}\big),
\end{equation}
and $\|w_{\mathrm{app}}\|_{H^{s_h+1}}\le\sum_jC^{(j)}\maj\le C_Q\majst\le\frac12$ on
$[0,T]$.

\emph{Continuity argument.} Let $t_*$ be the largest $t\in[0,T]$ such that
$\delta_{\mathrm r}(t)\le a/8$ and $\|w(\tau)\|_{H^{s_h}}\le1$ for $\tau\in[0,t]$; it
exists and is positive since $\delta_{\mathrm r}(0)=0$, $w(0)=0$, $\delta_{\mathrm r}$
is continuous (see its definition at the start of this section), and
$\|w\|_{H^{s_h}}$ is continuous by Theorem~\ref{thm:smooth}\textup{(1)}. Because $\supp\Phi\subset\{a\le|y|\le1\}$ and
$1+\frac a8\le\frac54$, the support of $\omega_{\mathrm r}(\tau)$ lies in
$\NNm\subset N$ for $\tau\le t_*$. So Lemma~\ref{lem:reduction}\textup{(3)} and
Lemma~\ref{lem:energy} with $s=s_h$ apply on $[0,t_*]$; the conditions
$\|w_{\mathrm{app}}\|_{H^{s_h+1}}\le1$ and $\|h\|_{C^{s_h+1}(N')}\le\eta\le1$ hold.
Since $v=w-w_{\mathrm{app}}$ has $v(0)=0$, Equation~\eqref{eq:Q-res} gives, for
$t\le t_*$,
\[
\|v(t)\|_{H^{s_h}}\le C_{\mathrm{res}}\int_0^te^{\Lam(s_h)(t-\tau)}\maj(\tau)^{N_Q}\big(\maj(\tau)+\eta(\tau)\big)\dd\tau .
\]
For $0\le\tau\le t$, (M1) and $N_Q\sigma\ge N_Q\sigb\ge\Lam(s_h)$ give
$e^{\Lam(s_h)(t-\tau)}\maj(\tau)^{N_Q}\le e^{(\Lam(s_h)-N_Q\sigma)(t-\tau)}\maj(t)^{N_Q}$,
which is at most $\maj(t)^{N_Q}$. By (M2) and (M3) we conclude, for $t\le t_*$,
\begin{equation}
\label{eq:Q-v}
\|v(t)\|_{H^{s_h}}\le C_{\mathrm{res}}\maj(t)^{N_Q}\int_0^t(\maj+\eta)\le C_{\mathrm{res}}\big(1+\sigb^{-1}\big)\maj(t)^{N_Q+1}.
\end{equation}
Since $w=\wlin+\sum_{j\ge2}w^{[j]}+v$ and $\maj\le1$, Equations~\eqref{eq:Q-hier}
and \eqref{eq:Q-v} give, for $t\le t_*$,
\begin{equation}
\label{eq:Q-b}
\|w(t)-\wlin(t)\|_{H^{s_h}}\le C_Q\maj(t)^2,\qquad\|w(t)\|_{H^{s_h}}\le C_Q\maj(t)\le\tfrac12 .
\end{equation}
For $y\in\supp\Phi$ and $\tau\le t\le t_*$, the point $\flow_\tau(y)$ lies in $N$, so
by Lemma~\ref{lem:reduction}\textup{(2)}, Lemma~\ref{lem:bilinear}\textup{(4)},
Equation~\eqref{eq:Q-b}, (M2) and (M3),
\[
\big||\flow_\tau(y)|-|y|\big|\le\int_0^t\big(\eta+C_{\mathrm{BS}}\|w\|_{H^{s_h}}\big)\le\maj(t)+\frac{C_{\mathrm{BS}}C_Q}{\sigb}\maj(t)=C_\Delta\maj(t).
\]
So $\delta_{\mathrm r}\le C_\Delta\maj\le a/16$ and $\|w\|_{H^{s_h}}\le\frac12$ on
$[0,t_*]$. If $t_*<T$, continuity would give a larger $t$ with
$\delta_{\mathrm r}(t)\le a/8$ and $\|w\|_{H^{s_h}}\le1$ on $[0,t]$; hence $t_*=T$.
Then (a) follows from $\delta_{\mathrm r}\le C_\Delta\maj$ and Lemma~\ref{lem:annuli},
and (b), (c) are
Equation~\eqref{eq:Q-b}.
\end{proof}

\subsection{An Active Ring}

A ring is active once its deviation from $\Phi$ has a component along the
unstable eigenfunction that grows at the rate $\sigma(G)$. Theorem~\ref{thm:active}
shows that the deviation then follows this component, up to a decaying transient
and errors of relative size $\hat a$, $\kappa^{1/2}$ and $\eta_b$, until its
envelope $\kappa e^{\sigma(G)u}$ reaches a small size $\delta_0$.
Corollary~\ref{cor:active-moment} gives the outer moment of the ring, through
which it drives the rings outside it. The unstable component of the linear
response is explicit, by Proposition~I.\ref{I-prop:semigroup}; the next lemma
writes it with a running coefficient $c(u)$, which accounts for the forcing by $h$.

\begin{lemma}
\label{lem:running}
Let $G\in[0,\Gone]$, $\bar u>0$, $w_0\in X^0$ real, and let $h$ be continuous on
$[0,\bar u]$ with values in $C^0(N';\R^2)$ and $3$-fold rotationally symmetric. Put
$F(u)=-h(u)\cdot\nabla\Phi$ and
\begin{equation}
\label{eq:running}
c(u)=\ellst(w_0)+\int_0^ue^{-\lamst v}\,\ellst(F(v))\dd v\qquad(0\le u\le\bar u),
\end{equation}
the \textbf{running coefficient}. Then
$\wlin_0(u):=e^{u\Lin(G)}w_0+\int_0^ue^{(u-v)\Lin(G)}F(v)\dd v$ satisfies
$\ellst(\wlin_0(u))=c(u)e^{\lamst u}$ and
$\Riesz\,\wlin_0(u)=2\Re\big(c(u)e^{\lamst u}\phist\big)$ for $0\le u\le\bar u$.
\end{lemma}
\begin{proof}
The function $F$ is continuous on $[0,\bar u]$ with values in $X^0$, since
$\|h(u)\cdot\nabla\Phi-h(u')\cdot\nabla\Phi\|_{L^2}\le\|h(u)-h(u')\|_{C^0(N')}\|\nabla\Phi\|_{L^2}$,
and $F(u)$ is real, supported in $N$ and $3$-fold rotationally symmetric. The functional
$\ellst$ is bounded on $X^0$, so it commutes with the integral, and
Proposition~I.\ref{I-prop:semigroup} gives $\ellst(e^{\tau\Lin}f)=e^{\lamst\tau}\ellst(f)$
for $f\in X^0$. Hence
$\ellst(\wlin_0(u))=e^{\lamst u}\ellst(w_0)+\int_0^ue^{\lamst(u-v)}\ellst(F(v))\dd v=c(u)e^{\lamst u}$.
Since $\wlin_0(u)$ is real, Definition~I.\ref{I-def:eigendata} gives
$\Riesz\,\wlin_0(u)=2\Re\big(\ellst(\wlin_0(u))\phist\big)=2\Re\big(c(u)e^{\lamst u}\phist\big)$.
\end{proof}

On the range of $1-\Riesz$ the semigroup grows at most at a rate $\sigma'<\sigb$
(Proposition~I.\ref{I-prop:semigroup}). The next lemma uses this to bound the rest
$(1-\Riesz)\wlin_0$ of the linear response.

\begin{lemma}
\label{lem:linear-part}
Assume \textup{(N1)}--\textup{(N3)}. Let $\sigma'\in(\frac9{200},\sigb)$, and let
$p\ge0$ and $k\ge2$ be integers. Put
$C_{\mathrm{lin}}(k)=C_S(k,\sigma')\big(1+\Cnab(k)/\sigma'+\Cnab(k)/(\sigb-\sigma')\big)$,
with $C_S$ of Proposition~I.\ref{I-prop:semigroup}. Let $G\in[0,\Gone]$,
$\sigma=\sigma(G)$, $\lambda=\lamst(G)$, $\bar u>0$, $A\ge0$, $\kappa\ge0$ and
$\eta_b\ge0$. Let $w_0\in X^k$ be real with $\|w_0\|_{H^k}\le A$, and let $h$ be
continuous on $[0,\bar u]$ with values in $C^k(N';\R^2)$, $3$-fold rotationally symmetric, with
$\|h(u)\|_{C^k(N')}\le A\jap{u}^p+\eta_b\kappa e^{\sigma u}$. Let $c$ and $\wlin_0$ be as
in Lemma~\ref{lem:running}. Then for $u\in[0,\bar u]$,
\[
\begin{gathered}
\wlin_0(u)=2\Re\big(c(u)e^{\lambda u}\phist\big)+(1-\Riesz)\wlin_0(u),\\
\|(1-\Riesz)\wlin_0(u)\|_{H^k}\le C_{\mathrm{lin}}(k)\big(A\jap{u}^pe^{\sigma'u}+\eta_b\kappa e^{\sigma u}\big).
\end{gathered}
\]
\end{lemma}
\begin{proof}
The decomposition is Lemma~\ref{lem:running}. For real $f\in X^k$,
Proposition~I.\ref{I-prop:semigroup} gives
$\Riesz e^{\tau\Lin}f=2\Re\big(\ellst(e^{\tau\Lin}f)\phist\big)=2\Re\big(e^{\lambda\tau}\ellst(f)\phist\big)=e^{\tau\Lin}\Riesz f$.
Since $\Riesz$ is bounded on $X^k$, it commutes with the integral, and
$(1-\Riesz)\wlin_0(u)=e^{u\Lin}(1-\Riesz)w_0+\int_0^ue^{(u-v)\Lin}(1-\Riesz)F(v)\dd v$.
By Proposition~I.\ref{I-prop:semigroup} and the bound
$\|h\cdot\nabla\Phi\|_{H^k}\le\Cnab(k)\|h\|_{C^k(N')}$ of Lemma~I.\ref{I-lem:uniform}, its
norm in $H^k$ is at most
$C_S(k,\sigma')\big(e^{\sigma'u}A+\Cnab(k)\int_0^ue^{\sigma'(u-v)}(A\jap{v}^p+\eta_b\kappa e^{\sigma v})\dd v\big)$.
Finally
$\int_0^ue^{\sigma'(u-v)}\jap{v}^p\dd v\le\jap{u}^pe^{\sigma'u}/\sigma'$ and
$\int_0^ue^{\sigma'(u-v)}e^{\sigma v}\dd v\le e^{\sigma u}/(\sigma-\sigma')\le e^{\sigma u}/(\sigb-\sigma')$.
\end{proof}

For $x\in(0,1]$ we write $\varsigma(x)=1+\sigb^{-1}\log(1/x)$. Since
$\sigma(G)\ge\sigb$, this number bounds $1+u$ for $u\le\sigma(G)^{-1}\log(1/x)$,
the time in which a coefficient of size $x$ that grows at the rate $\sigma(G)$
reaches size $1$. For an integer $s\ge3$
and a number $\sigma'\in(\frac9{200},\sigb)$ we put
\begin{equation}
\label{eq:W-const}
\begin{gathered}
p=3,\qquad N_G=\min\big\{n\in\Z:n\ge2,\ n\ge\Lam(s)/\sigb+1\big\},\qquad
M=s+N_G+2,\\ N_W=(p+2)N_G+1,\qquad
\bar C_S=\max_{0\le k\le M}C_S(k,\sigma'),\qquad
C_{\mathrm{lin}}=\max_{2\le k\le M}C_{\mathrm{lin}}(k),\\
C_{\mathrm{bil}}=\max_{2\le k\le M}\max\big(C_B(k),C_h(k),\Cnab(k)\big),\\
C_\phi=\max_{G\in[0,\Gone]}\|\phist(G)\|_{H^M},\qquad
\bar\rho=\sup_{u\ge0}\jap{u}^pe^{-(\sigb-\sigma')u},
\end{gathered}
\end{equation}
with $C_S$ of Proposition~I.\ref{I-prop:semigroup}, $C_B$ and $C_h$ of
Lemma~\ref{lem:bilinear}, $\Cnab$ of Lemma~I.\ref{I-lem:uniform} and
$C_{\mathrm{lin}}(k)$ of Lemma~\ref{lem:linear-part}. The maximum defining $C_\phi$
exists and is finite by Lemma~I.\ref{I-lem:uniform}, and $\bar\rho<\infty$ because
$\sigb>\sigma'$. All these numbers depend only on $(\Phi,\Gone,s,\sigma')$. As in
Theorem~\ref{thm:quiet}, the order $N_G$ of the approximate solution is chosen
large, so that its residual, which is one order higher, is small enough on the
window of Lemma~\ref{lem:window} to absorb the growth $e^{\Lam(s)u}$ allowed by
Lemma~\ref{lem:energy} (Equation~\eqref{eq:W-exp} below).

Lemma~\ref{lem:window} is the main step in the proof of Theorem~\ref{thm:active}.
It treats an active ring whose initial deviation and running coefficient are
bounded by a fixed multiple $C_\natural$ of a number $\kappa_1$, and whose field is
bounded accordingly. It does so on the window of time in which a component of
initial size $\kappa_1$, growing at the rate $\sigma(G)$, stays below $\delta_0$.

\begin{lemma}
\label{lem:window}
Assume \textup{(N1)}--\textup{(N3)}. Let $s\ge3$ be an integer and
$\sigma'\in(\frac9{200},\sigb)$, let $p$, $N_G$, $M$ and $N_W$ be as in
Equation~\eqref{eq:W-const}, and let $C_\natural\ge1$. There are $\delta_W\in(0,\frac12]$, $\kappa_W\in(0,1]$ and
$C_W\ge1$, depending only on $(\Phi,\Gone,s,\sigma',C_\natural)$, such that the
following holds. Let $G\in[0,\Gone]$, $\sigma=\sigma(G)$, $\lambda=\lamst(G)$,
$\kappa_1\in(0,\kappa_W]$, $\delta_0\in(0,\delta_W]$, $\eta_b\in[0,1]$,
$\hat T_1=\sigma^{-1}\log(\delta_0/\kappa_1)$ and $\bar u\in(0,\hat T_1]$. Let
$\omega$, $\omega_{\mathrm r},\omega_{\mathrm i},\omega_{\mathrm o}$ and $h$ be as in
Lemma~\ref{lem:reduction} on $[0,\bar u]$ with this $G$, put
$w=\omega_{\mathrm r}-\Phi$, and let $c$ be the running coefficient of Equation~\eqref{eq:running}
with $w_0=w(0)$. Assume that $\|w(0)\|_{H^{M-2}}\le C_\natural\kappa_1$, that
$\supp(\Phi+w(0))\subset\{(\frac78-\frac1{64})a\le|y|\le\frac54+\frac a{64}\}$, and
that for $u\in[0,\bar u]$,
\[
\|h(u)\|_{C^M(N')}\le C_\natural\kappa_1\jap{u}^p+\eta_b\kappa_1e^{\sigma u},\qquad
|c(u)|\le C_\natural\kappa_1 .
\]
Put $\hat a_1(u)=\kappa_1e^{\sigma u}$ and $\rho_1(u)=\jap{u}^pe^{-(\sigma-\sigma')u}$.
Then for every $u\in[0,\bar u]$, every trajectory of $\omega_{\mathrm r}$ moves
radially by at most $C_W\delta_0+C_W\kappa_1\varsigma(\kappa_1)^{N_W}\le a/32$ on
$[0,u]$, and
\[
\big\|w(u)-2\Re\big(c(u)e^{\lambda u}\phist\big)\big\|_{H^s}\le C_W\hat a_1(u)\big[\rho_1(u)+\hat a_1(u)+\kappa_1\varsigma(\kappa_1)^{N_W}+\eta_b\big].
\]
\end{lemma}
\begin{proof}
Write $\Lam=\Lam(s)$ and $\Lin=\Lin(G)$. Since $N_G\ge\Lam/\sigb+1$ and
$\Lam\ge2\sigma\ge2\sigb$ (Lemma~\ref{lem:energy}),
\begin{equation}
\label{eq:W-exp}
(N_G+1)\sigma-\Lam\ge2\sigma,\qquad\Lam-\sigma\ge\sigma\ge\sigb,\qquad N_G+1-\Lam/\sigma\ge2 .
\end{equation}

\emph{Constants.} With the numbers of Equation~\eqref{eq:W-const}, put
$C^{(1)}=2C_\natural C_\phi+C_{\mathrm{lin}}(C_\natural\bar\rho+1)$ and, for
$2\le j\le N_G$,
\[
C^{(j)}=2^{j-1}\max(1,\sigb^{-1})\,\bar C_SC_{\mathrm{bil}}\Big(\sum_{i+l=j}C^{(i)}C^{(l)}+C^{(j-1)}\Big).
\]
Put $C_{\mathrm{res}}=C_{\mathrm{bil}}\big(\sum_{i,l\le N_G,\,i+l>N_G}C^{(i)}C^{(l)}+2C^{(N_G)}\big)$,
$C_{\mathrm{en}}=2^{N_G-1}C_{\mathrm{res}}/\sigb$, $S_W=\sum_{j=1}^{N_G}C^{(j)}+C_{\mathrm{en}}$,
\[
C_W=2\big(1+\sigb^{-1}\big)\big(1+C_{\mathrm{BS}}\big)C_\natural^{N_G}\big(C_{\mathrm{lin}}C_\natural+S_W\big),\qquad
\delta_W=\min\Big(\frac12,\frac1{S_W},\frac a{64C_W}\Big),
\]
and let $\kappa_W$ be the largest number in $(0,1]$ such that
$C_Wx\varsigma(x)^{N_W}\le a/64$ for every $x\in(0,\kappa_W]$. It exists by the
argument given for $\kappa_0(C_0)$ in the proof of Theorem~\ref{thm:active}, because
$x\varsigma(x)^{N_W}\to0$ as $x\to0^+$ and the inequality is non-strict between
continuous functions. Taking $t=0$ in
Proposition~I.\ref{I-prop:semigroup} gives $C_S\ge1$, so $C_{\mathrm{lin}}\ge1$ and
$C_W\ge C_\natural$.

Now fix the data, and put $\varepsilon_W=C_\natural\kappa_1$,
$\mathfrak b(u)=\hat a_1(u)+\varepsilon_W\jap{u}^{p+2}$ and $k_j=M-1-j$ for
$1\le j\le N_G$, so $k_1=M-2$, $k_{N_G}=s+1$, $k_j+1=k_{j-1}$. The function
$\mathfrak b$ plays the role of $\maj$ in the proof of Theorem~\ref{thm:quiet}: the
terms $w^{[j]}$ below are bounded by multiples of $\hat a_1\mathfrak b^{j-1}$. For
$u\le\bar u\le\hat T_1\le\sigma^{-1}\log(1/\kappa_1)$ we have $\jap{u}\le1+u\le\varsigma(\kappa_1)$,
$\hat a_1(u)\le\delta_0$ and $\kappa_1\le\hat a_1(u)$. By the choice of $\kappa_W$ and
$\delta_W$, for $u\in[0,\bar u]$,
\begin{equation}
\label{eq:W-small}
\mathfrak b(u)\le\delta_0+C_\natural\kappa_1\varsigma(\kappa_1)^{p+2}\le1,\qquad
\|h(u)\|_{C^{s+1}(N')}\le C_\natural\kappa_1\varsigma(\kappa_1)^p+\delta_0\le1 .
\end{equation}

\emph{Approximate solution.} Let $w^{[j]}$ be the functions of
Equation~\eqref{eq:hier} with $J=N_G$ and $w^{[1]}=\wlin_0$, the function of
Lemma~\ref{lem:running} with $w_0=w(0)$. This choice of $w^{[1]}$ is admissible in
Equation~\eqref{eq:hier} because $\supp(\Phi+w(0))\subset N$, and the functions
$w^{[j]}$ are regular by Lemmas~\ref{lem:duhamel} and \ref{lem:bilinear}. We
claim that
\begin{equation}
\label{eq:W-hier}
\|w^{[j]}(u)\|_{H^{k_j}}\le C^{(j)}\hat a_1(u)\,\mathfrak b(u)^{j-1}\qquad(u\in[0,\bar u],\ 1\le j\le N_G).
\end{equation}
For $j=1$, Lemma~\ref{lem:linear-part} with $k=M-2$, $A=C_\natural\kappa_1$ and
$\kappa=\kappa_1$ gives
$\|(1-\Riesz)w^{[1]}(u)\|_{H^{M-2}}\le C_{\mathrm{lin}}\hat a_1(u)\big(C_\natural\rho_1(u)+\eta_b\big)$,
since $C_\natural\kappa_1\jap{u}^pe^{\sigma'u}=C_\natural\hat a_1(u)\rho_1(u)$; and
$\|2\Re(c(u)e^{\lambda u}\phist)\|_{H^{M-2}}\le2C_\phi C_\natural\hat a_1(u)$. As
$\rho_1\le\bar\rho$ and $\eta_b\le1$, this is Equation~\eqref{eq:W-hier} for $j=1$.
Let $j\ge2$. By Lemma~\ref{lem:bilinear}, Equation~\eqref{eq:W-hier} for the
indices less than $j$, and
$\|h(v)\|_{C^{k_j}(N')}\le\varepsilon_W\jap{v}^p+\hat a_1(v)$,
\[
\|g^{[j]}(v)\|_{H^{k_j}}\le C_{\mathrm{bil}}\Big(\sum_{i+l=j}C^{(i)}C^{(l)}+C^{(j-1)}\Big)\hat a_1(v)\big(\hat a_1(v)+\varepsilon_W\jap{v}^p\big)\mathfrak b(v)^{j-2}.
\]
Since $e^{\sigma(u-v)}\hat a_1(v)=\hat a_1(u)$ and
$\|e^{(u-v)\Lin}\|\le\bar C_Se^{\sigma(u-v)}$, it suffices to show that
\[
\int_0^u\big(\hat a_1+\varepsilon_W\jap{v}^p\big)\mathfrak b^{j-2}\dd v\le2^{j-1}\max(1,\sigb^{-1})\,\mathfrak b(u)^{j-1}.
\]
Expand $\mathfrak b^{j-2}=\sum_{i+q=j-2}\binom{j-2}i\hat a_1^i(\varepsilon_W\jap{v}^{p+2})^q$,
and use that $\hat a_1$ and $\jap{\cdot}$ are nondecreasing. Then
\[
\begin{gathered}
\int_0^u\hat a_1^{i+1}(\varepsilon_W\jap{v}^{p+2})^q\dd v\le\frac{\hat a_1(u)^{i+1}(\varepsilon_W\jap{u}^{p+2})^q}{(i+1)\sigma},\\
\int_0^u\varepsilon_W\jap{v}^p\hat a_1^i(\varepsilon_W\jap{v}^{p+2})^q\dd v\le\frac{\hat a_1(u)^i(\varepsilon_W\jap{u}^{p+2})^{q+1}}{i\sigma}\quad(i\ge1),\\
\int_0^u\varepsilon_W^{j-1}\jap{v}^{p+(p+2)(j-2)}\dd v\le(\varepsilon_W\jap{u}^{p+2})^{j-1},
\end{gathered}
\]
where the second line uses $\jap{u}^p\le\jap{u}^{p+2}$ and the third uses
$u\le\jap{u}$. Each right side is a monomial of degree $j-1$ in
$\hat a_1(u)$ and $\varepsilon_W\jap{u}^{p+2}$, hence at most
$\max(1,\sigb^{-1})\mathfrak b(u)^{j-1}$, and there are $2\cdot2^{j-2}$ terms counted with
the binomial weights. This proves Equation~\eqref{eq:W-hier}.

The residual is given by Equation~\eqref{eq:hier-res} with $J=N_G$. Since
$\mathfrak b\le1$ and $\hat a_1\le\mathfrak b$, we have
$\hat a_1^2\mathfrak b^{i+l-2}\le\hat a_1\mathfrak b^{N_G}$ for $i+l>N_G$, and
$\|h\|_{C^s(N')}\le\varepsilon_W\jap{u}^p+\hat a_1\le2\mathfrak b$. So
\begin{equation}
\label{eq:W-res}
\|R_{\mathrm{app}}(u)\|_{H^s}\le C_{\mathrm{res}}\hat a_1(u)\mathfrak b(u)^{N_G},\qquad
\|w_{\mathrm{app}}(u)\|_{H^{s+1}}\le\Big(\sum_jC^{(j)}\Big)\hat a_1(u)\le S_W\delta_0\le1 .
\end{equation}

\emph{Continuity argument.} Let $u_*$ be the largest $u\in[0,\bar u]$ with
$\delta_{\mathrm r}(u)\le a/16$ and $\|w(v)\|_{H^s}\le1$ for $v\in[0,u]$; at $u=0$
both hold, since $\delta_{\mathrm r}(0)=0$ and, as $s\le M-2$,
$\|w(0)\|_{H^s}\le C_\natural\kappa_1\le a/64$ (recall $C_W\ge C_\natural$, $\varsigma\ge1$
and the choice of $\kappa_W$). For $u\le u_*$ the
support of $\omega_{\mathrm r}(u)$ lies in
$\{\frac{51}{64}a\le|y|\le\frac54+\frac5{64}a\}\subset\Nflat\subset N$
(Lemma~\ref{lem:annuli}). So
Lemma~\ref{lem:reduction}\textup{(3)} and Lemma~\ref{lem:energy} apply on $[0,u_*]$,
by Equations~\eqref{eq:W-small} and \eqref{eq:W-res}. As $v=w-w_{\mathrm{app}}$
vanishes at $u=0$, and
$\mathfrak b^{N_G}\le2^{N_G-1}\big(\hat a_1^{N_G}+(\varepsilon_W\jap{\cdot}^{p+2})^{N_G}\big)$,
\[
\begin{aligned}
\|v(u)\|_{H^s}\le2^{N_G-1}C_{\mathrm{res}}\Big(&\int_0^ue^{\Lam(u-v)}\hat a_1(v)^{N_G+1}\dd v\\
&+(\varepsilon_W\jap{u}^{p+2})^{N_G}\int_0^ue^{\Lam(u-v)}\hat a_1(v)\dd v\Big).
\end{aligned}
\]
By Equation~\eqref{eq:W-exp}, the first integral is at most
$\hat a_1(u)^{N_G+1}/((N_G+1)\sigma-\Lam)\le\hat a_1(u)^{N_G+1}/(2\sigb)$, and the second is
at most $\kappa_1e^{\Lam u}/(\Lam-\sigma)\le\kappa_1e^{\Lam u}/\sigb$. For $u\le\hat T_1$,
$\kappa_1e^{\Lam u}\le\kappa_1^{1-\Lam/\sigma}\delta_0^{\Lam/\sigma}$. Hence, for $u\le u_*$,
\begin{equation}
\label{eq:W-v}
\|v(u)\|_{H^s}\le C_{\mathrm{en}}\big(\hat a_1(u)^{N_G+1}+\Theta_W(u)\big),\qquad
\Theta_W(u)=(\varepsilon_W\jap{u}^{p+2})^{N_G}\kappa_1^{1-\Lam/\sigma}\delta_0^{\Lam/\sigma},
\end{equation}
and, by Equation~\eqref{eq:W-exp}, $\kappa_1\le1$ and $\delta_0\le1$,
\begin{equation}
\label{eq:W-theta}
\Theta_W(u)\le C_\natural^{N_G}\varsigma(\kappa_1)^{(p+2)N_G}\kappa_1^2\le C_\natural^{N_G}\kappa_1\varsigma(\kappa_1)^{N_W}\cdot\kappa_1 .
\end{equation}

We collect the bounds for $u\le u_*$. Write
$w-2\Re(ce^{\lambda u}\phist)=(1-\Riesz)w^{[1]}+\sum_{j\ge2}w^{[j]}+v$. By the bound on
$(1-\Riesz)w^{[1]}$ above, Equation~\eqref{eq:W-hier} with $\mathfrak b\le1$,
Equations~\eqref{eq:W-v} and \eqref{eq:W-theta}, $\hat a_1^{N_G+1}\le\hat a_1^2$ and $\Theta_W\le\hat a_1\Theta_W/\kappa_1$,
\[
\begin{aligned}
\big\|w-2\Re(ce^{\lambda u}\phist)\big\|_{H^s}\le{}&C_{\mathrm{lin}}\hat a_1(C_\natural\rho_1+\eta_b)\\
&+S_W\hat a_1\big(2\hat a_1+C_\natural\kappa_1\varsigma(\kappa_1)^{p+2}+C_\natural^{N_G}\kappa_1\varsigma(\kappa_1)^{N_W}\big),
\end{aligned}
\]
which is at most $C_W\hat a_1\big[\rho_1+\hat a_1+\kappa_1\varsigma(\kappa_1)^{N_W}+\eta_b\big]$.
Similarly,
$\|w(u)\|_{H^s}\le S_W\big(\hat a_1(u)+\Theta_W(u)\big)\le C_W\big(\delta_0+\kappa_1\varsigma(\kappa_1)^{N_W}\big)\le\tfrac a{32}$.
For the displacement, Lemma~\ref{lem:reduction}\textup{(2)} and
Lemma~\ref{lem:bilinear}\textup{(4)} give, for $u\le u_*$,
\[
\begin{aligned}
\delta_{\mathrm r}(u)&\le\int_0^u\big(\|h\|_{C^0(N')}+C_{\mathrm{BS}}\|w\|_{H^s}\big)\\
&\le C_\natural\kappa_1\varsigma(\kappa_1)^{p+1}+\frac{\delta_0}{\sigb}+C_{\mathrm{BS}}\Big(\frac{S_W\delta_0}{\sigb}+C_{\mathrm{en}}C_\natural^{N_G}\varsigma(\kappa_1)^{N_W}\kappa_1^2\Big),
\end{aligned}
\]
where we used $\|w\|_{H^s}\le\sum_j\|w^{[j]}\|_{H^s}+\|v\|_{H^s}\le S_W\hat a_1+C_{\mathrm{en}}\Theta_W$,
$\int_0^u\hat a_1\le\hat a_1(u)/\sigma\le\delta_0/\sigb$, and
$\int_0^u\Theta_W\le\bar u\,C_\natural^{N_G}\varsigma(\kappa_1)^{(p+2)N_G}\kappa_1^2$ with
$\bar u\le\varsigma(\kappa_1)$. The right side is at most
$C_W\delta_0+C_W\kappa_1\varsigma(\kappa_1)^{N_W}\le a/32$. If $u_*<\bar u$, continuity
would give a larger $u$ with $\delta_{\mathrm r}(u)\le a/16$ and $\|w\|_{H^s}\le1$ on $[0,u]$, a
contradiction. So $u_*=\bar u$, and the bounds just proved hold on $[0,\bar u]$.
\end{proof}

\begin{theorem}
\label{thm:active}
Assume \textup{(N1)}--\textup{(N3)}. Fix an integer $s\ge3$, a number
$\sigma'\in(\frac9{200},\sigb)$, and put $p=3$ $(=m)$. Let $N_G$ be the least integer
with $N_G\ge2$ and $N_G\ge\Lam(s)/\sigb+1$, put $M=s+N_G+2$ (so $M\ge s+4$), and
put $\etmax=1$. There are $\dH\in(0,1)$, $C_2\ge1$ and $C_3>0$, depending only on
$(\Phi,\Gone,s,\sigma')$, and for every $C_0\ge1$ a number $\kappa_0(C_0)>0$,
depending only on $C_0$ and $(\Phi,\Gone,s,\sigma')$, such that the following holds.

Let $G\in[0,\Gone]$, $\sigma=\sigma(G)$, $\lambda=\lamst(G)$, $\kappa>0$,
$\eta_b\in[0,\etmax]$, $\delta_0\in(0,\dH]$, $\hat a(u)=\kappa e^{\sigma u}$ and
$\hat T=\sigma^{-1}\log(\delta_0/\kappa)$. Let $t_1=t_1(C_0,G)$ be the least
$t\ge0$ such that $2^{p/2}C_0\jap{t'}^pe^{-(\sigma-\sigma')t'}\le1$ for all $t'\ge t$,
and define the \textbf{transient}
\[
\rhohat(u)=C_0\jap{u}^pe^{-(\sigma-\sigma')u}\ (0\le u<t_1),\qquad
\rhohat(u)=\jap{u-t_1}^pe^{-(\sigma-\sigma')(u-t_1)}\ (u\ge t_1).
\]
Let $\bar u>0$, let $\omega$, $\omega_{\mathrm r},\omega_{\mathrm i},\omega_{\mathrm o}$
and $h$ be as in Lemma~\ref{lem:reduction} on $[0,\bar u]$ with this $G$, put
$w(u)=\omega_{\mathrm r}(u)-\Phi$, and let $c$ be the running coefficient of
Equation~\eqref{eq:running} with $w_0=w(0)$. Assume:
\begin{description}[font=\normalfont]
\item[\textup{(Init)}] $\|w(0)\|_{H^{M-1}}\le C_0\kappa$ and $\supp(\Phi+w(0))\subset\NNm$;
\item[\textup{(Field)}] $h=h_1+h_2$ on $[0,\bar u]$ with
  $\|h_1(u)\|_{C^M(N')}\le C_0\kappa\jap{u}^p$ and
  $\|h_2(u)\|_{C^M(N')}\le\eta_b\,\hat a(u)$;
\item[\textup{(Coef)}] $|c(u)|\le\kappa\big(2+C_0\jap{u}^pe^{-\sigma u}\big)$ on $[0,\bar u]$;
\item[\textup{(Small)}] $\kappa\le\kappa_0(C_0)$.
\end{description}
Then for every $u\in[0,\min(\bar u,\hat T)]$:
\begin{enumerate}[label=\textup{(\alph*)}]
\item $\supp(\Phi+w(u))\subset\NN$;
\item $\big\|w(u)-2\Re\big(c(u)e^{\lambda u}\phist\big)\big\|_{H^s}\le
  C_2\hat a(u)\big[\rhohat(u)+\hat a(u)+\kappa^{1/2}+\eta_b\big]$ and
  $\|w(u)\|_{H^s}\le C_2\hat a(u)(1+\rhohat(u))$.
\end{enumerate}
Moreover, let $\sigma_{\max}=\max_{G\in[0,\Gone]}\sigma(G)$, and let $\bar t_1(C_0)$ be the
least $t\ge0$ such that
\[
2^{p/2}C_0\jap{t'}^pe^{-(\sigb-\sigma')t'}\le1\qquad\text{for all }t'\ge t .
\]
Then $t_1(C_0,G)\le\bar t_1(C_0)$, and for all $u\ge0$,
\[
\hat a(u)\rhohat(u)\le C(C_0)\,\kappa\jap{u}^pe^{\sigma'u},\qquad
C(C_0)=\max\big(C_0,e^{\sigma_{\max}\bar t_1(C_0)}\big);
\]
the number $C(C_0)$ depends only on $C_0$, $\sigma'$, $p$, $\sigb$ and $\sigma_{\max}$.
\end{theorem}
\begin{proof}
Throughout,
$\Lin=\Lin(G)$, $F=-h\cdot\nabla\Phi$, and $\wlin$ is the function $\wlin_0$ of
Lemma~\ref{lem:running} with $w_0=w(0)$; it is an admissible choice of $w^{[1]}$
in Equation~\eqref{eq:hier}, because $\supp(\Phi+w(0))\subset\NNm\subset N$ by
(Init). The proof has two stages. On
the initial layer $[0,\min(\bar u,t_1)]$, whose length is at most $\bar t_1(C_0)$,
we compare $w$ with $\wlin$ by Lemma~\ref{lem:energy}. At the time $t_1$ the
transient has decayed, and we restart the problem there and apply
Lemma~\ref{lem:window} on the rest of the interval.

\emph{Constants.} Let $C_\phi$, $\bar C_S$, $C_{\mathrm{bil}}$ and
$C_{\mathrm{lin}}$ be the numbers of Equation~\eqref{eq:W-const}, $\|\ellst\|$ and
$\Cp$ as in Lemma~I.\ref{I-lem:uniform}, and put
$C_\natural=\max(4,\,6C_\phi+2C_{\mathrm{lin}}+1)$; this choice makes the restarted
problem below satisfy the assumptions of Lemma~\ref{lem:window}. Let $\delta_W,\kappa_W,C_W$ be the
constants of Lemma~\ref{lem:window} for this $C_\natural$, and put
\[
C_2=6C_\phi+3\big(C_W+C_{\mathrm{lin}}+1\big),\qquad C_3=(1+\Cp)C_2,\qquad
\dH=\min\Big(\delta_W,\frac a{32C_2}\Big).
\]
These depend only on $(\Phi,\Gone,s,\sigma')$. Now let $C_0\ge1$, write
$\bar t_1=\bar t_1(C_0)$, $\Lambda'=\Lam(M-1)$, $\Lambda''=\Lam(M-2)$,
$C_{\mathrm f}=C_0\jap{\bar t_1}^p+e^{\sigma_{\max}\bar t_1}$, and define the
constants of the initial layer by
\[
\begin{gathered}
C^{\mathrm{lay}}_1=e^{\Lambda'\bar t_1}\big(C_0+C_{\mathrm{bil}}C_{\mathrm f}\bar t_1\big),\qquad
C^{\mathrm{lay}}_2=\bar t_1\big(C_{\mathrm f}+C_{\mathrm{BS}}C^{\mathrm{lay}}_1\big),\\
C^{\mathrm{lay}}_3=\bar C_Se^{\sigma_{\max}\bar t_1}\big(C_0+C_{\mathrm{bil}}C_{\mathrm f}\bar t_1\big),\qquad
C^{\mathrm{lay}}_4=\bar t_1e^{\Lambda''\bar t_1}C_{\mathrm{bil}}\big(C^{\mathrm{lay}}_3+C_{\mathrm f}\big)C^{\mathrm{lay}}_3 .
\end{gathered}
\]
The smallness of $\kappa$ is imposed by three groups of conditions: (L1) is used on
the initial layer, (L2) to bound the terms that contain the error $v_1=w-\wlin$ of
the initial layer, and (L3) to apply Lemma~\ref{lem:window} after the restart.
Let $\kappa_0(C_0)$ be the largest $x\in(0,(a/32)^2]$ such that every
$\kappa\in(0,x]$ satisfies:
\begin{enumerate}[label=\textup{(L\arabic*)}]
\item $C_{\mathrm f}\kappa\le1$, $C^{\mathrm{lay}}_1\kappa\le\frac12$, $C^{\mathrm{lay}}_3\kappa\le1$ and
  $C^{\mathrm{lay}}_2\kappa\le\min(\frac a{128},\frac12\kappa^{1/2})$;
\item $(1+\|\ellst\|)C^{\mathrm{lay}}_4\kappa\le1$ and
  $4(1+C_\phi\|\ellst\|)C^{\mathrm{lay}}_4\kappa^{1/2}\le1$;
\item $e^{\sigma_{\max}\bar t_1}\kappa\le\kappa_W$ and
  $C_We^{\sigma_{\max}\bar t_1}\kappa\,\varsigma(\kappa)^{N_W}\le\frac12\kappa^{1/2}$.
\end{enumerate}
This number exists. Each inequality in (L1)--(L3) is non-strict between continuous
functions of $\kappa\in(0,1]$, so the set $S$ of $\kappa\in(0,(a/32)^2]$ satisfying all of
them is closed in $(0,(a/32)^2]$; and $S\supset(0,\epsilon]$ for some $\epsilon>0$,
because in each inequality the ratio of the left side to the right side tends to $0$
as $\kappa\to0^+$. So the supremum $x_0$ of the $x$ with $(0,x]\subset S$ is positive,
$(0,x_0)\subset S$, and $x_0\in S$ by closedness. The number $\kappa_0(C_0)$ depends
only on $C_0$ and $(\Phi,\Gone,s,\sigma')$, and by (Small) the $\kappa$ of the theorem
satisfies (L1)--(L3) and $\kappa\le(a/32)^2$.

\emph{The time $t_1$.} Since $\sigma\ge\sigb$,
$2^{p/2}C_0\jap{t'}^pe^{-(\sigma-\sigma')t'}\le2^{p/2}C_0\jap{t'}^pe^{-(\sigb-\sigma')t'}$, so
$t_1\le\bar t_1$. For $u\ge t_1$, the definition of $t_1$ with $t'=u$ gives
\begin{equation}
\label{eq:A-t1}
C_0\jap{u}^pe^{-\sigma u}\le2^{-p/2}e^{-\sigma'u}\le1,\qquad\text{hence }|c(u)|\le3\kappa\text{ by (Coef)}.
\end{equation}
For $u<t_1$,
$\hat a(u)\rhohat(u)=C_0\kappa\jap{u}^pe^{\sigma'u}$. For $u\ge t_1$,
$\hat a(u)\rhohat(u)=\kappa e^{\sigma t_1}\jap{u-t_1}^pe^{\sigma'(u-t_1)}\le e^{\sigma_{\max}\bar t_1}\kappa\jap{u}^pe^{\sigma'u}$.
This proves the last assertion of the theorem.

\emph{The initial layer $[0,\bar u_1]$, $\bar u_1=\min(\bar u,t_1)$.} Since
$t_1\le\bar t_1$ and $\eta_b\le1$, (Field) gives
$\|h(u)\|_{C^M(N')}\le C_{\mathrm f}\kappa\le1$ there. Let $u_*$ be the largest
$u\in[0,\bar u_1]$ such that $\delta_{\mathrm r}(u)\le a/64$ and $\|w(v)\|_{H^{M-1}}\le1$ for
$v\in[0,u]$; both hold at $u=0$ by (Init) and (L1). For $u\le u_*$ the support of
$\omega_{\mathrm r}(u)$ lies in $\{(\frac78-\frac1{64})a\le|y|\le\frac54+\frac a{64}\}\subset N$
(Lemma~\ref{lem:annuli}), so Lemma~\ref{lem:reduction}\textup{(3)} applies on $[0,u_*]$. Lemma~\ref{lem:energy}
at level $M-1$ with $w_{\mathrm{app}}=0$, for which $R_{\mathrm{app}}=h\cdot\nabla\Phi$, and
Lemma~\ref{lem:bilinear}\textup{(3)} give
\[
\|w(u)\|_{H^{M-1}}\le e^{\Lambda'u}C_0\kappa+\int_0^ue^{\Lambda'(u-v)}C_{\mathrm{bil}}C_{\mathrm f}\kappa\dd v\le C^{\mathrm{lay}}_1\kappa\le\tfrac12 ,
\]
and Lemma~\ref{lem:reduction}\textup{(2)} and Lemma~\ref{lem:bilinear}\textup{(4)} give
$\delta_{\mathrm r}(u)\le\int_0^u\big(\|h\|_{C^0(N')}+C_{\mathrm{BS}}\|w\|_{H^1}\big)\le C^{\mathrm{lay}}_2\kappa\le a/128$.
If $u_*<\bar u_1$, continuity would give a larger admissible $u$; so $u_*=\bar u_1$.
Next, $\wlin\in C([0,\bar u];X^{M-1})\cap C^1([0,\bar u];X^{M-2})$ is real with
$\partial_u\wlin=\Lin\wlin+F$ (Lemma~\ref{lem:duhamel}), and
Proposition~I.\ref{I-prop:semigroup} with Lemma~\ref{lem:bilinear} gives
$\|\wlin(u)\|_{H^{M-1}}\le C^{\mathrm{lay}}_3\kappa\le1$ for $u\le\bar u_1$. So
Lemma~\ref{lem:energy} at level $M-2$ applies to $w_{\mathrm{app}}=\wlin$ on
$[0,\bar u_1]$. Its residual, Equation~\eqref{eq:hier-res} with $J=1$, is
$R_{\mathrm{app}}=-B(\wlin,\wlin)+h\cdot\nabla\wlin$, and
$\|R_{\mathrm{app}}\|_{H^{M-2}}\le C_{\mathrm{bil}}(C^{\mathrm{lay}}_3+C_{\mathrm f})C^{\mathrm{lay}}_3\kappa^2$.
Since $v_1:=w-\wlin$ vanishes at $u=0$,
\begin{equation}
\label{eq:A-layer}
\|v_1(u)\|_{H^{M-2}}\le C^{\mathrm{lay}}_4\kappa^2,\qquad \delta_{\mathrm r}(u)\le C^{\mathrm{lay}}_2\kappa\qquad(0\le u\le\bar u_1).
\end{equation}

\emph{Proof of (b) for $u\le\min(\bar u_1,\hat T)$.} By
Lemma~\ref{lem:linear-part} with $k=s$ and $A=C_0\kappa$, and since
$C_0\kappa\jap{u}^pe^{\sigma'u}\le\hat a(u)\rhohat(u)$ for $u\le t_1$ (with
equality for $u<t_1$, while at $u=t_1$ the left side is at most
$2^{-p/2}\hat a(t_1)$ by the definition of $t_1$ and $\rhohat(t_1)=1$),
\[
\begin{gathered}
\wlin(u)=2\Re\big(c(u)e^{\lambda u}\phist\big)+(1-\Riesz)\wlin(u),\\
\|(1-\Riesz)\wlin(u)\|_{H^s}\le C_{\mathrm{lin}}\hat a(u)\big(\rhohat(u)+\eta_b\big).
\end{gathered}
\]
By Equation~\eqref{eq:A-layer} and (L2),
$\|v_1(u)\|_{H^s}\le C^{\mathrm{lay}}_4\kappa^2\le\kappa^{1/2}\kappa\le\kappa^{1/2}\hat a(u)$. Hence
$\|w-2\Re(ce^{\lambda u}\phist)\|_{H^s}\le(C_{\mathrm{lin}}+1)\hat a\big[\rhohat+\kappa^{1/2}+\eta_b\big]$.
By (Coef), $|c(u)|e^{\sigma u}\le2\hat a(u)+C_0\kappa\jap{u}^p\le\hat a(u)(2+\rhohat(u))$,
so $\|w(u)\|_{H^s}\le2C_\phi(2+\rhohat)\hat a+(C_{\mathrm{lin}}+1)(\rhohat+2)\hat a\le C_2\hat a(1+\rhohat)$.
The displacement is at most $C^{\mathrm{lay}}_2\kappa\le\frac12\kappa^{1/2}$ by (L1). This
proves (b), and bounds the displacement by $\frac12\kappa^{1/2}$, when
$\min(\bar u,\hat T)\le t_1$.

\emph{Restart.} Assume now that $t_1<u_e:=\min(\bar u,\hat T)$. The function
$\tilde\omega(u')=\omega(t_1+u')$ is the solution of Theorem~\ref{thm:smooth} with the
$3$-fold rotationally symmetric datum $\omega(t_1)$, by Theorem~\ref{thm:smooth}\textup{(5)}, and its
flow is $\tilde\flow_{u'}=\flow_{t_1+u'}\circ\flow_{t_1}^{-1}$. With
$\tilde\omega_\bullet(0)=\omega_\bullet(t_1)$ we get $\tilde\omega_\bullet(u')=\omega_\bullet(t_1+u')$,
the same $G$ (since $\int\omega_{\mathrm i}(t_1)=2\pi G$), and the field
$\tilde h(u')=h(t_1+u')$. So the restarted configuration satisfies the assumptions of
Lemma~\ref{lem:reduction} on $[0,\bar u-t_1]$, and its perturbation is
$\tilde w(u')=w(t_1+u')$. Its running coefficient is
$\tilde c(u')=\ellst(w(t_1))+\int_0^{u'}e^{-\lambda v}\ellst(F(t_1+v))\dd v$. By
Equation~\eqref{eq:running} and Lemma~\ref{lem:running},
$e^{\lambda t_1}c(t_1+u')=\ellst(\wlin(t_1))+\int_0^{u'}e^{-\lambda v}\ellst(F(t_1+v))\dd v$, so
\begin{equation}
\label{eq:A-restart}
\tilde c(u')=e^{\lambda t_1}c(t_1+u')+\ellst(v_1(t_1)).
\end{equation}

We check the assumptions of Lemma~\ref{lem:window} for the restarted problem,
with $C_\natural$ as above, $\kappa_1=\hat a(t_1)=\kappa e^{\sigma t_1}$, the same
$\delta_0$ and $\eta_b$, and $\bar u'=u_e-t_1$. First,
$\kappa\le\kappa_1\le e^{\sigma_{\max}\bar t_1}\kappa\le\kappa_W$ by (L3),
$\delta_0\le\dH\le\delta_W$, and $\sigma^{-1}\log(\delta_0/\kappa_1)=\hat T-t_1\ge\bar u'$.
By Equation~\eqref{eq:A-t1} at $u=t_1$, $|c(t_1)|e^{\sigma t_1}\le3\kappa_1$ and
$C_0\kappa\jap{t_1}^pe^{\sigma't_1}\le\kappa_1$, and
$\eta_b\kappa e^{\sigma t_1}\le\kappa_1$. So Lemma~\ref{lem:linear-part} at $k=M-2$,
Equation~\eqref{eq:A-layer} and (L2) give
\[
\|w(t_1)\|_{H^{M-2}}\le6C_\phi\kappa_1+C_{\mathrm{lin}}(\kappa_1+\kappa_1)+C^{\mathrm{lay}}_4\kappa^2\le(6C_\phi+2C_{\mathrm{lin}}+1)\kappa_1\le C_\natural\kappa_1 .
\]
Second, by Equation~\eqref{eq:A-layer} the support of $\Phi+w(t_1)$ lies in
$\{(\frac78-\frac1{64})a\le|y|\le\frac54+\frac a{64}\}$. Third, by Peetre's inequality
$\jap{t_1+u'}\le\sqrt2\jap{t_1}\jap{u'}$ and the definition of $t_1$,
\[
C_0\kappa\jap{t_1+u'}^p\le\big[2^{p/2}C_0\jap{t_1}^pe^{-(\sigma-\sigma')t_1}\big]e^{-\sigma't_1}\kappa_1\jap{u'}^p\le\kappa_1\jap{u'}^p,
\]
and $\eta_b\hat a(t_1+u')=\eta_b\kappa_1e^{\sigma u'}$; so
$\|\tilde h(u')\|_{C^M(N')}\le C_\natural\kappa_1\jap{u'}^p+\eta_b\kappa_1e^{\sigma u'}$. Fourth,
by Equations~\eqref{eq:A-restart}, \eqref{eq:A-t1} and \eqref{eq:A-layer} and (L2),
$|\tilde c(u')|\le3e^{\sigma t_1}\kappa+\|\ellst\|C^{\mathrm{lay}}_4\kappa^2\le4\kappa_1\le C_\natural\kappa_1$.

\emph{Proof of (b) for $u\in[t_1,u_e]$.} Lemma~\ref{lem:window} applies; put
$u'=u-t_1$. Then $\kappa_1e^{\sigma u'}=\hat a(u)$ and
$\jap{u'}^pe^{-(\sigma-\sigma')u'}=\rhohat(u)$. Since $\varsigma$ is decreasing and
$\kappa_1\ge\kappa$, (L3) gives $C_W\kappa_1\varsigma(\kappa_1)^{N_W}\le\frac12\kappa^{1/2}$.
By Equation~\eqref{eq:A-restart},
$\tilde c(u')e^{\lambda u'}=c(u)e^{\lambda u}+\ellst(v_1(t_1))e^{\lambda u'}$, and
$|\ellst(v_1(t_1))e^{\lambda u'}|\le\|\ellst\|C^{\mathrm{lay}}_4\kappa^2e^{\sigma u'}\le\|\ellst\|C^{\mathrm{lay}}_4\kappa\,\hat a(u)$.
Since $\|2\Re(X\phist)\|_{H^s}\le2C_\phi|X|$ for $X\in\C$, (L2) gives
\[
\begin{aligned}
\big\|w(u)-2\Re\big(c(u)e^{\lambda u}\phist\big)\big\|_{H^s}
&\le C_W\hat a(u)\big[\rhohat(u)+\hat a(u)+\tfrac12\kappa^{1/2}+\eta_b\big]\\
&\quad+\tfrac12\kappa^{1/2}\hat a(u),
\end{aligned}
\]
which gives the first bound of (b), since $C_W+1\le C_2$. By
Equation~\eqref{eq:A-t1}, the main term $2\Re(c(u)e^{\lambda u}\phist)$ has $H^s$ norm
at most $6C_\phi\hat a(u)$. Since $\hat a+\kappa^{1/2}+\eta_b\le3$, we get
\[
\|w(u)\|_{H^s}\le6C_\phi\hat a+(C_W+1)\hat a(\rhohat+3)\le C_2\hat a(1+\rhohat).
\]
For the displacement, a trajectory of $\omega_{\mathrm r}$ satisfies, with $z=\flow_{t_1}(y)$,
\[
|\flow_u(y)|-|y|=\big(|\tilde\flow_{u'}(z)|-|z|\big)+\big(|\flow_{t_1}(y)|-|y|\big),
\]
so by Lemma~\ref{lem:window} and Equation~\eqref{eq:A-layer} it moves radially by at
most $C_W\delta_0+\frac12\kappa^{1/2}+C^{\mathrm{lay}}_2\kappa\le C_2\delta_0+\kappa^{1/2}$.

\emph{Proof of (a).} In both cases the displacement is at most
$C_2\delta_0+\kappa^{1/2}\le\frac a{32}+\frac a{32}=\frac a{16}$, because
$\delta_0\le\dH\le a/(32C_2)$ and $\kappa\le(a/32)^2$. Since
$\supp\omega_{\mathrm r}(0)\subset\NNm$ by (Init), Lemma~\ref{lem:annuli} gives
$\supp(\Phi+w(u))\subset\nobreak\NN$.
\end{proof}

Recall from the expansion after Definition~I.\ref{I-def:multipole} that the outer
moment $\pout_3$ of a function determines the part in the angular modes $\pm3$ of its
field outside its support. So the outer moment of the deviation of a ring
determines the part in these modes of the field that the ring exerts on the rings
outside it. Recall also from Definition~I.\ref{I-def:eigendata} and the discussion
before it that $\pout_*(G)=\pout_3[\phist(G)]$ is the outer moment of the unstable
eigenfunction.
The next corollary gives this moment for an active ring.

\begin{corollary}
\label{cor:active-moment}
Under the assumptions of Theorem~\ref{thm:active}, put
$q(u)=2\pi\,\pout_3[w(u)]$. Then for every $u\in[0,\min(\bar u,\hat T)]$,
\[
\big|q(u)-2\pi\pout_*(G)\,c(u)e^{\lambda u}\big|\le2\pi C_3\,\hat a(u)\big[\rhohat(u)+\hat a(u)+\kappa^{1/2}+\eta_b\big].
\]
\end{corollary}
\begin{proof}
For $X\in\C$, $\pout_3[2\Re(X\phist)]=X\pout_3[\phist]+\bar X\pout_3[\overline{\phist}]$. The
second term vanishes: $\overline{\phist}$ lies in the angular mode $-3$, and
$\int_0^{2\pi}e^{-3i\theta}e^{-3i\theta}\dd\theta=0$. Since
$\pout_*(G)=\pout_3[\phist(G)]$, linearity gives
$\pout_3[w(u)]-\pout_*(G)c(u)e^{\lambda u}=\pout_3\big[w(u)-2\Re(c(u)e^{\lambda u}\phist)\big]$.
The function in brackets is supported in $N$, so Lemma~I.\ref{I-lem:uniform} and
Theorem~\ref{thm:active}\textup{(b)} bound its modulus by
$\Cp C_2\hat a(u)\big[\rhohat(u)+\hat a(u)+\kappa^{1/2}+\eta_b\big]$. Multiplying by $2\pi$
and using $\Cp C_2\le C_3$ proves the corollary.
\end{proof}

\subsection{The Seeded Ring}

The approximate datum of Definition~\ref{def:approx} replaces its innermost ring by
a rearrangement $\Phi^\kappa$ of $\Phi$ whose deviation from $\Phi$ has an unstable
component of size about $\kappa$. We obtain $\Phi^\kappa$ by composing $\Phi$ with
the time-$\kappa$ flow of a divergence-free field $V$. This flow preserves area, so
$\Phi^\kappa$ is equimeasurable with $\Phi$, and $V$ is chosen with
$V\cdot\nabla\Phi=2\Re\phist(0)$, so that $\Phi^\kappa-\Phi$ is $2\kappa\Re\phist(0)$ to
first order in $\kappa$.

\begin{lemma}
\label{lem:seed}
Assume \textup{(N1)}--\textup{(N3)}, let $s\ge3$ be an integer, and let
$M=s+N_G+2$ be the integer of Theorem~\ref{thm:active} for this $s$. There are a
number $\kL>0$ and functions $\Phi^\kappa$, $\kappa\in(0,\kL]$, all depending only
on $\Phi$, and a number $\Cseed<\infty$, depending only on $\Phi$, $\Gone$ and $s$,
such that for every
$\kappa\in(0,\kL]$: $\Phi^\kappa\in C_c^\infty(\R^2)$, $\Phi^\kappa\ge0$,
$\Phi^\kappa$ is $3$-fold rotationally symmetric and equimeasurable with $\Phi$;
$\supp\Phi^\kappa\subset\NNm$;
$\|\Phi^\kappa-\Phi\|_{H^{M-1}}\le\Cseed\kappa$;
and $|\ellst(0)(\Phi^\kappa-\Phi)|\in[\frac\kappa2,2\kappa]$.
\end{lemma}
\begin{proof}
Let $G=0$, $\lambda=\lamst(0)$, and let $\psi_*$ and $n_*=n_*(0)$ be
as in Theorem~I.\ref{I-thm:evans-simple} and Definition~I.\ref{I-def:eigendata}, for the
ring in ring units and $\ell=3$. By these definitions,
$\phist(0)=n_*^{-1}g_*(r)e^{3i\theta}$ with $g_*=\Mres_{\lambda}\kcoef_3\,\psi_*(\log r)$,
where $\kcoef_3(r)=3i\Phi'(r)/r$ and $\Mres_\lambda=(\lambda+3i\Omega(\cdot;0))^{-1}$
(Definition~I.\ref{I-def:sector-setting}). So the mode-$3$ profile of
$\phist(0)$ is
\[
\vphist(r)=n_*^{-1}\,\frac{3i\,\Phi'(r)}{r}\,\frac{\psi_*(\log r)}{\lambda+3i\Omega(r;0)} .
\]
Put $g_{\mathrm s}(r)=n_*^{-1}\psi_*(\log r)/(\lambda+3i\Omega(r;0))$ for $r>0$. It is
smooth: $\psi_*$ solves a linear equation with smooth coefficients
(Lemma~I.\ref{I-lem:sector}\textup{(3)}), $\Omega(\cdot;0)$ is smooth on $(0,\infty)$, and
$|\lambda+3i\Omega|\ge\Re\lambda>0$. Let $\chiseed\in C_c^\infty(\{a/2<|y|<2\})$ be
radial with $\chiseed=1$ on $\{a\le|y|\le1\}$, and put
$\psi_{\mathrm{seed}}=-2\Re\big(\chiseed(r)\,g_{\mathrm s}(r)\,e^{3i\theta}\big)$ and
$V=\nabla^\perp\psi_{\mathrm{seed}}$. Then $\psi_{\mathrm{seed}}$ is real, smooth,
$3$-fold rotationally symmetric and supported in the compact set
$K_V=\supp\chiseed\subset\{a/2<|y|<2\}$, and $V\in C_c^\infty(\R^2;\R^2)$ is
divergence-free and $3$-fold rotationally symmetric. Let $Y_\kappa=\flow_V(\kappa,0,\cdot)$ be the
time-$\kappa$ flow of $V$, with $\flow_V$ as in Lemma~\ref{lem:flow} for the
time-independent bounded field $V$. By Lemma~\ref{lem:flow}\textup{(a)}, \textup{(b)},
\textup{(d)}, $Y$ is defined for all $\kappa\in\R$ and is $C^\infty$ in $(\kappa,y)$,
each $Y_\kappa$ is a $C^\infty$ diffeomorphism of $\R^2$, and $\det DY_\kappa=1$
because $\divop V=0$. Since $V$ does not depend on time, uniqueness for the flow
gives $Y_{\kappa+\kappa'}=Y_\kappa\circ Y_{\kappa'}$ and, since
$V(\Orot y)=\Orot V(y)$, $Y_\kappa\circ\Orot=\Orot\circ Y_\kappa$. For $y\notin K_V$
the constant path $y$ solves $\dot x=V(x)$, since $V(y)=0$; so $Y_\kappa$ is the identity outside
$K_V$.

Put $\Phi^\kappa=\Phi\circ Y_\kappa$. Then $\Phi^\kappa\in C_c^\infty(\R^2)$ is nonnegative and $3$-fold
rotationally symmetric. It is equimeasurable with $\Phi$: for $c>0$,
$\{\Phi^\kappa>c\}=Y_\kappa^{-1}(\{\Phi>c\})$ has the measure of $\{\Phi>c\}$ because
$Y_\kappa$ preserves area, and $\{\Phi^\kappa<-c\}=\{\Phi<-c\}=\emptyset$. Since
$\Phi$ is radial, $V\cdot\nabla\Phi=(V\cdot\er)\Phi'$, and
$V\cdot\er=-r^{-1}\partial_\theta\psi_{\mathrm{seed}}=2\Re\big(\frac{3i}r\chiseed g_{\mathrm s}e^{3i\theta}\big)$.
On $\supp\Phi'\subset\{a\le|y|\le1\}$ we have $\chiseed=1$, so
$(V\cdot\er)\Phi'=2\Re(\vphist e^{3i\theta})=2\Re\phist(0)$; off $\supp\Phi'$ both
sides vanish, because $\vphist$ is a multiple of $\Phi'$. Thus
\begin{equation}
\label{eq:seed-first}
V\cdot\nabla\Phi=2\Re\phist(0).
\end{equation}

We now expand in $\kappa$. Fix $y$ and put $f_y(t)=\Phi(Y_t(y))$. Then
$f_y'(t)=(V\cdot\nabla\Phi)(Y_t(y))$ and $f_y''(t)=(V\cdot\nabla(V\cdot\nabla\Phi))(Y_t(y))$,
so with $\Psi_{\mathrm s}=V\cdot\nabla(V\cdot\nabla\Phi)\in C_c^\infty(\R^2)$,
\[
\Phi^\kappa-\Phi=\int_0^\kappa(V\cdot\nabla\Phi)\circ Y_t\dd t,\qquad
\Phi^\kappa-\Phi-\kappa V\cdot\nabla\Phi=\int_0^\kappa(\kappa-t)\,\Psi_{\mathrm s}\circ Y_t\dd t .
\]
The derivatives of $Y_t(y)-y$ of every order vanish outside $K_V$ and are
continuous on $[-1,1]\times K_V$, hence bounded there. By the chain rule and
$\det DY_t=1$, for every $k$ there is $C^Y_k$ with $\|g\circ Y_t\|_{H^k}\le C^Y_k\|g\|_{H^k}$
for $g\in C_c^\infty(\R^2)$ and $|t|\le1$. Put $C'_0=\frac12C^Y_0\|\Psi_{\mathrm s}\|_{L^2}$, which depends only on $\Phi$, and
$\Cseed=C^Y_{M-1}\|2\Re\phist(0)\|_{H^{M-1}}$. For $0<\kappa\le1$,
Equation~\eqref{eq:seed-first} and the two identities give
$\|\Phi^\kappa-\Phi-2\kappa\Re\phist(0)\|_{L^2}\le C'_0\kappa^2$ and
$\|\Phi^\kappa-\Phi\|_{H^{M-1}}\le\Cseed\kappa$.

Finally put
$\kL=\min\big(1,\,a/(8(1+\|V\|_{L^\infty})),\,1/(2(1+\|\ellst(0)\|_{(X^0)^*}C'_0))\big)$,
and let $\kappa\in(0,\kL]$. Since $|Y_\kappa^{-1}(x)-x|\le\kappa\|V\|_{L^\infty}\le a/8$ by
Lemma~\ref{lem:flow}\textup{(b)}, \textup{(c)}, and
$\supp\Phi^\kappa=Y_\kappa^{-1}(\supp\Phi)$ with $\supp\Phi\subset\{a\le|y|\le1\}$, we get
$\supp\Phi^\kappa\subset\{\frac78a\le|y|\le1+\frac a8\}\subset\NNm$. The function
$\Phi^\kappa-\Phi$ is supported in $N$ and $3$-fold rotationally symmetric, so it lies in $X^0$.
By Definition~I.\ref{I-def:eigendata}, $\ellst(0)(2\Re\phist(0))=\ellst(0)(\phist(0))+\ellst(0)(\overline{\phist(0)})=1$.
Hence
$|\ellst(0)(\Phi^\kappa-\Phi)-\kappa|\le\|\ellst(0)\|_{(X^0)^*}C'_0\kappa^2\le\frac\kappa2$, and
$|\ellst(0)(\Phi^\kappa-\Phi)|\in[\frac\kappa2,\frac{3\kappa}2]\subset[\frac\kappa2,2\kappa]$.
\end{proof}

\section{The Datum, the Approximants and the Phase Clock}
\label{sec:datum}

We define the datum $f$ of Theorem~\ref{thm:main} and its approximants
$f^{(n)}$; the solutions $\omega^{(n)}$ with these data form the
$2$-admissible sequence to which Corollary~I.\ref{I-cor:defect} is applied in
the proof of Proposition~\ref{prop:reduction}. We then split each $\omega^{(n)}$
into its rings, describe each ring in its own units by its perturbation and by
the field that the other rings exert on it, and define the phase clock, which
fixes, from $\omega^{(n)}$ itself, the time at which each ring becomes active and
the time at which its unstable coefficient is expected to reach the size
$\delta_0$.

\begin{definition}
\label{def:datum}
Let $A_1>0$, $\mu\in(1,2]$ and $\sbar\in(0,a/16]$, and define $A_k$, $s_k$,
$r_k$ ($k\ge1$) and $f$ by Equation~\eqref{eq:datum}.
\end{definition}

\begin{proposition}
\label{prop:integrability}
In Definition~\ref{def:datum}, the sets $\{ar_k\le|x|\le r_k\}$, $k\ge1$, are
pairwise disjoint.
The function $f$ is radial and nonnegative, $f\in (L^1\cap L^q)(\R^2)$ for every
$q\in[1,\infty)$, and $f\notin L^\infty(\R^2)$. Moreover
$\sum_{k\ge1}A_k^qr_k^2<\infty$ for every $q\in[1,\infty)$.
\end{proposition}
\begin{proof}
Since $\mu>1$ and $\sbar\le a/16$, $s_k\le\sbar/2<a$, so $r_{k+1}<ar_k$ and every
annulus of index $i>k$ lies in the open disk $B_{ar_k}$, because $r_i\le r_{k+1}$;
this disk does not meet the annulus of index $k$, so the annuli are pairwise
disjoint. Moreover
$r_k=\prod_{j<k}s_j=\sbar^{k-1}(2\mu)^{-k(k-1)/2}$.

By Lemma~I.\ref{I-lem:ring-facts}, $\Phi$ is radial and continuous, $0\le\Phi<1$,
$\Phi\not\equiv0$ and $\supp\Phi\subset\{a\le|y|\le1\}$. So the $k$-th term of
the series is radial, nonnegative and supported in the $k$-th annulus; since the
annuli are disjoint, at every point at most one term is nonzero, $f$ is radial and nonnegative, and, by
monotone convergence and the formula for $r_k$, for $q\in[1,\infty)$,
\begin{gather*}
\int f^q\dd x=\sum_{k\ge1}A_k^qr_k^2\|\Phi\|_{L^q}^q,\qquad
\|\Phi\|_{L^q}^q\le|\{a\le|y|\le1\}|\le\pi,\\
A_k^qr_k^2=A_1^q\mu^{q(k-1)}\sbar^{2(k-1)}(2\mu)^{-k(k-1)}=A_1^q\big(\mu^q\sbar^2(2\mu)^{-k}\big)^{k-1}.
\end{gather*}
If $(2\mu)^k\ge2\mu^q$, which holds for all large $k$ because $2\mu>1$, the
bracket is at most $\frac12$, because $\sbar<1$, so
$A_k^qr_k^2\le A_1^q2^{1-k}$ for all large $k$. Hence $\sum_kA_k^qr_k^2<\infty$
and $f\in L^q$. Finally, $f\ge A_k\Phi(\cdot/r_k)$, and, $\Phi$ being continuous,
$A_k\Phi(\cdot/r_k)>\frac12A_k\max\Phi>0$ on a nonempty open set; since $\mu>1$
gives $A_k\to\infty$, $f\notin L^\infty$.
\end{proof}

The approximants keep the first $n$ rings of $f$ and replace ring $n$, the
innermost of them, whose time scale $1/A_n$ is the shortest, by the seeded ring
$\Phi^{\kseed}$ of Lemma~\ref{lem:seed}. This is a nonnegative, $3$-fold
rotationally symmetric function, equimeasurable with $\Phi$, whose coefficient
$\ellst(0)(\Phi^{\kseed}-\Phi)$ along the unstable eigenfunction $\phist(0)$ has
modulus between $\kseed/2$ and $2\kseed$. The relay described in the
Introduction starts from this seed.

\begin{definition}
\label{def:approx}
Assume \textup{(N1)}--\textup{(N3)}, and let $\kL$ and $\Phi^\kappa$ be as in
Lemma~\ref{lem:seed}. For an integer $n\ge2$ and $\kseed\in(0,\kL]$ put
$f^{(n)}=\sum_{k<n}A_k\Phi(\cdot/r_k)+A_n\Phi^{\kseed}(\cdot/r_n)$, let $\omega^{(n)}$ be the solution of Theorem~\ref{thm:smooth} with datum
$f^{(n)}$, let $v^{(n)}=\KBS\ast\omega^{(n)}$, and let $\flow_t$ be the flow of
$v^{(n)}$.
\end{definition}

\begin{lemma}
\label{lem:adm}
In Definition~\ref{def:approx}, $f^{(n)}\in C_c^\infty(\R^2)$,
$f^{(n)}$ and $\omega^{(n)}(t)$ are $3$-fold rotationally symmetric for every
$t$, and, for fixed $\kseed$, $f^{(n)}\to f$ in $(L^1\cap L^q)(\R^2)$ for every
$q\in[1,\infty)$ as $n\to\infty$. Moreover $(f^{(n)},\omega^{(n)})_{n\ge2}$ is a
$2$-admissible sequence for $f$ on $[0,\infty)$ in the sense of
Definition~I.\ref{I-def:adm} (the sequence starts at
$n=2$; the shift $n\mapsto n+1$ gives the indexing of Definition~I.\ref{I-def:adm}).
\end{lemma}
\begin{proof}
By Lemmas~I.\ref{I-lem:ring-facts} and \ref{lem:seed}, $\Phi$ and $\Phi^{\kseed}$ lie
in $C_c^\infty(\R^2)$, are $3$-fold rotationally symmetric ($\Phi$ is radial)
and equimeasurable, so
they have the same $L^q$ norms. Hence $f^{(n)}\in C_c^\infty(\R^2)$ is
$3$-fold rotationally symmetric, and every $\omega^{(n)}(t)$ is $3$-fold
rotationally symmetric, by Theorem~\ref{thm:smooth}\textup{(6)} with the rotation $\Orot$. Let
$q\in[1,\infty)$. In
$f-f^{(n)}=A_n(\Phi-\Phi^{\kseed})(\cdot/r_n)+\sum_{k>n}A_k\Phi(\cdot/r_k)$, the
first term has $L^q$ norm at most $2A_nr_n^{2/q}\|\Phi\|_{L^q}$, by the triangle
inequality and the equality of the $L^q$ norms of $\Phi$ and $\Phi^{\kseed}$; the
terms of the sum have disjoint supports, by Proposition~\ref{prop:integrability}, so
\[
\|f-f^{(n)}\|_{L^q}\le\Big(2A_nr_n^{2/q}
+\Big(\sum_{k>n}A_k^qr_k^2\Big)^{1/q}\Big)\|\Phi\|_{L^q},
\]
and both terms on the right tend to $0$ as $n\to\infty$ because
$\sum_kA_k^qr_k^2<\infty$ (Proposition~\ref{prop:integrability}): $A_n^qr_n^2$ is
a term, and $\sum_{k>n}A_k^qr_k^2$ a tail, of this convergent series. Finally, by
Theorem~\ref{thm:smooth}\textup{(7)}, for $q=2$ and every
$T>0$ the restriction of $\omega^{(n)}$ to $[0,T)$ is a renormalized weak
solution in $X_2(T)$ with datum $f^{(n)}$; so $\omega^{(n)}$ is a renormalized
weak solution on $[0,\infty)$, and with the convergence just proved this is the
$2$-admissibility of Definition~I.\ref{I-def:adm}.
\end{proof}

Each ring is studied in its own scale, in which the $k$-th term of $f^{(n)}$
becomes $\Phi$, or $\Phi^{\kseed}$ if $k=n$, so that the results of
Section~\ref{sec:single}, stated in ring units, can be applied to it. For $A,r>0$ the function
$\tilde\omega(\tau,y)=A^{-1}\omega^{(n)}(\tau/A,ry)$ is the solution of
Theorem~\ref{thm:smooth} with datum $A^{-1}f^{(n)}(r\,\cdot)$, with velocity
$(\KBS\ast\tilde\omega)(\tau,y)=(Ar)^{-1}v^{(n)}(\tau/A,ry)$ and flow
$y\mapsto r^{-1}\flow_{\tau/A}(ry)$. Indeed, $\KBS$ is homogeneous of degree
$-1$, which gives the velocity; then
$\partial_\tau\tilde\omega+(\KBS\ast\tilde\omega)\cdot\nabla\tilde\omega
=A^{-2}(\partial_t\omega^{(n)}+v^{(n)}\cdot\nabla\omega^{(n)})(\tau/A,ry)=0$, and
Theorem~\ref{thm:smooth}\textup{(5)} identifies $\tilde\omega$ with that solution.
With $A=A_k$ and $r=r_k$, $\tilde\omega$ is $\omega^{(n)}$ in \textbf{ring-$k$
units}: points $x$ become $y=x/r_k$, vorticities are divided by $A_k$, and
$\tau=A_kt$ is the \textbf{ring-$k$ time}. For a velocity field $V$ we write
$V_{[k]}(y)=V(r_ky)/(A_kr_k)$ for $V$ in ring-$k$ units.

The next definition splits $\omega^{(n)}$ into its rings and names the
quantities that the later sections estimate. The set $S_k(t)$ carries ring $k$
along the flow, and $\omega_{\ge k}$ is the part of $\omega^{(n)}$ formed by ring
$k$ and the rings inside it. In ring-$k$ units, the perturbation $w_k$ is the
deviation of ring $k$ from $\Phi$; $2\pi G_k$ is the circulation of the point
vortex by which the linearization $\Lin(G_k)$ of Theorem~I.\ref{I-thm:linear}
represents the rings inside ring $k$; and the external field $h_k$ is the field
of all the other rings on $N'$ minus the field of this point vortex, that is,
the field $h$ of Lemma~\ref{lem:reduction}.

The number $\varepsilon_k$ is the small factor in the bounds of
Lemma~\ref{lem:transfer}\textup{(1)}--\textup{(3)} on the field, minus its
monopole, that the material inside ring $k$ induces on ring $k$. The
normalized moment $q_{k+1}$ is the mode-$3$ moment through which this material
drives the unstable coefficient of ring $k$ (Lemma~\ref{lem:coupling}), and the
Duhamel majorant $\maj_k$ bounds, up to a constant factor, the $H^{s_h}$ norm
of $w_k$ while ring $k$ is quiet and $\maj_k$ is small
(Theorem~\ref{thm:quiet}\textup{(c)}).

\begin{definition}
\label{def:regions}
Fix $n\ge2$ and the objects of Definition~\ref{def:approx}. For $1\le k\le n$:
\begin{itemize}
\item $f^{(n)}_k$ is the $k$-th term of $f^{(n)}$, that is,
  $f^{(n)}_k=A_k\Phi(\cdot/r_k)$ for $k<n$ and
  $f^{(n)}_n=A_n\Phi^{\kseed}(\cdot/r_n)$;
  \[
  S_k=\supp f^{(n)}_k,\quad S_k(t)=\flow_t(S_k),\quad D_k(t)=\bigcup_{i=k}^nS_i(t),\quad
  D_{n+1}(t)=\emptyset,
  \]
  and $\omega_{\ge k}(t)=\omega^{(n)}(t)\mathbf 1_{D_k(t)}$;
\item $N_k=r_kN$, $N'_k=r_kN'$, $\Nflat_k=r_k\Nflat$, $\NN_k=r_k\NN$,
  $\NNm_k=r_k\NNm$ (the annuli of Section~I.\ref{I-sec:ring-def} and of Section~\ref{sec:II-prelim} dilated by $r_k$);
\item while $S_k(t)\subset N_k$, the \textbf{perturbation} of ring $k$ at the
  time $t$, in ring-$k$ units, is
  $w_k(t,y)=A_k^{-1}\omega^{(n)}(t,r_ky)\mathbf 1_{S_k(t)}(r_ky)-\Phi(y)$;
\item the \textbf{inner circulation} $G_k$ is defined by
  $2\pi G_kA_kr_k^2=\int\omega_{\ge k+1}(0)$; we put
  $\varepsilon_k=\mu s_k^{m+2}$, and, if $G_k\le\Gone$, $\sigma_k=\sigma(G_k)$,
  $\lambda_k=\lamst(G_k)$, $x_k=\sigma_k/\mu$ and
  $\beta_k=\beta(G_k)s_k^{m+2}/(2\pi)$;
\item for $k<n$, the \textbf{normalized moment} of the inner region is
  \[
  q_{k+1}(t)=\overline{\Qm[\omega_{\ge k+1}(t)]}/(A_{k+1}r_{k+1}^{m+2})
  \]
  (the support of $\omega_{\ge k+1}(t)$ is compact and does not contain $0$, by
  Lemma~\ref{lem:regions-basic}\textup{(2)} below, so $\Qm$ applies);
\item at a time $t$ with $D_{k+1}(t)\subset B_{ar_k/8}$, the \textbf{inner field}
  of ring $k$ at ring-$k$ time $\tau=A_kt$ is
  \[
  \hin_k(\tau)=\KBS\ast\nu^{\mathrm{in}}_k-G_k\frac{y^\perp}{|y|^2}\quad\text{on }N',\qquad
  \nu^{\mathrm{in}}_k(y)=A_k^{-1}\omega_{\ge k+1}(t,r_ky);
  \]
  at a time $t$ with $\big(S_1(t)\cup\dots\cup S_{k-1}(t)\big)\cap B_{6r_k}=\emptyset$
  (the rings $i<k$ lie outside $B_{6r_k}$), the \textbf{outer field} is
  $\hout_k(\tau)=\KBS\ast\nu^{\mathrm{out}}_k$ on $N'$, with
  $\nu^{\mathrm{out}}_k(y)=A_k^{-1}\omega^{(n)}(t,r_ky)\mathbf 1_{S_1(t)\cup\dots\cup S_{k-1}(t)}(r_ky)$;
  when both are defined, $h_k=\hin_k+\hout_k$ is the \textbf{external field} of
  ring $k$;
\item for an integer $M_h\ge1$ (from Section~\ref{sec:relay} on, $M_h$ is the
  integer of \textup{(SA2)} in Definition~\ref{def:standing}) we put
  $\etin_k(\tau)=\|\hin_k(\tau)\|_{C^{M_h}(N')}$,
  $\etout_k(\tau)=\|\hout_k(\tau)\|_{C^{M_h}(N')}$ and $\eta_k=\etin_k+\etout_k$,
  as functions of ring-$k$ time;
\item if $G_k\le\Gone$ and $\eta_k$ is defined on $[0,\tau]$ (it is then
  continuous there, as we show after Lemma~\ref{lem:regions-basic}), the
  \textbf{Duhamel majorant} of ring $k$ is
  \[
  \maj_k(\tau)=\int_0^\tau e^{\sigma_k(\tau-v)}\eta_k(v)\dd v ,
  \]
  and $\maj^{\mathrm{in}}_k$ is defined in the same way from $\etin_k$.
\end{itemize}
\end{definition}

\begin{lemma}
\label{lem:regions-basic}
In Definition~\ref{def:regions} the following hold for every $t\ge0$.
\begin{enumerate}[label=\textup{(\arabic*)},beginpenalty=10000]
\item The sets $S_1,\dots,S_n$ are compact, pairwise disjoint and invariant
  under $\Orot$, and none contains $0$. Moreover
  $S_k\subset\{ar_k\le|x|\le r_k\}$ for $k<n$, $S_n\subset\NNm_n$, and
  $f^{(n)}\mathbf 1_{S_k}=f^{(n)}_k$.
\item $\flow_t\circ\Orot=\Orot\circ\flow_t$ and $\flow_t(0)=0$. The sets
  $S_1(t),\dots,S_n(t)$ are compact, pairwise disjoint and invariant under
  $\Orot$, none contains $0$, and
  $\supp\omega^{(n)}(t)\subset S_1(t)\cup\dots\cup S_n(t)$.
\item For $1\le k\le n$,
  \begin{equation}
  \label{eq:omega-ge}
  \omega^{(n)}(t)\mathbf 1_{S_k(t)}=f^{(n)}_k\circ\flow_t^{-1},\qquad
  \omega_{\ge k}(t)=\Big(\sum_{i=k}^nf^{(n)}_i\Big)\circ\flow_t^{-1}
  =\sum_{i=k}^n\omega^{(n)}(t)\mathbf 1_{S_i(t)},
  \end{equation}
  and $\int\omega^{(n)}(t)\mathbf 1_{S_k(t)}=A_kr_k^2\|\Phi\|_{L^1}$. These
  functions lie in $C_c^\infty(\R^2)$, are nonnegative and are $3$-fold
  rotationally symmetric.
\item If $S_k(t)\subset N_k$, then $w_k(t)\in C_c^\infty(\R^2)$ is real valued,
  $3$-fold rotationally symmetric and supported in $N$; so $w_k(t)\in X^s$ for every integer
  $s\ge0$. Moreover $\int w_k(t)=0$, $\|w_k(t)\|_{L^1}\le2\|\Phi\|_{L^1}$, and
  $\|w_k(t)\|_{L^1}\le|N|^{1/2}\|w_k(t)\|_{L^2}$.
\end{enumerate}
\end{lemma}

\begin{proof}
(1) For $k<n$, $f^{(n)}_k$ is radial and supported in $\{ar_k\le|x|\le r_k\}$
(Lemma~I.\ref{I-lem:ring-facts}); $f^{(n)}_n$ is $3$-fold rotationally symmetric and supported in
$\NNm_n$ (Lemma~\ref{lem:seed}). So each $S_k$ is compact, $\Orot$-invariant and
avoids $0$. Let $i<k$. By these support properties, $S_k\subset\overline B_{5r_k/4}$, and
$\frac54r_k\le\frac54s_ir_i<ar_i$, because $r_k\le r_{i+1}=s_ir_i$ and $s_i\le\sbar/2\le a/32$, while
$S_i\subset\{|x|\ge ar_i\}$; so $S_k\cap S_i=\emptyset$, and
$f^{(n)}\mathbf 1_{S_k}=f^{(n)}_k$.

(2) By Lemma~\ref{lem:adm}, $\omega^{(n)}(t)\circ\Orot=\omega^{(n)}(t)$; with
$\KBS(\Orot x)=\Orot\KBS(x)$ and the change of variables $y=\Orot y'$ this gives
$v^{(n)}(t,\Orot x)=\Orot v^{(n)}(t,x)$. So $t\mapsto\Orot^{-1}\flow_t(\Orot y)$
solves the equation of $\flow_t(y)$ with the same initial value; since
$v^{(n)}$ is $C^1$ and $\nabla v^{(n)}$ is bounded on $[0,t]\times\R^2$
(Theorem~\ref{thm:smooth}\textup{(2)}), Gronwall's inequality gives uniqueness,
and $\flow_t\circ\Orot=\Orot\circ\flow_t$. At $x=0$, $v^{(n)}(t,0)=\Orot
v^{(n)}(t,0)$, so $v^{(n)}(t,0)=0$ and $\flow_t(0)=0$. As $\flow_t$ is a
homeomorphism (Theorem~\ref{thm:smooth}\textup{(4)}), the sets
$S_k(t)=\flow_t(S_k)$ inherit compactness, disjointness and $\Orot$-invariance
from (1), $0\notin S_k(t)$, and
$\supp\omega^{(n)}(t)=\flow_t(\supp f^{(n)})\subset\bigcup_kS_k(t)$, because
$\omega^{(n)}(t)=f^{(n)}\circ\flow_t^{-1}$.

(3) Since $\omega^{(n)}(t)=f^{(n)}\circ\flow_t^{-1}$ and
$\mathbf 1_{S_k(t)}=\mathbf 1_{S_k}\circ\flow_t^{-1}$, (1) gives the first
identity, and (2), by which $D_k(t)$ is the disjoint union of
$S_k(t),\dots,S_n(t)$, gives the other two. Each $f^{(n)}_k$ lies in
$C_c^\infty(\R^2)$, is nonnegative and $3$-fold rotationally symmetric (Lemmas~I.\ref{I-lem:ring-facts}
and \ref{lem:seed}), and $\flow_t^{-1}$ is
a $C^\infty$ diffeomorphism commuting with $\Orot$; this gives the regularity
and symmetry. Since $\det D\flow_t=1$,
$\int\omega^{(n)}(t)\mathbf 1_{S_k(t)}=\int f^{(n)}_k=A_kr_k^2\int\Phi
=A_kr_k^2\|\Phi\|_{L^1}$, as $\Phi\ge0$; for $k=n$ we use that the nonnegative equimeasurable
functions $\Phi^{\kseed}$ and $\Phi$ have the same integral.

(4) The two terms of
$w_k(t,y)=A_k^{-1}(\omega^{(n)}(t)\mathbf 1_{S_k(t)})(r_ky)-\Phi(y)$ are real,
lie in $C_c^\infty(\R^2)$, are $3$-fold rotationally symmetric (by (3) and
Lemma~I.\ref{I-lem:ring-facts}) and are supported in $N$, because
$S_k(t)\subset N_k$ and $\supp\Phi\subset\{a\le|y|\le1\}\subset N$
(Lemma~\ref{lem:annuli}). By (3) and the change of variables $x=r_ky$, both are
nonnegative with integral, hence $L^1$ norm, $\|\Phi\|_{L^1}$; so $\int w_k(t)=0$ and $\|w_k(t)\|_{L^1}\le2\|\Phi\|_{L^1}$.
The last inequality is the Cauchy-Schwarz inequality on $N$.
\end{proof}

The fields of Definition~\ref{def:regions} depend continuously on time: if
$J\subset[0,\infty)$ is a compact interval with $D_{k+1}(t)\subset B_{ar_k/8}$
for $t\in J$, then $\tau\mapsto\hin_k(\tau)$ is continuous from $A_kJ$ into
$C^M(N')$ for every integer $M\ge0$. Indeed, $\hin_n=-G_ny^\perp/|y|^2$ is
constant, because $D_{n+1}(t)=\emptyset$. Let $k<n$ and $g=\sum_{i>k}f^{(n)}_i$, supported in the compact set
$D_{k+1}(0)$. By Equation~\eqref{eq:omega-ge}, $\omega_{\ge k+1}(t)=g\circ\flow_t^{-1}$,
so the change of variables $y=\flow_t(y')$, with $\det D\flow_t=1$, gives
$(\KBS\ast\omega_{\ge k+1}(t))(x)=\int\KBS(x-\flow_t(y'))g(y')\dd y'$. As
$(t,y')\mapsto\flow_t(y')$ is continuous (Theorem~\ref{thm:smooth}\textup{(4)}),
$\rho=\max\{|\flow_t(y')|:t\in J,\ y'\in D_{k+1}(0)\}$ exists, and $\rho<\frac a8r_k$,
because the maximum is attained at a point of one of the sets
$\flow_t(D_{k+1}(0))=D_{k+1}(t)\subset B_{ar_k/8}$.
Let $U$ be the set of the $x$ at distance less than $\frac a{16}r_k$ from $N'_k$;
as the inner radius of $N'_k$ is $\frac a4r_k$, every $x\in U$ has
$|x|>\frac3{16}ar_k$. So, for $x\in U$, $t\in J$ and $y'\in D_{k+1}(0)$,
$|x-\flow_t(y')|\ge\frac3{16}ar_k-\rho>\frac a{16}r_k$. Since
$|\partial^\beta\KBS(z)|\le C_\beta|z|^{-1-|\beta|}$, the derivatives
$\partial^\beta_x$, $|\beta|\le M$, of the integrand are bounded, for $x\in U$
and $t\in J$, by a constant times $|g(y')|$, which is integrable. So we may
differentiate under the integral sign, and each of these derivatives of the field is
continuous on $J\times U$ by dominated convergence, hence uniformly continuous on
the compact set $J\times N'_k$; since $G_ky^\perp/|y|^2$ does not depend on
$\tau$, passing to ring-$k$ units gives the claim. In the same way, if the rings $i<k$ lie outside $B_{6r_k}$
for $t\in J$, then $\tau\mapsto\hout_k(\tau)\in C^M(N')$ is continuous on $A_kJ$:
replace $g$ by $\sum_{i<k}f^{(n)}_i$ and $U$ by $B_{4r_k}$, which contains $N'_k$,
and use $|x-\flow_t(y')|\ge2r_k$ for
$x\in B_{4r_k}$ and $y'\in S_1\cup\dots\cup S_{k-1}$, which holds because
$\flow_t(y')$ then lies in $S_1(t)\cup\dots\cup S_{k-1}(t)$, outside $B_{6r_k}$. Consequently $\etin_k$, $\etout_k$ and
$\eta_k$ are continuous on every interval on which they are defined, so the
integrals that define $\maj_k$ and $\maj^{\mathrm{in}}_k$ exist.

If on $[0,T]$ both $D_{k+1}(t)\subset B_{ar_k/8}$ and the rings $i<k$ lie outside
$B_{6r_k}$, then, by Lemma~\ref{lem:regions-basic}, the solution in ring-$k$
units, split into its parts on $S_k(t)$, on $D_{k+1}(t)$ and on
$S_1(t)\cup\dots\cup S_{k-1}(t)$ (rescaled), satisfies the hypotheses of
Lemma~\ref{lem:reduction} on $[0,A_kT]$ with $G=G_k$, and $h_k$ is the field $h$
there. Indeed, by Lemma~\ref{lem:regions-basic}\textup{(2)} and \textup{(3)}, the
three parts have pairwise disjoint supports, are $3$-fold rotationally symmetric
and are transported by the flow of the solution in ring-$k$ units. The part on
$D_{k+1}(t)$ is $\nu^{\mathrm{in}}_k$, whose integral at $t=0$ is $2\pi G_k$ by
the definition of $G_k$, so $G=G_k$. The two conditions on $[0,T]$, read in ring-$k$
units, imply the support conditions of Lemma~\ref{lem:reduction}. While
$S_k(t)\subset N_k$, the function $w=\omega_{\mathrm r}-\Phi$ of
Lemma~\ref{lem:reduction}\textup{(3)} at ring-$k$ time $A_kt$ is $w_k(t)$. This is how Lemma~\ref{lem:reduction} and Section~\ref{sec:single} are
applied to ring $k$.

Part~I and Section~\ref{sec:single} treat the linearization $\Lin(G)$ for
$G\in[0,\Gone]$. The next lemma shows that every inner circulation $G_k$ lies in
this range when $\sbar$ is small, with a bound independent of $n$.

\begin{lemma}
\label{lem:inner-circ}
In Definition~\ref{def:regions}, $\int\omega_{\ge k+1}(t)$ is independent of
$t$, and $0\le G_k\le C_G\mu s_k^2\le C_G\mu\sbar^2$ for $1\le k\le n$, with
a constant $C_G$ depending only on $\Phi$; moreover $G_n=0$. In particular, if
$C_G\mu\sbar^2\le\Gone$, then $G_k\in[0,\Gone]$ for $1\le k\le n$.
\end{lemma}
\begin{proof}
Let $F\colon\R\to\R$ be continuous with $F(0)=0$. Since $\flow_t$ preserves
area (Theorem~\ref{thm:smooth}\textup{(4)}), Equation~\eqref{eq:omega-ge} gives
$\int F(\omega_{\ge k}(t))\dd x=\int F\big(\sum_{i=k}^nf^{(n)}_i\big)\dd x$. The
$f^{(n)}_i$ are supported in the pairwise disjoint sets $S_i$
(Lemma~\ref{lem:regions-basic}\textup{(1)}), so, as $F(0)=0$, the function
$F\big(\sum_{i=k}^nf^{(n)}_i\big)$ is the sum of the $F(f^{(n)}_i)$. Hence, for
every such $F$ and every $1\le k\le n+1$,
\begin{equation}
\label{eq:conserved}
\int F\big(\omega_{\ge k}(t)\big)\dd x=\sum_{i=k}^n\int F\big(f^{(n)}_i\big)\dd x .
\end{equation}
With $F(s)=s$, $\int\omega_{\ge k+1}(t)$ does not depend on $t$, and, as
$\int f^{(n)}_i=A_ir_i^2\|\Phi\|_{L^1}$ (Lemma~\ref{lem:regions-basic}\textup{(3)}
at $t=0$), $2\pi G_kA_kr_k^2=\|\Phi\|_{L^1}\sum_{i=k+1}^nA_ir_i^2$. So $G_k\ge0$
and $G_n=0$, as the sum is then empty. Since $A_{i+1}r_{i+1}^2=\mu s_i^2\,A_ir_i^2$ and
$\mu s_i^2\le2(a/16)^2<\frac12$, each term of the sum is less than half the
previous one, and comparison with a geometric series shows that the sum is at most
$2A_{k+1}r_{k+1}^2=2\mu s_k^2A_kr_k^2$; so
$G_k\le\|\Phi\|_{L^1}\mu s_k^2/\pi\le C_G\mu s_k^2$, and
$s_k\le\sbar$.
\end{proof}

The next lemma bounds the part $\omega_{\ge k}$ of the solution formed by ring
$k$ and the rings inside it, uniformly in $n$ and $t$. Its velocity bound controls the
velocity field of $\omega_{\ge k}$ itself; in the proof of
Proposition~\ref{prop:confine} it is one of the two terms of the speed bound for
the material inside a ring, and its constant $\Cstar$ enters the confinement
radius there.

\begin{lemma}
\label{lem:speed}
In Definition~\ref{def:regions}, there are absolute constants $\Mvort$
and $\Cstar$ such that for all $1\le k\le n$, $t\ge0$ and $q\in\{1,4\}$,
\[
\begin{gathered}
\|\omega_{\ge k}(t)\|_{L^q}\le\Mvort^{1/q}A_kr_k^{2/q},\qquad
\|\KBS\ast\omega_{\ge k}(t)\|_{L^\infty}\le\Cstar A_kr_k,\\
\|\omega_{\ge k}(t)\|_{L^1}=\|\Phi\|_{L^1}\sum_{i=k}^nA_ir_i^2\le2\|\Phi\|_{L^1}A_kr_k^2 .
\end{gathered}
\]
\end{lemma}
\begin{proof}
By Equation~\eqref{eq:conserved} with $F(s)=|s|^q$ and the equimeasurability of
$\Phi^{\kseed}$ and $\Phi$ (Lemma~\ref{lem:seed}),
$\|\omega_{\ge k}(t)\|_{L^q}^q=\|\Phi\|_{L^q}^q\sum_{i=k}^nA_i^qr_i^2$. For
$q\in\{1,4\}$ and $i\ge k$ we have
$A_{i+1}^qr_{i+1}^2=\mu^qs_i^2A_i^qr_i^2\le\mu^4\sbar^2A_i^qr_i^2\le\frac18A_i^qr_i^2$,
because $\mu\le2$ and $\sbar\le a/16<\frac1{16}$ give $\mu^4\sbar^2<\frac1{16}$;
so $\sum_{i=k}^nA_i^qr_i^2\le\frac87A_k^qr_k^2$. With $q=1$ this gives the
identity and the bound for $\|\omega_{\ge k}(t)\|_{L^1}$. As $0\le\Phi<1$ and
$\supp\Phi\subset\{a\le|y|\le1\}$ (Lemma~I.\ref{I-lem:ring-facts}), $\|\Phi\|_{L^q}^q\le\pi$, and
$\|\omega_{\ge k}(t)\|_{L^q}^q\le\frac{8\pi}7A_k^qr_k^2\le\Mvort A_k^qr_k^2$.
Finally, let $g=\omega_{\ge k}(t)\in (L^1\cap L^4)(\R^2)$; we may assume $g\ne0$. Let
$x\in\R^2$ and $R>0$; by Lemma~I.\ref{I-lem:BS-basic}(a), $2\pi|\KBS(x-y)|=|x-y|^{-1}$.
We split the integral of $|g(y)||x-y|^{-1}$ over $y$ into the near field and the
far field.
On the near field, H\"older's inequality with exponents $4,4/3$ and
integration in polar coordinates give
\[
\int_{|x-y|<R}|g(y)||x-y|^{-1}\dd y\lesssim\|g\|_{L^4}R^{1/2}.
\]
The far-field integral is at most $\|g\|_{L^1}R^{-1}$. Choosing
$R=(\|g\|_{L^1}/\|g\|_{L^4})^{2/3}$ balances these bounds and yields
\[
\|\KBS\ast g\|_{L^\infty}
\lesssim\|g\|_{L^1}^{1/3}\|g\|_{L^4}^{2/3}
\lesssim A_kr_k.
\]
The implicit constants are absolute; choose $\Cstar$ to dominate the last one.
\end{proof}

The phase clock turns the relay described in the Introduction into definite
times. It is defined recursively, from the innermost ring outward. Ring $n$
becomes active at time $0$. Ring $k$ becomes active at the time $\tsh_k$, and
$T_k$ is the time at which its unstable coefficient is expected to reach the
size $\delta_0$. Ring $k-1$ then becomes active after a pinning phase
of length $\Tpin$ in ring-$k$ time. At the end of this pinning phase the kick
time of ring $k-1$, defined below, equals $\Tpin$; there
Lemma~\ref{lem:coupling}\textup{(2)} evaluates the kick that ring $k$ and the
rings inside it give to ring $k-1$.

\begin{definition}
\label{def:clock}
Let $\delta_0,\Tpin>0$. In Definition~\ref{def:regions}, assume $G_k\le\Gone$
for $1\le k\le n$. Define recursively, downward in $k$:
\begin{itemize}
\item $\tsh_n=0$;
\item if $\tsh_k$ is defined and $S_k(\tsh_k)\subset N_k$:
  $c_k(0)=\ellst(G_k)\big(w_k(\tsh_k)\big)$ and $\ksh_k=|c_k(0)|$;
\item if moreover $\ksh_k\in(0,\delta_0)$:
  $\hat T_k=\sigma_k^{-1}\log(\delta_0/\ksh_k)$, $T_k=\tsh_k+\hat T_k/A_k$, and, if
  $k\ge2$, $\tsh_{k-1}=T_k+\Tpin/A_k$.
\end{itemize}
The clock is \textbf{defined up to level $k$} if $S_i(\tsh_i)\subset N_i$ and
$\ksh_i\in(0,\delta_0)$ for every $i\in[k,n]$, so that $T_i$ is defined for
these $i$. Put $T_{n+1}=0$, $\tash_k=A_k\tsh_k$, $W_k=[T_{k+1},T_k]$, the
\textbf{active phase} $\Phase_k=[\tsh_k,T_k]$ and the \textbf{pinning phase}
$\Pin_k=[T_{k+1},\tsh_k]$ ($k<n$). Ring $k$ is \textbf{quiet} on $[0,\tsh_k]$,
\textbf{active} on $\Phase_k$, and \textbf{debris} after $T_k$. On $\Phase_k$ we
use the \textbf{active time} $u=A_k(t-\tsh_k)$; if $h_k$ is defined on
$[\tsh_k,t]$, the running coefficient of ring $k$ is, for $0\le u\le A_k(t-\tsh_k)$,
\begin{equation}
\label{eq:running-k}
c_k(u)=c_k(0)+\int_0^ue^{-\lambda_kv}\,\ellst(G_k)\big(F_k(v)\big)\dd v ,
\end{equation}
where $F_k(v)=-h_k(\tash_k+v)\cdot\nabla\Phi$ on $N'$ and $F_k(v)=0$ elsewhere.
The \textbf{kick time} of ring $k<n$ is $\tauK=A_{k+1}(t-T_{k+1})$; it equals
$\Tpin+\mu u$ on $\Phase_k$, and $\tau_0=-A_{k+1}T_{k+1}$ is the kick time of
$t=0$.
\end{definition}

The window $W_k$ is the union of the pinning phase $\Pin_k$ and the active phase
$\Phase_k$ of ring $k$ (for $k=n$, of $\Phase_n$ alone).

The integral in Equation~\eqref{eq:running-k} exists, so $c_k(u)$ is well
defined: its integrand is continuous in $v$, since the map
$v\mapsto h_k(\tash_k+v)\in C^0(N')$ is continuous (as shown after
Lemma~\ref{lem:regions-basic}) and
$|\ellst(G_k)(-h\cdot\nabla\Phi)|\le\Cls\|h\|_{C^0(N')}$ for continuous $h$
(Lemma~I.\ref{I-lem:uniform}).

The relay lemmas of Section~\ref{sec:relay} are causal, and the bootstrap of
Section~\ref{sec:bootstrap}, a continuity argument in time, applies them at a
time $t$ under hypotheses on $[0,t]$. So the phase boundaries that enter them at
time $t$ must be determined by $\omega^{(n)}$ on $[0,t]$, in the following sense, and
Lemma~\ref{lem:clock} below shows that they are. A quantity built from $\omega^{(n)}$ by
Definitions~\ref{def:regions} and \ref{def:clock} is
\textbf{determined by $\omega^{(n)}$ on $[0,\bar t]$} if it keeps its value when
$\omega^{(n)}$ is replaced in these definitions by any
$\tilde\omega\in C^\infty([0,\infty)\times\R^2)$ such that $\tilde\omega=\omega^{(n)}$
on $[0,\bar t]\times\R^2$, every $\tilde\omega(t)$ has compact support, and
$\KBS\ast\tilde\omega$ is smooth with every derivative bounded on $[0,T]\times\R^2$
for every $T$; the flow $\flow$ is then replaced by the flow of
$\KBS\ast\tilde\omega$, and $f^{(n)}$, the sets $S_k$ and the parameters are kept.

\begin{lemma}
\label{lem:clock}
In Definition~\ref{def:clock}, assume that the clock is defined up to level $k$.
\begin{enumerate}[label=\textup{(\arabic*)}]
\item $0=\tsh_n<T_n<\tsh_{n-1}<T_{n-1}<\dots<\tsh_k<T_k$, and $T_k<\tsh_{k-1}$
  if $k\ge2$.
\item For every $i\in[k,n]$, each of $\tsh_i$, $\ksh_i$, $T_i$ and, if $i\ge2$,
  $\tsh_{i-1}$ is determined by $\omega^{(n)}$ on $[0,\tsh_i]$.
\item Hence, for every $t\ge0$, the set $\mathcal J_t$ of the $i\in[k,n]$ with
  $\tsh_i\le t$, and the numbers $\tsh_i$ and $T_i$ for $i\in\mathcal J_t$, are
  determined by $\omega^{(n)}$ on $[0,t]$. In particular, so are the phase
  boundaries $\tsh_i$, $T_i$ ($k\le i\le n$) not exceeding $t$, and the phase of
  each ring $i\in[k,n]$ at time $t$ in the sense of Definition~\ref{def:clock}.
\end{enumerate}
\end{lemma}
\begin{proof}
(1) Let $i\in[k,n]$. By Lemma~\ref{lem:inner-circ}, $G_i\ge0$, and $G_i\le\Gone$ by
the assumption of Definition~\ref{def:clock}; so Proposition~I.\ref{I-prop:SU} gives
$\sigma_i=\sigma(G_i)>\frac{211}{2500}$. Moreover $\ksh_i<\delta_0$, as the clock
is defined up to level $k$, so $\hat T_i>0$
and $T_i>\tsh_i$; if $i\ge2$, then $\tsh_{i-1}-T_i=\Tpin/A_i>0$. These inequalities, for $i=n,n-1,\dots,k$,
give (1).

(2) Let $\bar t\ge\tsh_i$, let $\tilde\omega$ be as above, and mark with a tilde
the objects built from $\tilde\omega$. The velocities $\KBS\ast\tilde\omega$ and
$v^{(n)}$ agree on $[0,\bar t]\times\R^2$ and are Lipschitz in $x$, so, by
uniqueness for the equation of the flow, the two
flows agree on $[0,\bar t]$; and $\tilde G_j=G_j$, because
$\tilde\omega(0)=f^{(n)}$, so the numbers $\sigma_j$ do not change either. We show by downward induction on $j\in[i,n]$ that
$\tilde t^\sharp_j=\tsh_j$, $\tilde\kappa^\sharp_j=\ksh_j$, $\tilde T_j=T_j$, and
$\tilde t^\sharp_{j-1}=\tsh_{j-1}$ if $j\ge2$. First, $\tilde t^\sharp_n=0=\tsh_n$.
If $j\in[i,n]$ and $\tilde t^\sharp_j=\tsh_j$, then $\tsh_j\le\tsh_i\le\bar t$ by (1), so
$\tilde S_j(\tsh_j)=S_j(\tsh_j)\subset N_j$ and $\tilde w_j(\tsh_j)=w_j(\tsh_j)$;
hence $\tilde\kappa^\sharp_j=\ksh_j\in(0,\delta_0)$, and
Definition~\ref{def:clock} gives $\tilde T_j=T_j$ and
$\tilde t^\sharp_{j-1}=\tsh_{j-1}$. This completes the induction and proves (2).

(3) Let $\tilde\omega$ agree with $\omega^{(n)}$ on $[0,t]\times\R^2$, and let
$\tilde{\mathcal J}_t$ be the set of the $i\in[k,n]$ for which
$\tilde t^\sharp_i$ is defined and at most $t$. Let $i_0=\min\mathcal J_t$, which
exists because $\tsh_n=0$. By (1), $\tsh_i$ decreases as $i$ increases, so
$\mathcal J_t=[i_0,n]$, and $\tsh_i,T_i>t$
for $i\notin\mathcal J_t$. By (2) with $\bar t=t$ and $i=i_0$, the numbers
$\tsh_i$, $T_i$, $i\ge i_0$, are the same for $\tilde\omega$, so
$[i_0,n]\subset\tilde{\mathcal J}_t$, and, if $i_0>k$,
$\tilde t^\sharp_{i_0-1}=\tsh_{i_0-1}>t$; the proof of (1) applies to the clock
of $\tilde\omega$ wherever it is defined, so every defined $\tilde t^\sharp_i$ with $i<i_0$ exceeds
$t$. Hence $\tilde{\mathcal J}_t=\mathcal J_t$. These data decide the phase of
ring $i\in[k,n]$ at $t$. If $i\notin\mathcal J_t$, then $t<\tsh_i<T_i$ and ring
$i$ is quiet and neither active nor debris, and for $\tilde\omega$ the number
$\tilde t^\sharp_i$ is either undefined or greater than $t$, so ring $i$ is
neither active nor debris and is quiet at $t$ in the reading of
Definition~\ref{def:loc} in Section~\ref{sec:several}, which takes $t\le\tsh_i$ to
be true when $\tsh_i$ is not defined. If $i\in\mathcal J_t$, ring $i$ is
quiet if and only if $t=\tsh_i$, active if and only if $t\le T_i$, and debris if
and only if $t>T_i$.
\end{proof}

\section{Several Rings: Transfer of Moments and Confinement}

We assume \textup{(N1)}--\textup{(N3)}, fix $n\ge2$, and use
Definitions~\ref{def:datum}, \ref{def:approx}, \ref{def:regions} and
\ref{def:clock} with $\mu\in(1,2]$ and $\sbar\in(0,a/16]$; we write
$\omega=\omega^{(n)}$, $v=\KBS\ast\omega$, and $\flow_t$ for the flow of $v$.
Since $s_k=\sbar(2\mu)^{-k}\le\sbar/2$ and $r_k/r_i=s_is_{i+1}\cdots s_{k-1}$,
\begin{equation}
\label{eq:sk-small}
s_k\le\frac a{32}\quad(k\ge1),\qquad
\frac{r_k}{r_i}\le s_{k-1}\,\sbar^{\,k-i-1}\quad(1\le i<k).
\end{equation}
Moreover $\sbar^{m-2}=\sbar<\frac12$, because $m=3$, and
$\mu^4\sbar^2\le16\cdot2^{-8}<\frac18$, because $a<1$. The conditions
$\sbar\le a/16$, $\sbar^{m-2}\le\frac12$ and $\mu^4\sbar^2\le\frac18$ appear
among the hypotheses of some statements below so that these can be quoted on
their own; under the standing assumptions of this section they always hold.

In this section $q$, without an index, is the order of a derivative, and $i$ is a ring index or
the imaginary unit. For a velocity field $V$, $V_{[j]}$ is $V$ in ring-$j$ units
(defined before Definition~\ref{def:regions}), and $V$ is \emph{bounded on
$N'_j$ in $C^k$, in ring-$j$ units, by} a number if and only if
$\|V_{[j]}\|_{C^k(N')}$ is at most that number. We use the \textbf{scaling
rule}: if $c\in\R$, $c'>0$ and $W(y)=cV(c'y)$, then
$\nabla^qW(y)=c\,c'^{\,q}(\nabla^qV)(c'y)$.

We write $F=V_1-iV_2$ for the complex form of a vector field $V$, and we use
Equation~\eqref{eq:BS-complex} for the complex form of $\KBS\ast g$
(Section~\ref{sec:II-prelim}). The monopole
$\frac1{2\pi}\big(\int g\big)x^\perp/|x|^2$ of $g$ has complex form $\big(\int g\big)/(2\pi iz)$,
and $\Orot$ is multiplication by $e^{2\pi i/m}$. For real $g$, the complex forms
of $\KBS\ast g$ and of $\KBS\ast g$ minus its monopole are holomorphic on
$U=\C\setminus(\supp g\cup\{0\})$, so Equation~\eqref{eq:BS-complex-deriv}
applies to them there.

\subsection{Fields Induced by One Ring on Another}

\begin{lemma}
\label{lem:transfer}
Assume $\sbar\le a/16$, and let $M_h\ge1$ and $s'\ge0$ be integers. There are
constants $C^{\mathrm T}_M$ and $C^{\mathrm T}_O$, depending only on $a$, $m$ and
$M_h$, such that with $C_{L1}=|N|^{1/2}$, where $|N|$ is the area of $N$,
\[
C_M=C^{\mathrm T}_M\max(1,C_{L1}),\qquad C_O=C^{\mathrm T}_O\max(1,C_{L1}),
\]
the following hold at every time $t\ge0$, for $1\le j<n$.
\begin{enumerate}[label=\textup{(\arabic*)}]
\item Let $g\in L^1(\R^2)$ be real valued, $3$-fold rotationally symmetric and supported in
  $\overline B_{ar_{j+1}/8}$, and let $E_g$ be its velocity field minus its
  monopole. Then
  $\|(E_g)_{[j]}\|_{C^{M_h}(N')}\le C_M\varepsilon_j\sup_{|y|=a/4}|(E_g)_{[j+1]}(y)|$.
\item If $S_{j+1}(t)\subset N_{j+1}$, the velocity field induced on $N'_j$ by
  $A_{j+1}w_{j+1}(t,\cdot/r_{j+1})$ is bounded in $C^{M_h}(N'_j)$, in
  ring-$j$ units, by $C_M\varepsilon_j\|w_{j+1}(t)\|_{H^{s'}}$.
\item If $D_{j+1}(t)\subset\overline B_R$ with $R\le ar_j/8$, the non-monopole
  field of $\omega_{\ge j+1}(t)$ on $N'_j$ is bounded in $C^{M_h}(N'_j)$, in ring-$j$
  units, by $C_M\varepsilon_j(R/r_{j+1})^m$.
\item For $1\le i<k\le n$ with $S_i(t)\subset N_i$, the velocity field induced on
  $N'_k$ by $A_iw_i(t,\cdot/r_i)$ is bounded in $C^{M_h}(N'_k)$, in ring-$k$
  units, by $C_Ob_{i,k}\|w_i(t)\|_{L^1}$ and by $C_Ob_{i,k}\|w_i(t)\|_{H^{s'}}$,
  where $b_{i,k}=\mu^{i-k}(r_k/r_i)^{m-2}$. If $S_i(t)\subset N_i$ for every
  $i<k$, the velocity field of $\omega(t)\mathbf 1_{S_1(t)\cup\dots\cup S_{k-1}(t)}$
  on $N'_k$ is the sum of these fields over $i<k$, and its norm in
  $C^{M_h}(N'_k)$, in ring-$k$ units, is at most
  $C_O\sum_{i<k}b_{i,k}\|w_i(t)\|_{L^1}$.
\end{enumerate}
\end{lemma}

\begin{proof}
With $C_{m,q}$ the constant of Lemma~\ref{lem:ext-multipole}, put
\[
\begin{gathered}
\mathsf S_{m,q}=\sum_{p\ge1}(pm+q)^q2^{-(p-1)m}<\infty,\qquad
C^{(1)}_{\mathrm T}=\sum_{q=0}^{M_h}2^{q/2}\Big(\frac4a\Big)^q\mathsf S_{m,q},\\
C^{(2)}_{\mathrm T}=\frac2{\pi a}\Big(\frac8a\Big)^mC^{(1)}_{\mathrm T},\qquad
C^{(3)}_{\mathrm T}=2\pi\sum_{q=0}^{M_h}C_{m,q}\Big(\frac4a\Big)^{m+1+q},\\
C^{(4)}_{\mathrm T}=\frac1{2\pi}\sum_{q=0}^{M_h}2^{q/2}\Big(\frac2a\Big)^{q+1}\Big(\frac6a\Big)^{m-1-q}\mathsf S_{m,q}.
\end{gathered}
\]
We prove (1)--(3) with $C^{(1)}_{\mathrm T},C^{(2)}_{\mathrm T},C^{(3)}_{\mathrm T}$,
respectively, in place of $C_M$ and (4) with $C^{(4)}_{\mathrm T}$ in place of
$C_O$, for the $L^1$ norms. Then
$C^{\mathrm T}_M=\max(C^{(1)}_{\mathrm T},C^{(2)}_{\mathrm T},C^{(3)}_{\mathrm T})$
and $C^{\mathrm T}_O=C^{(4)}_{\mathrm T}$
depend only on $a$, $m$ and $M_h$, and the $L^1$ bounds follow from
$C_M\ge C^{\mathrm T}_M$ and $C_O\ge C^{\mathrm T}_O$. The $H^{s'}$ bounds follow,
because $\|w\|_{L^1}\le C_{L1}\|w\|_{L^2}\le C_{L1}\|w\|_{H^{s'}}$ for
$w=w_{j+1}(t)$ and $w=w_i(t)$ (Lemma~\ref{lem:regions-basic}(4)). Put
$\rho_0=a/4$.

\emph{Proof of (1).} Put $V=(E_g)_{[j+1]}$,
$\mathcal M=\sup_{|z|=\rho_0}|V(z)|$ and
$\tilde g(\zeta)=g(r_{j+1}\zeta)/A_{j+1}$. Since $\KBS$ is homogeneous of degree
$-1$, $V$ is the field of $\tilde g$ minus its monopole, and $\tilde g$ is
$3$-fold rotationally symmetric and supported in $\overline B_{a/8}$. Inserting
$(z-\zeta)^{-1}=\sum_{l\ge0}\zeta^lz^{-l-1}$, uniformly convergent for
$|\zeta|\le a/8$ and $z$ in compact subsets of $\{|z|>a/8\}$, in
Equation~\eqref{eq:BS-complex}, the complex form of $V$ is
\[
F(z)=\sum_{l\ge1}\alpha_lz^{-l-1},\qquad \alpha_l=\frac1{2\pi i}\int\tilde g(\zeta)\zeta^l\dd A(\zeta)\qquad(|z|>a/8),
\]
the term $l=0$ being the monopole. The substitution $\zeta=\Orot\zeta'$ gives
$\alpha_l=e^{2\pi il/m}\alpha_l$, so $\alpha_l=0$ unless $m\mid l$. On
$|z|=\rho_0$ the series converges uniformly, and $\alpha_l\rho_0^{-l-1}$ is the
Fourier coefficient of $\theta\mapsto F(\rho_0e^{i\theta})$ at the frequency
$-(l+1)$; since $|F|=|V|$, $|\alpha_l|\le\mathcal M\rho_0^{l+1}$. For
$|z|\ge\rho_0/s_j$, since $(l+1)\cdots(l+q)\le(l+q)^q$ and $s_j\le\frac12$
(Equation~\eqref{eq:sk-small}),
\begin{align*}
|F^{(q)}(z)|&\le\mathcal M\sum_{l\in m\N}(l+q)^q\rho_0^{l+1}\Big(\frac{s_j}{\rho_0}\Big)^{l+1+q}
=\mathcal M\rho_0^{-q}s_j^{m+1+q}\sum_{p\ge1}(pm+q)^qs_j^{(p-1)m}\\
&\le\mathcal M\Big(\frac4a\Big)^q\mathsf S_{m,q}\,s_j^{m+1+q}.
\end{align*}
In ring-$j$ units
$(E_g)_{[j]}(y)=\mu s_jV(y/s_j)$, since $A_{j+1}r_{j+1}=\mu s_jA_jr_j$, and
$|y/s_j|\ge\rho_0/s_j$ for $y\in N'$. By the scaling rule and
Equation~\eqref{eq:BS-complex-deriv}, on $N'$,
$|\nabla^q(E_g)_{[j]}(y)|=\mu s_j^{1-q}2^{q/2}|F^{(q)}(y/s_j)|$, which is at most
$\mu s_j^{m+2}\mathcal M\,2^{q/2}(4/a)^q\mathsf S_{m,q}$; summing over
$q\le M_h$ gives $\|(E_g)_{[j]}\|_{C^{M_h}(N')}\le C^{(1)}_{\mathrm T}\varepsilon_j\mathcal M$,
which is (1).

\emph{Proof of (2).} Let $w=w_{j+1}(t)$; by Lemma~\ref{lem:regions-basic}(4), $w$
is $3$-fold rotationally symmetric, supported in $N\subset\overline B_2$, and $\int w=0$. The
field $V$ of $A_{j+1}w(\cdot/r_{j+1})$ satisfies $V_{[j+1]}=\KBS\ast w$, by
homogeneity, so as in (1) its complex form in ring-$(j+1)$ units is
$F(z)=\sum_{l\in m\N}\alpha_lz^{-l-1}$ for $|z|>2$, with
$|\alpha_l|\le\frac1{2\pi}2^l\|w\|_{L^1}$ and no term $l=0$, since $\int w=0$.
For $|z|\ge\rho_0/s_j$ we have $2/|z|\le8s_j/a\le\frac14$ by
Equation~\eqref{eq:sk-small}, and
\[
|F^{(q)}(z)|\le\frac{\|w\|_{L^1}}{2\pi}\sum_{l\in m\N}(l+q)^q2^l\Big(\frac{4s_j}a\Big)^{l+1+q}
\le\frac{\|w\|_{L^1}}{2\pi}\Big(\frac4a\Big)^{1+q}\Big(\frac8a\Big)^m\mathsf S_{m,q}\,s_j^{m+1+q}.
\]
Passing to ring-$j$ units as in (1),
$|\nabla^qV_{[j]}(y)|\le\varepsilon_j\|w\|_{L^1}\frac1{2\pi}2^{q/2}(4/a)^{1+q}(8/a)^m\mathsf S_{m,q}$
on $N'$, and the sum over $q\le M_h$ is $C^{(2)}_{\mathrm T}\varepsilon_j\|w\|_{L^1}$.

\emph{Proof of (3).} By Lemma~\ref{lem:regions-basic}(3), $g=\omega_{\ge j+1}(t)$
is $3$-fold rotationally symmetric and supported in $D_{j+1}(t)\subset\overline B_R$, and by
Lemma~\ref{lem:speed}, $\|g\|_{L^1}\le2\|\Phi\|_{L^1}A_{j+1}r_{j+1}^2\le2\pi
A_{j+1}r_{j+1}^2$, since $0\le\Phi<1$ and $\supp\Phi\subset\overline B_1$. Every
$x\in N'_j$ has $|x|\ge ar_j/4\ge2R$. For the field $E$ of $g$ minus its
monopole, Lemma~\ref{lem:ext-multipole} with $\varrho=R$ and the scaling rule
give, for $y\in N'$,
\begin{align*}
|\nabla^qE_{[j]}(y)|&=\frac{r_j^{q}}{A_jr_j}|\nabla^qE(r_jy)|
\le C_{m,q}\frac{2\pi A_{j+1}r_{j+1}^2R^m}{A_jr_j^{m+2}}|y|^{-m-1-q}\\
&\le2\pi C_{m,q}\Big(\frac4a\Big)^{m+1+q}\mu s_j^2\Big(\frac R{r_j}\Big)^m .
\end{align*}
Since $\mu s_j^2(R/r_j)^m=\mu s_j^{m+2}(R/r_{j+1})^m=\varepsilon_j(R/r_{j+1})^m$,
the sum over $q\le M_h$ is at most $C^{(3)}_{\mathrm T}\varepsilon_j(R/r_{j+1})^m$.

\emph{Proof of (4).} Let $w=w_i(t)$ and $s_{i,k}=r_k/r_i\le a/32$
(Equation~\eqref{eq:sk-small}). The field $V$ of $A_iw(\cdot/r_i)$ satisfies
$V_{[i]}=\KBS\ast w$, and $w$ is $3$-fold rotationally symmetric and supported in
$N=\{a/2\le|\zeta|\le2\}$. Inserting
$(z-\zeta)^{-1}=-\sum_{l\ge0}z^l\zeta^{-l-1}$, uniformly convergent for
$\zeta\in N$ and $z$ in compact subsets of $B_{a/2}$, in
Equation~\eqref{eq:BS-complex} gives, for $|z|<a/2$,
\[
F(z)=-\sum_{l\ge0}\alpha_lz^l,\qquad \alpha_l=\frac1{2\pi i}\int w(\zeta)\zeta^{-l-1}\dd A(\zeta),\qquad
|\alpha_l|\le\frac{\|w\|_{L^1}}{2\pi}\Big(\frac2a\Big)^{l+1},
\]
and the substitution $\zeta=\Orot\zeta'$ shows $\alpha_l=0$ unless $m\mid l+1$.
Let $|z|\le3s_{i,k}$, so that $6s_{i,k}/a\le\frac14$. Writing $l=pm-1$ with
$p\ge1$, and using $l(l-1)\cdots(l-q+1)\le(pm)^q$, we get
\[
|F^{(q)}(z)|\le\frac{\|w\|_{L^1}}{2\pi}\sum_{p\ge1}(pm)^q\Big(\frac2a\Big)^{pm}(3s_{i,k})^{pm-1-q}
\le\frac{\|w\|_{L^1}}{2\pi}\Big(\frac2a\Big)^{q+1}\Big(\frac{6s_{i,k}}a\Big)^{m-1-q}\mathsf S_{m,q};
\]
the terms with $pm-1<q$ vanish on the left and are nonnegative on the right,
and
\[
(2/a)^{pm}(3s_{i,k})^{pm-1-q}=(2/a)^{q+1}(6s_{i,k}/a)^{m-1-q}(6s_{i,k}/a)^{(p-1)m}.
\]
In ring-$k$ units,
$V_{[k]}(y)=\frac{A_ir_i}{A_kr_k}V_{[i]}(s_{i,k}y)=\mu^{i-k}s_{i,k}^{-1}V_{[i]}(s_{i,k}y)$,
and $|s_{i,k}y|\le3s_{i,k}$ for $y\in N'$. Hence, on $N'$, by the scaling rule,
Equation~\eqref{eq:BS-complex-deriv} and $\mu^{i-k}s_{i,k}^{m-2}=b_{i,k}$,
\[
|\nabla^qV_{[k]}(y)|=\mu^{i-k}s_{i,k}^{q-1}2^{q/2}|F^{(q)}(s_{i,k}y)|
\le b_{i,k}\|w\|_{L^1}\frac{2^{q/2}}{2\pi}\Big(\frac2a\Big)^{q+1}\Big(\frac6a\Big)^{m-1-q}\mathsf S_{m,q},
\]
and the sum over $q\le M_h$ is $C^{(4)}_{\mathrm T}b_{i,k}\|w\|_{L^1}$. For the last
assertion, Lemma~\ref{lem:regions-basic}(2) and the definition of $w_i$ give
$\omega(t)\mathbf 1_{S_1(t)\cup\dots\cup S_{k-1}(t)}=\sum_{i<k}\big(A_i\Phi(\cdot/r_i)+A_iw_i(t,\cdot/r_i)\big)$.
Each $A_i\Phi(\cdot/r_i)$ is radial and supported in $\{|x|\ge ar_i\}$, so its
field vanishes on $B_{ar_i}$ by Lemma~\ref{lem:hollow}, and
$N'_k\subset\overline B_{3r_k}\subset B_{ar_i}$ because $3r_k=3s_{i,k}r_i<ar_i$.
Summing the bounds over $i<k$ completes the proof.
\end{proof}

\begin{lemma}
\label{lem:inner-field}
Let $M_h\ge1$ be an integer. There is $\Cin$, depending only on $a$, $m$ and
$M_h$, such that the following holds. Let $1\le k<n$ and $t\ge0$ be such that
$\omega_{\ge k+1}(t)$ is supported in $\overline B_R$ with $R\le\frac a8r_k$ and
$\|\omega_{\ge k+1}(t)\|_{L^1}\le\Mvort A_{k+1}r_{k+1}^2$. Let $E$ be the
velocity field of $\omega_{\ge k+1}(t)$ minus its monopole. Then
$\|E_{[k]}\|_{C^{M_h}(N')}\le\Cin\,\mu s_k^2\,(R/r_k)^m$. If moreover
$D_{k+1}(t)\subset B_{ar_k/8}$, then $E_{[k]}=\hin_k(A_kt)$ on $N'$.
\end{lemma}

\begin{proof}
Put $\Cin=\Mvort\sum_{q=0}^{M_h}C_{m,q}(4/a)^{m+1+q}$, with $C_{m,q}$ of
Lemma~\ref{lem:ext-multipole}. The function $g=\omega_{\ge k+1}(t)$ is $3$-fold
rotationally symmetric by Lemma~\ref{lem:regions-basic}(3), and every $x\in N'_k$ has
$|x|\ge ar_k/4\ge2R$. The computation in the proof of
Lemma~\ref{lem:transfer}(3), with $\|g\|_{L^1}\le\Mvort A_{k+1}r_{k+1}^2$ in
place of $2\pi A_{k+1}r_{k+1}^2$, gives
\[
|\nabla^qE_{[k]}(y)|\le\Mvort C_{m,q}(4/a)^{m+1+q}\mu s_k^2(R/r_k)^m\qquad(y\in N'),
\]
and the sum over $q\le M_h$ is the claimed bound. Now let
$D_{k+1}(t)\subset B_{ar_k/8}$, so that $\hin_k(A_kt)$ is defined. By homogeneity
of $\KBS$, $(\KBS\ast\omega_{\ge k+1}(t))_{[k]}=\KBS\ast\nu^{\mathrm{in}}_k$, and the
monopole of $\omega_{\ge k+1}(t)$ is, in ring-$k$ units,
$\frac1{2\pi}\big(\int\nu^{\mathrm{in}}_k\big)y^\perp/|y|^2$. By
Lemma~\ref{lem:inner-circ}, $\int\omega_{\ge k+1}(t)=\int\omega_{\ge k+1}(0)$, so
$\int\nu^{\mathrm{in}}_k=A_k^{-1}r_k^{-2}\int\omega_{\ge k+1}(0)=2\pi G_k$. Hence
$E_{[k]}=\KBS\ast\nu^{\mathrm{in}}_k-G_ky^\perp/|y|^2=\hin_k(A_kt)$.
\end{proof}

\subsection{Confinement on a Window}

\begin{proposition}
\label{prop:confine}
Let $1\le k<n$, assume that the clock is defined up to level $k+1$, put
$t_0=T_{k+1}$, and let $t_1\ge t_0$. Put $\tauK=A_{k+1}(t-t_0)$ and
$\Lambda=2+2\Cstar A_{k+1}(t_1-t_0)$. Assume:
\begin{enumerate}[label=\textup{(\roman*)}]
\item $D_{k+1}(t_0)\subset\overline{B}_{2r_{k+1}}$;
\item for every $t\in[t_0,t_1]$ and every $j\le k$,
  $S_j(t)\subset\{|x|\ge\frac a2r_j\}$;
\item $\sbar\le a/16$, $\sbar^{m-2}\le\frac12$ and $\mu^4\sbar^2\le\frac18$;
\item $\Lambda s_k\le a/8$ and
  $\frac83\big(\frac2a\big)^m\Lambda^{m-1}s_k^{m-2}\le\mu\Cstar$.
\end{enumerate}
Then for every $t\in[t_0,t_1]$:
\begin{enumerate}[label=\textup{(\alph*)}]
\item $D_{k+1}(t)\subset\overline{B}_{\varrho(t)}$ with
  $\varrho(t)=(2+2\Cstar\tauK)r_{k+1}\le\frac a8r_k$;
\item $|q_{k+1}(t)|\le\Cmom(1+\tauK)^m$, with $\Cmom=\Mvort(2\max(1,\Cstar))^m$.
\end{enumerate}
None of the constants depends on $n$, $k$, $t_0$, $\mu$, $\sbar$, $\delta_0$ or
$\kseed$.
\end{proposition}

\begin{proof}
The only constants, $\Cstar$, $\Mvort$ and $\Cmom$, are absolute numbers by
Lemma~\ref{lem:speed}, which gives the last sentence.

\emph{Speed bound.} Fix $t\in[t_0,t_1]$ and $x\in\overline B_{\Lambda r_{k+1}}$;
by (iv), $|x|\le\Lambda s_kr_k\le\frac a8r_k$. By
Lemma~\ref{lem:regions-basic}(2),(3), $\omega(t)=\omega_{\ge k+1}(t)+\sum_{j\le k}g_j$
with $g_j=\omega(t)\mathbf 1_{S_j(t)}$, so
$v(t,x)=\KBS\ast\omega_{\ge k+1}(t)(x)+\sum_{j\le k}\KBS\ast g_j(x)$, and
$|\KBS\ast\omega_{\ge k+1}(t)(x)|\le\Cstar A_{k+1}r_{k+1}$ by Lemma~\ref{lem:speed}.
Let $j\le k$. By Lemma~\ref{lem:regions-basic}(3), $g_j\in C_c^\infty$ is
$3$-fold rotationally symmetric and $\|g_j\|_{L^1}=A_jr_j^2\|\Phi\|_{L^1}\le\pi A_jr_j^2$, since
$0\le\Phi<1$ and $\supp\Phi\subset\overline B_1$. By (ii), $g_j$ is supported in
$\{|y|\ge R_j\}$ with $R_j=\frac a2r_j$, and $|x|\le\frac a8r_k\le R_j/2$.
Lemma~\ref{lem:hollow} with $m_0=m$, $|x|\le\Lambda r_{k+1}$ and
$A_j=\mu^{j-k-1}A_{k+1}$ give
\begin{align*}
|\KBS\ast g_j(x)|&\le\frac4{3\pi}\pi A_jr_j^2|x|^{m-1}R_j^{-m}
=\frac43\Big(\frac2a\Big)^mA_jr_j\Big(\frac{|x|}{r_j}\Big)^{m-1}\\
&\le\frac43\Big(\frac2a\Big)^m\Lambda^{m-1}A_{k+1}r_{k+1}\,\mu^{j-k-1}\Big(\frac{r_{k+1}}{r_j}\Big)^{m-2}.
\end{align*}
By Equation~\eqref{eq:sk-small}, $r_{k+1}/r_j\le s_k\sbar^{k-j}$, and by (iii),
$\sbar^{(m-2)(k-j)}\le2^{j-k}$; with $\mu^{j-k}\le1$,
$\mu^{j-k-1}(r_{k+1}/r_j)^{m-2}\le\mu^{-1}s_k^{m-2}2^{j-k}$, whose sum over
$j\le k$ is less than $2\mu^{-1}s_k^{m-2}$. By (iv),
$\sum_{j\le k}|\KBS\ast g_j(x)|\le\frac83(\frac2a)^m\Lambda^{m-1}s_k^{m-2}\mu^{-1}A_{k+1}r_{k+1}\le\Cstar A_{k+1}r_{k+1}$,
so
\[
|v(t,x)|\le2\Cstar A_{k+1}r_{k+1}\qquad(t\in[t_0,t_1],\ x\in\overline B_{\Lambda r_{k+1}}).\tag{$*$}
\]

\emph{Proof of (a).} By definition $D_{k+1}(t)=\flow_t(D_{k+1}(0))$, and
$D_{k+1}(0)=S_{k+1}\cup\dots\cup S_n$ is compact. Let $t^*$ be the largest
$t\in[t_0,t_1]$ such that $D_{k+1}(t')\subset\overline B_{\Lambda r_{k+1}}$ for
every $t'\in[t_0,t]$; it exists, because the set of such $t$ contains $t_0$, by
(i) and $\Lambda\ge2$, and is closed, because $(t,y)\mapsto\flow_t(y)$ is
continuous (Theorem~\ref{thm:smooth}\textup{(4)}). For $t\in[t_0,t^*]$ and
$y\in D_{k+1}(0)$, the point $\flow_{t'}(y)$ lies in $\overline B_{\Lambda r_{k+1}}$
for $t'\in[t_0,t]$, so \textup{($*$)} gives
$|\partial_{t'}\flow_{t'}(y)|=|v(t',\flow_{t'}(y))|\le2\Cstar A_{k+1}r_{k+1}$, and by (i),
$|\flow_t(y)|\le|\flow_{t_0}(y)|+2\Cstar A_{k+1}r_{k+1}(t-t_0)\le\varrho(t)$. If
$t^*<t_1$, then $\varrho(t^*)<\varrho(t_1)=\Lambda r_{k+1}$, since $\Cstar>0$, and
the uniform continuity of $(t,y)\mapsto|\flow_t(y)|$ on
$[t_0,t_1]\times D_{k+1}(0)$ gives $\delta\in(0,t_1-t^*]$ with
$|\flow_t(y)|<\Lambda r_{k+1}$ for $t\in[t^*,t^*+\delta]$ and $y\in D_{k+1}(0)$,
contradicting the maximality of $t^*$. Hence $t^*=t_1$, and (a) holds, because
$\varrho(t)\le\varrho(t_1)=\Lambda s_kr_k\le\frac a8r_k$ by (iv).

\emph{Proof of (b).} By (a) and Lemma~\ref{lem:speed},
$|\Qm[\omega_{\ge k+1}(t)]|\le\varrho(t)^m\|\omega_{\ge k+1}(t)\|_{L^1}\le\Mvort A_{k+1}r_{k+1}^2\allowbreak\varrho(t)^m$,
so by the definition of $q_{k+1}$,
$|q_{k+1}(t)|\le\Mvort(\varrho(t)/r_{k+1})^m=\Mvort(2+2\Cstar\tauK)^m$, and
$2+2\Cstar\tauK\le2\max(1,\Cstar)(1+\tauK)$.
\end{proof}

\begin{definition}
\label{def:loc}
Let $t\ge0$. In this definition we read $t\le T_i$ and $t\le\tsh_i$ as true
when $T_i$, respectively $\tsh_i$, is not defined by Definition~\ref{def:clock},
and ring $j$ is \textbf{quiet at $t$} if and only if $t\le\tsh_j$. The
\textbf{localization hypothesis} $(\mathrm{Loc})$ holds at $t$ if and only if:
\begin{enumerate}[label=\textup{(\arabic*)}]
\item for every ring $i$ with $t\le T_i$, $S_i(t)\subset N_i$;
\item for every ring $j$ that is quiet at $t$: $S_i(t)\cap B_{6r_j}=\emptyset$ for
  every $i<j$; $D_{j+1}(t)\subset B_{ar_j/8}$ (the open disk); and, if $T_{j+1}$ is defined
  and $t\ge T_{j+1}$, $D_{j+1}(t)\subset\overline{B}_{\varrho_j(t)}$ with
  $\varrho_j(t)=\big(2+2\Cstar A_{j+1}(t-T_{j+1})\big)r_{j+1}$, the radius of
  Proposition~\ref{prop:confine}\textup{(a)} for $k=j$.
\end{enumerate}
\end{definition}

\begin{lemma}
\label{lem:field-bound}
Let $M_h\ge1$ be an integer. There is $C_\eta$, depending only on $a$, $m$, $M_h$
and $\|\Phi\|_{L^1}$, such that the following holds. Assume
$\sbar\le a/16$ and $\sbar^{m-2}\le\frac12$, let $1\le j<n$ and $t\ge0$, and
assume that $S_i(t)\subset N_i$ for every $i<j$, that $S_i(t)\cap B_{6r_j}=\emptyset$
for every $i<j$, and that $D_{j+1}(t)\subset B_{ar_j/8}$. Then
$\eta_j(A_jt)\le C_\eta\big(\mu\sbar^2+\sbar^{m-2}\big)$.
\end{lemma}

\begin{proof}
Put $C_\eta=\max\{(a/8)^mC_M,\,4\|\Phi\|_{L^1}C_O\}$, with $C_M$ and $C_O$ of
Lemma~\ref{lem:transfer} at order $M_h$; it depends only on $a$, $m$, $M_h$ and
$\|\Phi\|_{L^1}$. The hypotheses are those under which
Definition~\ref{def:regions} defines $h_j$ at time $t$, and
$\eta_j(A_jt)=\etin_j(A_jt)+\etout_j(A_jt)$, where $\etin_j(A_jt)$ is the norm in
$C^{M_h}(N')$ of the non-monopole field of $\omega_{\ge j+1}(t)$ in ring-$j$ units
(Lemma~\ref{lem:inner-field}, whose hypotheses hold with $R=ar_j/8$ by
Lemma~\ref{lem:speed}), and $\etout_j(A_jt)$ that of the field of
$\omega(t)\mathbf 1_{S_1(t)\cup\dots\cup S_{j-1}(t)}$. By
Lemma~\ref{lem:transfer}(3) with $R=ar_j/8$,
$\etin_j(A_jt)\le C_M\varepsilon_j(ar_j/(8r_{j+1}))^m=(a/8)^mC_M\mu s_j^2\le(a/8)^mC_M\mu\sbar^2$.
If $j=1$, then $\etout_1=0$. If $j\ge2$, Lemma~\ref{lem:transfer}(4) with $k=j$ and
Lemma~\ref{lem:regions-basic}(4) give
$\etout_j(A_jt)\le C_O\sum_{i<j}b_{i,j}\|w_i(t)\|_{L^1}\le2\|\Phi\|_{L^1}C_O\sum_{i<j}b_{i,j}$.
By Equation~\eqref{eq:sk-small}, $r_j/r_i\le\sbar^{j-i}$, and $\mu^{i-j}\le1$, so
$b_{i,j}\le(\sbar^{m-2})^{j-i}$ and, as $\sbar^{m-2}\le\frac12$,
$\sum_{i<j}b_{i,j}\le\sum_{l\ge1}(\sbar^{m-2})^l\le2\sbar^{m-2}$. Hence
$\etout_j(A_jt)\le4\|\Phi\|_{L^1}C_O\sbar^{m-2}$, and adding the two bounds proves
the lemma.
\end{proof}

\section{The Kick}
\label{sec:kick}

We compute the kick that ring $k$ receives from the rings inside it,
for $1\le k<n$. While ring $k$ is quiet, the inner field $\hin_k$ of
Definition~\ref{def:regions} drives the unstable coefficient of ring $k$ exactly
through the mode-$3$ moment $q_{k+1}$ of the region $D_{k+1}(t)$: the unstable
coefficient of the part of the linear response of ring $k$ driven by $\hin_k$
solves the scalar linear Equation~\eqref{eq:kick}, with drive $q_{k+1}$
(Lemma~\ref{lem:coupling}). Proposition~\ref{prop:kick} bounds the modulus of
the solution at the final time, with its free growth factored out, from above
and below, and Section~\ref{sec:relay} uses these bounds to start ring $k$
(Lemmas~\ref{lem:handoff} and~\ref{lem:kicktrue}).

\begin{lemma}
\label{lem:coupling}
Assume \textup{(N1)}--\textup{(N3)}, fix $n\ge2$ and the objects of
Definitions~\ref{def:regions} and \ref{def:clock}, and let $1\le k<n$.
\begin{enumerate}[label=\textup{(\arabic*)}]
\item Let $t\ge0$ with $D_{k+1}(t)\subset B_{ar_k/8}$, and put $\tau=A_kt$. Then
  $-\hin_k(\tau)\cdot\nabla\Phi$, extended by $0$ outside $N'$, lies in $X^s$
  for every $s\ge0$, and
  \[
  \ellst(G_k)\big(-\hin_k(\tau)\cdot\nabla\Phi\big)=\mu\,\beta_k\,q_{k+1}(t).
  \]
  Moreover, with $\nu^{\mathrm{in}}_k$ of Definition~\ref{def:regions} at the
  time $t$, the mode-$3$ profile of $\hin_k(\tau)\cdot\er$ is
  $\frac i2\pout_3[\nu^{\mathrm{in}}_k]\,r^{-4}$ for $a/4\le r\le3$.
\item Assume that the clock is defined up to level $k+1$ and that
  $D_{k+1}(t)\subset B_{ar_k/8}$ for every $t\in[0,\tsh_k]$. Then
  $\tau\mapsto-\hin_k(\tau)\cdot\nabla\Phi$ is continuous from $[0,\tash_k]$ into
  $X^s$ for every $s$, and we may put, in ring-$k$ time,
  \[
  \wlin_{\mathrm{in},k}(\tau)=\int_0^\tau e^{(\tau-v)\Lin(G_k)}\big(-\hin_k(v)\cdot\nabla\Phi\big)\dd v
  \qquad(0\le\tau\le\tash_k).
  \]
  Regard $q_{k+1}$ and $\cin=\ellst(G_k)\big(\wlin_{\mathrm{in},k}\big)$ as functions
  of the kick time $\tauK=A_{k+1}(t-T_{k+1})\in[\tau_0,\Tpin]$. Then
  $\cin\in C^1([\tau_0,\Tpin])$ and, exactly,
  \begin{equation}
  \label{eq:kick}
  \frac{d\cin}{d\tauK}=\zeta_0\,\cin+\beta_k\,q_{k+1},\qquad \zeta_0=\frac{\lambda_k}{\mu},\qquad
  \cin(\tau_0)=0 .
  \end{equation}
  The \textbf{kick amplitude} $\Ain(\tauK)=e^{-\zeta_0\tauK}\cin(\tauK)$ satisfies
  \[
  \cin_k:=\ellst(G_k)\big(\wlin_{\mathrm{in},k}(\tash_k)\big)=e^{\zeta_0\Tpin}\Ain(\Tpin).
  \]
\item If, in addition to the hypotheses of \textup{(2)}, the rings $i<k$ lie
  outside $B_{6r_k}$ for every $t\in[0,\tsh_k]$, then
  $\tau\mapsto-\hout_k(\tau)\cdot\nabla\Phi$ is continuous from $[0,\tash_k]$ into
  $X^s$ for every $s$, and $\wlin_{\mathrm{out},k}$, defined as
  $\wlin_{\mathrm{in},k}$ with $\hout_k$ in place of $\hin_k$, satisfies $\wlin_{\mathrm{in},k}+\wlin_{\mathrm{out},k}=\wlin$ on
  $[0,\tash_k]$, where $\wlin$ is the linear response of
  Theorem~\ref{thm:quiet} for ring $k$, with $G=G_k$ and $h=h_k$.
\end{enumerate}
\end{lemma}
\begin{proof}
\emph{(1).} This part makes precise the computation made before
Definition~I.\ref{I-def:eigendata}, where $b$ is introduced as the mode-$3$
forcing of the ring by an interior field of unit outer moment, and $\beta(G)$ as
the component of that forcing along the unstable eigenfunction. By
Lemma~\ref{lem:regions-basic}(2),(3), the function
$\nu^{\mathrm{in}}_k$ of Definition~\ref{def:regions} at time $t$, the vorticity of
$D_{k+1}(t)$ in ring-$k$ units, is real,
$3$-fold rotationally symmetric, in $C_c^\infty(\R^2)$, and its support avoids $0$; by the
hypothesis $D_{k+1}(t)\subset B_{ar_k/8}$ the support is a compact subset of
$B_{a/8}$, so the largest modulus $\rho$ of a point of the support satisfies $\rho<a/8$. By
Lemma~I.\ref{I-lem:BS-basic}(d), $\KBS\ast\nu^{\mathrm{in}}_k=\nabla^\perp\psi$ with
$\psi=\frac1{2\pi}\log|\cdot|\ast\nu^{\mathrm{in}}_k$ smooth on $\{|y|>\rho\}$.
The monopole $G_ky^\perp/|y|^2$ is orthogonal to $\er$, so the polar form
$v\cdot\er=-r^{-1}\partial_\theta\psi$ of the Biot-Savart law
(Section~I.\ref{I-sec:prelim}) gives $\hin_k(\tau)\cdot\er=-r^{-1}\partial_\theta\psi$
on $N'$. Put $p_3=\pout_3[\nu^{\mathrm{in}}_k]$, the outer mode-$3$ moment of
Definition~I.\ref{I-def:multipole}. By the expansion of the stream function stated
before Lemma~\ref{lem:ext-multipole}, with $m=3$ and $\varrho=\rho$, the angular mode $3$ of $\psi$ is
$-\frac16p_3r^{-3}$ for $r>\rho$, and by integration by parts in $\theta$ that of
$\partial_\theta\psi$ is $3i$ times that of $\psi$. Since
$N'=\{a/4\le|y|\le3\}\subset\{|y|>\rho\}$, the mode-$3$ profile of
$\hin_k(\tau)\cdot\er$ is $\frac i2p_3r^{-4}$ for $a/4\le r\le3$, the last claim
of (1). As $\Phi$ is radial and $\Phi'$ vanishes outside $[a,1]$,
$\hin_k(\tau)\cdot\nabla\Phi=\Phi'(r)\,\hin_k(\tau)\cdot\er$, so the mode-$3$
profile of $-\hin_k(\tau)\cdot\nabla\Phi$ is
$-\frac i2p_3r^{-4}\Phi'(r)=p_3\,b(r)$, with $b$ of Definition~I.\ref{I-def:eigendata}.
By Equation~\eqref{eq:BS-complex}, $\hin_k(\tau)$ is the restriction to
$N'$ of a field smooth on $\{|y|>\rho\}$, and it is $3$-fold rotationally symmetric because
$\nu^{\mathrm{in}}_k$ is; so $-\hin_k(\tau)\cdot\nabla\Phi$ is real, smooth,
$3$-fold rotationally symmetric and supported in $\{a\le|y|\le1\}\subset N$, hence in $X^s$
for every $s$. By Definition~I.\ref{I-def:eigendata}, $\ellst(G_k)(w)=n_*(G_k)\,\ellst(w_3)$
depends linearly on the mode-$3$ profile $w_3$ of $w$ only, and
$\beta(G_k)=\ellst(G_k)(b\,e^{3i\theta})=n_*(G_k)\ellst(b)$; so
$\ellst(G_k)(-\hin_k(\tau)\cdot\nabla\Phi)=p_3\,\beta(G_k)$. It remains to
express $p_3$ through the normalized moment $q_{k+1}(t)$ of
Definition~\ref{def:regions}. The change
of variables $x=r_ky$ and the reality of $\omega_{\ge k+1}(t)$ give
\[
p_3=\frac1{2\pi A_k}\int\omega_{\ge k+1}(t,r_ky)\,\bar y^{\,3}\dd y
=\frac{\overline{\Qm[\omega_{\ge k+1}(t)]}}{2\pi A_kr_k^5}
=\frac{A_{k+1}r_{k+1}^5}{2\pi A_kr_k^5}\,q_{k+1}(t)=\frac{\mu s_k^5}{2\pi}\,q_{k+1}(t),
\]
and (1) follows from $m+2=5$ (recall $m=3$) and $\beta_k=\beta(G_k)s_k^5/(2\pi)$
(Definition~\ref{def:regions}).

\emph{(2).} By the hypothesis and the paragraph after
Lemma~\ref{lem:regions-basic}, $\tau\mapsto\hin_k(\tau)$ is continuous from
$[0,\tash_k]$ into $C^M(N')$ for every integer $M\ge0$. Put
$g(\tau)=-\hin_k(\tau)\cdot\nabla\Phi$. By (1), $g(\tau)\in X^s$ for every $s$,
and by Lemma~I.\ref{I-lem:uniform},
$\|g(\tau)-g(\tau')\|_{H^M}\le\Cnab(M)\|\hin_k(\tau)-\hin_k(\tau')\|_{C^M(N')}$;
so $g$ is continuous into $X^s$ for every $s$. Since $e^{t\Lin(G_k)}$ is a
$C_0$-semigroup on $X^s$ (Lemma~I.\ref{I-lem:semigroup-construct}), the integrand
$v\mapsto e^{(\tau-v)\Lin(G_k)}g(v)$ is continuous into $X^s$, so
$\wlin_{\mathrm{in},k}$ is well defined.

Next we change from ring-$k$ time to kick time. Recall from
Definition~\ref{def:clock} that $\tau_0=-A_{k+1}T_{k+1}$ is the kick time of $t=0$.
Since $A_{k+1}=\mu A_k$ by Equation~\eqref{eq:datum}, we have
$A_kT_{k+1}=-\tau_0/\mu$, so ring-$k$
time and kick time are related by $\tau=(\tauK-\tau_0)/\mu$. As
$\tsh_k=T_{k+1}+\Tpin/A_{k+1}$ (Definition~\ref{def:clock}), $\tauK$ runs over
$[\tau_0,\Tpin]$ as $t$ runs over $[0,\tsh_k]$. As $\ellst(G_k)$ is bounded on
$X^0$ (Lemma~I.\ref{I-lem:uniform}), it commutes with the integral that defines
$\wlin_{\mathrm{in},k}$. In the display below, the first equality is then the
last identity of Proposition~I.\ref{I-prop:semigroup}, by which $\ellst(G_k)$
turns $e^{t\Lin(G_k)}$ into multiplication by $e^{\lambda_kt}$; the second is
(1); and the third is the substitution
$v=(\tauK'-\tau_0)/\mu$. This gives, for $\tau_0\le\tauK\le\Tpin$ and
$\tau=(\tauK-\tau_0)/\mu$, with $q_{k+1}$ written as a function of kick time in
the last line,
\begin{align*}
\cin(\tauK)&=\int_0^\tau e^{\lambda_k(\tau-v)}\,\ellst(G_k)\big(g(v)\big)\dd v
=\mu\beta_k\int_0^\tau e^{\lambda_k(\tau-v)}\,q_{k+1}(v/A_k)\dd v\\
&=\beta_k\int_{\tau_0}^{\tauK}e^{\zeta_0(\tauK-\tauK')}\,q_{k+1}(\tauK')\dd\tauK'.
\end{align*}
By Equation~\eqref{eq:omega-ge} and area preservation,
$\Qm[\omega_{\ge k+1}(t)]=\int\sum_{i>k}f^{(n)}_i(x')\,\flow_t(x')^m\dd x'$. The
integrand is supported in the compact set $D_{k+1}(0)$, it is bounded uniformly
for $t$ in compact intervals, and it is continuous in $t$, as
$(t,x')\mapsto\flow_t(x')$ is continuous (Theorem~\ref{thm:smooth}\textup{(4)});
so, by dominated convergence, $q_{k+1}$ is continuous. Hence, taking the factor
$e^{\zeta_0\tauK}$ out of the last integral and using the fundamental theorem of
calculus, $\cin\in C^1([\tau_0,\Tpin])$, $\cin(\tau_0)=0$, and differentiation gives
Equation~\eqref{eq:kick}. At $t=\tsh_k$ we have $\tau=\tash_k$ and
$\tauK=A_{k+1}(\tsh_k-T_{k+1})=\Tpin$, so $\cin_k=\cin(\Tpin)=e^{\zeta_0\Tpin}\Ain(\Tpin)$.

\emph{(3).} By the hypotheses and the paragraph after
Lemma~\ref{lem:regions-basic}, $\tau\mapsto\hout_k(\tau)$ is continuous from
$[0,\tash_k]$ into $C^M(N')$ for every $M$. By Lemma~\ref{lem:regions-basic}(3),
$\nu^{\mathrm{out}}_k$ is real and $3$-fold rotationally symmetric, and, as the
rings $i<k$ lie outside $B_{6r_k}$, it vanishes on $B_6$; so
$\hout_k(\tau)$ is real, $3$-fold rotationally symmetric and smooth on $B_6$, by
Equation~\eqref{eq:BS-complex}, and $B_6$ contains $N'$. Hence, as in (1),
$-\hout_k(\tau)\cdot\nabla\Phi\in X^s$
for every $s$, and it is continuous in $\tau$ into $X^s$ by
Lemma~I.\ref{I-lem:uniform}, as in (2). By the discussion after
Lemma~\ref{lem:regions-basic}, ring $k$ in ring-$k$ units satisfies the
hypotheses of Lemma~\ref{lem:reduction} on $[0,\tash_k]$ with $G=G_k$ and
$h=h_k=\hin_k+\hout_k$, so the linear response $\wlin$ of
Theorem~\ref{thm:quiet}\textup{(b)} is defined with this $h$; the identity for
$\wlin$ is then the linearity of the integral that defines it.
\end{proof}

Proposition~\ref{prop:kick} works in the following setting, which isolates
Equation~\eqref{eq:kick}. Let $\zeta_0,\beta\in\C$, put $x=\Re\zeta_0$, let
$\tau_0<\tau_1$ be real, put $I=[\tau_0,\tau_1]$, let
$q\colon I\to\C$ be continuous, let $c$ solve $c'=\zeta_0c+\beta q$ on $I$ with
$c(\tau_0)=0$, and put $A(\tau)=e^{-\zeta_0\tau}c(\tau)$. The factor
$e^{-\zeta_0\tau}$ removes the free growth of $c$, so that $A$ is a weighted
integral of the drive $q$, as the next display shows. Since
$A'=e^{-\zeta_0\tau}(c'-\zeta_0c)=\beta e^{-\zeta_0\tau}q$ and $A(\tau_0)=0$,
\begin{equation}
\label{eq:kick-A}
A(\tau)=\beta\int_{\tau_0}^\tau e^{-\zeta_0s}\,q(s)\dd s\qquad(\tau\in I).
\end{equation}
Under the hypotheses of Lemma~\ref{lem:coupling}\textup{(2)}, the functions
$\cin$, $q_{k+1}$ and $\Ain$, with $\zeta_0=\lambda_k/\mu$, $\beta=\beta_k$ and
$\tau_1=\Tpin$, are an instance of this setting; here
$\tau_0=-A_{k+1}T_{k+1}<0<\Tpin$ by Lemma~\ref{lem:clock}(1) and
Definition~\ref{def:clock}.

Recall $z=\lamst(0)$ and $\sigma=\Re z>0$ (Definition~I.\ref{I-def:eigendata},
Proposition~I.\ref{I-prop:U}(a)), and the integral $\Ip{p}{x}$ of
Equation~(I.\ref{I-eq:Ip}), a decreasing function of $x$.

\begin{proposition}
\label{prop:kick}
Assume \textup{(N1)}--\textup{(N3)}, work in the setting above, and let
$1<\mu\le2$ and $G,G'\in[0,\Gone]$ with $\zeta_0=\lamst(G)/\mu$ and
$|\lamst(G)-z|,|\lamst(G')-z|\le\sigma(\mu-1)/(8\mu)$. Put $z'=\lamst(G')$ and
$w=z'-\zeta_0$. Let $T_g>0$ with $\tau_0<-T_g$, assume $\tau_1>0$, let $\bar Q>0$
and let $p\ge0$ be an integer. Assume:
\begin{description}[font=\normalfont]
\item[\textup{(Q1)}] there are $\qs\in\C\setminus\{0\}$ and a measurable
  $\err\colon[-T_g,0]\to[0,\infty)$ with $|q(\tau)-\qs e^{z'\tau}|\le\err(\tau)$ on
  $[-T_g,0]$;
\item[\textup{(Q2)}] there is a measurable function
  $Q_{\mathrm{pre}}\colon[\tau_0,-T_g]\to[0,\infty)$ with
  $|q(\tau)|\le Q_{\mathrm{pre}}(\tau)$ on $[\tau_0,-T_g]$;
\item[\textup{(Q3)}] $|q(\tau)|\le\bar Q(1+\tau)^p$ for $\tau\in[0,\tau_1]$.
\end{description}
Put $\Errg=\int_{-T_g}^0e^{-x\tau}\err(\tau)\dd\tau$ and
$\Errpre=\int_{\tau_0}^{-T_g}e^{-x\tau}Q_{\mathrm{pre}}(\tau)\dd\tau$, and assume
\begin{description}[font=\normalfont]
\item[\textup{(Kick1)}] $\mu-1\le|\qs|/\big(10|z|\,\bar Q\,\Ip{p}{\sigma/4}\big)$;
\item[\textup{(Kick2)}] $e^{-\Re w\,T_g}\le\frac12$;
\item[\textup{(Kick3)}] $\Errg\le|\qs|/\big(20|z|(\mu-1)\big)$;
\item[\textup{(Kick4)}] $\Errpre\le|\qs|/\big(20|z|(\mu-1)\big)$.
\end{description}
Then at $\tau=\tau_1$,
\[
\frac{|\beta|\,|\qs|}{10|z|(\mu-1)}\le|A(\tau)|\le\frac{3|\beta|\,|\qs|\,\mu}{|z|(\mu-1)} .
\]
\end{proposition}

The hypotheses describe the drive $q$ on three consecutive intervals:
\textup{(Q2)} bounds it on $[\tau_0,-T_g]$, \textup{(Q1)} says that on $[-T_g,0]$
it is close to the exponential $\qs e^{z'\tau}$, and \textup{(Q3)} bounds its
growth on $[0,\tau_1]$. In Lemma~\ref{lem:kicktrue}, where $G'=G_{k+1}$ and
$T_g=\hat T_{k+1}$, these intervals are, for $k+1<n$, the quiet phase and the
active phase of ring $k+1$ and the pinning phase $\Pin_k$ of ring $k$, in kick
time; at the seeded level $k+1=n$ the proof of that lemma applies
Proposition~\ref{prop:kick} with a slightly smaller $T_g$. The proof below shows
that $\Re w>0$ and that $|w|$ is comparable to $|z|(\mu-1)/\mu$. By
\textup{(Q1)}, the integrand of Equation~\eqref{eq:kick-A} on $[-T_g,0]$ is
close to $\qs e^{ws}$. \textup{(Kick2)} makes $|1-e^{-wT_g}|\ge\frac12$, so
this interval contributes a term of modulus comparable to
$|\beta||\qs|/(|z|(\mu-1))$, and \textup{(Kick1)}, \textup{(Kick3)} and
\textup{(Kick4)} make the contributions of $[0,\tau_1]$, of the error in
\textup{(Q1)} and of $[\tau_0,-T_g]$ small against it.

\begin{proof}
We first compare $w$ with $z$ and bound $x$ from below. Put
$e_\mu=\frac{\mu-1}{8\mu}\le\frac1{16}$ and write $z'=z+d_1$,
$\lamst(G)=z+d_2$ with $|d_1|,|d_2|\le e_\mu\sigma$. Then
$w=z\frac{\mu-1}\mu+d_1-\frac{d_2}\mu$ with
$|d_1-\frac{d_2}\mu|\le2e_\mu\sigma=\frac{\sigma(\mu-1)}{4\mu}$, so, since
$\Re z=\sigma\le|z|$,
\[
\Re w\ge\frac{\sigma(\mu-1)}\mu-\frac{\sigma(\mu-1)}{4\mu}=\frac34\sigma\frac{\mu-1}\mu>0,\qquad
\frac34\,\frac{|z|(\mu-1)}\mu\le|w|\le\frac54\,\frac{|z|(\mu-1)}\mu .
\]
As $\mu\le2$, $x=\Re\lamst(G)/\mu\ge(1-e_\mu)\sigma/\mu\ge\frac{15}{32}\sigma\ge\frac\sigma4>0$,
and as $\Ip{p}{\cdot}$ is decreasing, $\Ip{p}{x}\le\Ip{p}{\sigma/4}$.

Next we split the integral of Equation~\eqref{eq:kick-A} at $-T_g$ and $0$,
according to the three hypotheses on $q$. Put $\mathsf K=\beta\qs(1-e^{-wT_g})/w$,
the main term, and let $\tau=\tau_1$. Then
\[
A(\tau)=\beta\int_{\tau_0}^{-T_g}e^{-\zeta_0s}q(s)\dd s+\beta\int_{-T_g}^0e^{-\zeta_0s}q(s)\dd s
+\beta\int_0^\tau e^{-\zeta_0s}q(s)\dd s ,
\]
and $|e^{-\zeta_0s}|=e^{-xs}$. By (Q2), the
first term on the right has modulus at most $|\beta|\Errpre$. Since
\[
\int_{-T_g}^0e^{-\zeta_0s}\qs e^{z's}\dd s=\qs\int_{-T_g}^0e^{ws}\dd s=\qs(1-e^{-wT_g})/w,
\]
the second term is $\mathsf K$ plus a remainder of modulus at most
$|\beta|\Errg$, by (Q1). By (Q3), $x>0$ and Equation~(I.\ref{I-eq:Ip}), the third
has modulus at most
$|\beta|\bar Q\int_0^\infty e^{-xs}(1+s)^p\dd s=|\beta|\bar Q\,\Ip{p}{x}$.

Finally we compare the main term and the three remainders with
$\mathsf Z=\frac25|\beta||\qs|/(|z|(\mu-1))$.
By (Kick2) and $|e^{-wT_g}|=e^{-\Re w\,T_g}$, $|1-e^{-wT_g}|\in[\frac12,\frac32]$,
and the bounds on $|w|$ from the first paragraph give
\[
|\mathsf K|\ge\frac12\cdot\frac45\cdot\frac{|\beta||\qs|\mu}{|z|(\mu-1)}=\mu\,\mathsf Z,\qquad
|\mathsf K|\le\frac32\cdot\frac43\cdot\frac{|\beta||\qs|\mu}{|z|(\mu-1)}=\frac{2|\beta||\qs|\mu}{|z|(\mu-1)} .
\]
By (Kick1) and $\Ip{p}{x}\le\Ip{p}{\sigma/4}$,
$|\beta|\bar Q\,\Ip{p}{x}\le|\beta||\qs|/(10|z|(\mu-1))=\frac14\mathsf Z$, and by
(Kick3) and (Kick4), $|\beta|\Errg,|\beta|\Errpre\le\frac18\mathsf Z$; hence
$|A(\tau)-\mathsf K|\le\frac12\mathsf Z$. Since $\mu>1$,
\begin{gather*}
|A(\tau)|\ge|\mathsf K|-\tfrac12\mathsf Z\ge\big(\mu-\tfrac12\big)\mathsf Z\ge\tfrac12\mathsf Z=\frac{|\beta||\qs|}{5|z|(\mu-1)},\\
|A(\tau)|\le|\mathsf K|+\tfrac12\mathsf Z\le\Big(2\mu+\frac15\Big)\frac{|\beta||\qs|}{|z|(\mu-1)}
\le\frac{3|\beta||\qs|\mu}{|z|(\mu-1)} .\qedhere
\end{gather*}
\end{proof}

\section{The Relay Lemmas}
\label{sec:relay}

We prove the estimates on the rings that the joint bootstrap of
Section~\ref{sec:bootstrap} combines. They are
organized by the phase clock of Definition~\ref{def:clock}: ring $k$ is quiet on
$[0,\tsh_k]$, active on $\Phase_k=[\tsh_k,T_k]$, and debris after $T_k$.
Section~\ref{sec:param} chooses the parameters so that the conditions on them
hold.

Every statement lists its hypotheses explicitly, as bounds on the rings and
conditions on the parameters, and is causal: its
conclusions at a time $t$ use its hypotheses on $[0,t]$ only, and every phase
boundary entering a bound at time $t$ is at most $t$ or determined by
$\omega^{(n)}$ on $[0,t]$ (Lemma~\ref{lem:clock}), as the continuity argument in
time of Section~\ref{sec:bootstrap} requires.

\subsection{Standing Assumptions}

\begin{definition}
\label{def:standing}
We say that the \textbf{standing assumptions (SA)} hold if \textup{(N1)}--\textup{(N3)}
hold and the following objects are fixed:
\begin{enumerate}[label=\textup{(SA\arabic*)},leftmargin=*]
\item an integer $s\ge3$ and $\sigma'\in\big(\frac9{200},\min(\sigb,\frac34\sigma)\big)$;
  with $p=m=3$, the constants $\dH,C_2,C_3,\etmax$, the integer $M$, the
  functions $\kappa_0(\cdot)$ and $C(\cdot)$ of Theorem~\ref{thm:active};
\item the integers $s_h=M-1$, $N_Q=\max(2,\lceil\Lam(s_h)/\sigb\rceil)$ and
  $M_h=s_h+N_Q+1$, with $\Lam$ of Lemma~\ref{lem:energy}; the constants $\majst,C_Q$ of
  Theorem~\ref{thm:quiet} for these $s_h$ and $N_Q$; $C_{L1}$, $C_M$, $C_O$ of
  Lemma~\ref{lem:transfer} and $\Cin$ of Lemma~\ref{lem:inner-field}, at order $M_h$;
\item numbers $\delta_0\in(0,\dH]$, $\mu\in(1,\frac{17}{16}]$ and $\Tpin>0$; a
  number $\kseed\in(0,\kL]$, with $\kL$ of Lemma~\ref{lem:seed}; a number
  $\sbar\in(0,a/16]$ with $C_G\mu\sbar^2\le\Gone$, where $C_G$ is the constant of
  Lemma~\ref{lem:inner-circ}; an integer $n\ge2$; and the objects of
  Definitions~\ref{def:datum}--\ref{def:clock} with these parameters.
\end{enumerate}
Under (SA) we use the following derived quantities:
\begin{gather*}
c_\gamma=1-\frac{16\sigma'}{15\sigma},\qquad
\vartheta_\mu=\frac{\mu-1}3,\qquad \gamma_\mu=\min(\vartheta_\mu,c_\gamma),\\
\gamma'_\mu=\min\Big(\frac{13}{20}\cdot\frac{\mu-1}\mu,c_\gamma\Big),\qquad
\CDuh=\frac{2\mu}{\sigb(\mu-1)},\\
\rT(T)=(1+T)^me^{-\sigma T/4}\quad(T\ge0),
\end{gather*}
for every level $j<n$ the
exponent
\[
\gamma_j=\frac{\max(\sigma_j,\mu\sigma')}{\mu\sigma_{j+1}} ;
\]
and, for every level $k$ for which $\ksh_k$ is defined, the envelope $\hat a_k(u)=\ksh_ke^{\sigma_ku}$
($u\ge0$) and the transient $\rhohat_k$ of Theorem~\ref{thm:active} with $G=G_k$
and a constant $C_0\ge1$ named in each statement. We say that the
\textbf{detuning hypothesis (Det)} holds if
$\sup_{0\le G\le C_G\mu\sbar^2}|\lamst(G)-z|\le\sigma(\mu-1)/(8\mu)$.
\end{definition}

By Lemma~\ref{lem:detuning}, $c_\gamma$ bounds the exponents $1-\sigma'/\sigma_k$ from below. Recall that
$x_k=\sigma_k/\mu$ (Definition~\ref{def:regions}). The functions $\hat a_k$ and
$\rhohat_k$ are the envelope and the transient of Theorem~\ref{thm:active} for
ring $k$. Hypothesis (Det) keeps the unstable eigenvalues $\lamst(G)$ of all
rings close to $z=\lamst(0)$ (Definition~I.\ref{I-def:eigendata}). In the time of
ring $k$, ring $k+1$ grows at the rate $\mu\sigma_{k+1}$ and ring $k$ at the rate
$\sigma_k$; Lemma~\ref{lem:detuning} derives from (Det) that the first rate
exceeds the second by at least $\frac12\sigb(\mu-1)$.

Levels are integers in $[1,n]$; ring $1$ is the outermost ring and ring $n$, the
seeded ring, the innermost. We apply Lemma~\ref{lem:reduction} and
Theorems~\ref{thm:quiet} and \ref{thm:active} to ring $k$ in ring-$k$ units, as
explained before Definition~\ref{def:regions} and after
Lemma~\ref{lem:regions-basic}, without further comment; Theorem~\ref{thm:quiet}
for ring $k$ is applied with $G=G_k$, $h=h_k$, $\eta=\eta_k$, $\maj=\maj_k$ and
perturbation $w_k$, so that its conclusion (c) reads
$\|w_k(t)\|_{H^{s_h}}\le C_Q\maj_k(A_kt)$. By the paragraph after
Lemma~\ref{lem:regions-basic}, $\etin_k$, $\etout_k$ and $\eta_k$ are continuous
on every compact interval on which they are defined, so $\maj_k$ and
$\maj^{\mathrm{in}}_k$ are continuous on every $[0,\tau]$ on which $h_k$ is
defined. By Lemma~\ref{lem:inner-circ} and (SA3), $0\le G_k\le C_G\mu\sbar^2\le\Gone$
for every level $k$, so $\sigma_k$, $\lambda_k$, $\beta_k$ and $\rhohat_k$ are
defined and $\sigma_k\ge\sigb$. Since $A_i=\mu^{i-j}A_j$, the time of ring $i$ is
$\mu^{i-j}$ times the time of ring $j$.

\subsection{Quiet Rings}

A quiet ring is driven by the velocity field of the region inside it and by
the velocity field of the rings outside it. Lemma~\ref{lem:transfer-step} bounds these two fields by
the perturbations of the other rings and, for the inner field, by the inner field
of the next ring inward. With it, Lemma~\ref{lem:O} shows that,
while ring $j$ and the rings outside it are quiet, ring $j$ dominates them:
their Duhamel majorants are smaller than that of ring $j$ by a product of small
factors.

\begin{lemma}
\label{lem:detuning}
Assume (SA) and (Det). Put $e_\mu=\frac{\mu-1}{8\mu}$. Then for all levels $i,l$:
$|\sigma_i-\sigma|\le e_\mu\sigma$, $e_\mu\le\frac1{136}$, $\sigma_i\ge\frac{15}{16}\sigma>\sigma'$,
$\frac\sigma4\le x_i\le\frac{17}{16}\sigma$, and
\[
\mu\sigma_i-\sigma_l\ge\sigma\big((\mu-1)-e_\mu(\mu+1)\big)\ge\tfrac34\sigma(\mu-1)\ge\tfrac12\sigb(\mu-1);
\]
in particular $\sigma_i\mu^{i-j}\le\sigma_j-\sigb(\mu-1)/(2\mu)$ for $i<j$.
\end{lemma}
\begin{proof}
By Lemma~\ref{lem:inner-circ} and (SA3), $0\le G_i\le C_G\mu\sbar^2$, so (Det)
gives $|\sigma_i-\sigma|=|\Re(\lamst(G_i)-z)|\le e_\mu\sigma$. As
$1<\mu\le\frac{17}{16}$, $e_\mu=\frac18(1-\frac1\mu)\le\frac1{136}$, so
$\sigma_i\ge(1-e_\mu)\sigma\ge\frac{15}{16}\sigma>\frac34\sigma>\sigma'$ by
(SA1) and $\sigma=\sigma(0)\ge\sigb>0$ (Proposition~I.\ref{I-prop:U}(a)); and
$x_i=\sigma_i/\mu$ lies between $\frac{15}{16}\cdot\frac{16}{17}\sigma\ge\frac\sigma4$
and $(1+e_\mu)\sigma\le\frac{17}{16}\sigma$. For levels $i,l$,
$\mu\sigma_i-\sigma_l\ge\mu(1-e_\mu)\sigma-(1+e_\mu)\sigma=\sigma((\mu-1)-e_\mu(\mu+1))$,
and $e_\mu(\mu+1)=\frac{(\mu-1)(\mu+1)}{8\mu}\le\frac14(\mu-1)$ because
$\mu+1\le2\mu$; with $\sigma\ge\sigb$ this gives the display. For $i<j$,
$\mu^{i-j}\le\mu^{-1}$, and the display with $(j,i)$ in place of $(i,l)$ gives
$\sigma_i\mu^{i-j}\le\sigma_i/\mu\le\sigma_j-\sigb(\mu-1)/(2\mu)$.
\end{proof}

\begin{lemma}
\label{lem:transfer-step}
Assume (SA), and let $t\ge0$.
\begin{enumerate}[label=\textup{(\arabic*)}]
\item Let $1\le i<n$. Assume that $S_{i+1}(t)\subset N_{i+1}$ and that the fields
  $h_i$ and $h_{i+1}$ are defined at $t$ (Definition~\ref{def:regions}). Then for
  every integer $s'\ge0$,
  \[
  \etin_i(A_it)\le C_M\varepsilon_i\big(\|w_{i+1}(t)\|_{H^{s'}}+\etin_{i+1}(A_{i+1}t)\big).
  \]
\item Let $1\le k\le n$. Assume that $h_k$ is defined at $t$ and that
  $S_i(t)\subset N_i$ for every $i<k$. Then for every integer $s'\ge0$,
  \[
  \etout_k(A_kt)\le C_O\sum_{i<k}b_{i,k}\|w_i(t)\|_{H^{s'}},\qquad b_{i,k}=\mu^{i-k}(r_k/r_i)^{m-2};
  \]
  in particular $\etout_1=0$.
\end{enumerate}
\end{lemma}
\begin{proof}
We use three facts. (a) $\omega_{\ge l}(t)$ is real and $3$-fold rotationally symmetric
(Lemma~\ref{lem:regions-basic}(3)). (b) $\int w_l(t)=0$ if $S_l(t)\subset N_l$
(Lemma~\ref{lem:regions-basic}(4)). (c) If $g\in L^1(\R^2)$ is real and radial
with $\supp g\subset\overline B_\rho$, then $g$ is $m'$-fold rotationally symmetric for every
integer $m'\ge1$, so Lemma~\ref{lem:ext-multipole} without derivatives bounds the
velocity of $g$ minus its monopole at $|x|\ge2\rho$ by
$C_{m',0}\|g\|_{L^1}\rho^{m'}|x|^{-m'-1}$, with
$C_{m',0}=\frac1{2\pi}\sup_{0<r\le1/2}(1-r^{m'})^{-1}\le\frac1\pi$. For
$|x|\ge2\rho$ this bound tends to $0$ as $m'\to\infty$, since $\rho/|x|\le\frac12$;
hence
$\KBS\ast g(x)=\frac1{2\pi}(\int g)\frac{x^\perp}{|x|^2}$ for $|x|\ge2\rho$.

(1) We split the vorticity inside ring $i$ into ring $i+1$ and the region inside
ring $i+1$; in physical units,
$\omega_{\ge i+1}(t)=A_{i+1}\Phi(\cdot/r_{i+1})+A_{i+1}w_{i+1}(t,\cdot/r_{i+1})+\omega_{\ge i+2}(t)$.
The first term is radial and supported in $\overline B_{r_{i+1}}$, and
$N'_i\subset\{|x|\ge\frac a4r_i\}$ with $2r_{i+1}=2s_ir_i\le\frac a8r_i$, so by (c)
its velocity on $N'_i$ is a monopole; the second term has zero integral by (b).
Hence on $N'_i$ the velocity of $\omega_{\ge i+1}(t)$ minus its monopole is the
velocity of the second term, bounded in $C^{M_h}(N'_i)$, in ring-$i$ units, by
$C_M\varepsilon_i\|w_{i+1}(t)\|_{H^{s'}}$ (Lemma~\ref{lem:transfer}(2)), plus
$E_g$, the velocity of $g=\omega_{\ge i+2}(t)$ minus its monopole. The function
$g$ is $3$-fold rotationally symmetric by (a) and supported in $D_{i+2}(t)\subset B_{ar_{i+1}/8}$,
as $h_{i+1}$ is defined at $t$; so by Lemma~\ref{lem:transfer}(1), $E_g$ is
bounded in $C^{M_h}(N'_i)$, in ring-$i$ units, by $C_M\varepsilon_i$ times
$\sup|E_g|$ on $\{|x|=ar_{i+1}/4\}$ in ring-$(i+1)$ units. In these units $E_g$ is
$\hin_{i+1}$ (its monopole is $G_{i+1}x^\perp/|x|^2$ by the definition of
$G_{i+1}$ and Lemma~\ref{lem:inner-circ}), and $\{|y|=a/4\}\subset N'$, so the
supremum is at most $\etin_{i+1}(A_{i+1}t)$.

(2) By definition, $\hout_k$ is the velocity of
$\sum_{i<k}\omega^{(n)}(t)\mathbf 1_{S_i(t)}$ on $N'_k$, in ring-$k$ units. By the
last assertion of Lemma~\ref{lem:transfer}(4), this is the sum over $i<k$ of the
fields induced on $N'_k$ by $A_iw_i(t,\cdot/r_i)$, and by its first assertion each
of them is bounded in $C^{M_h}(N'_k)$, in ring-$k$ units, by
$C_Ob_{i,k}\|w_i(t)\|_{H^{s'}}$. Sum over $i<k$; for $k=1$ the sum is empty.
\end{proof}

\begin{lemma}
\label{lem:O}
Assume (SA) and (Det), and put
\[
\KO=8C_M(C_Q\CDuh+1)+32C_MC_Q\CDuh .
\]
Let $1\le j<n$, assume that the clock is
defined up to level $j+1$, and let $\tau\in(0,\tash_j]$ (ring-$j$ time). Assume
that for every $t\in[0,\tau/A_j]$ and every $i\le j$: $S_l(t)\cap B_{6r_i}=\emptyset$
for every $l<i$, $D_{i+1}(t)\subset B_{ar_i/8}$, and $\eta_i(A_it)\le1$. Assume
moreover $\maj_j\le\majst$ on $[0,\tau]$, $\KO\sbar^{2m}\le\frac14$,
$16C_OC_Q\CDuh\KO\sbar^{2m}\le1$ and $\KO\mu\sbar^{m+2}\le\frac12$. Then for every
$t\in[0,\tau/A_j]$ and every $i<j$:
\[
\maj_i(A_it)\le\tfrac34\Ptil_{i,j}\,\maj_j(A_jt),\qquad \Ptil_{i,j}=\prod_{p=i}^{j-1}\KO\varepsilon_p ,
\]
and $\etout_j(A_jt)\le\theta_j\maj_j(A_jt)$ with
$\theta_j=4C_OC_Q\KO s_{j-1}^{2m}$ for $j\ge2$ and $\theta_1=0$. Moreover, for
every $t\in(0,\tau/A_j]$ and every $i\le j$, the hypotheses of
Theorem~\ref{thm:quiet} hold for ring $i$ on $[0,A_it]$, in ring-$i$ units, with
$G=G_i$, $h=h_i$, $\eta=\eta_i$ and $\maj=\maj_i$; in particular its conclusions
\textup{(a)}--\textup{(c)} hold there.
\end{lemma}

\begin{proof}
We use ring-$j$ time $\tau'=A_jt$; ring $i$ is then at its own time
$\mu^{i-j}\tau'$. Put $\hat\sigma_i=\mu^{i-j}\sigma_i$ and
$\Pprod_{i,l}=\prod_{p=i}^{l-1}C_M\varepsilon_p$ for $i<l$, so that
$\Pprod_{i,l}=(C_M/\KO)^{l-i}\Ptil_{i,l}$ and $\Ptil_{i,l}\Ptil_{l,j}=\Ptil_{i,j}$.
Under the geometric hypotheses Definition~\ref{def:regions} defines $h_i$,
$i\le j$, on $[0,\tau/A_j]$, and the substitution $v=\mu^{i-j}v'$ there gives
\begin{equation}
\label{eq:O-duh}
\maj_i(\mu^{i-j}\tau')=\mu^{i-j}\int_0^{\tau'}e^{\hat\sigma_i(\tau'-v')}\eta_i(\mu^{i-j}v')\dd v' .
\end{equation}
For $i<j$, Lemma~\ref{lem:detuning} gives $\hat\sigma_i\le\sigma_j-\sigb(\mu-1)/(2\mu)$;
that is, measured in ring-$j$ time, ring $i$ grows more slowly than ring $j$.
Hence, if $F\ge0$ on $[0,\tau]$ satisfies $e^{\sigma_j(\tau'-v')}F(v')\le F(\tau')$
for $0\le v'\le\tau'\le\tau$, then, as $\mu^{i-j}\le1$,
\begin{equation}
\label{eq:O-fast}
\mu^{i-j}\int_0^{\tau'}e^{\hat\sigma_i(\tau'-v')}F(v')\dd v'\le F(\tau')\int_0^{\tau'}e^{-(\sigma_j-\hat\sigma_i)(\tau'-v')}\dd v'\le\CDuh F(\tau') ;
\end{equation}
and, since $\hat\sigma_i\le\sigma_j$,
\begin{equation}
\label{eq:O-same}
\mu^{i-j}\int_0^{\tau'}e^{\hat\sigma_i(\tau'-v')}\etin_j(v')\dd v'\le\int_0^{\tau'}e^{\sigma_j(\tau'-v')}\eta_j(v')\dd v'=\maj_j(\tau') .
\end{equation}

We prove the bound on $\maj_i$ by a continuity argument in $t$.
Equation~\eqref{eq:O-BH} below is this bound with the factor $1$ in place of
$\frac34$ and with a small term $\epsilon e^{\sigma_j\tau'}$ added to $\maj_j$, so
that its right side is positive at $t=0$, and we improve the factor to $\frac14$.

Fix $\epsilon\in(0,\majst e^{-\sigma_j\tau}]$ (unrelated to the numbers
$\varepsilon_p$) and put $\maj^\epsilon_j(\tau')=\maj_j(\tau')+\epsilon e^{\sigma_j\tau'}$.
By (M1) of Lemma~\ref{lem:majorant}, $F=\maj_j$ and $F=\maj^\epsilon_j$ satisfy
the hypothesis of Equation~\eqref{eq:O-fast}. Let $t_\epsilon$ be the supremum of
the $t_1\in[0,\tau/A_j]$ such that
\begin{equation}
\label{eq:O-BH}
\maj_i(A_it)\le\Ptil_{i,j}\,\maj^\epsilon_j(A_jt)\qquad\text{for every }t\in[0,t_1]\text{ and every }i<j .
\end{equation}
Both sides are continuous in $t$, since $h_i$, $i\le j$, is defined on
$[0,\tau/A_j]$; at $t=0$ the left side vanishes and the right side is
$\Ptil_{i,j}\epsilon>0$. So Equation~\eqref{eq:O-BH} holds on $[0,t_\epsilon]$, and
$t_\epsilon>0$.

We first verify the hypotheses of Theorem~\ref{thm:quiet} on $[0,t_\epsilon]$.
Every factor of $\Ptil_{i,j}$ satisfies
$\KO\varepsilon_p=\KO\mu s_p^{m+2}\le\KO\mu\sbar^{m+2}\le\frac12$, so
$\Ptil_{i,j}\le\frac12$, and Equation~\eqref{eq:O-BH} with $\maj_j\le\majst$ gives
$\maj_i(A_it)\le\frac12(\majst+\epsilon e^{\sigma_j\tau})\le\majst$ for
$t\in[0,t_\epsilon]$ and $i<j$. Ring $i<n$ starts at $\Phi$ in ring-$i$ units,
$G_i\in[0,\Gone]$, the geometric hypotheses are those of
Lemma~\ref{lem:reduction} for ring $i$, and
$\|h_i\|_{C^{M_h}(N')}\le\etin_i+\etout_i=\eta_i\le1$ with $\eta_i$ continuous. So
Theorem~\ref{thm:quiet} applies to ring $i<j$ on $[0,A_it_\epsilon]$ and to ring
$j$ on $[0,\tau]$ by the hypothesis $\maj_j\le\majst$; in particular, for $t\in[0,t_\epsilon]$ and $i\le j$,
$S_i(t)\subset\NNm_i\subset N_i$ and
\begin{equation}
\label{eq:O-Q}
\|w_i(t)\|_{H^{s_h}}\le C_Q\maj_i(A_it).
\end{equation}

Let $t\in[0,t_\epsilon]$, $\tau'=A_jt$ and $i<j$. Lemma~\ref{lem:transfer-step}(1)
at the levels $i,\dots,j-1$ with $s'=s_h$ gives by induction
$\etin_i(A_it)\le\sum_{l=i+1}^{j}\Pprod_{i,l}\|w_l(t)\|_{H^{s_h}}+\Pprod_{i,j}\etin_j(\tau')$,
and Lemma~\ref{lem:transfer-step}(2) gives
$\etout_i(A_it)\le C_O\sum_{l<i}b_{l,i}\|w_l(t)\|_{H^{s_h}}$. We insert
Equations~\eqref{eq:O-Q} and \eqref{eq:O-BH} into these two bounds and group the
terms of $\eta_i(A_it)=\etin_i(A_it)+\etout_i(A_it)$ by the ring they come from.
Ring $j$ and the region inside it (the terms with $\|w_j(t)\|_{H^{s_h}}$ and
$\etin_j(\tau')$) give
at most $\frac{C_M}{\KO}\Ptil_{i,j}(C_Q\maj^\epsilon_j(\tau')+\etin_j(\tau'))$, and the
rings $l\in(i,j)$ give
$C_Q\Ptil_{i,j}\maj^\epsilon_j(\tau')\sum_{l=i+1}^{j-1}(C_M/\KO)^{l-i}\le\frac87C_Q\frac{C_M}{\KO}\Ptil_{i,j}\maj^\epsilon_j(\tau')$,
as $C_M/\KO\le\frac18$. For the rings $l<i$ we have
$b_{l,i}\Ptil_{l,j}=b_{l,i}\Ptil_{l,i}\Ptil_{i,j}$ and
\[
b_{l,i}\Ptil_{l,i}=\mu^{l-i}\Big(\frac{r_i}{r_l}\Big)^{m-2}\KO^{i-l}\mu^{i-l}\prod_{p=l}^{i-1}s_p^{m+2}
=\KO^{i-l}\Big(\frac{r_i}{r_l}\Big)^{2m}\le\big(\KO\sbar^{2m}\big)^{i-l}\le4^{l-i}
\]
because $s_l\cdots s_{i-1}=r_i/r_l\le\sbar^{i-l}$ and $\KO\sbar^{2m}\le\frac14$;
summing the geometric series, these rings give at most $\frac43C_OC_Q\KO\sbar^{2m}$ times $\Ptil_{i,j}\maj^\epsilon_j(\tau')$. We insert
these bounds into Equation~\eqref{eq:O-duh} and use Equation~\eqref{eq:O-fast}
with $F=\maj^\epsilon_j$ and Equation~\eqref{eq:O-same}:
\[
\maj_i(A_it)\le\Big[\frac{C_M}{\KO}\big(C_Q\CDuh+1\big)+\frac87\cdot\frac{C_MC_Q\CDuh}{\KO}+\frac43C_OC_Q\CDuh\KO\sbar^{2m}\Big]\Ptil_{i,j}\,\maj^\epsilon_j(A_jt).
\]
As $\KO\ge8C_M(C_Q\CDuh+1)$, $\KO\ge32C_MC_Q\CDuh$ and
$16C_OC_Q\CDuh\KO\sbar^{2m}\le1$, the bracket is at most
$\frac18+\frac1{28}+\frac1{12}<\frac14$, so
\begin{equation}
\label{eq:O-improved}
\maj_i(A_it)\le\tfrac14\Ptil_{i,j}\,\maj^\epsilon_j(A_jt)\qquad(t\in[0,t_\epsilon],\ i<j).
\end{equation}
The right side is positive, so Equation~\eqref{eq:O-BH} is strict at
$t=t_\epsilon$ and, by continuity, holds beyond $t_\epsilon$ unless
$t_\epsilon=\tau/A_j$. Hence $t_\epsilon=\tau/A_j$ for every admissible $\epsilon$.

So, for every $t\in(0,\tau/A_j]$, the hypotheses of Theorem~\ref{thm:quiet} hold
for ring $i\le j$ on $[0,A_it]$, since they hold on $[0,A_i\tau/A_j]$; and
$\epsilon\to0$ in Equation~\eqref{eq:O-improved} gives
$\maj_i(A_it)\le\frac14\Ptil_{i,j}\maj_j(A_jt)\le\frac34\Ptil_{i,j}\maj_j(A_jt)$.
Finally, Lemma~\ref{lem:transfer-step}(2) at level $j$, Equation~\eqref{eq:O-Q},
this bound and the identity for $b_{i,j}\Ptil_{i,j}$ give
\[
\etout_j(A_jt)\le\tfrac14C_OC_Q\sum_{i<j}\KO^{j-i}\Big(\frac{r_j}{r_i}\Big)^{2m}\maj_j(A_jt).
\]
The term $i=j-1$ is $\KO s_{j-1}^{2m}$, and the term $i<j-1$ is
\[
\KO s_{j-1}^{2m}\KO^{j-1-i}(r_{j-1}/r_i)^{2m}\le\KO s_{j-1}^{2m}(\KO\sbar^{2m})^{j-1-i}\le\KO s_{j-1}^{2m}4^{-(j-1-i)} .
\]
So the sum is at most $\frac43\KO s_{j-1}^{2m}$, and
$\etout_j(A_jt)\le\frac13C_OC_Q\KO s_{j-1}^{2m}\maj_j(A_jt)\le\theta_j\maj_j(A_jt)$.
For $j=1$ there is no outer ring and $\etout_1=0$.
\end{proof}

The next proposition bounds the majorant of a quiet ring $j$ in each phase of ring
$j+1$, the next ring inward. Its hypothesis (H1) is the localization hypothesis
(Loc) of Definition~\ref{def:loc}. While ring $j+1$ is quiet, (o) and (i) bound
$\maj^{\mathrm{in}}_j$ and $\maj_j$ by a multiple of $\maj_{j+1}$. After ring $j+1$ has become active, (ii) and (iii)
bound $\maj_j$ up to the activation time $\tash_j$ of ring $j$; in
Equation~\eqref{eq:relay}, the first two terms in the bracket come from the
envelope $\hat a_{j+1}$ of the active ring $j+1$ and from its debris after
saturation, and $\Xi_{j+1}$ collects the quiet phase of ring $j+1$ and the
remaining terms of its active phase.

\begin{proposition}
\label{prop:relay}
Assume (SA) and (Det), and let $C_0\ge1$ and $K_1>0$.
There is a constant $C_\Xi=C_\Xi(C_0,K_1)$ for the estimates below. It
depends only on $C_0$, $K_1$, $\sigb$ and the constants of (SA1)--(SA2), and not
on $\mu$, $\delta_0$,
$\Tpin$, $\sbar$, $\kseed$, $n$ or the level. Let $1\le j<n$, assume that the
clock is defined up to level $j+1$, let $\tau\in[0,\tash_j]$ (ring-$j$ time), and
assume:
\begin{description}[font=\normalfont]
\item[\textup{(H1)}] \textup{(Loc)} holds at every $t\in[0,\tau/A_j]$;
\item[\textup{(H2)}] $\etout_j\le\theta_j\maj_j$ on $[0,\tau]$, with $\theta_j$ of
  Lemma~\ref{lem:O}, and $\theta_j\tau\le\frac12$;
\item[\textup{(H3)}] if $j+1<n$, the conclusions of Theorem~\ref{thm:quiet} hold for
  ring $j+1$, in its own time, on the interval
  $\big[0,A_{j+1}\min(\tau/A_j,\tsh_{j+1})\big]$;
\item[\textup{(H4)}] if $\tau/A_j>\tsh_{j+1}$, then
  $\maj_{j+1}(\tash_{j+1})\le K_1\ksh_{j+1}$, and for
  $t\in[\tsh_{j+1},\min(\tau/A_j,T_{j+1})]$ and $u=A_{j+1}(t-\tsh_{j+1})$:
  $D_{j+2}(t)\subset B_{ar_{j+1}/8}$,
  $\|w_{j+1}(t)\|_{H^s}\le2C_2\hat a_{j+1}(u)(1+\rhohat_{j+1}(u))$ and
  $\|\hin_{j+1}(\tash_{j+1}+u)\|_{C^{M_h}(N')}\le C_0\ksh_{j+1}\jap{u}^m$.
\end{description}
Then $\maj_j\le2\maj^{\mathrm{in}}_j$ on $[0,\tau]$, and:
\begin{enumerate}[label=\textup{(\roman*)}]
\item[\textup{(o)}] if $j+1<n$, then for every $\tau'\in[0,\min(\tau,A_j\tsh_{j+1})]$,
  \[
  \maj^{\mathrm{in}}_j(\tau')\le C_M\Big(\frac{2C_Q}{\sigb(\mu-1)}+1\Big)\varepsilon_j\,\maj_{j+1}(\mu\tau');
  \]
\item if $\tau\le A_j\tsh_{j+1}$, then
  $\maj_j(\tau)\le2C_M\big(\frac{2C_Q}{\sigb(\mu-1)}+1\big)\varepsilon_j\,\maj_{j+1}(\mu\tau)$;
\item if $\tau=\tash_j$, then
  \begin{equation}
  \label{eq:relay}
  \maj_j(\tash_j)\le2C_M\varepsilon_je^{x_j\Tpin}\Big[\frac{4C_2\delta_0}{\sigb(\mu-1)}+C_D(\Tpin)+\Xi_{j+1}\Big],
  \end{equation}
  where $C_D(T)=(2+2\Cstar T)^mT$ and
  \[
  \Xi_{j+1}=C_\Xi\Big(1+\frac1{\sigb(\mu-1)}\Big)\delta_0\Big(\frac{\ksh_{j+1}}{\delta_0}\Big)^{1-\gamma_j}\Big\langle\log\frac{\delta_0}{\ksh_{j+1}}\Big\rangle^{m+1};
  \]
\item if $A_j\tsh_{j+1}<\tau<\tash_j$, then $\maj_j(\tau)$ is at most the right side
  of Equation~\eqref{eq:relay}.
\end{enumerate}
All times used are at most $\tau/A_j$ or are determined by $\omega^{(n)}$ on
$[0,\tau/A_j]$.
\end{proposition}
\begin{proof}
Write $\tau^{\mathrm{act}}=A_j\tsh_{j+1}$ and $\tau^{\mathrm{sat}}=A_jT_{j+1}$
(the ring-$j$ times at which ring $j+1$ becomes active and saturates),
$\kappa=\ksh_{j+1}$ and $\hat T=\hat T_{j+1}=\sigma_{j+1}^{-1}\log(\delta_0/\kappa)$.
As the clock is defined up to level $j+1$, $\kappa\in(0,\delta_0)$,
$\tau^{\mathrm{sat}}=\tau^{\mathrm{act}}+\hat T/\mu$ and
$\tash_j=\tau^{\mathrm{sat}}+\Tpin/\mu$; at a ring-$j$ time $v\ge\tau^{\mathrm{act}}$
the active time of ring $j+1$ is $\mu(v-\tau^{\mathrm{act}})$. Every ring-$j$ time
$v\le\tau$ corresponds to $v/A_j\le\tsh_j$, where ring $j$ is quiet; so by
(Loc)(2) the rings $i<j$ lie outside $B_{6r_j}$, $D_{j+1}(v/A_j)\subset B_{ar_j/8}$,
and $h_j$ is defined there.

We first reduce the claims to bounds on $\maj^{\mathrm{in}}_j$, the majorant of
the inner field. Let $\tau'\in[0,\tau]$. As $\eta_j=\etin_j+\etout_j$, (H2) and (M1) of
Lemma~\ref{lem:majorant} for $\maj_j$ give
$\maj_j(\tau')\le\maj^{\mathrm{in}}_j(\tau')+\theta_j\int_0^{\tau'}e^{\sigma_j(\tau'-v)}\maj_j(v)\dd v
\le\maj^{\mathrm{in}}_j(\tau')+\theta_j\tau'\maj_j(\tau')$,
and $\theta_j\tau'\le\frac12$, so $\maj_j\le2\maj^{\mathrm{in}}_j$. It remains to
bound $\maj^{\mathrm{in}}_j$ by half the right sides of (i)--(iii), and to prove (o).

\emph{Quiet phase of ring $j+1$.} If $j+1=n$, then $\tau^{\mathrm{act}}=0$, (o) is
void and (i) concerns only $\tau=0$, where $\maj_j(0)=0$. Let $j+1<n$,
$\tau'\le\min(\tau,\tau^{\mathrm{act}})$, $v\in[0,\tau']$ and $t=v/A_j$. Then
$t\le\tsh_{j+1}<T_{j+1}$, so ring $j+1$ is quiet at $t$: (Loc)(1) gives
$S_{j+1}(t)\subset N_{j+1}$, and (Loc)(2) for ring $j+1$ gives
$D_{j+2}(t)\subset B_{ar_{j+1}/8}$ and $S_i(t)\cap B_{6r_{j+1}}=\emptyset$ for $i\le j$;
so $h_{j+1}$ is defined at $t$. Lemma~\ref{lem:transfer-step}(1) with $s'=s_h$
and Theorem~\ref{thm:quiet}(c) for ring $j+1$, which holds by (H3), give
$\etin_j(v)\le C_M\varepsilon_j(C_Q\maj_{j+1}(\mu v)+\eta_{j+1}(\mu v))$. By (M1),
$\maj_{j+1}(\mu v)\le e^{-\mu\sigma_{j+1}(\tau'-v)}\maj_{j+1}(\mu\tau')$, and
$\mu\sigma_{j+1}-\sigma_j\ge\frac12\sigb(\mu-1)$ (Lemma~\ref{lem:detuning}), so
$\int_0^{\tau'}e^{\sigma_j(\tau'-v)}\maj_{j+1}(\mu v)\dd v\le2\maj_{j+1}(\mu\tau')/(\sigb(\mu-1))$.
As $\sigma_j/\mu\le\sigma_{j+1}$, the substitution $v'=\mu v$ gives
\[
\int_0^{\tau'}e^{\sigma_j(\tau'-v)}\eta_{j+1}(\mu v)\dd v=\frac1\mu\int_0^{\mu\tau'}e^{\frac{\sigma_j}\mu(\mu\tau'-v')}\eta_{j+1}(v')\dd v'\le\maj_{j+1}(\mu\tau') .
\]
Together,
\begin{equation}
\label{eq:relay-quiet}
\maj^{\mathrm{in}}_j(\tau')\le C_M\Big(\frac{2C_Q}{\sigb(\mu-1)}+1\Big)\varepsilon_j\,\maj_{j+1}(\mu\tau')\qquad\big(0\le\tau'\le\min(\tau,\tau^{\mathrm{act}})\big),
\end{equation}
which is (o). If $\tau\le\tau^{\mathrm{act}}$, (o) at $\tau'=\tau$ and $\maj_j\le2\maj^{\mathrm{in}}_j$
give (i).

Let now $\tau^{\mathrm{act}}<\tau\le\tash_j$, and split
$\maj^{\mathrm{in}}_j(\tau)=\int_0^\tau e^{\sigma_j(\tau-v)}\etin_j(v)\dd v$ over
$\Iint_1=[0,\tau^{\mathrm{act}}]$, $\Iint_2=[\tau^{\mathrm{act}},\tau^{\mathrm{sat}}]$ and
$\Iint_3=[\tau^{\mathrm{sat}},\tash_j]$ (intersected with $[0,\tau]$), on which ring
$j+1$ is quiet, active and debris, respectively. In each part we use
$e^{\sigma_j(\tau-v)}\le e^{\sigma_j(\tash_j-v)}$ and
\begin{equation}
\label{eq:relay-id}
e^{\sigma_j(\tash_j-\tau^{\mathrm{sat}})}=e^{x_j\Tpin},\qquad
\kappa\,e^{\sigma_j(\tau^{\mathrm{sat}}-\tau^{\mathrm{act}})}=\kappa\Big(\frac{\delta_0}\kappa\Big)^{\frac{\sigma_j}{\mu\sigma_{j+1}}}\le\delta_0\Big(\frac\kappa{\delta_0}\Big)^{1-\gamma_j},
\end{equation}
where the inequality holds as $\sigma_j/(\mu\sigma_{j+1})\le\gamma_j$ and $\kappa<\delta_0$.

\emph{The part over $\Iint_1$} vanishes if $j+1=n$. Otherwise it equals
$e^{\sigma_j(\tau-\tau^{\mathrm{act}})}\maj^{\mathrm{in}}_j(\tau^{\mathrm{act}})$. We
bound the exponential factor by its value at $\tau=\tash_j$ and split its exponent
at $\tau^{\mathrm{sat}}$, so that both parts of Equation~\eqref{eq:relay-id} apply. With
Equation~\eqref{eq:relay-quiet} at $\tau'=\tau^{\mathrm{act}}$, the bound
$\maj_{j+1}(\mu\tau^{\mathrm{act}})=\maj_{j+1}(\tash_{j+1})\le K_1\kappa$ of (H4) and
Equation~\eqref{eq:relay-id}, this bounds the part over $\Iint_1$ by
$C_M\varepsilon_je^{x_j\Tpin}K_1\big(\frac{2C_Q}{\sigb(\mu-1)}+1\big)\delta_0(\kappa/\delta_0)^{1-\gamma_j}$.

\emph{The part over $\Iint_2$.} Let $v\in[\tau^{\mathrm{act}},\min(\tau,\tau^{\mathrm{sat}})]$,
$t=v/A_j$ and $u=\mu(v-\tau^{\mathrm{act}})\in[0,\hat T]$. Then $t\le T_{j+1}$, so
$S_{j+1}(t)\subset N_{j+1}$ by (Loc)(1). The field $h_{j+1}$ is defined at $t$:
$D_{j+2}(t)\subset B_{ar_{j+1}/8}$ by (H4), the rings $i<j$ lie outside
$B_{6r_j}\supset B_{6r_{j+1}}$, and ring $j$ lies in $N_j\subset\{|x|\ge\frac a2r_j\}$
by (Loc)(1), which misses $B_{6r_{j+1}}$ as $6s_j\le\frac{6a}{16}<\frac a2$.
Lemma~\ref{lem:transfer-step}(1) with $s'=s$ and (H4) give
$\etin_j(v)\le C_M\varepsilon_j(2C_2\hat a_{j+1}(u)+2C_2\hat a_{j+1}(u)\rhohat_{j+1}(u)+C_0\kappa\jap{u}^m)$.
As $\hat a_{j+1}(u)=\kappa e^{\sigma_{j+1}u}=\delta_0e^{-\mu\sigma_{j+1}(\tau^{\mathrm{sat}}-v)}$,
Lemma~\ref{lem:detuning} bounds the contribution of the first term by
$2C_2C_M\varepsilon_je^{x_j\Tpin}\delta_0\int_{-\infty}^{\tau^{\mathrm{sat}}}e^{-(\mu\sigma_{j+1}-\sigma_j)(\tau^{\mathrm{sat}}-v)}\dd v
\le C_M\varepsilon_je^{x_j\Tpin}\cdot4C_2\delta_0/(\sigb(\mu-1))$.
For the second term, the last assertion of Theorem~\ref{thm:active} (with
$G=G_{j+1}$) gives $\hat a_{j+1}(u)\rhohat_{j+1}(u)\le C(C_0)\kappa\jap{u}^me^{\sigma'u}$;
moreover $e^{\sigma_j(\tash_j-v)}=e^{x_j\Tpin}e^{\frac{\sigma_j}\mu(\hat T-u)}$ and,
for $u\in[0,\hat T]$,
$\frac{\sigma_j}\mu(\hat T-u)+\sigma'u\le\max(\frac{\sigma_j}\mu,\sigma')\hat T=\gamma_j\sigma_{j+1}\hat T=\gamma_j\log\frac{\delta_0}\kappa$.
With $\dd v=\dd u/\mu\le\dd u$, the second term contributes at most $2C_2C(C_0)$
times the left side of the following display, and the third (in the same way,
with $e^{\sigma'u}$ replaced by $1$) at most $C_0$ times it:
\[
C_M\varepsilon_je^{x_j\Tpin}\kappa\Big(\frac{\delta_0}\kappa\Big)^{\gamma_j}\int_0^{\hat T}\jap{u}^m\dd u
\le C_M\varepsilon_je^{x_j\Tpin}\delta_0\Big(\frac\kappa{\delta_0}\Big)^{1-\gamma_j}\hat T\jap{\hat T}^m .
\]
Finally $\hat T\le\sigb^{-1}\log(\delta_0/\kappa)$, as $\sigma_{j+1}\ge\sigb$, and
$\jap{cx}\le\max(1,c)\jap{x}$ for $c>0$, so
$\hat T\jap{\hat T}^m\le\langle\log(\delta_0/\kappa)\rangle^{m+1}/(\sigb\min(1,\sigb)^m)$.

\emph{The part over $\Iint_3$} matters only if $\tau>\tau^{\mathrm{sat}}$. Let
$v\in[\tau^{\mathrm{sat}},\tau]$ and $t=v/A_j$, so that $T_{j+1}\le t\le\tsh_j$ and
$\tauK=A_{j+1}(t-T_{j+1})=\mu(v-\tau^{\mathrm{sat}})\le\mu(\tash_j-\tau^{\mathrm{sat}})=\Tpin$. Ring $j$ is quiet at
$t$, so by (Loc)(2) $D_{j+1}(t)\subset B_{ar_j/8}\cap\overline B_{\varrho_j(t)}$ with
$\varrho_j(t)=(2+2\Cstar\tauK)r_{j+1}$. Hence $R_t=\max\{|x|:x\in D_{j+1}(t)\}$
satisfies $R_t<ar_j/8$ and $R_t\le(2+2\Cstar\Tpin)r_{j+1}$, and
Lemma~\ref{lem:transfer}(3) with $R=R_t$ gives
$\etin_j(v)\le C_M\varepsilon_j(R_t/r_{j+1})^m\le C_M\varepsilon_j(2+2\Cstar\Tpin)^m$. As
$e^{\sigma_j(\tau-v)}\le e^{\sigma_j(\tash_j-\tau^{\mathrm{sat}})}=e^{x_j\Tpin}$ (as
$v\ge\tau^{\mathrm{sat}}$, $\tau\le\tash_j$, and by Equation~\eqref{eq:relay-id}) and the
interval has length at most $\Tpin/\mu\le\Tpin$, this part is at most
$C_M\varepsilon_je^{x_j\Tpin}C_D(\Tpin)$.

Adding the three parts, with $\langle\log(\delta_0/\kappa)\rangle\ge1$ for the
part over $\Iint_1$, $\maj^{\mathrm{in}}_j(\tau)$ is at most
$C_M\varepsilon_je^{x_j\Tpin}[\frac{4C_2\delta_0}{\sigb(\mu-1)}+C_D(\Tpin)
+C'_\Xi\delta_0(\frac\kappa{\delta_0})^{1-\gamma_j}\langle\log\frac{\delta_0}\kappa\rangle^{m+1}]$
where $C'_\Xi\lesssim1+(\sigb(\mu-1))^{-1}$, with a constant depending
only on $C_0$, $K_1$ and the fixed data of (SA1)--(SA2).
Choose $C_\Xi(C_0,K_1)$ large enough that
$C'_\Xi\le C_\Xi(1+(\sigb(\mu-1))^{-1})$.
So $\maj^{\mathrm{in}}_j(\tau)$ is at most half the right side of
Equation~\eqref{eq:relay}, and $\maj_j\le2\maj^{\mathrm{in}}_j$ gives (ii) and (iii).

Causality: if $\tau\le\tau^{\mathrm{act}}$, the proof used the solution on
$[0,\tau/A_j]$ and the fact that ring $j+1$ is quiet there, which is determined by
$\omega^{(n)}$ on $[0,\tau/A_j]$ (Lemma~\ref{lem:clock}); otherwise
$\tsh_{j+1}<\tau/A_j$, and $\ksh_{j+1}$, $T_{j+1}$ and $\tsh_j$ are determined by
$\omega^{(n)}$ on $[0,\tsh_{j+1}]$ (Lemma~\ref{lem:clock}).
\end{proof}

While ring $k$ is active, the quiet rings outside it are driven by its growing
perturbation and act back on it through the outer field $\hout_k$. The next lemma
bounds this field by a small multiple $\eta_{b,k}$ of the envelope $\hat a_k$; this
is the form of the second part $h_2$ of the field in hypothesis (Field) of
Theorem~\ref{thm:active}.

\begin{lemma}
\label{lem:Oprime}
Assume (SA) and (Det), and let $C_0\ge1$, $K_1>0$ and $C_F\in(0,C_0]$.
There is a constant $\Cop=\Cop(C_0,K_1)$ for the estimate below. It
depends only on $C_0$, $K_1$, $\sigma$, $\sigb$ and the constants of
(SA1)--(SA2), and not on $C_F$, $\mu$,
$\delta_0$, $\Tpin$, $\sbar$, $\kseed$, $n$ or the level. Let $1\le k\le n$, assume
that the clock is defined up to level $k$, and let $t\in\Phase_k$ with $t>0$. Assume on
$[0,t]$: \textup{(Loc)}; if $k\ge2$, the hypotheses and conclusions of
Lemma~\ref{lem:O} for $j=k-1$ and $\tau=A_{k-1}t$ (ring $k-1$ is quiet on $[0,t]$,
since $t\le T_k<\tsh_{k-1}$), and $\theta_{k-1}A_{k-1}t\le\frac12$; the active
bound $\|w_k(\tsh_k+u/A_k)\|_{H^s}\le2C_2\hat a_k(u)(1+\rhohat_k(u))$ and the inner bound
$\|\hin_k(\tash_k+u)\|_{C^{M_h}(N')}\le C_F\ksh_k\jap{u}^m$ for
$u\in[0,A_k(t-\tsh_k)]$; and, if $k<n$: the conclusions of
Theorem~\ref{thm:quiet} for ring $k$ on $[0,\tash_k]$ (in its own time),
$\maj_k(\tash_k)\le K_1\ksh_k$, and $D_{k+1}(t')\subset B_{ar_k/8}$ for
$t'\in[\tsh_k,t]$. Then on $[0,A_k(t-\tsh_k)]$,
\begin{equation}
\label{eq:etab}
\|\hout_k(\tash_k+u)\|_{C^{M_h}(N')}\le\eta_{b,k}\,\hat a_k(u),\qquad
\eta_{b,k}=\Cop\frac{s_{k-1}^{2m}}{\sigb(\mu-1)}\le\eta_b=\Cop\frac{\sbar^{2m}}{\sigb(\mu-1)}
\end{equation}
for $k\ge2$, and $\hout_1=0$.
\end{lemma}
\begin{proof}
For $k=1$, $\hout_1=0$ by Lemma~\ref{lem:transfer-step}(2). Let $k\ge2$. We use
ring-$(k-1)$ time $\tau'=A_{k-1}t'$ and put $\tau^{\mathrm{act}}=A_{k-1}\tsh_k$; at
$\tau'\ge\tau^{\mathrm{act}}$ the active time of ring $k$ is
$u=\mu(\tau'-\tau^{\mathrm{act}})$, and $\mu\tau^{\mathrm{act}}=\tash_k$.

\emph{Reduction to $\maj_{k-1}$.} Let $t'\in[\tsh_k,t]$. By (Loc) and the
hypotheses of Lemma~\ref{lem:O}, $h_k$ is defined at $t'$:
$D_{k+1}(t')\subset B_{ar_k/8}$ by hypothesis (trivially if $k=n$), the rings
$i<k-1$ lie outside $B_{6r_{k-1}}\supset B_{6r_k}$, and ring $k-1$ lies in
$N_{k-1}$, which misses $B_{6r_k}$. Lemma~\ref{lem:transfer-step}(2) with
$s'=s_h$, the bound $\|w_i(t')\|_{H^{s_h}}\le C_Q\maj_i(A_it')$ of
Theorem~\ref{thm:quiet}(c) for $i\le k-1$ (its hypotheses hold by
Lemma~\ref{lem:O} if $t'>0$, and $w_i(0)=0$ as these rings start at $\Phi$), and
$\maj_i\le\frac34\Ptil_{i,k-1}\maj_{k-1}$ (Lemma~\ref{lem:O}) give
$\etout_k(A_kt')\le C_OC_Q[b_{k-1,k}+\frac34\sum_{i<k-1}b_{i,k}\Ptil_{i,k-1}]\maj_{k-1}(\tau')$.
As $b_{i,k}=b_{k-1,k}\,b_{i,k-1}$ and $b_{i,k-1}\Ptil_{i,k-1}\le4^{-(k-1-i)}$ (proof
of Lemma~\ref{lem:O}), the bracket is at most $b_{k-1,k}(1+\frac34\cdot\frac13)=\frac54b_{k-1,k}$, and
\begin{equation}
\label{eq:Op-red}
\etout_k(A_kt')\le2C_OC_Q\,b_{k-1,k}\,\maj_{k-1}(\tau') .
\end{equation}

\emph{Bound on $\maj_{k-1}$.} Proposition~\ref{prop:relay} applies with $j=k-1$,
$\tau=A_{k-1}t$ and the present $C_0$, $K_1$: the clock is defined up to level
$k$, and $A_{k-1}t\le\tash_{k-1}$ as $t\le T_k<\tsh_{k-1}$; (H1) is assumed; (H2)
follows from $\etout_{k-1}\le\theta_{k-1}\maj_{k-1}$ on $[0,A_{k-1}t]$
(Lemma~\ref{lem:O}) and $\theta_{k-1}A_{k-1}t\le\frac12$; (H3) is assumed, as
$[0,A_k\min(t,\tsh_k)]=[0,\tash_k]$; for (H4), let $t>\tsh_k$: the bound
$\maj_k(\tash_k)\le K_1\ksh_k$ is assumed if $k<n$ and holds as
$\maj_n(\tash_n)=\maj_n(0)=0$ if $k=n$, and on $[\tsh_k,\min(t,T_k)]=[\tsh_k,t]$
the inclusion $D_{k+1}\subset B_{ar_k/8}$, the active bound and the inner bound
(with $C_F\le C_0$) are assumed. So $\maj_{k-1}\le2\maj^{\mathrm{in}}_{k-1}$ on
$[0,A_{k-1}t]$, and, if $k<n$, (o) holds at $\tau'=\tau^{\mathrm{act}}=\min(A_{k-1}t,A_{k-1}\tsh_k)$.

We split $\maj^{\mathrm{in}}_{k-1}$ at $\tau^{\mathrm{act}}$. For
$\tau'\in[\tau^{\mathrm{act}},A_{k-1}t]$,
$\maj^{\mathrm{in}}_{k-1}(\tau')=e^{\sigma_{k-1}(\tau'-\tau^{\mathrm{act}})}\maj^{\mathrm{in}}_{k-1}(\tau^{\mathrm{act}})
+\int_{\tau^{\mathrm{act}}}^{\tau'}e^{\sigma_{k-1}(\tau'-v)}\etin_{k-1}(v)\dd v$, and $e^{\sigma_{k-1}(\tau'-\tau^{\mathrm{act}})}\le e^{\sigma_ku}$, as
$\sigma_{k-1}\le\mu\sigma_k$ (Lemma~\ref{lem:detuning}); recall
$\ksh_ke^{\sigma_ku}=\hat a_k(u)$. The first term vanishes if $k=n$, as then
$\tau^{\mathrm{act}}=0$; if $k<n$, (o) at $\tau^{\mathrm{act}}$,
$\mu\tau^{\mathrm{act}}=\tash_k$ and $\maj_k(\tash_k)\le K_1\ksh_k$ bound it by
$C_M\big(\frac{2C_Q}{\sigb(\mu-1)}+1\big)K_1\varepsilon_{k-1}\hat a_k(u)$.

For the second term, let $v\in[\tau^{\mathrm{act}},\tau']$ and
$u_v=\mu(v-\tau^{\mathrm{act}})$. At $t'=v/A_{k-1}\in[\tsh_k,t]$, ring $k$ lies in
$N_k$ by (Loc)(1), $h_{k-1}$ is defined by (Loc)(2) for the quiet ring $k-1$, and
$h_k$ is defined, as shown above; so Lemma~\ref{lem:transfer-step}(1) with $s'=s$
and the hypotheses on ring $k$ give
$\etin_{k-1}(v)\le C_M\varepsilon_{k-1}(2C_2\hat a_k(u_v)+2C_2\hat a_k(u_v)\rhohat_k(u_v)+C_F\ksh_k\jap{u_v}^m)$.
As $\hat a_k(u_v)=\ksh_ke^{\mu\sigma_k(v-\tau^{\mathrm{act}})}$ and
$\mu\sigma_k-\sigma_{k-1}\ge\frac12\sigb(\mu-1)$, the first part contributes at most
$4C_2C_M\varepsilon_{k-1}\hat a_k(u)/(\sigb(\mu-1))$. For the other two we write
$e^{\sigma_{k-1}(\tau'-v)}=e^{\sigma_{k-1}(\tau'-\tau^{\mathrm{act}})}e^{-\sigma_{k-1}v'}$
with $v'=v-\tau^{\mathrm{act}}$, so $u_v=\mu v'$, and use
$\jap{\mu v'}\le\mu(1+v')$ and $\mu^m\le2$. By the last assertion of
Theorem~\ref{thm:active},
$\hat a_k\rhohat_k(u_v)\le C(C_0)\ksh_k\jap{\mu v'}^me^{\mu\sigma'v'}$, and
$\sigma_{k-1}-\mu\sigma'>\frac{15}{16}\sigma-\frac{17}{16}\cdot\frac34\sigma=\frac{9\sigma}{64}$
by Lemma~\ref{lem:detuning} and (SA1). Hence, by Equation~(I.\ref{I-eq:Ip}), the
transient part contributes at most
$2C_2C(C_0)C_M\varepsilon_{k-1}\ksh_ke^{\sigma_ku}\mu^m\int_0^\infty(1+v')^me^{-\frac{9\sigma}{64}v'}\dd v'
\le4C_2C(C_0)\Ip{m}{\frac{9\sigma}{64}}C_M\varepsilon_{k-1}\hat a_k(u)$,
and, as $\sigma_{k-1}\ge\frac{9\sigma}{64}$, $\Ip{m}{\cdot}$ is decreasing and
$C_F\le C_0$, the inner part contributes at most
$2C_0\Ip{m}{\frac{9\sigma}{64}}C_M\varepsilon_{k-1}\hat a_k(u)$.

Using $\sigb(\mu-1)\le\sigb$ and choosing $\Cop(C_0,K_1)$ sufficiently large, the coefficients of
$C_M\varepsilon_{k-1}\hat a_k(u)$ found above add up to at most
$\Cop/(4C_OC_QC_M\sigb(\mu-1))$; so, for $\tau'\in[\tau^{\mathrm{act}},A_{k-1}t]$,
\[
\maj_{k-1}(\tau')\le2\maj^{\mathrm{in}}_{k-1}(\tau')\le\frac{\Cop}{2C_OC_Q}\cdot\frac{\varepsilon_{k-1}\hat a_k(u)}{\sigb(\mu-1)} .
\]
Finally $b_{k-1,k}\varepsilon_{k-1}=\mu^{-1}s_{k-1}^{m-2}\cdot\mu s_{k-1}^{m+2}=s_{k-1}^{2m}$,
so Equation~\eqref{eq:Op-red} gives
\[
\etout_k(A_kt')\le\Cop s_{k-1}^{2m}\hat a_k(u)/(\sigb(\mu-1))=\eta_{b,k}\hat a_k(u)
\]
for $u=A_k(t'-\tsh_k)$; this is Equation~\eqref{eq:etab}, as $A_kt'=\tash_k+u$ and so
$\|\hout_k(\tash_k+u)\|_{C^{M_h}(N')}=\etout_k(A_kt')$.
Also $\eta_{b,k}\le\eta_b$ as $s_{k-1}\le\sbar$. Every bound used at $t'$ involves
the solution on $[0,t']$ only.
\end{proof}

\subsection{The Active Ring and the Hand-Off}

This subsection provides two of the hypotheses of Theorem~\ref{thm:active} for an
active ring $k$: hypothesis (Coef) on its running coefficient
(Lemma~\ref{lem:drift}), and, for $k<n$, the bounds at the activation time
$\tsh_k$ that give hypothesis (Init) (Lemma~\ref{lem:handoff}).

When we say that the conclusions of Theorem~\ref{thm:active}, or of
Corollary~\ref{cor:active-moment}, hold for ring $k$ on $\Phase_k$, we apply that
result with $G=G_k$ and $\kappa=\ksh_k$ to the solution in ring-$k$ units with
time origin $\tsh_k$, decomposed into the parts of the vorticity on $S_k$, on
$D_{k+1}$ and on $S_1\cup\dots\cup S_{k-1}$. Its time $u$ is the active time
$A_k(t-\tsh_k)$, and $w(u)=w_k(\tsh_k+u/A_k)$, $h(u)=h_k(\tash_k+u)$ and $c=c_k$
(the argument of $w_k$ is the time of $\omega^{(n)}$, that of $h_k$ is ring-$k$
time).

By Equation~\eqref{eq:running-k}, the running coefficient $c_k$ of an active ring
moves away from $c_k(0)$ only through the external field $h_k=\hin_k+\hout_k$. For
$k<n$, the inner field acts through the moment $q_{k+1}$ at kick times $\tauK\ge\Tpin$, that
is, after the pinning phase $\Pin_k$, and its contribution is small when
$\rT(\Tpin)$ is; hypothesis (iii) below on the outer field follows from the bound
of Lemma~\ref{lem:Oprime}. In (ii), $q'$ stands for the amplitude $\qs$ of
Lemma~\ref{lem:kicktrue}, whose conclusion gives the lower bound on $|\Ain(\Tpin)|$.

\begin{lemma}
\label{lem:drift}
Assume (SA) and (Det). Let $1\le k\le n$, assume that the clock is defined up to
level $k$, and let $0\le\bar u\le\hat T_k$ and $\eta\ge0$. Assume:
\begin{enumerate}[label=\textup{(\roman*)}]
\item if $k<n$: $|q_{k+1}|\le\Cmom(1+\tauK)^m$ for $\tauK\in[\Tpin,\Tpin+\mu\bar u]$,
  with $\Cmom$ of Proposition~\ref{prop:confine};
\item if $k<n$: the hypotheses of Lemma~\ref{lem:coupling}\textup{(2)} hold,
  $\ksh_k\ge\frac12e^{x_k\Tpin}|\Ain(\Tpin)|$, and there is
  $q'\in\C\setminus\{0\}$ with
  $|\Ain(\Tpin)|\ge|\beta_k||q'|/(10|z|(\mu-1))$ and
  $\mu-1\le|q'|/(10|z|\Cmom\Ip{m}{\sigma/4})$;
\item $\|\hout_k(\tash_k+u)\|_{C^0(N')}\le\eta\,\hat a_k(u)$ for $u\in[0,\bar u]$
  (the argument of $\hout_k$ is ring-$k$ time);
\item for every $t\in[\tsh_k,\tsh_k+\bar u/A_k]$, $D_{k+1}(t)\subset B_{ar_k/8}$
  and the rings $i<k$ lie outside $B_{6r_k}$, so that $h_k$ and $c_k$ are
  defined there.
\end{enumerate}
Then for $0\le u\le u'\le\bar u$,
\begin{equation}
\label{eq:drift-interval}
|c_k(u')-c_k(u)|\le2\rT(\Tpin+\mu u)\,\ksh_k+\Cls\,\eta\,(u'-u)\,\ksh_k,
\end{equation}
without the first term when $k=n$. If $\rT(\Tpin)\le\frac1{64}$ and
$\Cls\eta\hat T_k\le\frac1{32}$, then hypothesis \textup{(Coef)} of
Theorem~\ref{thm:active} holds on $[0,\bar u]$
with $G=G_k$ and $\kappa=\ksh_k$, and if $\bar u=\hat T_k$ then
$|c_k(\hat T_k)|e^{\sigma_k\hat T_k}\in[\frac{15}{16}\delta_0,\frac{17}{16}\delta_0]$.
\end{lemma}

\begin{proof}
Write $\kappa=\ksh_k$, $\lambda=\lambda_k$ and $t(v)=\tsh_k+v/A_k$. By (iv),
$h_k=\hin_k+\hout_k$ is defined at every ring-$k$ time $\tash_k+v$,
$v\in[0,\bar u]$, so $c_k$ is defined on $[0,\bar u]$, and by
Equation~\eqref{eq:running-k} and the linearity of $\ellst(G_k)$ (recall
$G_k\in[0,\Gone]$, by Lemma~\ref{lem:inner-circ} and (SA3)), $c_k(u')-c_k(u)$ is the sum of the inner and outer integrals
$\int_u^{u'}e^{-\lambda v}\,\ellst(G_k)\big(-h^\bullet_k(\tash_k+v)\cdot\nabla\Phi\big)\dd v$,
$\bullet\in\{\mathrm{in},\mathrm{out}\}$.

By Lemma~I.\ref{I-lem:uniform} and (iii),
$|\ellst(G_k)(-\hout_k(\tash_k+v)\cdot\nabla\Phi)|\le\Cls\|\hout_k(\tash_k+v)\|_{C^0(N')}\le\Cls\eta\,\kappa e^{\sigma_kv}$;
as $|e^{-\lambda v}|=e^{-\sigma_kv}$, the outer integral is at most
$\Cls\eta\kappa(u'-u)$.

If $k=n$, then $D_{n+1}=\emptyset$ and $G_n=0$, so $\hin_n=0$ and the inner
integral vanishes. Let $k<n$. By (iv) and Lemma~\ref{lem:coupling}\textup{(1)},
$\ellst(G_k)(-\hin_k(\tash_k+v)\cdot\nabla\Phi)=\mu\beta_k\,q_{k+1}(t(v))$, where
$q_{k+1}$, as a function of the kick time, is evaluated at $\Tpin+\mu v$
(Definition~\ref{def:clock}). With $\zeta_0=\lambda/\mu$, $\Re\zeta_0=x_k$, the
substitution $\tauK=\Tpin+\mu v$ gives
\[
\int_u^{u'}e^{-\lambda v}\mu\beta_kq_{k+1}\dd v=e^{\zeta_0\Tpin}\beta_k\int_{\tau_a}^{\Tpin+\mu u'}e^{-\zeta_0\tauK}q_{k+1}(\tauK)\dd\tauK,
\qquad \tau_a=\Tpin+\mu u .
\]
By (i) and $1+\tau_a+y\le(1+\tau_a)(1+y)$ for $y\ge0$, its modulus is at most
\begin{align*}
e^{x_k\Tpin}|\beta_k|\Cmom\int_{\tau_a}^\infty e^{-x_k\tauK}(1+\tauK)^m\dd\tauK
&\le e^{x_k\Tpin}|\beta_k|\Cmom(1+\tau_a)^me^{-x_k\tau_a}\Ip{m}{x_k}\\
&\le e^{x_k\Tpin}|\beta_k|\Cmom\Ip{m}{\sigma/4}\,\rT(\tau_a),
\end{align*}
because
$x_k\ge\sigma/4$ (Lemma~\ref{lem:detuning}) and $\Ip{m}{\cdot}$ is decreasing. By (ii),
\[
|\beta_k|\Cmom\Ip{m}{\sigma/4}\le\frac{|\beta_k||q'|}{10|z|(\mu-1)}\le|\Ain(\Tpin)|\le2\kappa e^{-x_k\Tpin}.
\]
So the inner integral is at most $2\kappa\rT(\tau_a)$, which proves
Equation~\eqref{eq:drift-interval}.

For $u\in[0,\bar u]$, Equation~\eqref{eq:drift-interval} with the pair
$(0,u)$, $u\le\hat T_k$ and the two smallness conditions give
$|c_k(u)-c_k(0)|\le(\frac2{64}+\frac1{32})\kappa=\frac1{16}\kappa$, and, as
$|c_k(0)|=\kappa$, $|c_k(u)|\le\frac{17}{16}\kappa\le2\kappa$. As $c_k$ is the
running coefficient of Equation~\eqref{eq:running} for ring $k$ in ring-$k$
units, with $w_0=w_k(\tsh_k)$ and the field $u\mapsto h_k(\tash_k+u)$, (Coef)
holds with $\kappa=\ksh_k$, since the right side of (Coef) is at least $2\kappa$.
Finally
$|c_k(\hat T_k)|\in[\frac{15}{16}\kappa,\frac{17}{16}\kappa]$ and
$\kappa e^{\sigma_k\hat T_k}=\delta_0$ (Definition~\ref{def:clock}).
\end{proof}

At the activation time $\tsh_k$ of ring $k<n$, its unstable coefficient has
modulus $\ksh_k$ (Definition~\ref{def:clock}). The next lemma compares $\ksh_k$
with the modulus of the coefficient $\cin_k$ produced by the kick from the region
inside ring $k$ (Lemma~\ref{lem:coupling}\textup{(2)}), and bounds the perturbation
of ring $k$ at $\tsh_k$ by $\ksh_k$.

\begin{lemma}
\label{lem:handoff}
Assume (SA) and (Det), let $K_1>0$, let $1\le k<n$, and assume that the clock is
defined up to level $k+1$. Assume:
\begin{enumerate}[label=\textup{(\roman*)}]
\item for every $t\in[0,\tsh_k]$, $D_{k+1}(t)\subset B_{ar_k/8}$ and the rings
  $i<k$ lie outside $B_{6r_k}$, so that $\eta_k$ is defined on $[0,\tash_k]$
  (Definition~\ref{def:regions}); and the hypotheses of Theorem~\ref{thm:quiet}
  hold for ring $k$ on $[0,\tash_k]$, with $\eta=\eta_k$ and $\maj=\maj_k$;
\item for some $\theta_k\ge0$, $\etout_k\le\theta_k\maj_k$ on $[0,\tash_k]$ and
  $\Cls\theta_k\tash_kK_1\le\frac14$;
\item $\maj_k(\tash_k)\le\frac12K_1|\cin_k|$, with $\cin_k$ of Lemma~\ref{lem:coupling};
\item $\|\ellst\|C_QK_1\maj_k(\tash_k)\le\frac12$.
\end{enumerate}
Then, with $\Chand=K_1C_Q$ and in ring-$k$ units,
\begin{equation}
\label{eq:handoff}
\begin{gathered}
\tfrac12|\cin_k|\le\ksh_k\le\tfrac32|\cin_k|,\qquad \maj_k(\tash_k)\le K_1\ksh_k,\\
\|w_k(\tsh_k)\|_{H^{s_h}}\le\Chand\ksh_k,\qquad \supp(\Phi+w_k(\tsh_k))\subset\NNm .
\end{gathered}
\end{equation}
\end{lemma}

\begin{proof}
Write $\tau^\sharp=\tash_k$, $\maj=\maj_k$ and $c=\cin_k$. By (i) and
Theorem~\ref{thm:quiet}(a), for $t\in[0,\tsh_k]$, $D_{k+1}(t)\subset B_{ar_k/8}$,
the rings $i<k$ lie outside $B_{6r_k}$ and $S_k(t)\subset N_k$. So the hypotheses
of Lemma~\ref{lem:coupling}\textup{(2)} and \textup{(3)} hold, $c$ is defined,
and by Lemma~\ref{lem:coupling}\textup{(3)} the linear response of
Theorem~\ref{thm:quiet} for ring $k$ is
$\wlin=\wlin_{\mathrm{in},k}+\wlin_{\mathrm{out},k}$; at ring-$k$ time $\tau^\sharp$,
\[
\ellst(G_k)\big(w_k(\tsh_k)\big)=c+\ellst(G_k)\big(\wlin_{\mathrm{out},k}(\tau^\sharp)\big)+\ellst(G_k)\big(w_k(\tsh_k)-\wlin(\tau^\sharp)\big),
\]
and $\ksh_k$ is the modulus of the left side (Definition~\ref{def:clock}). As
$G_k\in[0,\Gone]$ (Lemma~\ref{lem:inner-circ} and (SA3)), $\ellst(G_k)$ is bounded on $X^0$ (Lemma~I.\ref{I-lem:uniform}), so
it commutes with the integral defining $\wlin_{\mathrm{out},k}$, and
Proposition~I.\ref{I-prop:semigroup} gives
$\ellst(G_k)(\wlin_{\mathrm{out},k}(\tau^\sharp))=\int_0^{\tau^\sharp}e^{\lambda_k(\tau^\sharp-v)}\ellst(G_k)(-\hout_k(v)\cdot\nabla\Phi)\dd v$.
By Lemma~I.\ref{I-lem:uniform}, $\|\hout_k\|_{C^0(N')}\le\etout_k$, (ii), (M1) of
Lemma~\ref{lem:majorant} and (iii),
\[
\big|\ellst(G_k)\big(\wlin_{\mathrm{out},k}(\tau^\sharp)\big)\big|
\le\Cls\theta_k\int_0^{\tau^\sharp}e^{\sigma_k(\tau^\sharp-v)}\maj(v)\dd v
\le\Cls\theta_k\tau^\sharp\maj(\tau^\sharp)\le\frac1{4K_1}\cdot\frac{K_1|c|}2=\frac{|c|}8 .
\]
Since $\ellst(G_k)$ has norm at most $\|\ellst\|$ on $X^0$
(Lemma~I.\ref{I-lem:uniform}) and $\|\cdot\|_{L^2}\le\|\cdot\|_{H^{s_h}}$,
Theorem~\ref{thm:quiet}(b), (iv) and (iii) give
\[
\big|\ellst(G_k)\big(w_k(\tsh_k)-\wlin(\tau^\sharp)\big)\big|\le\|\ellst\|C_Q\maj(\tau^\sharp)^2
\le\frac{\maj(\tau^\sharp)}{2K_1}\le\frac{|c|}4 .
\]
Hence $|\ksh_k-|c||\le\frac38|c|$, which gives the first claim. Then (iii) and
$|c|\le2\ksh_k$ give $\maj(\tau^\sharp)\le K_1\ksh_k$, and
Theorem~\ref{thm:quiet}(c) and (a) (the ring part is $\Phi+w_k(\tsh_k)$) give
$\|w_k(\tsh_k)\|_{H^{s_h}}\le C_Q\maj(\tau^\sharp)\le\Chand\ksh_k$ and
$\supp(\Phi+w_k(\tsh_k))\subset\NNm$.
\end{proof}

The next definition fixes the constants with which Section~\ref{sec:bootstrap}
applies the statements of this section: $K_1$ in Proposition~\ref{prop:relay} and
Lemmas~\ref{lem:Oprime}, \ref{lem:handoff} and~\ref{lem:kicktrue}; $\Abar$ in
the bound $\maj_j\le2\Abar\varepsilon_j$ on a quiet ring $j$; $\Klo$ and $\Kup$
in the bounds $\Klo s_k^{m+2}\le\ksh_k\le\Kup s_k^{m+2}$, $k<n$, on the
amplitudes at activation; and $\Chand$ in Equation~\eqref{eq:handoff}.

\begin{definition}
\label{def:K1}
Under (SA), with $c_N,C_N$ of Lemma~I.\ref{I-lem:uniform}, put $c_K=\frac{3c_N}{32\mu}$,
\begin{equation}
\label{eq:K1}
\begin{gathered}
K_1=\frac{8C_M|z|}{c_K}\Big[\frac{4C_2+1}{\sigb}+\frac{C_D(\Tpin)(\mu-1)}{\delta_0}\Big],\\
\Abar=2C_Me^{17\sigma\Tpin/16}\Big[\frac{(4C_2+1)\delta_0}{\sigb(\mu-1)}+C_D(\Tpin)\Big],
\end{gathered}
\end{equation}
with $C_D$ of Proposition~\ref{prop:relay}, and
\[
\Klo=\frac{c_N\delta_0}{22|z|(\mu-1)},\qquad \Kup=\frac{5C_N\delta_0\mu}{|z|(\mu-1)}e^{17\sigma\Tpin/16},\qquad \Chand=K_1C_Q .
\]
\end{definition}

\subsection{The Kick on the True Field}

Proposition~\ref{prop:kick} bounds the kick amplitude $\Ain$ under hypotheses on
the moment $q_{k+1}$ of the region inside ring $k$. This subsection verifies these
hypotheses for the solution $\omega^{(n)}$: on the active phase of ring $k+1$ by
Theorem~\ref{thm:active}, Corollary~\ref{cor:active-moment} and
Lemma~\ref{lem:drift}, on its quiet phase by Theorem~\ref{thm:quiet}, and for the
region inside ring $k+1$ by Proposition~\ref{prop:confine}; the bound on the
pinning phase $\Pin_k$ is a hypothesis.

The sets $S_1,\dots,S_n$ are pairwise disjoint
(Lemma~\ref{lem:regions-basic}\textup{(1)}) and each $\flow_t$ is a bijection, so
$D_j(t)$ is the disjoint union of $S_j(t)$ and $D_{j+1}(t)$ for $j\le n$, $t\ge0$.

\begin{lemma}
\label{lem:kicktrue}
Assume (SA) and (Det), and let $C_0\ge1$ and $K_1>0$.
Choose constants $C_{\mathrm{id}}$, $C_{\mathrm{od}}$ and $C_{\mathrm{qr}}$
depending only on the fixed data of (SA1)--(SA2), and
$C_{\mathrm{tr}}(C_0)$ with the additional dependence on $C_0$, large enough
for the error bounds in the proof below. Also choose $\Cdeep>0$ depending
only on $\mu$, $\Tpin$, $\Klo$, $\sigma$ and $\Cmom$ for the deep-ring bound there,
with $\Cmom$ of
Proposition~\ref{prop:confine}. Let $1\le k<n$, assume that
the clock is defined up to level $k+1$, put $\kappa=\ksh_{k+1}$,
$T_g=\hat T_{k+1}$ and
\[
\qs=2\pi\pout_*(G_{k+1})\,c_{k+1}(\hat T_{k+1})\,e^{\lambda_{k+1}\hat T_{k+1}} .
\]
Assume the following.
\begin{itemize}
\item $D_{k+1}(t)\subset B_{ar_k/8}$ for every $t\in[0,\tsh_k]$ (the hypotheses
  of Lemma~\ref{lem:coupling}\textup{(2)}).
\item The conclusions of Theorem~\ref{thm:active} and
  Corollary~\ref{cor:active-moment} hold for ring $k+1$ on $\Phase_{k+1}$ with
  $G=G_{k+1}$, $\kappa=\ksh_{k+1}$, the constant $C_0$ and $\eta_b$ of
  Equation~\eqref{eq:etab}.
\item The hypotheses \textup{(i)}--\textup{(iv)} of Lemma~\ref{lem:drift} hold for
  ring $k+1$ with $\bar u=\hat T_{k+1}$ and some $\eta\in[0,\eta_b]$, and
  $\rT(\Tpin)\le\frac1{64}$ and $\Cls\eta\hat T_{k+1}\le\frac1{32}$.
\item If $k+1<n$: for every $t\in[0,\tsh_{k+1}]$, $D_{k+2}(t)\subset B_{ar_{k+1}/8}$
  and the rings $i\le k$ lie outside $B_{6r_{k+1}}$, so that $\eta_{k+1}$ is
  defined on $[0,\tash_{k+1}]$ (Definition~\ref{def:regions}); the hypotheses of
  Theorem~\ref{thm:quiet} hold for ring $k+1$ on $[0,\tash_{k+1}]$, with
  $\eta=\eta_{k+1}$ and $\maj=\maj_{k+1}$; $\maj_{k+1}(\tash_{k+1})\le K_1\kappa$;
  the hypotheses of Proposition~\ref{prop:confine} hold with $k+1$ in place of $k$
  and $t_1=T_{k+1}$; and $\kappa\ge\Klo s_{k+1}^{m+2}$.
\item $|q_{k+1}|\le\Cmom(1+\tauK)^m$ for $\tauK\in[0,\Tpin]$ (also if $k+1=n$).
\item The following inequalities hold:
\begin{enumerate}[label=\textup{(\alph*)}]
\item $\frac{16}{15}\cdot\frac{2\pi C_3\delta_0^2}{\sigma}\le\frac14\cdot\frac{15}{16}\cdot\frac{\cout\delta_0}{20|z|}$;
\item $\mu-1\le\frac{15}{16}\cdot\cout\delta_0/\big(10|z|\Cmom\Ip{m}{\sigma/4}\big)$;
\item $\kappa\le\delta_04^{-\mu/(\mu-1)}$;
\item \begin{multline*}
  \frac{2\pi C_3\kappa^{1/2}\mu}{\frac34\sigma(\mu-1)}+C_{\mathrm{tr}}(C_0)\Big(\frac\kappa{\delta_0}\Big)^{\gamma_\mu}\log^{m+1}\frac{\delta_0}\kappa\\
  +\big[\Cdeep+C_{\mathrm{id}}\rT(\Tpin)\Ip{m}{\tfrac\sigma8}\big]\Big(\frac\kappa{\delta_0}\Big)^{\vartheta_\mu}\log^{m}\frac{\delta_0}\kappa
  \le\frac12\cdot\frac{15}{16}\cdot\frac{\cout}{20|z|(\mu-1)};
  \end{multline*}
\item $C_{\mathrm{qr}}\big(\frac{4\mu C_{L1}C_Q}{3\sigma(\mu-1)}+1\big)K_1\big(\frac\kappa{\delta_0}\big)^{\vartheta_\mu}\le\frac{15}{16}\cdot\frac{\cout}{20|z|(\mu-1)}$;
\item $C_{\mathrm{od}}\eta_b\delta_0\frac{\mu^2}{\sigma^2(\mu-1)^2}+2\pi C_3\eta_b\frac{4\mu\delta_0}{3\sigma(\mu-1)}
  \le\frac14\cdot\frac{15}{16}\cdot\frac{\cout\delta_0}{20|z|(\mu-1)}$.
\end{enumerate}
\end{itemize}
Then $|\qs|\in[\frac{15}{16},\frac{17}{16}]\,2\pi|\pout_*(G_{k+1})|\delta_0$, and the
kick amplitude $\Ain$ of Lemma~\ref{lem:coupling}\textup{(2)} satisfies
\[
\frac{|\beta_k|\,|\qs|}{10|z|(\mu-1)}\le|\Ain(\Tpin)|\le\frac{3|\beta_k|\,|\qs|\,\mu}{|z|(\mu-1)} .
\]
\end{lemma}

\begin{proof}
We apply Proposition~\ref{prop:kick} to the instance described after
Lemma~\ref{lem:coupling} (an instance of its setting by
Lemma~\ref{lem:coupling}\textup{(2)}): $G=G_k$, $G'=G_{k+1}$, $\beta=\beta_k$,
$q=q_{k+1}$, $c=\cin$, $A=\Ain$, $\tau_0=-A_{k+1}T_{k+1}$, $\tau_1=\Tpin$ and
$(\bar Q,p)=(\Cmom,m)$. We verify its hypotheses (Q1)--(Q3) and (Kick1)--(Kick4),
and then the remaining ones; (a)--(f) refer to the inequalities in the statement
above. Write $\sigma_+=\sigma_{k+1}$, $x=x_k=\sigma_k/\mu$,
$L=\log(\delta_0/\kappa)$, $z'=\lambda_{k+1}$, $\zeta_0=\lambda_k/\mu$ and
$w=z'-\zeta_0$, so that $\Re w=\sigma_+-x$ and $T_g=L/\sigma_+$, and let
$\mathsf B=\frac{15}{16}\cout\delta_0/(20|z|(\mu-1))$ be the common budget of (d),
(e) and (f).

By Lemma~\ref{lem:detuning} and its proof,
$\frac{15}{16}\sigma\le\sigma_+\le\frac{17}{16}\sigma$, $x\ge\frac\sigma4$,
$\mu\sigma_+-\sigma_k\ge\frac34\sigma(\mu-1)$, $e_\mu(\mu+1)\le\frac14(\mu-1)$ and
$|\sigma_i-\sigma|\le e_\mu\sigma$ for $i\in\{k,k+1\}$, so
$\mu\sigma_+-\sigma_k\le\mu(1+e_\mu)\sigma-(1-e_\mu)\sigma=\sigma((\mu-1)+e_\mu(\mu+1))\le\frac54\sigma(\mu-1)$; as
$\mu-1\le\frac1{16}$,
\begin{equation}
\label{eq:Rew}
\frac34\sigma\,\frac{\mu-1}\mu\le\Re w=\frac{\mu\sigma_+-\sigma_k}\mu\le\frac54\sigma(\mu-1)\le\frac\sigma8 .
\end{equation}
By the bounds on $\mu\sigma_+-\sigma_k$ and $\sigma_+$ above and $\mu\le\frac{36}{17}$,
$1-x/\sigma_+=\Re w/\sigma_+=(\mu\sigma_+-\sigma_k)/(\mu\sigma_+)\ge\frac{12}{17}\cdot\frac{\mu-1}\mu\ge\frac{\mu-1}3=\vartheta_\mu$;
also $1-\sigma'/\sigma_+\ge1-\frac{16\sigma'}{15\sigma}=c_\gamma$, and
$\gamma_\mu=\min(\vartheta_\mu,c_\gamma)$. By (c), $L\ge\frac\mu{\mu-1}\log4\ge1$;
and $T_g\le\frac{16}{15\sigma}L$. As $e^{\sigma_+T_g}=\delta_0/\kappa$,
$\kappa e^{yT_g}=\delta_0(\kappa/\delta_0)^{1-y/\sigma_+}$ for every real $y$.

\emph{The amplitude $\qs$ and (Kick1).} By the hypotheses on ring $k+1$,
Lemma~\ref{lem:drift} applies to ring
$k+1$ with $\bar u=\hat T_{k+1}$ and the given $\eta$, and gives
$|c_{k+1}(\hat T_{k+1})|e^{\sigma_+\hat T_{k+1}}\in[\frac{15}{16}\delta_0,\frac{17}{16}\delta_0]$;
as $|e^{\lambda_{k+1}\hat T_{k+1}}|=e^{\sigma_+\hat T_{k+1}}$, this is the claimed
range of $|\qs|$. As
$G_{k+1}\in[0,\Gone]$ (Lemma~\ref{lem:inner-circ} and (SA3)), $2\pi|\pout_*(G_{k+1})|\ge\cout$ (Lemma~I.\ref{I-lem:uniform}),
so $|\qs|\ge\frac{15}{16}\cout\delta_0>0$, hence $\qs\ne0$, and $\mathsf B\le|\qs|/(20|z|(\mu-1))$. Then
(b) gives $\mu-1\le|\qs|/(10|z|\Cmom\Ip{m}{\sigma/4})$, which is (Kick1).

\emph{(Kick2) and (Q3).} By Equation~\eqref{eq:Rew}, $\sigma_+\le\frac{17}{16}\sigma$
and (c),
$\Re w\,T_g\ge\frac34\cdot\frac{16}{17}\cdot\frac{\mu-1}\mu\,L\ge\frac{12}{17}\log4\ge\log2$,
so $e^{-\Re w\,T_g}\le\frac12$. (Q3) on $[0,\Tpin]$ is the assumed
bound $|q_{k+1}|\le\Cmom(1+\tauK)^m$.

\emph{(Q1).} Let $\tau\in[-T_g,0]$ and $t=T_{k+1}+\tau/A_{k+1}$. By
Definition~\ref{def:clock}, $T_{k+1}-\tsh_{k+1}=T_g/A_{k+1}$, so
$t\in\Phase_{k+1}$, and $u=\tau+T_g=A_{k+1}(t-\tsh_{k+1})\in[0,T_g]$ is the
active time of ring $k+1$.
We split $q_{k+1}$ into the moment of ring $k+1$ and that of the region inside it.
As $D_{k+1}(t)$ is the disjoint union of $S_{k+1}(t)$ and $D_{k+2}(t)$ and $\Qm$ is
linear,
$\Qm[\omega_{\ge k+1}(t)]=\Qm[\omega^{(n)}(t)\mathbf 1_{S_{k+1}(t)}]+\Qm[\omega_{\ge k+2}(t)]$;
the second term vanishes if $k+1=n$, as $D_{n+1}=\emptyset$. Otherwise, dividing by
$A_{k+1}r_{k+1}^{m+2}$ and using
$A_{k+2}r_{k+2}^{m+2}=\mu s_{k+1}^{m+2}A_{k+1}r_{k+1}^{m+2}$, the second term
becomes $\mu s_{k+1}^{m+2}q_{k+2}(t)$. As $T_{k+1}-T_{k+2}=\Tpin/A_{k+2}+T_g/A_{k+1}$
(Definition~\ref{def:clock}) and $A_{k+2}=\mu A_{k+1}$, the kick time of ring $k+1$ is
$\tau'=A_{k+2}(t-T_{k+2})=\mu A_{k+1}(t-T_{k+1})+A_{k+2}(T_{k+1}-T_{k+2})=\mu\tau+\Tpin+\mu T_g$.
So $q_{k+1}(\tau)=q^{\mathrm{ring}}(\tau)+q^{\mathrm{deep}}(\tau)$ with
\[
q^{\mathrm{ring}}(\tau)=\frac{\overline{\Qm[\omega^{(n)}(t)\mathbf 1_{S_{k+1}(t)}]}}{A_{k+1}r_{k+1}^{m+2}},\qquad
q^{\mathrm{deep}}(\tau)=\mu s_{k+1}^{m+2}\,q_{k+2}\big(\Tpin+\mu(\tau+T_g)\big),
\]
the argument of $q_{k+2}$ being the kick time of ring $k+1$, and
$q^{\mathrm{deep}}\equiv0$ if $k+1=n$.
By the assumed conclusions of Theorem~\ref{thm:active}, $w_{k+1}(t)$ is defined and
$\supp(\Phi+w_{k+1}(t))\subset\NN$, and in ring-$(k+1)$ units
$\omega^{(n)}(t)\mathbf 1_{S_{k+1}(t)}=A_{k+1}(\Phi+w_{k+1}(t))$
(Definition~\ref{def:regions}). As $\Phi$ is radial and $w_{k+1}(t)$ is real,
$q^{\mathrm{ring}}(\tau)=\overline{\Qm[w_{k+1}(t)]}=2\pi\pout_3[w_{k+1}(t)]$, which
is the function $q(u)$ of Corollary~\ref{cor:active-moment} for ring $k+1$, as
$w(u)=w_{k+1}(t)$. As $\qs e^{z'\tau}=2\pi\pout_*(G_{k+1})c_{k+1}(T_g)e^{z'u}$,
\begin{align*}
q_{k+1}(\tau)-\qs e^{z'\tau}&=\big[q^{\mathrm{ring}}(\tau)-2\pi\pout_*(G_{k+1})c_{k+1}(u)e^{z'u}\big]\\
&\quad+2\pi\pout_*(G_{k+1})\big(c_{k+1}(u)-c_{k+1}(T_g)\big)e^{z'u}+q^{\mathrm{deep}}(\tau).
\end{align*}
Corollary~\ref{cor:active-moment} bounds the bracket, and
Equation~\eqref{eq:drift-interval} for ring $k+1$, with the pair $(u,T_g)$ and
$\eta=\eta_b$, bounds the drift (hypothesis (iii) of Lemma~\ref{lem:drift} with
$\eta\le\eta_b$ implies (iii) with $\eta_b$, and
Equation~\eqref{eq:drift-interval} does not use the bound on
$\Cls\eta\hat T_{k+1}$). With $\hat a(u)=\kappa e^{\sigma_+u}$,
$\rhohat=\rhohat_{k+1}$ and $P=\max_{G\in[0,\Gone]}|\pout_*(G)|$, (Q1) holds with
\begin{align*}
\err(\tau)&=2\pi C_3\hat a(u)\big[\rhohat(u)+\hat a(u)+\kappa^{1/2}+\eta_b\big]\\
&\quad+2\pi P\big(2\kappa\,\rT(\Tpin+\mu u)\,\mathbf 1_{k+1<n}+\Cls\eta_b\kappa(T_g-u)\big)e^{\sigma_+u}
+|q^{\mathrm{deep}}(\tau)| .
\end{align*}

\emph{(Kick3).} We bound
$\Errg=\int_{-T_g}^0e^{-x\tau}\err(\tau)\dd\tau=\int_0^{T_g}e^{x(T_g-u)}\err(u-T_g)\dd u$
term by term. In the order in which they appear in $\err$, the terms are the
transient term, the $\hat a^2$-term, the $\kappa^{1/2}$-term and the $\eta_b$-term,
which come from Corollary~\ref{cor:active-moment}; the inner and the outer drift
terms, which come from Equation~\eqref{eq:drift-interval}; and the deep term
$|q^{\mathrm{deep}}|$. As $x<\sigma_+$, $2\sigma_+-x\ge\sigma_+\ge\frac{15}{16}\sigma$, so by
(a) and $\mu-1\le1$ the $\hat a^2$-term contributes
$2\pi C_3\kappa^2e^{xT_g}\int_0^{T_g}e^{(2\sigma_+-x)u}\dd u\le\frac{2\pi C_3\kappa^2e^{2\sigma_+T_g}}{2\sigma_+-x}
\le\frac{16}{15}\cdot\frac{2\pi C_3\delta_0^2}{\sigma}\le\frac{\mathsf B}4$.
In the same way, by Equation~\eqref{eq:Rew}, the $\kappa^{1/2}$-term and the
$\eta_b$-term contribute at most
$2\pi C_3\kappa^{1/2}\kappa e^{\sigma_+T_g}/\Re w\le\delta_0\cdot2\pi C_3\kappa^{1/2}\mu/(\frac34\sigma(\mu-1))$
and $2\pi C_3\eta_b\delta_0/\Re w\le2\pi C_3\eta_b\cdot4\mu\delta_0/(3\sigma(\mu-1))$.
For the transient term, Theorem~\ref{thm:active} gives
$\hat a(u)\rhohat(u)\le C(C_0)\kappa\jap{u}^me^{\sigma'u}$; as
$x(T_g-u)+\sigma'u\le\max(x,\sigma')T_g$ on $[0,T_g]$,
$\jap{T_g}\le(1+\frac{16}{15\sigma})L$, $1-\max(x,\sigma')/\sigma_+\ge\gamma_\mu$ and
$\kappa<\delta_0$, it contributes at most
\[
2\pi C_3C(C_0)\,T_g\jap{T_g}^m\kappa e^{\max(x,\sigma')T_g}
\le C_{\mathrm{tr}}(C_0)\,\delta_0
\Big(\frac\kappa{\delta_0}\Big)^{\gamma_\mu}L^{m+1},
\]
where $C_{\mathrm{tr}}(C_0)$ depends only on $C_0$ and the fixed data of
(SA1)--(SA2).
The inner drift term is absent if $k+1=n$. If $k+1<n$, we use
$\rT(\Tpin+\mu u)\le\rT(\Tpin)(1+\mu u)^me^{-\sigma\mu u/4}$,
$(1+\mu u)^m\le\mu^m(1+u)^m$, $\mu\le\frac{17}{16}$, $1-x/\sigma_+\ge\vartheta_\mu$ and
$\sigma_+-x-\frac{\mu\sigma}4\le\frac\sigma8-\frac\sigma4=-\frac\sigma8$
(Equation~\eqref{eq:Rew}); it contributes at most
\[
4\pi P\kappa e^{xT_g}\rT(\Tpin)\mu^m\int_0^{\infty}(1+u)^me^{-\sigma u/8}\dd u
\le C_{\mathrm{id}}\,\rT(\Tpin)\,\Ip{m}{\tfrac\sigma8}\,\delta_0\Big(\frac\kappa{\delta_0}\Big)^{\vartheta_\mu}.
\]
With $u=T_g-v$ and Equation~\eqref{eq:Rew}, the outer drift term contributes at
most
$2\pi P\Cls\eta_b\kappa e^{\sigma_+T_g}\allowbreak\int_0^{T_g}v\,e^{-\Re w\,v}\dd v
\le\frac{2\pi P\Cls\eta_b\delta_0}{(\Re w)^2}
\le C_{\mathrm{od}}\eta_b\delta_0\frac{\mu^2}{\sigma^2(\mu-1)^2}$.
The deep term vanishes if $k+1=n$. Let $k+1<n$. As
$\tsh_{k+1}=T_{k+2}+\Tpin/A_{k+2}\ge T_{k+2}$, $\Phase_{k+1}$ lies in
$[T_{k+2},T_{k+1}]$, and Proposition~\ref{prop:confine}(b) at level $k+1$,
with $t_0=T_{k+2}$ and $t_1=T_{k+1}$, gives $|q_{k+2}(\tau')|\le\Cmom(1+\tau')^m$
there; with $\tau'=\Tpin+\mu(\tau+T_g)$,
$|q^{\mathrm{deep}}(\tau)|\le\mu s_{k+1}^{m+2}\Cmom\big(1+\Tpin+\mu(\tau+T_g)\big)^m$
on $[-T_g,0]$. As $x\ge\frac\sigma4$,
$\int_{-T_g}^0e^{-x\tau}\dd\tau\le4e^{xT_g}/\sigma$. On $[-T_g,0]$,
$1+\Tpin+\mu(\tau+T_g)\le1+\Tpin+2T_g\le(1+\Tpin+\frac{32}{15\sigma})L$, as
$\mu\le2$, $T_g\le\frac{16}{15\sigma}L$ and $L\ge1$. With
$s_{k+1}^{m+2}\le\kappa/\Klo$, $\kappa e^{xT_g}=\delta_0(\kappa/\delta_0)^{1-x/\sigma_+}$,
$1-x/\sigma_+\ge\vartheta_\mu$ and $\kappa<\delta_0$, the deep term contributes at most
\[
\frac{4\mu\Cmom}{\sigma\Klo}\Big(1+\Tpin+\frac{32}{15\sigma}\Big)^mL^m\,\kappa e^{xT_g}
\le\Cdeep\,\delta_0\Big(\frac\kappa{\delta_0}\Big)^{1-x/\sigma_+}L^m
\le\Cdeep\,\delta_0\Big(\frac\kappa{\delta_0}\Big)^{\vartheta_\mu}L^m .
\]
As $L\ge1$, the $\kappa^{1/2}$-, transient,
inner drift and deep terms add up to at most $\delta_0$ times the left side of
(d), hence to at most $\frac{\mathsf B}2$; the $\eta_b$-term and the outer drift
term add up to at most $\frac{\mathsf B}4$ by (f). With the $\hat a^2$-term,
$\Errg\le\frac{\mathsf B}4+\frac{\mathsf B}2+\frac{\mathsf B}4=\mathsf B\le|\qs|/(20|z|(\mu-1))$, which is (Kick3).

\emph{(Q2) and (Kick4) when $k+1<n$.} By Definition~\ref{def:clock},
$\tau_0=-A_{k+1}T_{k+1}=-\tash_{k+1}-T_g<-T_g$, as $\tash_{k+1}>0$
(Lemma~\ref{lem:clock}(1)). On $[\tau_0,-T_g]$ we take $Q_{\mathrm{pre}}=|q_{k+1}|$,
so (Q2) holds. In the ring-$(k+1)$ time $v=\tau-\tau_0\in[0,\tash_{k+1}]$,
$-x\tau=xT_g+x(\tash_{k+1}-v)$, and the time of $\omega^{(n)}$ is $v/A_{k+1}$, the
argument of $q_{k+1}$ in Definition~\ref{def:regions}.

To bound $\Errpre$ we bound $q_{k+1}$ on the quiet phase of ring $k+1$ by the
perturbation and the inner field of that ring.
Put $j=k+1$, let $t\in[0,\tsh_j]$ and work in ring-$j$ units. As
$D_{j+1}(t)\subset B_{ar_j/8}$, $\hin_j(A_jt)$ is defined, and
Theorem~\ref{thm:quiet}(a) gives $S_j(t)\subset N_j$, so $w_j(t)$ is defined. As
$D_j(t)$ is the disjoint union of $S_j(t)$ and $D_{j+1}(t)$,
$\Qm[\omega_{\ge j}(t)]=A_jr_j^{m+2}(\Qm[\Phi+w_j(t)]+\Qm[g])$ with
$g(y)=A_j^{-1}\omega_{\ge j+1}(t,r_jy)$. Here $\Qm[\Phi]=0$ as $\Phi$ is radial,
and $\Phi$ and $\Phi+w_j(t)$ are supported in $\NNm\subset N\subset\overline B_2$
(Lemmas~I.\ref{I-lem:ring-facts} and~\ref{lem:annuli} and Theorem~\ref{thm:quiet}(a)), so
$|\Qm[w_j(t)]|\le2^m\|w_j(t)\|_{L^1}$. The real function $g$ is
$\nu^{\mathrm{in}}_j$ at ring-$j$ time $A_jt$ (Definition~\ref{def:regions}); put
$p_m=\pout_m[g]=\frac1{2\pi}\int g(y')\,\bar y'^{\,m}\dd y'$
(Definition~I.\ref{I-def:multipole}). As $D_{j+1}(t)\subset B_{ar_j/8}$,
Lemma~\ref{lem:coupling}\textup{(1)} with $j$ in place of $k$ shows that the
mode-$m$ profile of $\hin_j(A_jt)\cdot\er$ is $\frac i2p_mr^{-m-1}$ for
$a/4\le r\le3$ (recall $m=3$); at $r=1$,
$\frac1{2\pi}\int_0^{2\pi}\hin_j(A_jt)(e^{i\theta})\cdot\er\,e^{-im\theta}\dd\theta=\frac i2p_m$,
so $|p_m|\le2\max_{|y|=1}|\hin_j(A_jt)(y)|\le2\etin_j(A_jt)$, the unit circle
lying in $N'$. As $g$ is real, $\overline{\Qm[g]}=\int g(y)\,\bar y^{\,m}\dd y=2\pi p_m$,
so $|\Qm[g]|\le4\pi\etin_j(A_jt)$. Dividing by $A_jr_j^{m+2}$ gives
$|q_j(t)|\le C_{\mathrm{qr}}(\|w_j(t)\|_{L^1}+\etin_j(A_jt))$.

For $v\in[0,\tash_j]$, $w_j(v/A_j)$ is supported in $N$, so the
Cauchy-Schwarz inequality, Theorem~\ref{thm:quiet}(c) and (M1) of
Lemma~\ref{lem:majorant} give, with $C_{L1}=|N|^{1/2}$ of (SA2),
\[
\|w_j(v/A_j)\|_{L^1}\le C_{L1}\|w_j(v/A_j)\|_{H^{s_h}}\le C_{L1}C_Q\maj_j(v)\le C_{L1}C_Qe^{-\sigma_+(\tash_j-v)}\maj_j(\tash_j).
\]
Since $\int_0^{\tash_j}e^{-(\sigma_+-x)(\tash_j-v)}\dd v\le(\sigma_+-x)^{-1}$, the
part of $\int_0^{\tash_j}e^{x(\tash_j-v)}|q_j(v/A_j)|\dd v$ coming from
$\|w_j\|_{L^1}$ is at most $C_{\mathrm{qr}}C_{L1}C_Q\maj_j(\tash_j)/(\sigma_+-x)$.
As $x<\sigma_+$ and $\etin_j\le\eta_j$,
$\int_0^{\tash_j}e^{x(\tash_j-v)}\etin_j(v)\dd v\le\maj_j(\tash_j)$.
These two bounds, $\maj_{k+1}(\tash_{k+1})\le K_1\kappa$,
Equation~\eqref{eq:Rew} for $\sigma_+-x=\Re w$, and $1-x/\sigma_+\ge\vartheta_\mu$ give
\[
\Errpre=e^{xT_g}\int_0^{\tash_{k+1}}e^{x(\tash_{k+1}-v)}|q_{k+1}(v/A_{k+1})|\dd v
\le C_{\mathrm{qr}}\Big(\frac{4\mu C_{L1}C_Q}{3\sigma(\mu-1)}+1\Big)K_1\,\delta_0\Big(\frac\kappa{\delta_0}\Big)^{\vartheta_\mu},
\]
which is at most $\mathsf B\le|\qs|/(20|z|(\mu-1))$ by (e); this is (Kick4). The
remaining hypotheses of Proposition~\ref{prop:kick} hold: $G_k,G_{k+1}\in[0,\Gone]$
(Lemma~\ref{lem:inner-circ} and (SA3)),
$|\lamst(G_i)-z|\le\sigma(\mu-1)/(8\mu)$ for $i\in\{k,k+1\}$ by (Det),
$1<\mu\le2$, $T_g>0$ and $\tau_1=\Tpin>0$. So Proposition~\ref{prop:kick} applies,
and its conclusion at $\tau_1$ is the claimed bound on $|\Ain|$.

\emph{The seeded level $k+1=n$.} Here $\tsh_n=0$, so $\tau_0=-T_g$, while
Proposition~\ref{prop:kick} asks for $\tau_0<-T_g$. We apply it with $T_g$
replaced by some $T'_g\in[\frac{17}{24}T_g,T_g)$ and the same $\qs$, $z'$ and
$\err$. (Q1) and $\Errg\le\mathsf B$ hold on the shorter interval $[-T'_g,0]$, so
(Kick3) holds; (Kick1) and (Q3) do not involve $T_g$; and
$\Re w\,T'_g\ge\frac{17}{24}\Re w\,T_g\ge\log2$, because the computation for (Kick2)
gives $\Re w\,T_g\ge\frac{12}{17}\log4=\frac{24}{17}\log2$.
On $[\tau_0,-T'_g]=[-T_g,-T'_g]$ we take $Q_{\mathrm{pre}}=|q_n|$; then (Q2) holds
and $\Errpre\le(T_g-T'_g)e^{xT_g}\max_{[-T_g,0]}|q_n|$, as $q_n$ is continuous.
This bound tends to $0$ as $T'_g\to T_g$, so $T'_g$ close enough to $T_g$ gives
(Kick4). The remaining hypotheses hold as in
the case $k+1<n$, and Proposition~\ref{prop:kick} gives the claimed bound.
\end{proof}

\subsection{Time}

The relay ends when the outermost ring saturates, at the time $T_1$.
Lemma~\ref{lem:time} bounds $T_1$ by a number $\Tstar$ independent of $n$, given
lower bounds on the amplitudes $\ksh_k$: the active phase of ring $k$ lasts a
ring-$k$ time $\hat T_k$ of at most order $\log(1/s_k)$, which grows linearly in
$k$, while $A_k$ grows geometrically. The resulting bounds on the relay times
are summable over infinitely many rings, so the total time remains bounded
as the number of rings tends to infinity.
Lemma~\ref{lem:quiet-time} bounds in the same way the ring-$j$ time that a quiet
ring $j$ spends before its activation; Section~\ref{sec:bootstrap} uses it, with
Lemma~\ref{lem:theta-quiet} of Section~\ref{sec:param}, to verify the conditions on $\theta_j\tau$ in
Proposition~\ref{prop:relay} and on $\theta_k\tash_k$ in Lemma~\ref{lem:handoff}.

\begin{lemma}
\label{lem:time}
Assume (SA) and (Det). Put $\Tz=\Tpin+1+\sigma'^{-1}\log_+(\delta_0/\Klo)$ and
\[
\Tstar=\frac1{A_1}\Big[\frac{\mu}{\mu-1}\Big(\Tz+\frac{m+2}{\sigma'}\log\frac1{\sbar}\Big)
+\frac{(m+2)\log(2\mu)}{\sigma'}\cdot\frac{\mu^2}{(\mu-1)^2}+1+\frac{1}{\sigma'}\log\frac{2\delta_0}{\kseed}\Big].
\]
If the clock is defined up to level $1$, $\ksh_k\ge\Klo s_k^{m+2}$ for $k<n$,
and $\ksh_n\ge\frac12\kseed$, then $T_1\le\Tstar$.
\end{lemma}

\begin{proof}
By Lemma~\ref{lem:detuning} and (SA1), $\sigma_k\ge\frac{15}{16}\sigma>\frac34\sigma>\sigma'$
for every $k$. Since the clock is defined up to level $1$, $\ksh_k\in(0,\delta_0)$,
so $\log(\delta_0/\ksh_k)>0$ and $\hat T_k\le\sigma'^{-1}\log(\delta_0/\ksh_k)$. For
$k<n$, the hypothesis gives
$\hat T_k\le\sigma'^{-1}\big[\log_+(\delta_0/\Klo)+(m+2)\log(1/s_k)\big]$, and by
Definition~\ref{def:clock}, with $A_{k+1}=\mu A_k\ge A_k$,
\[
T_k-T_{k+1}=\frac{\Tpin}{A_{k+1}}+\frac{\hat T_k}{A_k}\le\frac1{A_k}\Big[\Tz+\frac{m+2}{\sigma'}\log\frac1{s_k}\Big].
\]
For $k=n$, $\tsh_n=0$ and $\ksh_n\ge\frac12\kseed$ give
$T_n=\hat T_n/A_n\le\sigma'^{-1}\log(2\delta_0/\kseed)/A_1$. We sum over $k$, using
$T_1=T_n+\sum_{k=1}^{n-1}(T_k-T_{k+1})$, $\log(1/s_k)=\log(1/\sbar)+k\log(2\mu)$,
$A_k=A_1\mu^{k-1}$, and
\[
\sum_{k\ge1}\mu^{1-k}=\frac\mu{\mu-1},\qquad\sum_{k\ge1}k\mu^{1-k}=\frac{\mu^2}{(\mu-1)^2}.
\]
All terms are nonnegative, so $T_1$ is at most $\Tstar$ minus $1/A_1$.
\end{proof}

\begin{lemma}
\label{lem:quiet-time}
Assume (SA), (Det) and $\Kup\sbar^{m+2}\le\kseed$, and let $\Tz$ be as
in Lemma~\ref{lem:time}. Let $1\le j<n$ and put
\[
\bar\tau^\sharp_j=\frac{3\Tz+(2m+5)\sigma'^{-1}\log(1/s_j)}{\mu-1}+\frac{(m+2)\mu\log(2\mu)}{\sigma'(\mu-1)^2}.
\]
Let $t\ge0$, and assume: for every $i\in(j,n]$ such that $\tsh_i$ is defined and
$\tsh_i\le t$, the clock is defined up to level $i$, and moreover
$\ksh_i\ge\Klo s_i^{m+2}$ if $i<n$; $t\le\tsh_j$ whenever $\tsh_j$ is defined;
and $\ksh_n\ge\frac12\kseed$. Then $A_jt\le\bar\tau^\sharp_j$.
\end{lemma}

\begin{proof}
Recall (Definition~\ref{def:clock}) that $\tsh_i$ is defined once the clock is
defined up to level $i+1$, and then $\tsh_{i+1}\le T_{i+1}\le\tsh_i$, as
$\hat T_{i+1}\ge0$. Let $i_*$ be the least $i\in(j,n]$ such that $\tsh_i$ is
defined and $\tsh_i\le t$; it exists as $\tsh_n=0$. For $i\in[i_*,n]$, $\tsh_i$ is
defined and $\tsh_i\le\tsh_{i_*}\le t$, so by hypothesis the clock is defined up to level $i$,
and $\ksh_i\ge\Klo s_i^{m+2}$ if $i<n$; in particular $T_{i_*}$ is defined. If
$i_*-1>j$, then $\tsh_{i_*-1}=T_{i_*}+\Tpin/A_{i_*}$ is defined and, by the
minimality of $i_*$, larger than $t$; if $i_*-1=j$, then
$t\le\tsh_j=T_{j+1}+\Tpin/A_{j+1}$ by hypothesis. In both cases, by
Definition~\ref{def:clock},
\[
t\le T_{i_*}+\frac{\Tpin}{A_{i_*}}=\frac{\hat T_n}{A_n}+\sum_{i=i_*}^{n-1}\Big(\frac{\Tpin}{A_{i+1}}+\frac{\hat T_i}{A_i}\Big)+\frac{\Tpin}{A_{i_*}}.
\]
We multiply by $A_j$ and use $A_j/A_i=\mu^{j-i}$. The pinning terms contribute at
most $\Tpin\sum_{i>j}\mu^{j-i}=\Tpin/(\mu-1)$. For $j<i<n$ with $i\ge i_*$, as in
the proof of Lemma~\ref{lem:time},
$\hat T_i\le\sigma'^{-1}[(m+2)\log(1/s_i)+\log_+(\delta_0/\Klo)]$ with
$\log(1/s_i)=\log(1/s_j)+(i-j)\log(2\mu)$; with $\sum_{i>j}\mu^{j-i}=\frac1{\mu-1}$
and $\sum_{i>j}(i-j)\mu^{j-i}=\frac\mu{(\mu-1)^2}$, these active phases
contribute at most
\[
\frac{\sigma'^{-1}\log_+(\delta_0/\Klo)+(m+2)\sigma'^{-1}\log(1/s_j)}{\mu-1}+\frac{(m+2)\mu\log(2\mu)}{\sigma'(\mu-1)^2}.
\]
The seeded phase contributes $\mu^{j-n}\hat T_n\le\hat T_n$. As $c_N\le C_N$ and
$\frac1{22}\le5\mu e^{17\sigma\Tpin/16}$, $\Klo\le\Kup$, so
$\ksh_n\ge\frac12\kseed\ge\frac12\Kup\sbar^{m+2}\ge\frac12\Klo\sbar^{m+2}$ and
\[
\hat T_n\le\frac1{\sigma'}\Big[\log2+\log_+\frac{\delta_0}{\Klo}+(m+2)\log\frac1{\sbar}\Big]
\le\Tz+\frac{m+3}{\sigma'}\log\frac1{s_j},
\]
because $\log2\le\log(1/\sbar)\le\log(1/s_j)$, by $\sbar\le a/16<\frac12$ and
$s_j\le\sbar$. Adding the three groups and using $\mu-1\le1$,
\[
A_jt\le\frac{2\Tz+(2m+5)\sigma'^{-1}\log(1/s_j)}{\mu-1}+\frac{(m+2)\mu\log(2\mu)}{\sigma'(\mu-1)^2}\le\bar\tau^\sharp_j .\qedhere
\]
\end{proof}

\section{Choice of Parameters}
\label{sec:param}

We choose the parameters of the construction so that the hypotheses of the
statements of the earlier sections that Section~\ref{sec:bootstrap} applies, in
particular the lemmas of Section~\ref{sec:relay}, hold at once. These statements
give their hypotheses explicitly, as inequalities between the parameters and the
constants of the earlier sections and of Part~I, and Section~\ref{sec:bootstrap}
assumes the choices made here when it applies them.

Recall that $\mu$ fixes the ratio of the amplitudes of consecutive rings and,
with $\sbar$, the ratios $s_k$ of their radii in Equation~\eqref{eq:datum}, and
that $\kseed$ is the seed size of Definition~\ref{def:approx}. The numbers
$\delta_0$ and $\Tpin$ are those of the clock of Definition~\ref{def:clock}. The choices are made in the
order recorded in Table~\ref{tab:param}, and Lemma~\ref{lem:nonempty} shows that
they can be made. Throughout this section we assume
\textup{(N1)}--\textup{(N3)}.

\begin{definition}
\label{def:P0-P4}
\emergencystretch=3em
The \textbf{structural list} consists of the following choices; the objects of
(SA1), (SA2), and the quantities that Definition~\ref{def:standing} derives from
them, are defined there; the list only fixes the order in which they are chosen.
\begin{description}[font=\normalfont]
\item[\textup{(P0.1)}] $m=3$; $\Phi$, $a$, $\Gone$ of Definition~I.\ref{I-def:ring}; $z$, $\sigma$,
  $\sigb$ of Definition~I.\ref{I-def:eigendata}.
\item[\textup{(P0.2)}] An integer $s$ and a number $\sigma'$ as in (SA1); then
  $c_\gamma$ of Definition~\ref{def:standing}.
\item[\textup{(P0.3)}] $p=m$; the integers $N_G$, $M$, the constants
  $\dH,C_2,C_3,\etmax$ and $\sigma_{\max}$, and the functions
  $\kappa_0(\cdot)$, $\bar t_1(\cdot)$, $C(\cdot)$ of Theorem~\ref{thm:active};
  $\kL$ and $\Cseed$ of Lemma~\ref{lem:seed}; $\|\ellst\|$, $\Cls$ and $\Cp$ of
  Lemma~I.\ref{I-lem:uniform}.
\item[\textup{(P0.4)}] The integers $s_h$, $N_Q$, $M_h$ and the constants
  $\majst,C_Q$ of (SA2).
\item[\textup{(P0.5)}] At order $M_h$: the constants $C_{L1}$, $C_M$, $C_O$,
  $\Cin$ of (SA2), and $C_\eta$ of Lemma~\ref{lem:field-bound}. Further $\Mvort$,
  $\Cstar$ of Lemma~\ref{lem:speed}; $C_G$ of Lemma~\ref{lem:inner-circ}; $\Cmom$
  of Proposition~\ref{prop:confine}.
\item[\textup{(P0.6)}] $c_N$, $C_N$, $\cout$ of Lemma~I.\ref{I-lem:uniform};
  $C_{\mathrm{qr}}$, $C_{\mathrm{id}}$, $C_{\mathrm{od}}$
  and the function $C_{\mathrm{tr}}(\cdot)$ of Lemma~\ref{lem:kicktrue}.
\item[\textup{(P1)}] Choose $\delta_0\in(0,\dH]$ sufficiently small that
  $\frac{16}{15}\cdot\frac{2\pi C_3\delta_0^2}{\sigma}\le\frac14\cdot\frac{15}{16}\cdot\frac{\cout\delta_0}{20|z|}$
  and $\Cp C_2\delta_0\le\frac{\cout}{16\pi}$.
\item[\textup{(P2)}] $\mu\in(1,\mu_{\max}]$, with
  $\mu_{\max}=1+\min\big(c'\delta_0,\frac1{16}\big)$ and
  a fixed $c'>0$ sufficiently small that
  $c'\le\frac{15}{16}\cdot\frac{\cout}{10|z|\,\Cmom\,\Ip{m}{\sigma/4}}$; then
  $\vartheta_\mu$, $\gamma_\mu$, $\gamma'_\mu$, $\CDuh$ of Definition~\ref{def:standing},
  and $\KO$ of Lemma~\ref{lem:O}.
\item[\textup{(P3)}] $\Tpin>0$ with $\rT(\Tpin)\le\frac1{64}$, with the function
  $\rT$ of Definition~\ref{def:standing}. Then $C_D(\Tpin)$
  (Proposition~\ref{prop:relay}); $c_K$, $K_1$, $\Abar$, $\Klo$, $\Kup$, $\Chand$
  (Definition~\ref{def:K1});
  $C_F\ge\Cin\,\mu\,(2+2\Cstar\Tpin)^m(1+2\Cstar)^m/\Klo$; $\Cdeep$
  (Lemma~\ref{lem:kicktrue}); $\Tz$ (Lemma~\ref{lem:time}) and
  choose $C_\tau\ge\max\{3,\frac{2m+5}{\sigma'},\frac{3(m+2)}{\sigma'}\}$.
\item[\textup{(P4)}] Choose $C_0\ge\max\{1,\,2\Cseed,\,\Chand,\,C_F\}$. Then $\kappa_0(C_0)$,
  $\bar t_1(C_0)$ and $C(C_0)$ (Theorem~\ref{thm:active}); $C_{\mathrm{tr}}(C_0)$
  (Lemma~\ref{lem:kicktrue}); $C_\Xi(C_0,K_1)$ (Proposition~\ref{prop:relay});
  $\Cop(C_0,K_1)$ (Lemma~\ref{lem:Oprime}).
\end{description}
\end{definition}

The number $C_0$ of (P4) is the constant of hypotheses (Init) and (Field) of
Theorem~\ref{thm:active}, which the proof of Proposition~\ref{prop:bootstrap}
applies to every active ring. That proof bounds the initial perturbation of the
seeded ring by $2\Cseed\ksh_n$, that of ring $k$ at its hand-off by
$\Chand\ksh_k$, and the inner field of the active ring $k$ by
$C_F\ksh_k\jap{u}^m$. Since (Init) and (Field) ask for bounds $C_0\kappa$ and
$C_0\kappa\jap{u}^p$, with $p=m$, the number $C_0$ must be at least $2\Cseed$,
$\Chand$ and $C_F$, and at least $1$ as Theorem~\ref{thm:active} requires. None of
$\Chand$, $C_F$, $\Cseed$ depends on $C_0$, so the choice of $C_0$ is not
circular. Likewise $\Tpin$ depends only on $\sigma$ and $m$. It is chosen before
$C_F$, whose formula in (P3) contains it. This order is not circular either: the
condition on $\Tpin$ in (P3) is the hypothesis $\rT(\Tpin)\le\frac1{64}$ of
Lemmas~\ref{lem:drift} and~\ref{lem:kicktrue}, and the drift bound of
Lemma~\ref{lem:drift} does not involve $C_F$.

\begin{definition}
\label{def:P5}
The \textbf{$\kappa$-list (P5)}: choose $\kseed>0$ so small that, for every
$\kappa\in(0,2\kseed]$, the following hold.
\begin{description}[font=\normalfont]
\item[\textup{($\kappa$1)}] $\kappa\le\delta_04^{-\mu/(\mu-1)}$.
\item[\textup{($\kappa$2)}] \begin{multline*}
  \frac{2\pi C_3\kappa^{1/2}\mu}{\frac34\sigma(\mu-1)}+C_{\mathrm{tr}}(C_0)\Big(\frac\kappa{\delta_0}\Big)^{\gamma_\mu}\log^{m+1}\frac{\delta_0}\kappa\\
  +\big[\Cdeep+C_{\mathrm{id}}\rT(\Tpin)\Ip{m}{\tfrac\sigma8}\big]\Big(\frac\kappa{\delta_0}\Big)^{\vartheta_\mu}\log^{m}\frac{\delta_0}\kappa
  \le\frac12\cdot\frac{15}{16}\cdot\frac{\cout}{20|z|(\mu-1)}.
  \end{multline*}
\item[\textup{($\kappa$3)}] $\Cpre(\kappa/\delta_0)^{\vartheta_\mu}\le\frac{15}{16}\cdot\frac{\cout}{20|z|(\mu-1)}$,
  where $\Cpre=C_{\mathrm{qr}}\big(\frac{4\mu C_{L1}C_Q}{3\sigma(\mu-1)}+1\big)K_1$.
\item[\textup{($\kappa$4)}] $\kappa\le\kappa_0(C_0)$.
\item[\textup{($\kappa$5)}] For every $G\in[0,\Gone]$, with $\hat T=\sigma(G)^{-1}\log(\delta_0/\kappa)$
  and $\rhohat$ the transient of Theorem~\ref{thm:active} for $(C_0,G)$:
  $\Cp C_2\big(\rhohat(\hat T)+\kappa^{1/2}\big)\le\frac{\cout}{16\pi}$.
\item[\textup{($\kappa$6)}] $\kappa\le\kL$.
\item[\textup{($\kappa$7)}] $C_\Xi(C_0,K_1)\big(1+\frac1{\sigb(\mu-1)}\big)(\kappa/\delta_0)^{\gamma'_\mu}\langle\log(\delta_0/\kappa)\rangle^{m+1}\le\frac1{\sigb(\mu-1)}$.
\end{description}
\end{definition}

In (P5), $\kappa$ stands for an amplitude $\ksh_k$ with which a ring starts its
active phase and, in ($\kappa$1) and ($\kappa$6), also for $2\kseed$.
Theorem~\ref{thm:active} is applied to ring $k$ with $\kappa=\ksh_k$, and
Lemma~\ref{lem:kicktrue}, at level $k$, with $\kappa=\ksh_{k+1}$.
Conditions ($\kappa$1)--($\kappa$3) are hypotheses (c)--(e) of
Lemma~\ref{lem:kicktrue}, ($\kappa$4) is hypothesis (Small) of
Theorem~\ref{thm:active}, ($\kappa$5) is used in the proof of
Proposition~\ref{prop:reduction}, and ($\kappa$6) makes the seed of
Definition~\ref{def:approx} admissible. Every amplitude $\ksh_k$ lies in
$(0,2\kseed]$ (Steps~0 and~3 of the proof of Proposition~\ref{prop:bootstrap}),
so the conditions are imposed on that whole interval.

Condition ($\kappa$7) is imposed after $C_0$ and $K_1$ are fixed. Here $K_1$
(Definition~\ref{def:K1}) is an explicit function of quantities fixed at
(P0)--(P3); it does not involve $C_0$, $\kseed$ or the deeper rings. Ring $j+1$
enters the hand-off of ring $j$ only through the quantity $\Xi_{j+1}$ of
Proposition~\ref{prop:relay}, and ($\kappa$7) at $\kappa=\ksh_{j+1}$ makes
$\Xi_{j+1}$ at most $\delta_0/(\sigb(\mu-1))$ at every level (Step~3 of the
proof of Proposition~\ref{prop:bootstrap}). So the bound on ring $j$ at its
hand-off holds with the same constants $K_1$ and $\Abar$ at every level.

The number $\sbar$, which fixes the ratios $s_k$ of consecutive radii, is chosen
last, because the conditions (B1)--(B11) on it involve constants fixed at all the
earlier steps; each of them holds for every
sufficiently small $\sbar$ (proof of Lemma~\ref{lem:nonempty}).

\begin{definition}
\label{def:P6}
The \textbf{$\sbar$-list (P6)}: choose $\sbar>0$ such that the following hold,
where $s_k=\sbar(2\mu)^{-k}$, $\eta_{b,1}=0$ and $\eta_{b,k}$, $\eta_b$ are as in
Equation~\eqref{eq:etab}.
\begin{description}[font=\normalfont]
\item[\textup{(B1)}] $\sbar\le a/16$ and $\sbar^{m-2}\le\frac12$.
\item[\textup{(B2)}] $\mu^4\sbar^2\le\frac18$.
\item[\textup{(B3)}] $C_G\mu\sbar^2\le\Gone$.
\item[\textup{(B4)}] $\sup_{0\le G\le C_G\mu\sbar^2}|\lamst(G)-z|\le\sigma(\mu-1)/(8\mu)$
  (this is (Det)).
\item[\textup{(B5)}] For every $k\ge1$: $\Lconf_ks_k\le a/8$ and
  $\frac83(\frac2a)^m(\Lconf_k)^{m-1}s_k^{m-2}\le\mu\Cstar$, where
  \[
  \Lconf_k=2+2\Cstar\mu\Big[\Tpin+1+\frac1{\sigma'}\Big((m+2)\log\frac1{s_k}+\log_+\frac{\delta_0}{\Klo}\Big)\Big].
  \]
\item[\textup{(B6)}] $\Kup\sbar^{m+2}\le\kseed$.
\item[\textup{(B7)}] $16C_OC_Q\CDuh\KO\sbar^{2m}\le1$, $\KO\sbar^{2m}\le\frac14$,
  $\KO\mu\sbar^{m+2}\le\frac12$, and
  \begin{equation}
  \label{eq:E22}
  \KO\sbar^{2m}\Big[\Tz+2\log\frac1{\sbar}+\frac1{\mu-1}\Big]\le\frac{\mu-1}{4C_OC_QC_\tau}\min\Big\{\frac12,\frac1{4\Cls K_1}\Big\}.
  \end{equation}
\item[\textup{(B8)}] $\eta_b\le\frac12\etmax$; for every $k\ge2$,
  \[
  \frac{\Cls\eta_{b,k}}{\sigma'}\Big[(m+2)\log\frac1{s_k}+\log_+\frac{\delta_0}{\Klo}\Big]\le\frac1{32};
  \]
  and $\Cls\eta_b\sigma'^{-1}\log(2\delta_0/\kseed)\le\frac1{32}$.
\item[\textup{(B9)}] $2\Abar\mu\sbar^{m+2}\le\majst$ and $C_\eta(\mu\sbar^2+\sbar^{m-2})\le1$.
\item[\textup{(B10)}] $4C_M\big(\frac{2C_Q}{\sigb(\mu-1)}+1\big)\mu\sbar^{m+2}\le\frac12$.
\item[\textup{(B11)}] $\|\ellst\|C_QK_1\Abar\mu\sbar^{m+2}\le\frac12$, and
  \[
  C_{\mathrm{od}}\eta_b\delta_0\frac{\mu^2}{\sigma^2(\mu-1)^2}+2\pi C_3\eta_b\frac{4\mu\delta_0}{3\sigma(\mu-1)}
  \le\frac14\cdot\frac{15}{16}\cdot\frac{\cout\delta_0}{20|z|(\mu-1)} .
  \]
\end{description}
\end{definition}

The second part of (B1), and (B2), hold automatically (see the remark after
Equation~\eqref{eq:sk-small}); they are listed because Proposition~\ref{prop:confine} and
Lemma~\ref{lem:field-bound} state them as hypotheses.

\begin{table}[!tb]
\centering
\small
\begin{tabular}{@{}p{0.40\textwidth}lp{0.42\textwidth}@{}}
\hline
\raggedright Quantity & Fixed at & \raggedright Determined by \tabularnewline
\hline
\raggedright $s$, $\sigma'$, $c_\gamma$ & (P0.2) & \raggedright $\sigb$, $\sigma$ ($s\ge3$ free) \tabularnewline
\raggedright $N_G$, $M$, $\dH$, $C_2$, $C_3$, $\etmax$, $\sigma_{\max}$, $\kappa_0(\cdot)$, $\bar t_1(\cdot)$, $C(\cdot)$ & (P0.3) & \raggedright $\Phi$, $\Gone$, $s$, $\sigma'$ \tabularnewline
\raggedright $\kL$; $\Cseed$; $\|\ellst\|$, $\Cls$, $\Cp$ & (P0.3) & \raggedright $\Phi$; $\Phi$, $\Gone$, $s$; $\Phi$, $\Gone$ \tabularnewline
\raggedright $s_h$, $N_Q$, $M_h$, $\majst$, $C_Q$ & (P0.4) & \raggedright $M$, $\Lam$, $\sigb$ \tabularnewline
\raggedright $C_{L1}$, $C_M$, $C_O$, $\Cin$, $C_\eta$ & (P0.5) & \raggedright $a$, $m$, $M_h$, $\|\Phi\|_{L^1}$ \tabularnewline
\raggedright $\Mvort$, $\Cstar$, $C_G$, $\Cmom$ & (P0.5) & \raggedright absolute; $\|\Phi\|_{L^1}$ \tabularnewline
\raggedright $c_N$, $C_N$, $\cout$, $C_{\mathrm{qr}}$, $C_{\mathrm{id}}$, $C_{\mathrm{od}}$, $C_{\mathrm{tr}}(\cdot)$ & (P0.6) & \raggedright $\Phi$, $\Gone$, $m$; $C_3$, $\sigma$, $C(\cdot)$ \tabularnewline
\raggedright $\delta_0$ & (P1) & \raggedright $\dH$, $C_2$, $C_3$, $\Cp$, $\cout$, $\sigma$, $|z|$ \tabularnewline
\raggedright $c'$, $\mu_{\max}$, $\mu$; $\vartheta_\mu$, $\gamma_\mu$, $\gamma'_\mu$, $\CDuh$, $\KO$ & (P2) & \raggedright $\delta_0$, $\cout$, $|z|$, $\Cmom$, $\sigma$; $c_\gamma$, $\sigb$, $C_M$, $C_Q$ \tabularnewline
\raggedright $\Tpin$ & (P3) & \raggedright $\sigma$, $m$ \tabularnewline
\raggedright $C_D(\Tpin)$, $c_K$, $K_1$, $\Abar$, $\Klo$, $\Kup$, $\Chand$ & (P3) & \raggedright (P0)--(P2), $\Tpin$ \tabularnewline
\raggedright $C_F$, $\Cdeep$, $\Tz$, $C_\tau$ & (P3) & \raggedright (P0)--(P2), $\Tpin$, $\Klo$ \tabularnewline
\raggedright $C_0$ & (P4) & \raggedright $\Cseed$, $\Chand$, $C_F$ \tabularnewline
\raggedright $\kappa_0(C_0)$, $\bar t_1(C_0)$, $C(C_0)$, $C_{\mathrm{tr}}(C_0)$, $C_\Xi(C_0,K_1)$, $\Cop(C_0,K_1)$ & (P4) & \raggedright $C_0$, $K_1$ and (P0) \tabularnewline
\raggedright $\kseed$ & (P5) & \raggedright (P0)--(P4) \tabularnewline
\raggedright $\sbar$; $\eta_b$, $\eta_{b,k}$, $\theta_j$ & (P6) & \raggedright (P0)--(P5); $\Cop$, $\sbar$, $\mu$, $C_O$, $C_Q$, $\KO$ \tabularnewline
\hline
\end{tabular}
\caption{Where each constant of the construction is fixed, and the quantities
that determine it. Every entry of the last column is fixed at an earlier row or
at the same step. The depth $n$, the amplitude $A_1$ and the solution
$\omega^{(n)}$ appear in no row; a level enters only as the index of
$\eta_{b,k}$ and $\theta_j$.}
\label{tab:param}
\end{table}

The last inequality of (B7), Equation~\eqref{eq:E22}, is used only through the
following lemma. Recall that $\theta_j$ is the factor of Lemma~\ref{lem:O} that
bounds the field of the outer rings on ring $j$ by the Duhamel majorant
$\maj_j$, and that $\bar\tau^\sharp_j$ of Lemma~\ref{lem:quiet-time} bounds the
length of the quiet phase of ring $j$ in its own time. The lemma bounds their
product, which must be at most $\frac12$ for hypothesis \textup{(H2)} of
Proposition~\ref{prop:relay} and for the corresponding hypothesis of
Lemma~\ref{lem:Oprime}, and at most $1/(4\Cls K_1)$ for hypothesis
\textup{(ii)} of Lemma~\ref{lem:handoff}.

\begin{lemma}
\label{lem:theta-quiet}
Assume (SA) and \textup{(B7)}, and let $K_1$ be as in Definition~\ref{def:K1},
$\Tz$ and $C_\tau$ as in \textup{(P3)}, $\theta_j$ as in Lemma~\ref{lem:O} and
$\bar\tau^\sharp_j$ as in Lemma~\ref{lem:quiet-time}. Then $\theta_j\bar\tau^\sharp_j\le\min\{\frac12,\frac1{4\Cls K_1}\}$
for every level $1\le j<n$.
\end{lemma}

\begin{proof}
For $j=1$, $\theta_1=0$ by Lemma~\ref{lem:O}, and the claim holds. Let
$2\le j<n$. By Lemma~\ref{lem:O}, $\theta_j=4C_OC_Q\KO s_{j-1}^{2m}$.

We first bound $\bar\tau^\sharp_j$ by a multiple of $C_\tau$. By (SA3),
$1<\mu\le\frac{17}{16}$, so $\mu\log(2\mu)\le\frac{17}{16}\log\frac{17}8<1$.
The three terms of $\bar\tau^\sharp_j$ (Lemma~\ref{lem:quiet-time}) are those in
$\Tz$, in $\log(1/s_j)$ and in $(\mu-1)^{-2}$; we compare them with
those of $C_\tau[\Tz+\log(1/s_j)+(\mu-1)^{-1}]/(\mu-1)$. By the definition of
$C_\tau$ in \textup{(P3)} and the bound on $\mu\log(2\mu)$, the coefficients
satisfy $3\le C_\tau$,
$(2m+5)/\sigma'\le C_\tau$ and $(m+2)\mu\log(2\mu)/\sigma'\le(m+2)/\sigma'\le C_\tau$;
this gives $\bar\tau^\sharp_j\le C_\tau\big[\Tz+\log(1/s_j)+(\mu-1)^{-1}\big]/(\mu-1)$.
Next we pass from $s_j$ to $s_{j-1}$. Since
$s_{j-1}\le\sbar\le a/16<\frac1{2\mu}$, we have
$\log\frac1{s_j}=\log\frac1{s_{j-1}}+\log(2\mu)\le2\log\frac1{s_{j-1}}$, so
\[
\theta_j\bar\tau^\sharp_j\le\frac{4C_OC_QC_\tau}{\mu-1}\,\KO\,g(s_{j-1}),\qquad
g(x)=x^{2m}\Big[\Tz+2\log\frac1x+\frac1{\mu-1}\Big].
\]
It remains to replace $s_{j-1}$ by $\sbar$, since Equation~\eqref{eq:E22}
bounds $\KO g(\sbar)$; for this we show that $g$ is nondecreasing near $0$. For $0<x<1$,
$g'(x)=x^{2m-1}\big[2m\big(\Tz+2\log\frac1x+\frac1{\mu-1}\big)-2\big]$, which is
nonnegative when $\log(1/x)\ge1/(2m)$, because $\Tz>0$ and $\mu>1$. So $g$ is
nondecreasing on $(0,e^{-1/(2m)}]$. Since $s_{j-1}\le\sbar\le a/16<\frac1{16}<e^{-1/(2m)}$,
$g(s_{j-1})\le g(\sbar)$, and Equation~\eqref{eq:E22} gives
\[
\theta_j\bar\tau^\sharp_j\le\frac{4C_OC_QC_\tau}{\mu-1}\KO g(\sbar)\le\min\{\frac12,\frac1{4\Cls K_1}\}.\qedhere
\]
\end{proof}

\begin{lemma}
\label{lem:nonempty}
Assume \textup{(N1)}--\textup{(N3)}. The choices (P0)--(P6) can be made in this
order, with all quantities finite and positive (except $\eta_{b,1}=\theta_1=0$),
and each quantity is determined by quantities fixed at earlier steps or earlier in the same step, as in
Table~\ref{tab:param}. None of the choices depends on $n$, $A_1$ or
$\omega^{(n)}$. Finally, if (P0)--(P6) hold, then for every integer $n\ge2$ and
every $A_1>0$ the standing assumptions (SA) hold with these parameters, and
(Det) holds.
\end{lemma}

\begin{proof}
We take the steps in order. At each step the quantities of the earlier steps are
fixed, and we show that the conditions of the step can be met.

\emph{(P0)--(P4).} The constants of (P0) exist and are finite by the results
cited there; in particular, by the formulas of (SA2) and $M\ge s+4$
(Theorem~\ref{thm:active}), $s_h=M-1\ge s+3\ge3$ and $N_Q\sigb\ge\Lam(s_h)$, as
Theorem~\ref{thm:quiet} requires. The interval $\big(\frac9{200},\min(\sigb,\frac34\sigma)\big)$ for $\sigma'$
in (SA1) is nonempty: by Proposition~I.\ref{I-prop:U}\textup{(a)}, $\sigb$ is attained and
$\sigb>\frac{211}{2500}>\frac9{200}$; and
$\frac34\sigma\ge\frac34\sigb>\frac{633}{10000}>\frac9{200}$. Moreover
$c_\gamma>0$, because $\sigma'<\frac34\sigma<\frac{15}{16}\sigma$. The conditions in (P1) hold for sufficiently small $\delta_0>0$ because
$\sigma\ge\sigb>0$, $\cout>0$ and the other constants have already been fixed.
In (P2), $c'>0$ because $\Cmom>0$ and $0<\Ip{m}{\sigma/4}<\infty$, so
$\mu_{\max}>1$. In (P3), $\Tpin$ exists because $\rT(T)=(1+T)^me^{-\sigma T/4}$
(Definition~\ref{def:standing}) tends to $0$ as $T\to\infty$, as $\sigma>0$. The formulas of (P3) and (P4) have positive finite denominators:
$c_K=3c_N/(32\mu)>0$, $\delta_0>0$, $\sigb>0$, $\mu-1>0$ and $\Klo>0$. So all the
constants of (P0)--(P4) are finite and positive, and $C_0\ge1$ and
$0<C_F\le C_0$, as Theorem~\ref{thm:active} and Lemma~\ref{lem:Oprime} require.
Each can be chosen in terms of the quantities in its row of Table~\ref{tab:param}, which are fixed at
earlier rows or earlier in the same step.

\emph{(P5).} Since (P5) is imposed for every $\kappa\in(0,2\kseed]$, it suffices
to find $x_0>0$ such that each condition holds for every $\kappa\in(0,x_0]$; then every $\kseed\in(0,x_0/2]$ is admissible. For ($\kappa$1), ($\kappa$4), ($\kappa$6) this holds
because $\delta_04^{-\mu/(\mu-1)}$, $\kappa_0(C_0)$, $\kL$ are positive. For
($\kappa$2), ($\kappa$3), ($\kappa$7) it suffices that each left side tends to
$0$ as $\kappa\to0$. The left sides are finite sums of terms $C\kappa^{1/2}$ and
$C(\kappa/\delta_0)^\gamma\langle\log(\delta_0/\kappa)\rangle^N$ with $C$ fixed at
(P0)--(P4), integers $N\ge0$, and
$\gamma\in\{\gamma_\mu,\vartheta_\mu,\gamma'_\mu\}$, positive because $\mu>1$ and
$c_\gamma>0$; each term tends to $0$. For ($\kappa$5) we need a bound uniform in
$G$. Let $\kappa<\delta_0$ and $G\in[0,\Gone]$, and put $c=\sigb-\sigma'>0$, where
$\sigma'<\sigb$ by (SA1). By Theorem~\ref{thm:active},
$t_1(C_0,G)\le\bar t_1(C_0)$, and $\sigma(G)\in[\sigb,\sigma_{\max}]$, so
$\hat T\ge\sigma_{\max}^{-1}\log(\delta_0/\kappa)$. If moreover
$\kappa\le\delta_0e^{-\sigma_{\max}\bar t_1(C_0)}$, then $\hat T\ge\bar t_1(C_0)$, so
$y=\hat T-t_1(C_0,G)\ge\hat T-\bar t_1(C_0)\ge0$. Hence $\hat T$ lies in the range
$u\ge t_1$ of the second formula for the transient in Theorem~\ref{thm:active},
and, as $\sigma(G)-\sigma'\ge c$,
\[
\rhohat(\hat T)=\jap{y}^pe^{-(\sigma(G)-\sigma')y}\le\jap{y}^pe^{-cy}\le K_ce^{-cy/2}
\le K_ce^{c\bar t_1(C_0)/2}\Big(\frac\kappa{\delta_0}\Big)^{c/(2\sigma_{\max})},
\]
where $K_c=\sup_{y\ge0}\jap{y}^pe^{-cy/2}<\infty$, and the last step uses
$y\ge\hat T-\bar t_1(C_0)$ and the lower bound on $\hat T$. The right side does not depend on
$G$ and tends to $0$ as $\kappa\to0$, and so does $\kappa^{1/2}$; so ($\kappa$5)
holds for all small $\kappa$, uniformly in $G\in[0,\Gone]$. We take for $x_0$
the least of the thresholds found for ($\kappa$1)--($\kappa$7).

\emph{(P6)}, with (P0)--(P5) fixed. It suffices that each condition holds for
every $\sbar$ in some interval $(0,\sbar_i]$, $\sbar_i>0$; then
every $\sbar\in(0,\min_i\sbar_i]$ is admissible. The conditions (B1)--(B3), (B6), (B7), (B9),
(B10), (B11), the first part of (B8) and its last part (with $\kseed$ fixed)
have left sides that are finite sums of terms $C\sbar^\alpha\log^N(1/\sbar)$ with
$\alpha>0$, $N\in\{0,1\}$ and $C$ fixed at (P0)--(P5); here we use
$\eta_b=\Cop\sbar^{2m}/(\sigb(\mu-1))$, and $m=3$, so that the exponent $m-2$
in (B1) and (B9) is positive. Each such term tends to $0$ as
$\sbar\to0$, while the right sides are fixed positive numbers. For (B4), the
supremum tends to $0$ as $\sbar\to0$, because $\lamst$ is continuous at $G=0$
(Proposition~I.\ref{I-prop:U}\textup{(a)}) and $\lamst(0)=z$; the right side
$\sigma(\mu-1)/(8\mu)$ is positive. The conditions (B5) and (B8) are imposed at
every level $k$, and we reduce each of them to one condition on $\sbar$. For
(B5), put
$\alpha_0=2+2\Cstar\mu\big[\Tpin+1+\sigma'^{-1}\log_+(\delta_0/\Klo)\big]$ and
$\beta_0=2\Cstar\mu(m+2)/\sigma'$, so that $\Lconf_k=\alpha_0+\beta_0\log(1/s_k)$.
The functions $x\mapsto x(\alpha_0+\beta_0\log(1/x))$ and
$x\mapsto x^{m-2}(\alpha_0+\beta_0\log(1/x))^{m-1}=x(\alpha_0+\beta_0\log(1/x))^2$
tend to $0$ as $x\to0$, so there is $\sbar_{\mathrm{B5}}>0$ such that both inequalities of
(B5) hold at level $k$ whenever $s_k\le\sbar_{\mathrm{B5}}$. Since $s_k\le\sbar$ for every
$k$, (B5) holds for every $\sbar\le\sbar_{\mathrm{B5}}$. For (B8) and $k\ge2$, we have
$s_k=s_{k-1}/(2\mu)$, and, by the formula for $\eta_{b,k}$ in
Equation~\eqref{eq:etab}, the $k$-th inequality reads $F(s_{k-1})\le\frac1{32}$
with
\[
F(x)=\frac{\Cls\Cop}{\sigma'\sigb(\mu-1)}\,x^{2m}\Big[(m+2)\Big(\log\frac1x+\log(2\mu)\Big)+\log_+\frac{\delta_0}{\Klo}\Big].
\]
Since $F(x)\to0$ as $x\to0$, there is $\sbar_{\mathrm{B8}}>0$ with $F\le\frac1{32}$ on
$(0,\sbar_{\mathrm{B8}}]$, and, since $s_{k-1}\le\sbar$ for $k\ge2$, the
inequalities hold for all $k\ge2$ once $\sbar\le\sbar_{\mathrm{B8}}$.

No formula or condition of (P0)--(P6) contains $n$, $A_1$ or $\omega^{(n)}$; the
levels in (B5) and (B8) range over all integers $k\ge1$ and $k\ge2$, not only
up to $n$.

\emph{(SA) and (Det).} Assume (P0)--(P6) and let $n\ge2$ and $A_1>0$. (SA1) holds by (P0.2) and (P0.3), and (SA2) by (P0.4) and (P0.5). For
(SA3): $\mu\in(1,\frac{17}{16}]$ because
$\mu_{\max}\le1+\frac1{16}$; $\sbar\in(0,a/16]$ and $C_G\mu\sbar^2\le\Gone$ by (B1) and (B3);
$\delta_0\in(0,\dH]$ by (P1); $\Tpin>0$ by (P3); and $\kseed\in(0,\kL]$ by
($\kappa$6) at $\kappa=2\kseed$. The objects of
Definitions~\ref{def:datum}--\ref{def:clock} are then defined with these
parameters: Definition~\ref{def:datum} needs $\mu\in(1,2]$ and $\sbar\in(0,a/16]$,
Definition~\ref{def:approx} needs $\kseed\in(0,\kL]$, and
Definition~\ref{def:clock} needs $\delta_0,\Tpin>0$ and $G_k\le\Gone$ for
$1\le k\le n$, which holds by Lemma~\ref{lem:inner-circ} and (B3). Hypothesis
(Det) is (B4).
\end{proof}

\section{The Joint Bootstrap and the Proof of the Reduction}
\label{sec:bootstrap}

We prove Proposition~\ref{prop:reduction}. The main step is
Proposition~\ref{prop:bootstrap}: for every number $n\ge2$ of rings the clock is
defined up to level $1$, ring $1$ completes its active phase at a time
$T_1\le\Tstar$, where $\Tstar$ does not depend on $n$, and at $T_1$ the material
inside ring $1$ lies in $B_{a/8}$. From this the proof of Proposition~\ref{prop:reduction} obtains a
lower bound, independent of $n$, on a non-radial moment of $\omega^{(n)}(T_1)$,
and applies Corollary~I.\ref{I-cor:defect}.

Throughout this section we assume \textup{(N1)}--\textup{(N3)} and
\textup{(P0)}--\textup{(P6)}, fix $A_1>0$ and an integer $n\ge2$, write
$\omega=\omega^{(n)}$, and use the objects of
Definitions~\ref{def:datum}--\ref{def:clock} with the parameters of
Section~\ref{sec:param}. By Lemma~\ref{lem:nonempty}, (SA) and (Det) hold, so
the results of Section~\ref{sec:relay} apply. We use $\hat a_k$ and $\rhohat_k$
as in Definition~\ref{def:standing}, with the constant $C_0$ of (P4). The
conditions ($\kappa$1)--($\kappa$7) are those of the $\kappa$-list (P5)
(Definition~\ref{def:P5}), and (B1)--(B11) those of the $\sbar$-list (P6)
(Definition~\ref{def:P6}).

Proposition~\ref{prop:bootstrap} is proved by a continuity argument in time
under a hypothesis (BA) bounding every ring by twice what its phase allows; the
lemmas of Section~\ref{sec:relay} return the bounds without the factor $2$. The
clock of Definition~\ref{def:clock} is defined level by level only while the amplitudes $\ksh_k$ lie in
$(0,\delta_0)$, which is one of the conclusions. So we run the argument up to the
time at which the clock would stop, and show at the end that it does not stop.

\begin{definition}
\label{def:stop}
If the clock is defined up to level $1$, put $\Tend=T_1$. Otherwise let
$k_{\mathrm s}$ be the level $k\in\{1,\dots,n\}$ such that the clock is defined up
to level $k+1$ but not up to level $k$ (for $k=n$ the first condition is void),
and put $\Tend=\tsh_{k_{\mathrm s}}$; we then say that the clock \textbf{stops} at
level $k_{\mathrm s}$.
\end{definition}

If the clock is defined up to level $k$, it is defined up to every level
$k'\ge k$; so these levels form a set $\{k_0,\dots,n\}$ or the empty set, and if
$1$ is not in it, there is exactly one level $k_{\mathrm s}$ as in
Definition~\ref{def:stop}: $k_{\mathrm s}=k_0-1$ if the set is
$\{k_0,\dots,n\}$, and $k_{\mathrm s}=n$ if it is empty. By Definition~\ref{def:clock}, for $j<n$, $\tsh_j$
is defined if and only if the clock is defined up to level $j+1$, and $T_j$ if
and only if it is defined up to level $j$. So, if the clock stops, $\tsh_j$ is
defined exactly for $j\ge k_{\mathrm s}$, $T_j$ exactly for $j>k_{\mathrm s}$,
and $S_{k_{\mathrm s}}(\tsh_{k_{\mathrm s}})\not\subset N_{k_{\mathrm s}}$ or
$\ksh_{k_{\mathrm s}}\notin(0,\delta_0)$. In either case $\Tend<\infty$, and by
Lemma~\ref{lem:clock}(1), $[0,\Tend]$ is the union of at most $2n-1$
consecutive \textbf{phase intervals} $\Phase_n=[0,T_n]$, $\Pin_{n-1}$,
$\Phase_{n-1},\dots$, ending with $\Phase_1$ if the clock does not stop, with
$\Pin_{k_{\mathrm s}}=[T_{k_{\mathrm s}+1},\tsh_{k_{\mathrm s}}]$ if it stops at
$k_{\mathrm s}<n$; if $k_{\mathrm s}=n$, then $\Tend=0$.

As in Definition~\ref{def:loc}, ring $j$ is \textbf{quiet} at $t\in[0,\Tend]$ if
$\tsh_j$ is not defined or $t\le\tsh_j$, \textbf{active} at $t$ if $T_j$ is
defined and $\tsh_j\le t\le T_j$, and \textbf{debris} at $t$ if $T_j$ is defined
and $t>T_j$. Every ring has at least one phase at every $t\in[0,\Tend]$: if $T_j$
is not defined, then $j\le k_{\mathrm s}$, and either $\tsh_j$ is not defined or
$j=k_{\mathrm s}$ and $t\le\Tend=\tsh_j$. By Lemma~\ref{lem:clock}(1), if ring
$j$ is quiet or active at $t$, then every ring $i<j$ is quiet at $t$.

\begin{definition}
For $t\in[0,\Tend]$, \textup{(BA)}$(t)$ means that for every $t'\in[0,t]$:
\begin{description}[font=\normalfont]
\item[\textup{(BA1)}] for every ring $j$ that is not debris at $t'$: $S_j(t')\subset\Nflat_j$;
\item[\textup{(BA2)}] for every ring $j<n$ that is quiet at $t'$: the field $h_j$ is
  defined at every time in $[0,t']$, and $\maj_j(A_jt')\le2\Abar\varepsilon_j$;
\item[\textup{(BA3)}] for every ring $k$ that is active at $t'$, with
  $u=A_k(t'-\tsh_k)$: $\|w_k(t')\|_{H^s}\le2C_2\hat a_k(u)(1+\rhohat_k(u))$.
\end{description}
\textup{(BA$^*$)}$(t)$ is the same statement with $\Nflat_j$, $2\Abar$ and $2C_2$
replaced by $\NN_j$, $\Abar$ and $C_2$. \textup{(BA$^\circ$)}$(t)$ is
\textup{(BA)}$(t)$ without, for each level $k$, the instance of \textup{(BA3)}
for ring $k$ at $t'=\tsh_k$.
\end{definition}

In (BA3), $w_k(t')$ is defined, since (BA1) gives
$S_k(t')\subset\Nflat_k\subset N_k$. Since $\NN\subset\operatorname{int}\Nflat$
(Lemma~\ref{lem:annuli}) and $\hat a_k>0$, (BA$^*$)$(t)$ implies (BA)$(t)$ with
room in each inequality. The instance of (BA3) for ring $k$ at $t'=\tsh_k$ is
left out of (BA$^\circ$) because ring $k$ becomes active only at $\tsh_k$, so
this bound cannot be obtained by continuity from bounds at earlier times. The
proof of Proposition~\ref{prop:bootstrap} proves it directly: for $k<n$ in
Step~3, from the hand-off of ring $k$ (Equation~\eqref{eq:boot-switch}), and for
$k=n$ in Step~0 (Equation~\eqref{eq:boot-start}).

\begin{proposition}
\label{prop:bootstrap}
Assume \textup{(N1)}--\textup{(N3)} and \textup{(P0)}--\textup{(P6)}. Then for
every $A_1>0$ and every $n\ge2$:
\begin{enumerate}[label=\textup{(\arabic*)}]
\item $\Tend=T_1$; that is, the clock is defined up to level $1$. Moreover
  $\ksh_1\le\Kup s_1^{m+2}$;
\item the conclusions of Theorem~\ref{thm:active} hold for ring $1$ on
  $\Phase_1$, in ring-$1$ units and active time, with $G=G_1$,
  $\kappa=\ksh_1$, the constant $C_0$ and $\eta_b=\eta_{b,1}$; and
  $|c_1(\hat T_1)|e^{\sigma_1\hat T_1}\ge\frac{15}{16}\delta_0$;
\item $D_2(T_1)\subset B_{a/8}$;
\item $T_1\le\Tstar$, with $\Tstar$ of Lemma~\ref{lem:time}.
\end{enumerate}
\end{proposition}

The next lemma collects the geometric consequences of (BA1) and of lower bounds
on the amplitudes $\ksh_k$: for each ring $j$ that is not debris, the material
inside ring $j$ stays in the small disk $B_{ar_j/8}$ and the rings outside it stay
outside $B_{6r_j}$, so that its external field is defined; and (Loc) holds. Its proof
applies Proposition~\ref{prop:confine} level by level. Step~1 of the proof of
Proposition~\ref{prop:bootstrap} applies the lemma on each phase interval.

\begin{lemma}
\label{lem:geometry}
Let $t\in[0,\Tend]$. Assume that \textup{(BA1)} holds at every $t'\in[0,t]$, and
that $\ksh_k\ge\Klo s_k^{m+2}$ for every level $k<n$ with $\tsh_k<t$. For a level
$k<n$ such that $T_{k+1}$ is defined, put $e_k=T_k$ if $T_k$ is defined and
$e_k=\tsh_k$ otherwise. Then:
\begin{enumerate}[label=\textup{(\alph*)}]
\item for every level $k<n$ such that $T_{k+1}$ is defined and $T_{k+1}\le t$, the
  hypotheses of Proposition~\ref{prop:confine} hold with $t_0=T_{k+1}$ and
  $t_1=\min(t,e_k)$; and for every $t'\in[T_{k+1},\min(t,e_k)]$,
  \[
  D_{k+1}(t')\subset\overline B_{\varrho_k(t')}\subset B_{ar_k/8},\qquad
  |q_{k+1}(t')|\le\Cmom\big(1+A_{k+1}(t'-T_{k+1})\big)^m,
  \]
  where $\varrho_k(t')=\big(2+2\Cstar A_{k+1}(t'-T_{k+1})\big)r_{k+1}$;
\item for every $t'\in[0,t]$ and every level $j$ that is not debris at $t'$, every
  ring $i<j$ is quiet at $t'$ and satisfies
  $S_i(t')\subset\Nflat_i\subset\R^2\setminus B_{6r_j}$, and
  $D_{j+1}(t')\subset B_{ar_j/8}$; hence the field $h_j$ is defined at $t'$, the
  support hypotheses of Lemma~\ref{lem:reduction} hold for ring $j$ in ring-$j$
  units at $t'$, and, if $j<n$, $\eta_j(A_jt')\le1$;
\item \textup{(Loc)} holds at every $t'\in[0,t]$.
\end{enumerate}
\end{lemma}

\begin{proof}
Recall that $\Nflat_i=r_i\Nflat$, with the annulus $\Nflat$ of
Section~\ref{sec:II-prelim}. By (BA1), a ring $i$ that is not debris at $t'\in[0,t]$ has
$S_i(t')\subset\Nflat_i\subset\{\frac58ar_i\le|x|\le\frac74r_i\}$; and
$s_k\le a/32$ for every $k$ by Equation~\eqref{eq:sk-small}.

\emph{Outer rings in \textup{(b)}.} Let $j$ be not debris at $t'\in[0,t]$. Every
ring $i<j$ is quiet at $t'$ (Lemma~\ref{lem:clock}(1)), so
$S_i(t')\subset\{|x|\ge\frac58ar_i\}$, and
$\frac58ar_i\ge\frac58ar_{j-1}=\frac{5a}{8s_{j-1}}r_j\ge20r_j$; hence
$S_i(t')\cap B_{6r_j}=\emptyset$.

The material inside ring $j$ is controlled from the inside out: the confinement
of $D_{j+1}$ after the time $T_{j+1}$, at which the active phase of ring $j+1$
ends, starts from the confinement of $D_{j+2}$ at $T_{j+1}$. So we argue by
downward induction on $j=n,\dots,1$ and show that $D_{j+1}(t')\subset
B_{ar_j/8}$ at every $t'\in[0,t]$ at which ring $j$ is not debris, and, if
$j<n$, that (a) holds for $k=j$. For $j=n$, $D_{n+1}=\emptyset$. Let $j<n$ and
assume the claim for $j+1$.

\emph{\textup{(a)} for $k=j$.} Let $T_{j+1}\le t$ be defined and
$t_1=\min(t,e_j)$; the clock is defined up to level $j+1$. We check the
hypotheses of Proposition~\ref{prop:confine}.
(i): ring $j+1$ is not debris at $T_{j+1}$, so
$S_{j+1}(T_{j+1})\subset\overline B_{7r_{j+1}/4}$, and
$D_{j+2}(T_{j+1})\subset B_{ar_{j+1}/8}$ by induction; so
$D_{j+1}(T_{j+1})\subset\overline B_{2r_{j+1}}$. (ii): for
$t''\in[T_{j+1},t_1]$ and $i\le j$, ring $i$ is not debris at $t''$, since
$t''\le e_j$ and $T_i\ge T_j$ whenever $T_i$ is defined
(Lemma~\ref{lem:clock}(1)); so $S_i(t'')\subset\{|x|\ge\frac a2r_i\}$. (iii) is
(B1) and (B2). (iv): we bound the kick time $A_{j+1}(t_1-T_{j+1})$ of ring $j$
at $t_1$. If $t_1\le\tsh_j$, then $A_{j+1}(t_1-T_{j+1})\le\Tpin$, as
$\tsh_j=T_{j+1}+\Tpin/A_{j+1}$ (Definition~\ref{def:clock}). If $t_1>\tsh_j$,
then $e_j=T_j$, so the clock is defined up to level $j$, and $\tsh_j<t$. So
$\ksh_j\ge\Klo s_j^{m+2}$ by hypothesis, and $\sigma_j\ge\sigma'$ by
Lemma~\ref{lem:detuning}. As $t_1\le T_j=\tsh_j+\hat T_j/A_j$ and
$A_{j+1}=\mu A_j$, $A_{j+1}(t_1-T_{j+1})\le\Tpin+\mu\hat T_j$. It remains to
bound $\hat T_j$, the length of the active phase of ring $j$ in active time: as
$\ksh_j<\delta_0$ (the clock is defined up to level $j$), $\Klo>0$ (Definition~\ref{def:K1}; $c_N>0$ by Lemma~I.\ref{I-lem:uniform}) and
$\sigma_j\ge\sigma'>0$ (SA1),
$\hat T_j=\sigma_j^{-1}\log(\delta_0/\ksh_j)\le\sigma'^{-1}\log\big(\delta_0/(\Klo s_j^{m+2})\big)\le\sigma'^{-1}\big((m+2)\log\frac1{s_j}+\log_+\frac{\delta_0}{\Klo}\big)$.
In both cases, as $\Tpin<\mu(\Tpin+1)$,
\[
A_{j+1}(t_1-T_{j+1})<\mu\Big[\Tpin+1+\frac1{\sigma'}\Big((m+2)\log\frac1{s_j}+\log_+\frac{\delta_0}{\Klo}\Big)\Big],
\]
so $\Lambda=2+2\Cstar A_{j+1}(t_1-T_{j+1})<\Lconf_j$, with $\Lconf_j$ of (B5).
The two inequalities of (iv) increase with $\Lambda$, so they follow from (B5).
Proposition~\ref{prop:confine}(a), (b) give, for $t'\in[T_{j+1},t_1]$,
$D_{j+1}(t')\subset\overline B_{\varrho_j(t')}$ and the bound on $q_{j+1}(t')$,
and $\varrho_j(t')\le\Lambda r_{j+1}<\Lconf_js_jr_j\le\frac a8r_j$ by (B5).

\emph{The claim on $D_{j+1}$.} Let ring $j$ be not debris at $t'\in[0,t]$. If
ring $j+1$ is not debris at $t'$, then $S_{j+1}(t')\subset\overline
B_{7r_{j+1}/4}$, $D_{j+2}(t')\subset B_{ar_{j+1}/8}$ by induction, and
$2r_{j+1}=2s_jr_j\le\frac a{16}r_j$; so $D_{j+1}(t')\subset B_{ar_j/8}$.
Otherwise $T_{j+1}<t'\le t$, and $t'\le e_j$ (if $T_j$ is defined, this is
$t'\le T_j$; otherwise $j=k_{\mathrm s}$ and $t'\le\Tend=\tsh_j$); so (a) for
$k=j$ gives the claim. This completes the induction.

\emph{End of \textup{(b)}.} Let $j$ be not debris at $t'\in[0,t]$. Then
$D_{j+1}(t')\subset B_{ar_j/8}$ and the rings $i<j$ do not meet $B_{6r_j}$, so
$h_j$ is defined at $t'$ (Definition~\ref{def:regions}), and in ring-$j$ units
the inner part is supported in $D_{j+1}(t')/r_j\subset B_{a/8}$ and the outer part in
$\R^2\setminus B_6$, as Lemma~\ref{lem:reduction} requires. If $j<n$, the
hypotheses of Lemma~\ref{lem:field-bound} hold at $t'$, since (B1) holds and
$S_i(t')\subset\Nflat_i\subset N_i$ for $i<j$; so $\eta_j(A_jt')\le
C_\eta(\mu\sbar^2+\sbar^{m-2})\le1$ by (B9).

\emph{\textup{(c)}.} Item (1) of Definition~\ref{def:loc} holds because a ring
$i$ with $t'\le T_i$, or with $T_i$ not defined, is not debris at $t'$, so
$S_i(t')\subset\Nflat_i\subset N_i$. Item (2) concerns a ring $j$ quiet at
$t'$. Its first two conditions, on the rings $i<j$ and on $D_{j+1}(t')$, hold
by (b). Its last clause, the confinement $D_{j+1}(t')\subset\overline
B_{\varrho_j(t')}$ when $T_{j+1}$ is defined and $t'\ge T_{j+1}$, holds by (a)
with $k=j<n$, since a ring $j$ quiet at $t'$ is not debris, so $t'\le e_j$
as above; for $j=n$ it holds as $D_{n+1}=\emptyset$.
\end{proof}

\begin{proof}[Proof of Proposition~\ref{prop:bootstrap}]
Throughout, (SA) and (Det) hold by
Lemma~\ref{lem:nonempty}, and the phases are ordered as in
Lemma~\ref{lem:clock}\textup{(1)}. Step~0 checks the bootstrap assumption at
$t=0$, Steps~1--6 deduce (BA$^*$)$(t)$ from (BA$^\circ$)$(t)$, and Step~7 is the
continuity argument.

\emph{Step 0 (the initial time; it does not use \textup{(BA)}).} At $t=0$, $\tsh_n=0$, $G_n=0$ and, in
ring-$n$ units, $w_n(0)=\Phi^{\kseed}-\Phi$. By ($\kappa$6) at $\kappa=2\kseed$,
$\kseed\le\kL$, so Lemma~\ref{lem:seed} and (P4) give
$S_n(0)=S_n\subset\NNm_n\subset N_n$,
$\ksh_n=|\ellst(0)(\Phi^{\kseed}-\Phi)|\in[\frac12\kseed,2\kseed]$ and
$\|w_n(0)\|_{H^{M-1}}\le\Cseed\kseed\le2\Cseed\ksh_n\le C_0\ksh_n$. By
($\kappa$1) at $\kappa=2\kseed$, $2\kseed<\delta_0$. So $\ksh_n\in(0,\delta_0)$,
the clock is defined up to level $n$, and $k_{\mathrm s}<n$ if it stops.
Recall that $\rhohat_n$ is the transient of Theorem~\ref{thm:active} for
$(C_0,G_n)$; it equals $C_0\jap{u}^pe^{-(\sigma_n-\sigma')u}$ for
$u<t_1(C_0,G_n)$, and $t_1(C_0,G_n)>0$ because $2^{p/2}C_0>1$, so
$\rhohat_n(0)=C_0$. As $s\le M-1$ and $C_2\ge1$,
\begin{equation}
\label{eq:boot-start}
\|w_n(0)\|_{H^s}\le C_0\ksh_n<C_2\ksh_n(1+C_0)=C_2\hat a_n(0)\big(1+\rhohat_n(0)\big).
\end{equation}
For $j<n$, $S_j(0)=S_j\subset\{ar_j\le|x|\le r_j\}\subset r_j\NNm$ and
$\maj_j(0)=0$; $h_j$ is defined at $0$, since
$D_{j+1}(0)\subset\overline B_{5r_{j+1}/4}\subset B_{ar_j/8}$ and the rings $i<j$
lie in $\{|x|\ge ar_i\}\subset\R^2\setminus B_{6r_j}$. Hence (BA$^*$)$(0)$ holds.

Now let $t\in[0,\Tend]$ and assume (BA$^\circ$)$(t)$. We prove the conclusions of
Steps~1--6 for the times in $[0,t]$, treating the phase intervals of $[0,t]$ in
chronological order: Steps~1, 2, 4, 5, 6 on $\Phase_n\cap[0,t]$; then, for
$k=n-1,n-2,\dots$, Steps~1, 2 on $\Pin_k\cap[0,t]$, Step~3 at $\tsh_k$ if
$\tsh_k\le t$, Steps~5, 6 on $\Pin_k\cap[0,t]$, and Steps~1, 2, 4, 5, 6 on
$\Phase_k\cap[0,t]$. Each step uses only (BA$^\circ$)$(t)$ and the steps already
carried out. In particular, Step~3 at $\tsh_k$ uses, besides Step~0, only
Steps~1, 2 on $[0,\tsh_k]$, Step~4 on $[0,T_{k+1}]$ and Step~3 at the levels
$i>k$; and the bound $\ksh_k\ge\Klo s_k^{m+2}$, which Lemma~\ref{lem:quiet-time}
in Step~5 at the time $\tsh_k$ and Step~1 on $\Phase_k$ need, comes from Step~3
at $\tsh_k$.

\emph{Step 1 (geometry).} Let $I$ be the current phase interval and $\bar t\le t$
the right end of $I\cap[0,t]$. Lemma~\ref{lem:geometry} applies with $\bar t$
in place of $t$: (BA1) holds on $[0,\bar t]$, and every level $k<n$ with
$\tsh_k<\bar t$ has $\tsh_k$ at most the left end of $I$, so Step~3 at $\tsh_k$
has already been carried out and has given $\ksh_k\ge\Klo s_k^{m+2}$. Below we use Lemma~\ref{lem:geometry}(a)--(c)
on $[0,\bar t]$ without further mention.

\emph{Step 2 (quiet rings).} Let $t'$ be in the current interval and $j<n$ quiet
at $t'$. If $t'=0$, Equation~\eqref{eq:boot-quiet} below holds at $t''=0$ by
Step~0, as $\maj_j(0)=0$ and $w_j(0)=0$ (ring $j$ starts at $\Phi$ in ring-$j$
units); let $t'>0$ and $\tau=A_jt'$. Suppose first that $\tsh_j$ is defined;
then the clock is defined up to level $j+1$ and $\tau\le\tash_j$. We apply
Lemma~\ref{lem:O} to $j$ and $\tau$: at every $t''\in[0,t']$ every ring $i\le j$
is quiet, so Lemma~\ref{lem:geometry}(b) gives its geometric hypotheses and
$\eta_i(A_it'')\le1$; (BA2) and (B9) give, as $\varepsilon_j=\mu s_j^{m+2}$ and
$s_j\le\sbar$,
$\maj_j\le2\Abar\varepsilon_j\le2\Abar\mu\sbar^{m+2}\le\majst$ on $[0,\tau]$; and
(B7) gives the three smallness conditions on $\sbar$ in Lemma~\ref{lem:O}. So, for $t''\in[0,t']$ and $i<j$,
$\maj_i(A_it'')\le\frac34\Ptil_{i,j}\maj_j(A_jt'')$, the hypotheses of
Theorem~\ref{thm:quiet} hold for ring $i$ on $[0,A_it']$ (for $i=j$ as well), and
$\etout_j\le\theta_j\maj_j$ on $[0,\tau]$. If $\tsh_j$ is not defined, the clock stops,
$j<k_{\mathrm s}$, ring $k_{\mathrm s}$ is quiet on $[0,\Tend]$, and the case just
treated, applied to $k_{\mathrm s}$ at $t'$, shows that Theorem~\ref{thm:quiet}
applies to ring $j$ on $[0,A_jt']$ and
$\maj_j(A_jt')\le\frac34\Ptil_{j,k_{\mathrm s}}\maj_{k_{\mathrm s}}(A_{k_{\mathrm s}}t')$.
In both cases Theorem~\ref{thm:quiet} gives, for $t''\in[0,t']$,
\begin{equation}
\label{eq:boot-quiet}
S_j(t'')\subset\NNm_j,\qquad \|w_j(t'')\|_{H^{s_h}}\le C_Q\maj_j(A_jt'').
\end{equation}

\emph{Step 3 (hand-off of ring $k<n$ at $\tsh_k$, if $\tsh_k\le t$).} This step
starts ring $k$: it shows that the clock is defined up to level $k$, with
$\ksh_k$ between $\Klo s_k^{m+2}$ and $\Kup s_k^{m+2}$, and proves (BA3) for ring
$k$ at $\tsh_k$. Lemma~\ref{lem:kicktrue} bounds the kick amplitude
$|\Ain(\Tpin)|$ that ring $k+1$ and the rings inside it give to ring $k$
(Lemma~\ref{lem:coupling}\textup{(2)}),
Proposition~\ref{prop:relay} bounds the Duhamel majorant $\maj_k(\tash_k)$ from
above, and Lemma~\ref{lem:handoff} combines the two. The clock is
defined up to level $k+1$, Steps~1, 2 have been carried out on $[0,\tsh_k]$ and
Steps~4, 5, 6 on $[0,T_{k+1}]$, and no ring is active on $(T_{k+1},\tsh_k)$; so
Step~4 has been carried out on the whole of $\Phase_{k+1}$ and, if $k+1<n$,
Step~3 for ring $k+1$. By Step~3 for ring $k+1$ and
(B6), or by Step~0 if $k+1=n$,
\begin{equation}
\label{eq:boot-kappa}
0<\ksh_{k+1}\le2\kseed .
\end{equation}
We apply Lemma~\ref{lem:kicktrue} with the constants $C_0$ and $K_1$, and check
its hypotheses in the order in which it lists them. First,
$D_{k+1}\subset B_{ar_k/8}$ on $[0,\tsh_k]$ by
Lemma~\ref{lem:geometry}(b). Second, the conclusions of Theorem~\ref{thm:active} and
Corollary~\ref{cor:active-moment} for ring $k+1$ on $\Phase_{k+1}$ hold by
Step~4 with $\eta_{b,k+1}$, hence with $\eta_b\ge\eta_{b,k+1}$, as the bounds
increase with $\eta_b$. Third, Step~4 for ring $k+1$ verified hypotheses (i)--(iv) of
Lemma~\ref{lem:drift} for ring $k+1$ with $\bar u=\hat T_{k+1}$ and
$\eta=\eta_{b,k+1}\in[0,\eta_b]$, together with $\rT(\Tpin)\le\frac1{64}$ and
$\Cls\eta_{b,k+1}\hat T_{k+1}\le\frac1{32}$. Fourth, if $k+1<n$: on $[0,\tsh_{k+1}]$
$D_{k+2}\subset B_{ar_{k+1}/8}$ and the rings $i\le k$ lie outside $B_{6r_{k+1}}$
by Lemma~\ref{lem:geometry}(b); the hypotheses of Theorem~\ref{thm:quiet} hold for
ring $k+1$ on $[0,\tash_{k+1}]$ by Step~2; $\maj_{k+1}(\tash_{k+1})\le K_1\ksh_{k+1}$ is
Equation~\eqref{eq:handoff} for ring $k+1$; the hypotheses of
Proposition~\ref{prop:confine} with $k+1$ in place of $k$ and $t_1=T_{k+1}$ hold
by Lemma~\ref{lem:geometry}(a) for $k+1$; and
$\ksh_{k+1}\ge\Klo s_{k+1}^{m+2}$ by Step~3 for ring $k+1$. Fifth, the bound
$|q_{k+1}|\le\Cmom(1+\tauK)^m$ for $\tauK\in[0,\Tpin]$ is
Lemma~\ref{lem:geometry}(a) on $\Pin_k$. Finally, conditions (a)--(f) of
Lemma~\ref{lem:kicktrue} are (P1), (P2), ($\kappa$1), ($\kappa$2), ($\kappa$3) at
$\kappa=\ksh_{k+1}$, which Equation~\eqref{eq:boot-kappa} allows, and (B11). So
$|\qs|\in[\frac{15}{16},\frac{17}{16}]2\pi|\pout_*(G_{k+1})|\delta_0$ and
\begin{equation}
\label{eq:boot-kick}
\frac{|\beta_k|\,|\qs|}{10|z|(\mu-1)}\le|\Ain(\Tpin)|\le\frac{3|\beta_k|\,|\qs|\,\mu}{|z|(\mu-1)} .
\end{equation}

Next we apply Proposition~\ref{prop:relay} with the constants $C_0$ and $K_1$,
$j=k$ and $\tau=\tash_k$. By
Lemma~\ref{lem:quiet-time} at the time $\tsh_k$ (its hypothesis
$\Kup\sbar^{m+2}\le\kseed$ is (B6); its hypotheses on the levels
$i\in(k,n)$ hold by Step~3 for these levels, and $\ksh_n\ge\frac12\kseed$ by
Step~0) and Lemma~\ref{lem:theta-quiet},
\begin{equation}
\label{eq:boot-theta}
\theta_k\tash_k\le\theta_k\bar\tau^\sharp_k\le\min\Big\{\frac12,\frac1{4\Cls K_1}\Big\}.
\end{equation}
The hypotheses (H1)--(H4) of Proposition~\ref{prop:relay} hold. (H1): (Loc) holds on
$[0,\tsh_k]$ by Lemma~\ref{lem:geometry}(c). (H2): $\etout_k\le\theta_k\maj_k$ on
$[0,\tash_k]$ by Step~2, and $\theta_k\tash_k\le\frac12$ by
Equation~\eqref{eq:boot-theta}. (H3): if $k+1<n$, the
conclusions of Theorem~\ref{thm:quiet} for ring $k+1$ hold on
$[0,\tash_{k+1}]$ by Step~2. (H4): $\maj_{k+1}(\tash_{k+1})\le K_1\ksh_{k+1}$ by
Step~3 for ring $k+1$; if $k+1=n$, then $\tash_n=0$ and $\maj_n(0)=0$. The
other conditions of (H4), on $\Phase_{k+1}$, hold: $D_{k+2}\subset B_{ar_{k+1}/8}$ by Lemma~\ref{lem:geometry}(b);
(BA3) for ring $k+1$ holds by (BA$^\circ$)$(t)$, and at $\tsh_{k+1}$ by
Equation~\eqref{eq:boot-switch} for ring $k+1$ or by
Equation~\eqref{eq:boot-start}; and
$\|\hin_{k+1}\|_{C^{M_h}(N')}\le C_0\ksh_{k+1}\jap{u}^m$ by Step~4 for ring
$k+1$. Hence, by Proposition~\ref{prop:relay}(ii), Equation~\eqref{eq:relay}
holds at $\tash_k$.

We bound the right side of Equation~\eqref{eq:relay}. Its last term
$\Xi_{k+1}$ collects the quiet phase of ring $k+1$ and the remaining terms of its
active phase. We bound it by ($\kappa$7), after bounding the exponent
$1-\gamma_k$ from below. By
Lemma~\ref{lem:detuning} and its proof, $(1-e_\mu)\sigma\le\sigma_i\le(1+e_\mu)\sigma$
for $i\in\{k,k+1\}$, $e_\mu(\mu+1)\le\frac14(\mu-1)$, $\frac{15}{16}\le1-e_\mu\le1$
and $x_k\le\frac{17}{16}\sigma$. Hence
\begin{gather*}
1-\frac{\sigma_k}{\mu\sigma_{k+1}}\ge1-\frac{1+e_\mu}{\mu(1-e_\mu)}
=\frac{(\mu-1)-e_\mu(\mu+1)}{\mu(1-e_\mu)}\ge\frac34\cdot\frac{\mu-1}\mu,\\
1-\frac{\sigma'}{\sigma_{k+1}}\ge1-\frac{16\sigma'}{15\sigma}=c_\gamma ,
\end{gather*}
and, as $\gamma_k$ is the larger of $\sigma_k/(\mu\sigma_{k+1})$ and
$\sigma'/\sigma_{k+1}$, $1-\gamma_k\ge\min(\frac34\frac{\mu-1}\mu,c_\gamma)\ge\gamma'_\mu$.
As $\ksh_{k+1}<\delta_0$,
$(\ksh_{k+1}/\delta_0)^{1-\gamma_k}\le(\ksh_{k+1}/\delta_0)^{\gamma'_\mu}$, and
($\kappa$7) at $\kappa=\ksh_{k+1}$ gives $\Xi_{k+1}\le\delta_0/(\sigb(\mu-1))$. So
the bracket in Equation~\eqref{eq:relay} with $j=k$ is at most
$(4C_2+1)\delta_0/(\sigb(\mu-1))+C_D(\Tpin)$; with Equation~\eqref{eq:relay}
this gives the first inequality below, and the bound on $x_k$ above with Equation~\eqref{eq:K1} gives the second:
\[
\maj_k(\tash_k)\le2C_M\varepsilon_ke^{x_k\Tpin}\Big[\frac{(4C_2+1)\delta_0}{\sigb(\mu-1)}+C_D(\Tpin)\Big]\le\Abar\,\varepsilon_k .
\]
Next we bound from below the coefficient $\cin_k$ produced by the kick. By
Lemma~\ref{lem:coupling}\textup{(2)}, $\cin_k=e^{\zeta_0\Tpin}\Ain(\Tpin)$ with
$\Re\zeta_0=\sigma_k/\mu=x_k$, so that $|\cin_k|=e^{x_k\Tpin}|\Ain(\Tpin)|$. As
$G_k,G_{k+1}\in[0,\Gone]$, Lemma~I.\ref{I-lem:uniform} gives the first bounds below,
and they give the second with $\beta_k=\beta(G_k)s_k^{m+2}/(2\pi)$ and the bound
on $|\qs|$:
\[
c_N\le|\beta(G_k)||\pout_*(G_{k+1})|\le C_N,\qquad
\tfrac{15}{16}c_Ns_k^{m+2}\delta_0\le|\beta_k||\qs|\le\tfrac{17}{16}C_Ns_k^{m+2}\delta_0 .
\]
With the lower bound in Equation~\eqref{eq:boot-kick}, $\varepsilon_k=\mu s_k^{m+2}$
and $\frac{15}{16}\cdot\frac1{10}=\frac3{32}$,
\[
|\cin_k|\ge e^{x_k\Tpin}\frac{|\beta_k||\qs|}{10|z|(\mu-1)}\ge\frac{3c_N}{32\mu}\cdot\frac{\varepsilon_k\delta_0e^{x_k\Tpin}}{|z|(\mu-1)}
=\frac{c_K\varepsilon_k\delta_0e^{x_k\Tpin}}{|z|(\mu-1)} .
\]
Multiplying by $\frac12K_1$, Equation~\eqref{eq:K1} and the first inequality of
the bound on $\maj_k(\tash_k)$ above give
\[
\frac12K_1|\cin_k|\ge4C_M\varepsilon_ke^{x_k\Tpin}\Big[\frac{(4C_2+1)\delta_0}{\sigb(\mu-1)}+C_D(\Tpin)\Big]\ge2\maj_k(\tash_k).
\]

The hypotheses of Lemma~\ref{lem:handoff} hold: (i) by
Lemma~\ref{lem:geometry}(b) and Step~2; (ii), with the $\theta_k$ of
Lemma~\ref{lem:O}, by Step~2 and Equation~\eqref{eq:boot-theta}; (iii) by the
last inequality; (iv) by $\maj_k(\tash_k)\le\Abar\varepsilon_k\le\Abar\mu\sbar^{m+2}$
and (B11). So Equation~\eqref{eq:handoff} holds at $\tsh_k$. To bound $\ksh_k$
we use its first bound, the identity
$|\cin_k|=e^{x_k\Tpin}|\Ain(\Tpin)|$, Equation~\eqref{eq:boot-kick} and the bounds
on $|\beta_k||\qs|$. As $e^{x_k\Tpin}\ge1$ and $\frac{15}{320}\ge\frac1{22}$,
$\ksh_k\ge\frac12|\Ain(\Tpin)|\ge\frac12\cdot\frac1{10}\cdot\frac{15}{16}\cdot\frac{c_N\delta_0s_k^{m+2}}{|z|(\mu-1)}\ge\Klo s_k^{m+2}$;
and as $x_k\le\frac{17}{16}\sigma$ and
$\frac32\cdot3\cdot\frac{17}{16}=\frac{153}{32}\le5$,
\[
\ksh_k\le\frac32e^{x_k\Tpin}|\Ain(\Tpin)|\le\frac32\cdot3\cdot\frac{17}{16}\cdot\frac{C_N\delta_0\mu\,s_k^{m+2}}{|z|(\mu-1)}e^{17\sigma\Tpin/16}\le\Kup s_k^{m+2}.
\]
Hence
$0<\ksh_k\le\Kup\sbar^{m+2}\le\kseed<\delta_0$ by (B6) and ($\kappa$1), and
$S_k(\tsh_k)\subset\NNm_k\subset N_k$ by Equation~\eqref{eq:boot-quiet}; so the
clock is defined up to level $k$. Finally, as $s\le s_h=M-1$, $C_0\ge\Chand$,
$C_2\ge1$ and $\rhohat_k(0)=C_0$,
\begin{equation}
\label{eq:boot-switch}
\|w_k(\tsh_k)\|_{H^s}\le\Chand\ksh_k\le C_0\ksh_k<C_2\hat a_k(0)\big(1+\rhohat_k(0)\big),
\end{equation}
which is (BA3) for ring $k$ at $t'=\tsh_k$, with a margin. Step~3 uses neither
(BA3) for ring $k$ at $\tsh_k$ nor any time after $\tsh_k$.

\emph{Step 4 (the active phase of ring $k$).} This step verifies the hypotheses
(Init), (Field), (Coef) and (Small) of Theorem~\ref{thm:active} for ring $k$ on
$\Phase_k\cap[0,t]$. Let $t'\in\Phase_k\cap[0,t]$ and
$\bar u=A_k(t'-\tsh_k)\le\hat T_k$. If $t'=\tsh_k$, we need only the two bounds
at the end of this step for $t''=\tsh_k$, and they hold already: $S_k(\tsh_k)\subset\NNm_k\subset\NN_k$ by Step~0
if $k=n$ and by Equation~\eqref{eq:handoff} if $k<n$, and the bound on
$\|w_k(\tsh_k)\|_{H^s}$ is Equation~\eqref{eq:boot-start}, respectively
\eqref{eq:boot-switch}. So let $t'>\tsh_k$; then $\bar u>0$, and $T_k$ is defined
by Step~0 or Step~3. By the invariance of Equation~\eqref{eq:euler} under
translations in time and the scaling of ring-$k$ units, and by
Theorem~\ref{thm:smooth}(5), the solution in ring-$k$ units with time origin
$\tsh_k$ is the solution of Theorem~\ref{thm:smooth} with datum
$A_k^{-1}\omega(\tsh_k,r_k\cdot)$. We apply Lemma~\ref{lem:reduction} and
Theorem~\ref{thm:active} to it, with the parts of the vorticity on $S_k$, on
$D_{k+1}$ and on $S_1\cup\dots\cup S_{k-1}$, with $G=G_k$, $\kappa=\ksh_k$, the
constant $C_0$, $\eta_b=\eta_{b,k}$ and $\delta_0$; then
$w(u)=w_k(\tsh_k+u/A_k)$, $h(u)=h_k(\tash_k+u)$ and $c=c_k$. The support
hypotheses of Lemma~\ref{lem:reduction} on $[0,\bar u]$ hold by
Lemma~\ref{lem:geometry}(b). For Theorem~\ref{thm:active}: $G_k\in[0,\Gone]$,
$\delta_0\le\dH$ by (P1), $\eta_{b,k}\le\eta_b\le\etmax$ by (B8), and (Small)
holds: $\ksh_k\le2\kseed\le\kappa_0(C_0)$ by Step~0 if $k=n$, by Step~3 and (B6) if
$k<n$, and by ($\kappa$4). It remains to check (Init), (Field) and (Coef).

(Init): for $k<n$, Equation~\eqref{eq:handoff} gives
$\|w_k(\tsh_k)\|_{H^{M-1}}\le\Chand\ksh_k\le C_0\ksh_k$ and
$\supp(\Phi+w_k(\tsh_k))\subset\NNm$; for $k=n$ this is Step~0.

(Field), inner part $h_1(u)=\hin_k(\tash_k+u)$. For $k=n$, $h_1=0$, because
$D_{n+1}=\emptyset$ and $G_n=0$. Let $k<n$, $u\in[0,\bar u]$ and
$t''=\tsh_k+u/A_k$. By Lemma~\ref{lem:geometry}(a), $\omega_{\ge k+1}(t'')$ is
supported in $\overline B_R$ with $R=\varrho_k(t'')\le\frac a8r_k$,
$A_{k+1}(t''-T_{k+1})=\Tpin+\mu u$, and $D_{k+1}(t'')\subset B_{ar_k/8}$, so
$h_1(u)$ is the field $E_{[k]}$ of Lemma~\ref{lem:inner-field} at $t''$; and
$\|\omega_{\ge k+1}(t'')\|_{L^1}\le\Mvort A_{k+1}r_{k+1}^2$ by
Lemma~\ref{lem:speed}. Lemma~\ref{lem:inner-field}, the inequalities
$2+2\Cstar\Tpin+2\Cstar\mu u\le(2+2\Cstar\Tpin)(1+\Cstar\mu u)$ and
$1+\Cstar\mu u\le(1+\Cstar\mu)\max(1,u)\le(1+2\Cstar)\jap{u}$, the definition of
$C_F$ in (P3), and $\ksh_k\ge\Klo s_k^{m+2}$ (Step~3) give
\begin{equation}
\label{eq:boot-inner}
\begin{aligned}
\|h_1(u)\|_{C^{M_h}(N')}&\le\Cin\mu s_k^2(R/r_k)^m=\Cin\mu s_k^{m+2}\big(2+2\Cstar(\Tpin+\mu u)\big)^m\\
&\le\Cin\mu(2+2\Cstar\Tpin)^m(1+2\Cstar)^ms_k^{m+2}\jap{u}^m\\
&=C_F\Klo s_k^{m+2}\jap{u}^m\le C_F\ksh_k\jap{u}^m .
\end{aligned}
\end{equation}
As $M\le M_h$ and $C_F\le C_0$, this is the first bound of (Field).

(Field), outer part $h_2(u)=\hout_k(\tash_k+u)$. We apply
Lemma~\ref{lem:Oprime} with the constants $C_0$, $K_1$, $C_F$ at the time $t'$.
The clock is defined up to level $k$, and (Loc) holds on $[0,t']$ by
Lemma~\ref{lem:geometry}(c). If $k\ge2$, ring $k-1$ is quiet on $[0,t']$, so
Step~2 gives the hypotheses and conclusions of Lemma~\ref{lem:O} for $j=k-1$ and
$\tau=A_{k-1}t'$; and $\theta_{k-1}A_{k-1}t'\le\theta_{k-1}\tash_{k-1}\le\frac12$
(as ring $k-1$ is quiet at $t'$) by Lemma~\ref{lem:quiet-time} at the time
$\tsh_{k-1}$ and Lemma~\ref{lem:theta-quiet}; the hypotheses of
Lemma~\ref{lem:quiet-time} besides (B6) concern the levels $i\ge k$, whose amplitudes are
bounded below by Step~3 at the times $\tsh_i\le\tsh_k$ and, for $i=n$, by
Step~0. (BA3) for ring $k$ holds on $(\tsh_k,t']$ by (BA$^\circ$)$(t)$ and at
$\tsh_k$ by Equation~\eqref{eq:boot-switch} or \eqref{eq:boot-start}; the inner
bound is Equation~\eqref{eq:boot-inner}. If $k<n$, the conclusions of
Theorem~\ref{thm:quiet} for ring $k$ on $[0,\tash_k]$ hold by Step~2,
$\maj_k(\tash_k)\le K_1\ksh_k$ by Step~3, and $D_{k+1}\subset B_{ar_k/8}$ on
$[\tsh_k,t']$ by Lemma~\ref{lem:geometry}(a). So
$\|h_2(u)\|_{C^{M_h}(N')}\le\eta_{b,k}\hat a_k(u)$ on $[0,\bar u]$, which is the
second bound of (Field).

(Coef): we apply Lemma~\ref{lem:drift} to ring $k$ with this $\bar u$ and
$\eta=\eta_{b,k}$. Hypotheses (i) and (ii) concern only $k<n$. (i) is
Lemma~\ref{lem:geometry}(a) on $W_k$, since the kick time is $\tauK=\Tpin+\mu u$.
(ii), with $q'=\qs$ of Step~3: the hypotheses of Lemma~\ref{lem:coupling}(2) hold
by Lemma~\ref{lem:geometry}(b), so Equation~\eqref{eq:handoff} and
Lemma~\ref{lem:coupling}(2) give
$\ksh_k\ge\frac12|\cin_k|=\frac12e^{x_k\Tpin}|\Ain(\Tpin)|$; Step~3 gives
the lower bound on $|\Ain(\Tpin)|$. By Step~3 and Lemma~I.\ref{I-lem:uniform},
$|\qs|\ge\frac{15}{16}2\pi|\pout_*(G_{k+1})|\delta_0\ge\frac{15}{16}\cout\delta_0>0$,
so $\qs\ne0$, and (P2) gives
$\mu-1\le c'\delta_0\le|\qs|/(10|z|\Cmom\Ip{m}{\sigma/4})$. (iii) is the outer
bound of (Field), the $C^0$ norm being at most the $C^{M_h}$ norm, and (iv) is
Lemma~\ref{lem:geometry}(b). Moreover $\rT(\Tpin)\le\frac1{64}$ by (P3), and
$\Cls\eta_{b,k}\hat T_k\le\frac1{32}$: for $k=1$, $\eta_{b,1}=0$; for $2\le k<n$,
$\hat T_k\le\sigma'^{-1}\big((m+2)\log\frac1{s_k}+\log_+\frac{\delta_0}{\Klo}\big)$,
as $\sigma_k\ge\sigma'$ and $\ksh_k\ge\Klo s_k^{m+2}$, and (B8) applies; for
$k=n$, $\hat T_n\le\sigma'^{-1}\log(2\delta_0/\kseed)$ by Step~0 and
$\eta_{b,n}\le\eta_b$, and (B8) applies. So Lemma~\ref{lem:drift} gives (Coef) on
$[0,\bar u]$ and, if $\bar u=\hat T_k$,
$|c_k(\hat T_k)|e^{\sigma_k\hat T_k}\in[\frac{15}{16}\delta_0,\frac{17}{16}\delta_0]$.

Theorem~\ref{thm:active} and Corollary~\ref{cor:active-moment} therefore hold on
$[0,\bar u]$: in particular $S_k(t'')\subset\NN_k$ and
$\|w_k(t'')\|_{H^s}\le C_2\hat a_k(u)(1+\rhohat_k(u))$ for
$t''=\tsh_k+u/A_k\in[\tsh_k,t']$.

\emph{Step 5 (quiet size).} Let $t'$ be in the current interval, $j<n$ quiet at
$t'$ and $\tau=A_jt'$. We show $\maj_j(\tau)\le\Abar\varepsilon_j$; by
Lemma~\ref{lem:geometry}(b), $h_j$ is defined on $[0,t']$. If $t'=0$, then
$\maj_j(0)=0$ (Step~0); let $t'>0$. Suppose first that
$\tsh_j$ is defined, so that $\tau\le\tash_j$, and apply
Proposition~\ref{prop:relay} with the constants $C_0$ and $K_1$ to $j$ and
$\tau$. (Loc) holds on $[0,t']$ by
Lemma~\ref{lem:geometry}(c), and $\etout_j\le\theta_j\maj_j$ on $[0,\tau]$ by
Step~2. Lemma~\ref{lem:quiet-time} at the time $t'$, whose hypotheses on the
levels $i\in(j,n)$ with $\tsh_i\le t'$ hold by Step~3 at the times $\tsh_i\le t'$
and $\ksh_n\ge\frac12\kseed$ by Step~0, gives $\tau\le\bar\tau^\sharp_j$, and
Lemma~\ref{lem:theta-quiet} gives $\theta_j\bar\tau^\sharp_j\le\frac12$; so
$\theta_j\tau\le\frac12$. If $j+1<n$, the conclusions of Theorem~\ref{thm:quiet}
for ring $j+1$ on $[0,A_{j+1}\min(t',\tsh_{j+1})]$ hold by Step~2. If
$t'>\tsh_{j+1}$, the remaining hypotheses hold exactly as in Step~3 with $k=j$.

If $\tau\le A_j\tsh_{j+1}$, then $j+1<n$, as $\tsh_n=0<t'$, and ring $j+1$ is
quiet at $\tau/A_j$; so Proposition~\ref{prop:relay}(i), (BA2) for ring $j+1$ at the time $\tau/A_j$ and
(B10) give, as $\varepsilon_{j+1}\le\mu\sbar^{m+2}$,
\[
\maj_j(\tau)\le2C_M\Big(\frac{2C_Q}{\sigb(\mu-1)}+1\Big)\varepsilon_j\cdot2\Abar\varepsilon_{j+1}
\le\Abar\varepsilon_j .
\]
If $\tau>A_j\tsh_{j+1}$, Proposition~\ref{prop:relay}(ii) or (iii) bounds
$\maj_j(\tau)$ by the right side of Equation~\eqref{eq:relay}. As in Step~3 with
$k=j$ ($1-\gamma_j\ge\gamma'_\mu$ and $\ksh_{j+1}<\delta_0$), $\Xi_{j+1}\le\delta_0/(\sigb(\mu-1))$ by ($\kappa$7),
and $x_j\le\frac{17}{16}\sigma$ by Lemma~\ref{lem:detuning}; so that right side
is at most
\[
2C_M\varepsilon_je^{17\sigma\Tpin/16}\Big[\frac{(4C_2+1)\delta_0}{\sigb(\mu-1)}+C_D(\Tpin)\Big]=\Abar\varepsilon_j .
\]
Suppose now that $\tsh_j$ is not defined. Then $j<k_{\mathrm s}$, and by Step~2
and the case just treated for ring $k_{\mathrm s}$, which is quiet at $t'$,
\[
\maj_j(\tau)\le\tfrac34\Ptil_{j,k_{\mathrm s}}\maj_{k_{\mathrm s}}(A_{k_{\mathrm s}}t')
\le\tfrac34\Abar\,\KO\varepsilon_j\,\varepsilon_{k_{\mathrm s}}\prod_{p=j+1}^{k_{\mathrm s}-1}\KO\varepsilon_p\le\tfrac38\Abar\varepsilon_j,
\]
because $\KO\varepsilon_p\le\KO\mu\sbar^{m+2}\le\frac12$ for every $p$ by (B7).
So (BA2) holds with $\Abar$ in place of $2\Abar$.

\emph{Step 6 (localization).} Let $j$ be not debris at $t'\in[0,t]$. If $j$ is
quiet at $t'$, then $S_j(t')\subset\NNm_j\subset\NN_j$ by
Equation~\eqref{eq:boot-quiet} (for $j=n$, quiet only at $t'=0$, by Step~0); if
$j$ is active at $t'$, then $S_j(t')\subset\NN_j$ by Step~4. So (BA1) holds with
$\NN_j$ in place of $\Nflat_j$. With Step~5, and with Step~4 and
Equations~\eqref{eq:boot-switch}, \eqref{eq:boot-start} for (BA3): if
(BA$^\circ$)$(t)$ holds, then (BA)$(t)$ and (BA$^*$)$(t)$ hold.

\emph{Step 7 (continuity, and the clock does not stop).} Let $\mathcal B$ be the
set of $t\in[0,\Tend]$ such that (BA$^\circ$)$(t)$ holds. It contains $0$ by
Step~0, and with every $t$ it contains $[0,t]$. Let $t_*=\sup\mathcal B$. We use
three continuity properties. The flow $\flow$ is continuous on
$[0,\Tend]\times\R^2$ (Theorem~\ref{thm:smooth}(2),(4)), so for a compact set $S$
and an open set $U$ with $\flow_{t_0}(S)\subset U$ we have $\flow_{t'}(S)\subset U$
for $t'$ near $t_0$. On an interval on which $S_k(t')\subset N_k$, the map
$t'\mapsto w_k(t')\in H^s$ is continuous, because
$\omega\mathbf 1_{S_k(t')}=f^{(n)}_k\circ\flow_{t'}^{-1}$
(Equation~\eqref{eq:omega-ge}) is smooth in $(t',x)$ with support in a fixed
compact set. On an interval on which $h_j$ is defined, $\tau\mapsto
h_j(\tau)\in C^{M_h}(N')$ is continuous (Lemma~\ref{lem:reduction}(1)), so
$\eta_j$ and $\maj_j$ are continuous.

\emph{$t_*\in\mathcal B$.} The conditions of (BA$^\circ$)$(t_*)$ at times
$t'<t_*$ hold, because $\mathcal B\supset[0,t')$; so let $t_*>0$ and consider
$t'=t_*$. A ring $j$ that is not debris at $t_*$ is not debris on $[0,t_*]$, and
$S_j(t'')\subset\Nflat_j$ for $t''<t_*$; as $\Nflat_j$ is closed,
$S_j(t_*)\subset\Nflat_j$. So (BA1) holds on $[0,t_*]$. Every level $k<n$ with
$\tsh_k<t_*$ has $\tsh_k\in\mathcal B$, so Step~3 gives $\ksh_k\ge\Klo s_k^{m+2}$,
and Lemma~\ref{lem:geometry} applies with $t=t_*$; so, for a ring $j<n$ quiet at
$t_*$, $h_j$ is defined and $\maj_j$ is continuous on $[0,t_*]$, and
$\maj_j(A_jt_*)\le2\Abar\varepsilon_j$ follows from the same bound at earlier
times. Let $k$ be active at $t_*$ with $\tsh_k<t_*$. Then $w_k$ is continuous on
$[\tsh_k,t_*]$, and the right side $R_k(u)=2C_2\hat a_k(u)(1+\rhohat_k(u))$ of
(BA3) is continuous except at $u=t_1(C_0,G_k)$, where it is right-continuous
and jumps upward: the transient $\rhohat_k$ of Theorem~\ref{thm:active} takes the value
$1$ there, while its left limit there is the value of
$C_0\jap{u}^pe^{-(\sigma_k-\sigma')u}$ at $u=t_1(C_0,G_k)$, which is at most
$2^{-p/2}$ by the definition of $t_1(C_0,G_k)$.
So $R_k(u_*)\ge\limsup_{u\uparrow u_*}R_k(u)$ at $u_*=A_k(t_*-\tsh_k)$, and (BA3)
at $t_*$ follows from (BA3) at earlier times. Hence $t_*\in\mathcal B$.

\emph{$t_*=\Tend$.} Suppose $t_*<\Tend$. By Steps~1--6, (BA$^*$)$(t_*)$ holds.
Choose $\Delta>0$ such that $t_*+\Delta\le\Tend$ and no $\tsh_i$ or $T_i$ lies
in $(t_*,t_*+\Delta]$. Then a ring that is not debris (respectively quiet,
active) at a time in $(t_*,t_*+\Delta]$ is not debris (respectively quiet, active
with $\tsh_k\le t_*$) at $t_*$. Each $S_j(t_*)$ with $j$ not debris at $t_*$ is a
compact subset of $\NN_j\subset\operatorname{int}\Nflat_j$, so after decreasing
$\Delta$, (BA1) holds on $[0,t_*+\Delta]$. No $\tsh_k$ lies in
$(t_*,t_*+\Delta]$, so Lemma~\ref{lem:geometry} applies with $t=t_*+\Delta$, and
$h_j$ is defined on $[0,t_*+\Delta]$ for every ring $j<n$ quiet at $t_*+\Delta$;
as $\maj_j(A_jt_*)\le\Abar\varepsilon_j<2\Abar\varepsilon_j$ and $\maj_j$ is
continuous, after decreasing $\Delta$, (BA2) holds on $[0,t_*+\Delta]$. For a
ring $k$ active at times in $(t_*,t_*+\Delta]$,
$\|w_k(t_*)\|_{H^s}\le C_2\hat a_k(1+\rhohat_k)<R_k$ at $u_*$, because
$\hat a_k>0$; as $w_k$ is continuous and $R_k$ is right-continuous, after
decreasing $\Delta$, (BA3) holds on $(t_*,t_*+\Delta]$. So
$t_*+\Delta\in\mathcal B$, contradicting $t_*=\sup\mathcal B$. Hence $t_*=\Tend$
and $\mathcal B=[0,\Tend]$, and by
Steps~1--6, (BA)$(t)$ and (BA$^*$)$(t)$ hold for every $t\in[0,\Tend]$.

If the clock were not defined up to level $1$, then $k_{\mathrm s}<n$ by Step~0
and $\Tend=\tsh_{k_{\mathrm s}}\in\mathcal B$, so Step~3 for $k=k_{\mathrm s}$
would show that the clock is defined up to level $k_{\mathrm s}$, contrary to
Definition~\ref{def:stop}. Hence $\Tend=T_1$.

We read off the items. (1): $\Tend=T_1$; the bound on $\ksh_1$ is Step~3 for
$k=1$. (2): Step~4 at $t=T_1$ gives the conclusions of
Theorem~\ref{thm:active} on all of $\Phase_1$, and Lemma~\ref{lem:drift} with
$\bar u=\hat T_1$ gives the bound on $|c_1(\hat T_1)|e^{\sigma_1\hat T_1}$.
(3): by (BA)$(T_1)$ and Step~3, Lemma~\ref{lem:geometry} applies with $t=T_1$,
and its part (b) with $t'=T_1$ and $j=1$, ring $1$ being active at $T_1$, gives
$D_2(T_1)\subset B_{ar_1/8}=B_{a/8}$, as $r_1=1$. (4): the clock is defined up to
level $1$, $\ksh_k\ge\Klo s_k^{m+2}$ for $k<n$ (Step~3) and $\ksh_n\ge\frac12\kseed$ (Step~0); so
Lemma~\ref{lem:time} gives $T_1\le\Tstar$.
\end{proof}

\begin{proof}[Proof of Proposition~\ref{prop:reduction}]
Assume \textup{(N1)}--\textup{(N3)}. By Lemma~\ref{lem:nonempty}, we may choose
parameters satisfying (P0)--(P6). This fixes $\mu\in(1,\mu_{\max}]$ with
$\mu_{\max}\le\frac{17}{16}$ (by (P2)) and $\sbar\in(0,a/16]$ (by (P6) and
(B1)), and none of (P0)--(P6) involves $A_1$. Let $A_1>0$, let $f$ be the
function of Definition~\ref{def:datum} with these parameters, which is the
function $f$ of Equation~\eqref{eq:datum}, and let $\Tstar$ be the number of
Lemma~\ref{lem:time}. By Proposition~\ref{prop:integrability}, $f$ is radial
and nonnegative, $f\in (L^1\cap L^q)(\R^2)$ for every $q\in[1,\infty)$, and
$f\notin L^\infty(\R^2)$. The bracket in the definition of $\Tstar$ is a finite
sum of the term $1$ and of nonnegative terms, because $\Tz>0$, $\mu>1$,
$0<\sbar\le a/16<1$ and $\kseed<\delta_0$ by ($\kappa$1) at $\kappa=2\kseed$;
so $\Tstar\in[1/A_1,\infty)$, and in particular $\Tstar\in(0,\infty)$.

Fix $\chiout\in C_c^\infty((a/4,3))$ with $0\le\chiout\le1$ and $\chiout=1$ on
$[a/2,2]$, and put $\Psione(x)=\chiout(|x|)\bar x^{\,m}$; it does not depend on
$n$. The cutoff $\chiout$ is chosen so that, at the time $T_1$ below, it equals
$1$ on the support of ring $1$ and vanishes on the material inside ring $1$; so
the moment $\int\omega^{(n)}(T_1)\Psione$ sees only ring $1$.

For $n\ge2$ let $f^{(n)}$ and $\omega^{(n)}$ be as in Definition~\ref{def:approx},
with the seed size $\kseed$ of (P5), which lies in $(0,\kL]$ by ($\kappa$6) at
$\kappa=2\kseed$. By Lemma~\ref{lem:adm}, $(f^{(n)},\omega^{(n)})_{n\ge2}$ is a
$2$-admissible sequence for $f$ on $[0,\infty)$ (indexed from $n=2$), and
$\omega^{(n)}(t)$ is $3$-fold rotationally symmetric for every $n$ and $t$.
For $n\ge2$ put $T^{(n)}=T_1$, which is defined by
Proposition~\ref{prop:bootstrap}(1) and lies in $[0,\Tstar]$ by
Proposition~\ref{prop:bootstrap}(4). (SA) and (Det) hold by
Lemma~\ref{lem:nonempty}, and the clock is defined up to level $1$. By
Proposition~\ref{prop:bootstrap}(2), the conclusions of
Theorem~\ref{thm:active}(a)--(b) hold for ring $1$ on $\Phase_1$ with $G=G_1$,
$\kappa=\ksh_1$, the constant $C_0$ and $\eta_b=\eta_{b,1}=0$
(Definition~\ref{def:P6}). Ring $1$ has
$r_1=1$, so ring-$1$ units are the original lengths; at $t=T_1$ its active time
is $\hat T_1$, so $w(\hat T_1)=w_1(T_1)$ in Theorem~\ref{thm:active}. The
vorticity $\omega^{(n)}(T_1)$ is supported in $D_1(T_1)$, the disjoint union of
$S_1(T_1)$ and $D_2(T_1)$. By Proposition~\ref{prop:bootstrap}(3),
$D_2(T_1)\subset B_{a/8}$, so the part on $D_2(T_1)$ does not
contribute, as $\chiout$ vanishes on $[0,a/4]$. On $S_1(T_1)$,
$\omega^{(n)}(T_1)=A_1(\Phi+w_1(T_1))$ (Definition~\ref{def:regions}), which is
supported in $\NN\subset\{a/2\le|x|\le2\}$ by Theorem~\ref{thm:active}(a), where
$\chiout(|x|)=1$. As $\Phi$ is radial, $\pout_m[\Phi]=0$, with the
multipole moment $\pout_m$ of Definition~I.\ref{I-def:multipole}, so
$\int\omega^{(n)}(T_1)\Psione=A_1\int(\Phi+w_1(T_1))\,\bar y^{\,m}\dd y=2\pi A_1\,\pout_m[w_1(T_1)]$.

Write $c=c_1(\hat T_1)$, $\lambda=\lambda_1$ and $\phist=\phist(G_1)$, the
unstable eigenfunction of Definition~I.\ref{I-def:eigendata}. We bound
$|\pout_m[w_1(T_1)]|$ from below by comparing $w_1(T_1)$ with
$2\Re(ce^{\lambda\hat T_1}\phist)$, as in Theorem~\ref{thm:active}\textup{(b)}. As
$\pout_m[\phist]=\pout_*(G_1)$ and $\pout_m[\overline{\phist}]=0$ (its integrand
carries the factor $e^{-6i\theta}$),
$\pout_m[2\Re(ce^{\lambda\hat T_1}\phist)]=ce^{\lambda\hat T_1}\pout_*(G_1)$. The
function $w_1(T_1)-2\Re(ce^{\lambda\hat T_1}\phist)$ is supported in $N$, because
$\Phi$, $\Phi+w_1(T_1)$ and $\phist\in X^0$ are. (P1) gives
$\Cp C_2\delta_0\le\cout/(16\pi)$, and ($\kappa$5) with $G=G_1\in[0,\Gone]$
(Lemma~\ref{lem:inner-circ} and (SA3)) and
$\kappa=\ksh_1$, which lies in $(0,\kseed]$ by
Proposition~\ref{prop:bootstrap}(1) and (B6) (it is positive since the clock is
defined up to level $1$, Definition~\ref{def:clock}), gives
$\Cp C_2(\rhohat_1(\hat T_1)+\ksh_1{}^{1/2})\le\cout/(16\pi)$. As
$\hat a_1(\hat T_1)=\delta_0$, Lemma~I.\ref{I-lem:uniform},
Theorem~\ref{thm:active}(b) with $\eta_b=0$ and these two inequalities give
\begin{align*}
\big|\pout_m[w_1(T_1)]-ce^{\lambda\hat T_1}\pout_*(G_1)\big|
&\le\Cp\big\|w_1(T_1)-2\Re(ce^{\lambda\hat T_1}\phist)\big\|_{H^s}\\
&\le\Cp C_2\delta_0\big[\rhohat_1(\hat T_1)+\delta_0+\ksh_1{}^{1/2}\big]
\le\frac{\cout\delta_0}{8\pi}.
\end{align*}
As $G_1\in[0,\Gone]$, $|\pout_*(G_1)|\ge\cout/(2\pi)$ (Lemma~I.\ref{I-lem:uniform}), and
$|ce^{\lambda\hat T_1}|=|c_1(\hat T_1)|e^{\sigma_1\hat T_1}\ge\frac{15}{16}\delta_0$
by Proposition~\ref{prop:bootstrap}(2). Hence
$|\pout_m[w_1(T_1)]|\ge\frac{\cout}{2\pi}\cdot\frac{15}{16}\delta_0-\frac{\cout\delta_0}{8\pi}
=\frac{11}{16}\cdot\frac{\cout\delta_0}{2\pi}\ge\frac{\cout\delta_0}{4\pi}$, and
multiplying by $2\pi A_1$ gives
$|\int\omega^{(n)}(T^{(n)})\Psione|\ge c_0:=\frac12A_1\cout\delta_0>0$. The number $c_0$
involves only $A_1$, $\cout$ and $\delta_0$, and $\Tstar$ only $A_1$, $\mu$,
$\sbar$, $\sigma'$, $\Tz$, $\delta_0$ and $\kseed$; none of them depends on $n$.
Moreover $\Psione$ has the form required in
Corollary~I.\ref{I-cor:defect}, because $\chiout\in C_c^\infty((a/4,3))$ is real.
Corollary~I.\ref{I-cor:defect} with $\bar T=\Tstar$ gives a renormalized weak
solution $\omega_\infty$ on $[0,\infty)$ with datum $f$, in the sense of
Definition~I.\ref{I-def:weak} with $q=2$, such that
$\omega_\infty\in C([0,\infty);L^p(\R^2))$ for every $p\in[1,\infty)$ and, for
every $t\ge0$, $\omega_\infty(t)$ is equimeasurable with $f$ and $3$-fold
rotationally symmetric, and a time $t_0\in[0,\Tstar)$ such that $\omega_\infty$ and $f$ are
distinct elements of $X_q(T')$ for every $T'>t_0$ and every $q\in(4/3,\infty)$.

Let $T>t_0$ and $q\in(4/3,\infty)$; then $T>0$. By
Lemma~I.\ref{I-lem:radial-stationary}, the stationary function $t\mapsto f$ is a
renormalized weak solution in $X_q(T)$ with datum $f$. Also $\omega_\infty\in
X_q(T)$, since, by equimeasurability, $\|\omega_\infty(t)\|_{L^r}=\|f\|_{L^r}$ for every
$r\in[1,\infty)$ and $\omega_\infty\in C([0,T];L^1\cap L^q)$. By
Definition~I.\ref{I-def:weak}, the restriction of $\omega_\infty$ to $[0,T)$
satisfies Equations~(I.\ref{I-eq:weak}) and (I.\ref{I-eq:renorm}) for every test function
$\varphi\in\DT$ and every renormalization $\Ren$; these identities involve only
$\omega_\infty$, its velocity $\KBS\ast\omega_\infty$ and the datum $f$, and not
the exponent of the space. So $\omega_\infty$ is a renormalized weak solution in
$X_q(T)$ with datum $f$, and, as $T>t_0$, the two solutions are distinct elements
of $X_q(T)$. This holds in particular for every $T>\Tstar$, since $t_0<\Tstar$.
So $f$, $\Tstar$, the stationary solution, $\omega_\infty$ and $t_0$ have every
property stated in Theorem~\ref{thm:main}, for every $A_1>0$ with the same $\mu$
and $\sbar$, which were fixed before $A_1$.
\end{proof}


\appendix
\section{Verification of the Numerical Hypotheses}
\label{app:verify}

This appendix proves Theorem~\ref{thm:N-holds}.

Throughout, $a_0,d,\xiL,\xia,\xib,\xiR$, $H$, $\chi$, $\prof$, $\Omega_G$,
$\psi_\ell$, $\Evans_\ell$, $\Thl{\ell}$ and $\epsst$ are as in
Definition~I.\ref{I-def:ring} and Hypothesis~I.\ref{I-hyp:N}, and $G\ge0$. Only the
properties of $\chi$ listed in Definition~I.\ref{I-def:ring} are used. We identify $\C$ with $\R^2$, write $|Y|_w=\max(\ell|u|,|v|)$ for
$\ell\ge1$ and $Y=(u,v)\in\C^2$, and put
\[
x_a=\frac{\xia-a_0}{d},\quad y_b=-\frac{\xib}{d},\quad \Gmax=\frac1{50},\quad \Gmax'=d,\quad
\Rect=\Big[\frac9{200},3\Big]\times[-5,3].
\]
Definition~I.\ref{I-def:ring} gives exactly $\xia-a_0=-\frac12+2^{-54}$ and
$d=\frac1{50}(1+6\cdot2^{-58})$. So $d>\frac1{50}$, whence the range
$[0,\Gmax']$ of the computations contains the range $[0,\Gmax]$ of $G$ in
\textup{(N1)} and \textup{(N2)}, and $50(1-6\cdot2^{-58})<1/d<50$, so
$x_a=-(\frac12-2^{-54})/d$ and $y_b=-1/(2d)$ satisfy
\begin{equation}
\label{eq:xayb}
-25<x_a<-25+4\cdot10^{-15},\qquad -25<y_b<-25+10^{-15}.
\end{equation}
For a continuous closed curve $c\colon[0,1]\to\C\setminus\{0\}$,
$\wind(c,0)=(\theta(1)-\theta(0))/2\pi$, where $\theta$ is any continuous
function with $c=|c|e^{i\theta}$. If $\Gamma$ is the boundary of a closed
rectangle in $\{\Re\lambda>0\}$ and $\Evans_3(\cdot;G)\ne0$ on $\Gamma$,
$\wind(\Evans_3(\Gamma;G),0)$ is the winding number of
$t\mapsto\Evans_3(\gamma(t);G)$, where $\gamma$ traverses $\Gamma$ once
counterclockwise.

\subsection{Ball Arithmetic and the Trusted Base}
\label{app:arith}

The computations use \textbf{ball arithmetic}: a real number is enclosed in a
ball $[m\pm r]=[m-r,m+r]$ whose midpoint $m$ and radius $r\ge0$ are binary
floating-point numbers, and a complex number in a box $X+iY$, the product of two
real balls. We use the implementation of Johansson \cite{Johansson2017}, part of
FLINT \cite{FLINT} from version~3 on, in FLINT~3.6.0 as bundled with
version~0.9.0 of its Python interface python-flint \cite{pythonflint}, at a
working precision of $128$ bits. For a real ball $X$, $\mathrm{lo}(X)$ and
$\mathrm{up}(X)$ are its endpoints, which are dyadic rationals; for a box $Z$,
$|Z|^{\rm up}$ is a dyadic upper bound of $|z|$ over $z\in Z$; for a ball $\rho$,
$\mathrm{box}(\rho)=[-\mathrm{up}(\rho),\mathrm{up}(\rho)]
+i[-\mathrm{up}(\rho),\mathrm{up}(\rho)]$, which contains the closed disk of
radius $\rho'$ about $0$ for every $\rho'\in\rho$ with $\rho'\ge0$;
$\mathrm{hull}(X,X')$ is a ball computed from $\mathrm{lo}(X)$ and
$\mathrm{up}(X')$; and $X\cup X'$ is the computed union of two real balls, again
a ball. Every enclosure rests on the following \textbf{containment property}.
\begin{enumerate}[label=\textup{(CP)}, leftmargin=*]
\item For every arithmetic operation $f$ that we use (sums, differences, products
  and quotients of real balls and of boxes, integer powers of a real ball, the
  maximum of two real balls, and the exponential and the square root of a real
  ball), and for all input balls $X_1,\dots,X_k$, the output ball contains
  $f(x_1,\dots,x_k)$ for every choice of $x_j\in X_j$ at which $f$ is defined; a
  ball with infinite radius, or with an undefined midpoint, is regarded as the
  whole real line. A ball built from an integer contains that integer, a ball
  built from an exact midpoint $m$ and an exact radius $r$ is $[m-r,m+r]$, and a
  box built from two real balls is their product. Moreover
  $\mathrm{lo}(X)\le x\le\mathrm{up}(X)$ for every $x\in X$, $|z|\le|Z|^{\rm up}$ for
  every $z\in Z$, $\mathrm{hull}(X,X')$ contains every real number between
  $\mathrm{lo}(X)$ and $\mathrm{up}(X')$, $X\cup X'$ contains $X$, $X'$ and every
  real number between a point of $X$ and a point of $X'$, the computed
  intersection of a ball $X$ with $[0,\infty)$ contains $X\cap[0,\infty)$, and a
  comparison such as $X>0$ or $X\le c$ returns true only if every point of $X$
  satisfies it.
\end{enumerate}
The data of Definition~I.\ref{I-def:ring}, the contour points, and the endpoints of
the intervals of $G$ and the rectangles of Section~\ref{app:runs} are rational.
They are held exactly, as pairs of integers, in the arbitrary-precision integer
arithmetic of CPython, and every comparison between them is made there. A
rational $n/k$ in lowest terms, $k\ge1$, enters ball arithmetic as follows. If
$k=2^e$ and $|n|<2^{125}$, it enters as the product of the ball built from $n$ and
the power $-e$ of the ball built from $2$; the program checks that this product
has radius $0$, so by (CP) it is the point $n/k$. Otherwise it enters as the
quotient of the balls built from $n$ and $k$, which contains $n/k$ by (CP). No
hardware floating-point number of the programming language enters any enclosure:
every ball is built from integers or from other balls.

The proof of Proposition~\ref{prop:runs} is a computation. It is correct provided
the following \textbf{trusted base} holds.
\begin{enumerate}[label=\textup{(TB\arabic*)}, leftmargin=*]
\item The ball arithmetic of FLINT~3.6.0, as bundled with python-flint~0.9.0,
  satisfies (CP) for the operations listed there.
\item CPython integer arithmetic, and the rational arithmetic of its standard
  module \texttt{fractions}, are correct.
\item The program runs as written, with its assertions enabled, and its routines
  carry out the enclosure procedure of Section~\ref{app:procedure} and the
  certifications and runs of Section~\ref{app:runs}.
\end{enumerate}
We did not inspect FLINT's internal error analysis; if FLINT uses hardware
floating-point arithmetic internally, its correctness is part of (TB1). The
function \texttt{evans} of the module \texttt{ringcert.evans} follows steps
(i)--(iv) of Section~\ref{app:procedure} line by line, and the part of (TB3) that
concerns it is established by reading it against them; the tests of the package
detect some departures from the procedure (dropping the Taylor remainder in
step~(iii)), but not all (omitting the a priori enlargement that defines
$P^{\rm st}$ and $P'^{\rm st}$ in step~(iii)), and do not replace this reading. A certification (Section~\ref{app:runs})
is carried out by the function \texttt{certify\_cell} of the module
\texttt{ringcert.winding}, which follows the description of the certification in
Section~\ref{app:runs} line by line, and a run
by the function \texttt{certify\_cell\_adaptive} of the same module. Everything
else in this paper has a written proof that uses no machine computation.

\subsection{Two Elementary Lemmas}

\begin{lemma}
\label{lem:end-regions}
The following two estimates hold.
\begin{enumerate}[label=\textup{(\arabic*)}]
\item We have $\int_{-\infty}^{\xia}|\prof'|\dd\xi\le\frac94e^{2x_a}$ and
  $\int_{\xib}^{\infty}|\prof'|\dd\xi\le\frac94e^{2y_b}$, and both bounds are less
  than $3e^{-50}$.
\item Let $\ell\ge1$, $G\ge0$ and $\Re\lambda>0$, let $\psi=\psi_\ell(\cdot;\lambda,G)$,
  and let $[\alpha,\beta]$ be one of the intervals $[\xiL,\xia]$ and $[\xib,\xiR]$.
  Let $\psi_0$ solve $\psi_0''=\ell^2\psi_0$ with $\psi_0(\alpha)=\psi(\alpha)$ and
  $\psi_0'(\alpha)=\psi'(\alpha)$, and put $Y=(\psi,\psi')$, $Y_0=(\psi_0,\psi_0')$
  and $\epsend=3e^{-50}/\Re\lambda$. Then, for every $\xi\in[\alpha,\beta]$,
  \[
  |Y(\xi)-Y_0(\xi)|_w\le e^{\ell(\xi-\alpha)}\,|Y(\alpha)|_w\,\big(e^{\epsend}-1\big).
  \]
\end{enumerate}
\end{lemma}
\begin{proof}
(1) Write $A(\xi)=H(\frac{\xi-a_0}d)$, $B(\xi)=H(-\frac\xi d)$, so that
$\prof=\chi AB$, and $x=\frac{\xi-a_0}d$, $y=-\frac\xi d$. Since
$H'=\frac12\operatorname{sech}^2$, $A'=\frac1{2d}\operatorname{sech}^2x\ge0$ and
$B'=-\frac1{2d}\operatorname{sech}^2y\le0$; we use $H(x)\le e^{2x}$,
$\operatorname{sech}^2x\le4e^{2x}$, $0\le H\le1$ and $\operatorname{sech}^2\le1$. On
$(-\infty,\xia]$, $\prof'=\chi'AB+\chi(A'B+AB')$, $\chi$ is nondecreasing from
$\chi(\xiL)=0$ to $\chi(\xia)=1$, and $AB\le A\le A(\xia)=H(x_a)$ since $A$ is
increasing; so $\int_{-\infty}^{\xia}|\chi'|AB\le H(x_a)\le e^{2x_a}$. Next,
$0\le\chi\le1$, $|A'B+AB'|\le\frac1{2d}(\operatorname{sech}^2x+H(x)\operatorname{sech}^2y)
\le\frac5{2d}e^{2x}$ and $\int_{-\infty}^{\xia}e^{2(\xi-a_0)/d}\dd\xi=\frac d2e^{2x_a}$;
adding gives $\frac94e^{2x_a}$. On $[\xib,\infty)$ exchange the roles of $x$ and
$y$: $\chi$ is nonincreasing from $1$ to $0$, $AB\le B\le B(\xib)=H(y_b)\le e^{2y_b}$,
$|A'B+AB'|\le\frac1{2d}(\operatorname{sech}^2x\,H(y)+\operatorname{sech}^2y)\le\frac5{2d}e^{2y}$
and $\int_{\xib}^\infty e^{-2\xi/d}\dd\xi=\frac d2e^{2y_b}$. By \eqref{eq:xayb}, both
bounds are at most $\frac94e^{-50}e^{8\cdot10^{-15}}<3e^{-50}$, as
$e^{8\cdot10^{-15}}<\frac43$.

(2) The vector $Z=(\ell\psi,\psi')$ satisfies $Z'=A_0Z+\mathcal QZ$ with
$A_0=\begin{psmallmatrix}0&\ell\\ \ell&0\end{psmallmatrix}$ and
$\mathcal Q=\begin{psmallmatrix}0&0\\ \delta/\ell&0\end{psmallmatrix}$, where
$\delta=i\ell\prof'/(\lambda+i\ell\Omega_G)$, and $Z_0=(\ell\psi_0,\psi_0')$ satisfies
$Z_0'=A_0Z_0$; and $|Y|_w=\|Z\|_\infty$, the maximum of the moduli of the
components. Since $\Omega_G$ is real, $|\lambda+i\ell\Omega_G|\ge\Re\lambda$, so
the $\|\cdot\|_\infty$ operator norm of $\mathcal Q$ is
$b=|\prof'|/|\lambda+i\ell\Omega_G|\le|\prof'|/\Re\lambda$, and
$\int_\alpha^\beta b\le\epsend$ by (1). For $t\ge0$,
$e^{tA_0}=\begin{psmallmatrix}\cosh\ell t&\sinh\ell t\\ \sinh\ell t&\cosh\ell t\end{psmallmatrix}$
has norm $e^{\ell t}$. By variation of constants,
$Z(\xi)-Z_0(\xi)=\int_\alpha^\xi e^{(\xi-\eta)A_0}\mathcal Q(\eta)Z(\eta)\dd\eta$. With
$u(\xi)=e^{-\ell(\xi-\alpha)}\|Z(\xi)-Z_0(\xi)\|_\infty$, $c=\|Z(\alpha)\|_\infty$ and
$\|Z_0(\eta)\|_\infty\le e^{\ell(\eta-\alpha)}c$, this gives
$u(\xi)\le\int_\alpha^\xi b(\eta)(c+u(\eta))\dd\eta$, so by Gronwall's inequality
$c+u(\xi)\le c\exp(\int_\alpha^\xi b)$ and $u(\xi)\le c\,(e^{\epsend}-1)$.
\end{proof}

\begin{lemma}
\label{lem:apriori}
Let $0\le G\le\Gmax'$. Every zero $\lambda$ of $\Evans_3(\cdot;G)$ with
$\Re\lambda\ge\frac9{200}$ lies in $\Rect$.
\end{lemma}
\begin{proof}
Let $\ell=3$, let $\lambda$ be a zero of $\Evans_\ell(\cdot;G)$ with
$\Re\lambda\ge\frac9{200}$, so $\Re\lambda>0$, and put $TV=\int_{\R}|\prof'|\dd\xi$ and
$\Omega_{\max}=\frac12+\Gmax'e^{-2\xiL}$. On $[\xia,\xib]$, $\chi=1$, so, with $A,B$ as in the proof of
Lemma~\ref{lem:end-regions}, $\prof'=A'B+AB'$ with $A'\ge0\ge B'$ and $0\le A,B\le1$.
Hence $\int_{\xia}^{\xib}|\prof'|\le\int_{\xia}^{\xib}(A'-B')
=A(\xib)-A(\xia)+B(\xia)-B(\xib)<2$, and Lemma~\ref{lem:end-regions}(1) gives
$TV<2+6e^{-50}$. Since $e>2.7$ and $2.7^7>1000$, $e^{50}>2.7^{49}=(2.7^7)^7>10^{21}$, so
$6e^{-50}<10^{-20}$; and
$\log2=\sum_{k\ge1}\frac1{k2^k}>\sum_{k=1}^6\frac1{k2^k}>0.69$. So
$TV/\log2<(2+10^{-20})/0.69<2.9$.

For $\xi\in\R$, $0\le\prof\le1$ gives
$0\le\int_{-\infty}^\xi\prof(\xi')e^{2\xi'}\dd\xi'\le\frac12e^{2\xi}$, so
$0\le\Omega_G(\xi)\le\frac12+Ge^{-2\xi}\le\Omega_{\max}$ for $\xi\ge\xiL$. By
Definition~I.\ref{I-def:ring}, $-2\xiL=4528212286209965\cdot2^{-51}<2.011$; $e<2.7183$ gives $e^2<7.3892$, and
$e^x\le1+x+x^2$ ($0\le x\le1$) gives $e^{0.011}<1.01113$, so
$e^{-2\xiL}<7.3892\cdot1.01113<7.4715$. As $d<0.0200001$,
$\Omega_{\max}<\frac12+0.0200001\cdot7.4715<0.6495$.

Let $\psi=\psi_\ell(\cdot;\lambda,G)$ and
$\mu=\min_{\xi\in[\xiL,\xiR]}|\lambda+i\ell\Omega_G(\xi)|$, positive as $\Re\lambda>0$. The proof of
Lemma~\ref{lem:end-regions}(2), run on $[\xiL,\xiR]$ with
$\psi_0(\xi)=e^{\ell(\xi-\xiL)}$ (which has the same data at $\xiL$),
$\int_{\xiL}^{\xiR}b\le TV/\mu=:\varepsilon'$ and $|Y(\xiL)|_w=\ell$, gives
$|Y(\xiR)-Y_0(\xiR)|_w\le e^{\ell(\xiR-\xiL)}\ell\,(e^{\varepsilon'}-1)$. Since
$\psi_0'+\ell\psi_0=2\ell e^{\ell(\xi-\xiL)}$ and
$|(\psi'-\psi_0')+\ell(\psi-\psi_0)|\le2|Y-Y_0|_w$, we get
$|\Evans_\ell(\lambda;G)-2\ell|\le2\ell(e^{\varepsilon'}-1)$. If $\mu>TV/\log2$, then
$\varepsilon'<\log2$, the right side is less than $2\ell$, and
$\Evans_\ell(\lambda;G)\ne0$. As $\lambda$ is a zero, $\mu\le TV/\log2$, and $\mu$ is attained at some
$\xi\in[\xiL,\xiR]$, so $|\lambda+i\ell\,\Omega_G(\xi)|\le TV/\log2$. At that $\xi$,
$\Re\lambda\le TV/\log2<2.9$ and $|\Im\lambda+\ell\Omega_G(\xi)|\le TV/\log2$, so
$0\le\Omega_G(\xi)\le\Omega_{\max}$ gives $-\ell\,\Omega_{\max}-2.9<\Im\lambda<2.9$. As
$\ell=3$, $-3\Omega_{\max}-2.9>-3\cdot0.6495-2.9>-4.85$. Hence
$\lambda\in[\frac9{200},2.9)\times(-4.85,2.9)\subset\Rect$.
\end{proof}

\subsection{The Enclosure Procedure}
\label{app:procedure}

The procedure encloses $\Evans_\ell(\lambda;G)$ for all $\lambda$ in a box and
all $G$ in an interval. It crosses the end regions by
Lemma~\ref{lem:end-regions}, and $[\xia,\xib]$, where $\chi=1$, by a Taylor
method with coefficients and remainders enclosed in ball arithmetic. Fix the
Taylor order $\NT=16$, written $N$ in this subsection, and the points $s_0<s_1<\dots<s_n$, $n=360$, given by
$s_0=\xia$, $s_n=\xib$ and $s_j=\RN(\xia+\frac jn(\xib-\xia))$ for $0<j<n$, where
$\RN(q)$ is the binary64 number nearest to $q$ (ties to even). Only three facts are
used: the $s_j$ are dyadic rationals, they are strictly increasing, and the steps
$h_j=s_{j+1}-s_j$ are exact dyadic rationals.

The input is an integer $\ell\ge1$, a box $\Lambda^{\rm b}=\mathsf X+i\mathsf Y$ and a
real ball $\mathcal G^{\rm b}$. Each operation below is a ball operation; the
procedure \textbf{fails} if a comparison that it requires is not certified as
true, or if its output is not a box of finite radius.
\begin{enumerate}[label=\textup{(\roman*)}]
\item \emph{Coefficient data, independent of $\lambda$ and $G$.} For each $j<n$,
  and for $S=\{s_j\}$ (the \emph{point data}) and for $S=[s_j,s_{j+1}]$ (the
  \emph{step data}), compute
  \[
  \begin{gathered}
  \mathsf A_0=\mathrm{hull}\big(H(\tfrac{\min S-a_0}d),H(\tfrac{\max S-a_0}d)\big),\quad
  \bar{\mathsf A}_0=\mathrm{hull}\big(H(\tfrac{a_0-\max S}d),H(\tfrac{a_0-\min S}d)\big),\\
  \mathsf B_0=\mathrm{hull}\big(H(\tfrac{-\max S}d),H(\tfrac{-\min S}d)\big),\quad
  \bar{\mathsf B}_0=\mathrm{hull}\big(H(\tfrac{\min S}d),H(\tfrac{\max S}d)\big),
  \end{gathered}
  \]
  with $H(x)$ evaluated as $1/(1+\exp(-2x))$; then, for $0\le k\le N$,
  \[
  \begin{gathered}
  \mathsf A_{k+1}=\frac{2}{d(k+1)}\sum_{r=0}^k\mathsf A_r\bar{\mathsf A}_{k-r},\qquad
  \bar{\mathsf A}_{k+1}=-\mathsf A_{k+1},\\
  \mathsf B_{k+1}=-\frac{2}{d(k+1)}\sum_{r=0}^k\mathsf B_r\bar{\mathsf B}_{k-r},\qquad
  \bar{\mathsf B}_{k+1}=-\mathsf B_{k+1},
  \end{gathered}
  \]
  and $\mathsf w_k=\sum_{r\le k}\mathsf A_r\mathsf B_{k-r}$ for $k\le N+1$,
  $\mathsf f_k=(k+1)\mathsf w_{k+1}$, $\mathsf E^\pm_k=(\pm2)^ke^{\pm2S}/k!$ (with $e^{\pm2S}$ a
  ball containing $e^{\pm2\xi}$ for $\xi\in S$) and
  $\mathsf q_k=\sum_{r\le k}\mathsf w_r\mathsf E^+_{k-r}$ for $k\le N$. Then compute the
  balls
  \[
  \gamma_0=\mathrm{hull}\Big(0,\frac{e^{2x_a+2\xia}}{2/d+2}\Big),\qquad
  \gamma_{j+1}=\gamma_j+\Big(\sum_{k<N}\mathsf q_k^{\rm pt}\frac{h_j^{k+1}}{k+1}
  +\mathsf q_N^{\rm st}\frac{h_j^{N+1}}{N+1}\Big)\cap[0,\infty),
  \]
  where the superscripts ${\rm pt}$ and ${\rm st}$ refer to the point and the step
  data of step $j$, and, with $\mathsf g_0=\gamma_j$ for the point data,
  $\mathsf g_0=\mathrm{hull}(\gamma_j,\gamma_{j+1})$ for the step data, and
  $\mathsf g_r=\mathsf q_{r-1}/r$ for $1\le r<N$, put
  $\mathsf o_k=\sum_{r\le k}\mathsf g_r\mathsf E^-_{k-r}$ for $k<N$.
\item \emph{Left end.} Require $\mathrm{lo}(\mathsf X)>0$. Let $\varepsilon$ be an upper
  bound of $3e^{-50}/\mathrm{lo}(\mathsf X)$, and require $\varepsilon\le1$. With
  $L_1=\xia-\xiL$, put
  $P=e^{\ell L_1}+\mathrm{box}(2\varepsilon\,\mathrm{up}(e^{\ell L_1}))$ and $P'=\ell P$.
\item \emph{Middle steps.} For $j=0,1,\dots,n-1$ in turn, do the following, and then
  replace $(P,P')$ by the new values. For the point data and for the step data of
  step $j$, compute
  \[
  \begin{gathered}
  \mathsf D_0=\Lambda^{\rm b}+i\ell(\mathsf o_0+\mathcal G^{\rm b}\mathsf E^-_0),\qquad
  \mathsf D_k=i\ell(\mathsf o_k+\mathcal G^{\rm b}\mathsf E^-_k)\quad(1\le k<N),\\
  \mathsf R_0=1/\mathsf D_0,\qquad \mathsf R_k=-\mathsf R_0\sum_{r=1}^k\mathsf D_r\mathsf R_{k-r}\quad(1\le k<N),\\
  \mathsf V_k=i\ell\sum_{r\le k}\mathsf f_r\mathsf R_{k-r}+\ell^2[k=0]\quad(0\le k<N),
  \end{gathered}
  \]
  where $[k=0]$ is $1$ if $k=0$ and $0$ otherwise. Let $q_*=|\mathsf V^{\rm st}_0|^{\rm up}$,
  require that $q_*$ is finite, let $\kst$ be an upper bound of
  $(\max(q_*,\ell^2))^{1/2}$, require $\kst h_j\le700$, choose a dyadic number
  $M\ge\max(\kst|P|^{\rm up},|P'|^{\rm up})\,e^{\kst h_j}$, put
  $P^{\rm st}=P+\mathrm{box}(h_jM)$ and $P'^{\rm st}=P'+\mathrm{box}(h_jq_*M/\kst)$,
  and compute $\Psi_0=P$, $\Psi_1=P'$ and
  $\Psi_{k+2}=\frac1{(k+1)(k+2)}\sum_{r\le k}\mathsf V^{\rm pt}_r\Psi_{k-r}$, and in the same
  way $\Psi^{\rm st}_k$ from $P^{\rm st}$, $P'^{\rm st}$ and $\mathsf V^{\rm st}$, for
  $k\le N+1$. The new values are
  \[
  P_{\rm new}=\sum_{k<N}\Psi_kh_j^k+\Psi^{\rm st}_Nh_j^N,\qquad
  P'_{\rm new}=\sum_{k=1}^Nk\Psi_kh_j^{k-1}+(N+1)\Psi^{\rm st}_{N+1}h_j^N .
  \]
\item \emph{Right end and output.} With $L_2=\xiR-\xib$ and
  $W=\max(\ell|P|^{\rm up},|P'|^{\rm up})$, return
  \[
  \Evans^{\rm b}=e^{-\ell(\xiR-\xiL)}\Big(e^{\ell L_2}(P'+\ell P)
  +\mathrm{box}\big(4\varepsilon W\,\mathrm{up}(e^{\ell L_2})\big)\Big).
  \]
\end{enumerate}

\begin{proposition}
\label{prop:enclosure}
Let $\ell\ge1$ be an integer, let $\Lambda\subset\{\Re\lambda>0\}$ be a compact
axis-parallel box, and let $\Cell\subset[0,\infty)$ be a compact interval. Assume
that every ball operation satisfies \textup{(CP)}. Run the procedure of
Section~\ref{app:procedure} with $\ell$, a box $\Lambda^{\rm b}\supset\Lambda$ and a ball
$\mathcal G^{\rm b}\supset\Cell$. If it does not fail, it returns a box
$\Evans^{\rm b}$ of finite radius such that $\Evans_\ell(\lambda;G)\in\Evans^{\rm b}$ for every
$\lambda\in\Lambda$ and every $G\in\Cell$.
\end{proposition}
\begin{proof}
Fix $\lambda\in\Lambda$ and $G\in\Cell$, and let $\psi=\psi_\ell(\cdot;\lambda,G)$. We
show that every computed ball contains a corresponding exact quantity, defined
from $\lambda$, $G$ and $\chi$; the last one is $\Evans_\ell(\lambda;G)$. By (CP) it
suffices to check, at each step, that the exact quantity is the same function of
the earlier exact quantities, or else to justify the enclosure separately. For a
smooth function $F$ we write $F_k(\xi)=F^{(k)}(\xi)/k!$.

\emph{Step \textup{(i)}.} On $[\xia,\xib]$, $\chi=1$ and, $\chi$ being smooth, all its
derivatives vanish there, so $\prof^{(k)}=(AB)^{(k)}$, where
$A(\xi)=H(\frac{\xi-a_0}d)$, $B(\xi)=H(-\frac\xi d)$, $\bar A=1-A=H(\frac{a_0-\xi}d)$ and
$\bar B=1-B=H(\frac\xi d)$, as $1-H(x)=H(-x)$. Since $H'=2H(1-H)$, $A'=\frac2dA\bar A$, $\bar A'=-A'$,
$B'=-\frac2dB\bar B$ and $\bar B'=-B'$. The Leibniz rule gives
$A_{k+1}=\frac2{d(k+1)}\sum_{r\le k}A_r\bar A_{k-r}$, likewise for $B$, and
$\prof_k=\sum_{r\le k}A_rB_{k-r}$, $(\prof')_k=(k+1)\prof_{k+1}$,
$(e^{\pm2\xi})_k=(\pm2)^ke^{\pm2\xi}/k!$; the function $q=\prof e^{2\xi}$ has
$q_k=\sum_{r\le k}\prof_r(e^{2\xi})_{k-r}$. For $\xi\in S$, since $H$ is increasing,
$A(\xi)\in\mathsf A_0$, $\bar A(\xi)\in\bar{\mathsf A}_0$, $B(\xi)\in\mathsf B_0$ and
$\bar B(\xi)\in\bar{\mathsf B}_0$; by induction on $k$, each of
$\mathsf A_k,\bar{\mathsf A}_k,\mathsf B_k,\bar{\mathsf B}_k,\mathsf w_k,
\mathsf f_k,\mathsf E^\pm_k,\mathsf q_k$ contains the Taylor coefficient at $\xi$ of
$A,\bar A,B,\bar B,\prof,\prof',e^{\pm2\xi},q$ respectively. So the point data enclose
the Taylor coefficients at $s_j$, and the step data at every point of
$[s_j,s_{j+1}]$.

The function $g(\xi)=\int_{-\infty}^\xi\prof(\xi')e^{2\xi'}\dd\xi'$ has $g'=q\ge0$, so it is
nondecreasing. On $(-\infty,\xia]$, $0\le\prof\le A\le e^{2(\xi-a_0)/d}$, so
$0\le g(\xia)\le\int_{-\infty}^{\xia}e^{2(\xi-a_0)/d+2\xi}\dd\xi=e^{2x_a+2\xia}/(2/d+2)$,
that is, $g(s_0)\in\gamma_0$. Taylor's theorem with the Lagrange remainder,
applied to $t\mapsto\int_{s_j}^{s_j+t}q$, gives
\[
g(s_{j+1})-g(s_j)=\sum_{k<N}q_k(s_j)\frac{h_j^{k+1}}{k+1}+q_N(\eta_j)\frac{h_j^{N+1}}{N+1}
\]
for some $\eta_j\in[s_j,s_{j+1}]$. This increment is nonnegative and lies in the
ball in parentheses defining $\gamma_{j+1}$, so $g(s_j)\in\gamma_j$ for every $j$ by
induction, and $g(\xi)\in\mathrm{hull}(\gamma_j,\gamma_{j+1})$ for
$\xi\in[s_j,s_{j+1}]$ since $g$ is nondecreasing. As $g_r=q_{r-1}/r$ for $r\ge1$, the
balls $\mathsf o_k$ contain the Taylor coefficients of $ge^{-2\xi}$, at $s_j$ for the
point data and at every point of the step for the step data.

\emph{Step \textup{(ii)}.} Since $\lambda\in\Lambda^{\rm b}$,
$\Re\lambda\ge\mathrm{lo}(\mathsf X)>0$, so $\varepsilon\ge\epsend$ (as in
Lemma~\ref{lem:end-regions}) and, as $\varepsilon\le1$,
$e^{\epsend}-1\le e^\varepsilon-1\le(e-1)\varepsilon\le2\varepsilon$.
Lemma~\ref{lem:end-regions}(2) on $[\xiL,\xia]$, where $\psi_0(\xi)=e^{\ell(\xi-\xiL)}$
and $|Y(\xiL)|_w=\ell$, gives $|\psi(\xia)-e^{\ell L_1}|\le2\varepsilon e^{\ell L_1}$ and
$|\psi'(\xia)-\ell e^{\ell L_1}|\le2\ell\varepsilon e^{\ell L_1}$. Hence
$\psi(\xia)\in P$ and $\psi'(\xia)\in\ell P=P'$.

\emph{Step \textup{(iii)}.} Assume $\psi(s_j)\in P$ and $\psi'(s_j)\in P'$. On
$[\xia,\xib]$ the coefficient $V=\ell^2+i\ell\prof'/(\lambda+i\ell\Omega_G)$ is smooth, so
$\psi$ is smooth there. Put $D=\lambda+i\ell\Omega_G=\lambda+i\ell(g+G)e^{-2\xi}$ and
$R=1/D$, where $D\ne0$ as $\Re D=\Re\lambda>0$. The identity $DR=1$ gives
$R_0=1/D_0$ and $R_k=-R_0\sum_{r=1}^kD_rR_{k-r}$ for $k\ge1$, and
$V=\ell^2+i\ell\prof'R$ gives $V_k=i\ell\sum_{r\le k}(\prof')_rR_{k-r}+\ell^2[k=0]$. By
step (i) and (CP), $\mathsf V^{\rm pt}_k\ni V_k(s_j)$ and
$\mathsf V^{\rm st}_k\ni V_k(\xi)$ for every $\xi\in[s_j,s_{j+1}]$; in particular
$|V|\le q_*\le\kst^2$ on the step. For the a priori enclosure, put
$u(t)=\max(\kst|\psi(s_j+t)|,|\psi'(s_j+t)|)$ for $t\in[0,h_j]$. Then
$\kst|\psi(s_j+t)|\le\kst|\psi(s_j)|+\int_0^t\kst|\psi'|$ and
$|\psi'(s_j+t)|\le|\psi'(s_j)|+\int_0^t\kst\cdot\kst|\psi|$, so
$u(t)\le u(0)+\kst\int_0^tu$, and Gronwall's inequality gives
$u(t)\le u(0)e^{\kst t}\le M$. Hence, for $\xi\in[s_j,s_{j+1}]$,
$|\psi(\xi)-\psi(s_j)|\le h_jM$ and, since $|\psi|\le M/\kst$ on the step,
$|\psi'(\xi)-\psi'(s_j)|\le\int|V||\psi|\le h_jq_*M/\kst$; so $\psi(\xi)\in P^{\rm st}$
and $\psi'(\xi)\in P'^{\rm st}$. The equation $\psi''=V\psi$ gives
$(k+1)(k+2)\psi_{k+2}=\sum_{r\le k}V_r\psi_{k-r}$, so $\Psi_k\ni\psi_k(s_j)$ and
$\Psi^{\rm st}_k\ni\psi_k(\xi)$ for every $\xi\in[s_j,s_{j+1}]$ and $k\le N+1$.
Taylor's theorem with the integral remainder gives
\[
\psi(s_{j+1})=\sum_{k<N}\psi_k(s_j)h_j^k+h_j^N\int_0^1N(1-\tau)^{N-1}\psi_N(s_j+\tau h_j)\dd\tau .
\]
The weight $N(1-\tau)^{N-1}$ is a probability density on $[0,1]$, so the integral
lies in the closed convex hull of $\{\psi_N(\xi):\xi\in[s_j,s_{j+1}]\}$, which lies in
the closed convex box $\Psi^{\rm st}_N$. Hence $\psi(s_{j+1})\in P_{\rm new}$; the same
argument for $\psi'$, with $(\psi')_r=(r+1)\psi_{r+1}$, gives
$\psi'(s_{j+1})\in P'_{\rm new}$. By induction on $j$, $\psi(\xib)\in P$ and
$\psi'(\xib)\in P'$ at the end of step (iii).

\emph{Step \textup{(iv)}.} Apply Lemma~\ref{lem:end-regions}(2) on $[\xib,\xiR]$.
Since $(\psi_0'+\ell\psi_0)'=\ell(\psi_0'+\ell\psi_0)$,
$(\psi_0'+\ell\psi_0)(\xiR)=e^{\ell L_2}(\psi'+\ell\psi)(\xib)$, and
\[
|(\psi'+\ell\psi)(\xiR)-e^{\ell L_2}(\psi'+\ell\psi)(\xib)|\le2|Y(\xiR)-Y_0(\xiR)|_w
\le2e^{\ell L_2}|Y(\xib)|_w\cdot2\varepsilon\le4\varepsilon We^{\ell L_2}.
\]
Multiplying by $e^{-\ell(\xiR-\xiL)}$ gives $\Evans_\ell(\lambda;G)\in\Evans^{\rm b}$.
\end{proof}

\subsection{Certified Winding Numbers}

A set $B\subset\C$ is of \textbf{type P} if it does not meet $(-\infty,0]$, and
of \textbf{type N} if it does not meet $[0,\infty)$.

\begin{proposition}
\label{prop:winding}
Let $\Gamma$ be the boundary of a closed rectangle $\Rect'$ in $\{\Re\lambda>0\}$,
traversed counterclockwise and cut into segments $[p_i,p_{i+1}]$,
$i=0,\dots,I-1$, with $p_I=p_0$. Let $\Cell$ be a set, let $E(\cdot;G)$ be continuous
on $\Gamma$ for each $G\in\Cell$, and let $\mathcal B_i\subset\C$ be sets with
$E(\lambda;G)\in\mathcal B_i$ for all $\lambda\in[p_i,p_{i+1}]$ and $G\in\Cell$. Assume that
each $\mathcal B_i$ is of type P or of type N, and assign to each $i$ one type for which this
holds. Then $E(\cdot;G)\ne0$ on $\Gamma$ and, for every $G\in\Cell$,
\[
\wind\big(E(\Gamma;G),0\big)=\sum_{i\ \text{of type N}}\Big(\mathbf 1\big[\Im E(p_{i+1};G)<0\big]-\mathbf 1\big[\Im E(p_i;G)<0\big]\Big).
\]
In this sum the two terms at a point shared by two segments of type N cancel.
\end{proposition}
\begin{proof}
Fix $G\in\Cell$. Every $\mathcal B_i$ misses $0$, so $E(\cdot;G)\ne0$ on $\Gamma$. The
principal argument $\mathrm{Arg}\colon\C\setminus\{0\}\to(-\pi,\pi]$ is continuous on
$\C\setminus(-\infty,0]$, and
$\theta_{\rm N}(z)=\mathrm{Arg}(z)+2\pi\,\mathbf 1[\Im z<0]$ is continuous on
$\C\setminus[0,\infty)$, with values in $(0,2\pi)$. Since $E(\cdot;G)$ maps the
$i$-th segment into $\mathcal B_i$, $\mathrm{Arg}(E(\cdot;G))$ is a continuous
argument on a segment of type P, and $\theta_{\rm N}(E(\cdot;G))$ on one of type N. Each differs by a constant from a continuous argument along the
whole curve, so the winding number is $\frac1{2\pi}$ times the sum over the
segments of their increments. The increments of $\mathrm{Arg}$ telescope around the
closed curve to $0$, and what remains is the indicator sum. At a point shared by
two segments of type N the indicator enters once with each sign.
\end{proof}

\subsection{The Two Computations}
\label{app:runs}

\begin{definition}
\label{def:Nprime}
We consider the following two statements.
\begin{enumerate}[label=\textup{(N\arabic*$'$)}]
\item For every $G\in[0,\Gmax']$, $\Evans_3(\cdot;G)\ne0$ on $\partial\Rect$ and
  $\wind(\Evans_3(\partial\Rect;G),0)=1$.
\item There are $236$ closed intervals $\Cell_i$, $1\le i\le236$, with
  union $[0,\Gmax']$, and closed rectangles
  $\Rect_1,\dots,\Rect_{236}\subset\Rect$, the left edge of
  each of which lies on the line $\Re\lambda=\frac{211}{2500}$, such that for
  every $i$ and every $G\in\Cell_i$, $\Evans_3(\cdot;G)\ne0$ on $\partial\Rect_i$
  and $\wind(\Evans_3(\partial\Rect_i;G),0)=1$.
\end{enumerate}
\end{definition}

\begin{lemma}
\label{lem:Nprime-N}
\textup{(N1$'$)} implies \textup{(N1)}. \textup{(N1)} and \textup{(N2$'$)} imply
\textup{(N2)}.
\end{lemma}
\begin{proof}
Let $G\in[0,\Gmax]\subset[0,\Gmax']$. By
Lemma~I.\ref{I-lem:evans-basic}\textup{(1)} and its last assertion, $\Evans_3(\cdot;G)$ is holomorphic on
$\{\Re\lambda>0\}$, which contains $\Rect$ and every $\Rect_i$.

Assume \textup{(N1$'$)}. By the argument principle, $\Evans_3(\cdot;G)$ has exactly
one zero, counted with multiplicity, in the interior of $\Rect$. By
Lemma~\ref{lem:apriori}, every zero with $\Re\lambda\ge\frac9{200}$
lies in $\Rect$, and by \textup{(N1$'$)} none lies on $\partial\Rect$; so the zeros in
$\{\Re\lambda\ge\frac9{200}\}$ are exactly those in the interior of $\Rect$. This is
\textup{(N1)}.

Assume \textup{(N1)} and \textup{(N2$'$)}, let $\lamst$ be the zero of
\textup{(N1)}, and choose $i$ with $G\in\Cell_i$. By \textup{(N2$'$)} and the argument
principle, $\Evans_3(\cdot;G)$ has a zero $\lambda'$ in the interior of $\Rect_i$, whose left edge
lies on $\{\Re\lambda=\frac{211}{2500}\}$, so $\Re\lambda'>\frac{211}{2500}>\frac9{200}$, and $\lambda'=\lamst$ by \textup{(N1)}.
So the only zero of $\Evans_3(\cdot;G)$ in $\{\Re\lambda\ge\frac9{200}\}$ has real
part greater than $\frac{211}{2500}$, and there is none in
$\{\frac9{200}\le\Re\lambda\le\frac{211}{2500}\}$. This is \textup{(N2)}.
\end{proof}

Both computations use the procedure of Section~\ref{app:procedure} with $\ell=3$.
For a closed interval $\Cell$ with rational endpoints, a closed rectangle
$\Rect'=[x_0,x_1]\times[y_0,y_1]\subset\{\Re\lambda>0\}$ with rational vertices, a
number $\nu>0$ and integers $J,K\ge1$, the \textbf{certification} of
$(\Cell,\Rect')$ with density $\nu$, at most $J$ rounds and at most $K$ segments
is the following computation. The boundary of $\Rect'$ is traversed
counterclockwise from $(x_0,y_0)$, and each edge, of length $L$, is cut into
$\max(2,\lceil\nu L\rceil)$ segments of equal length; all cut points are rational.
For a segment $[p,q]$, the procedure is run with
$\Lambda^{\rm b}=(X_p\cup X_q)+i(Y_p\cup Y_q)$, where $X_p,X_q,Y_p,Y_q$ are the balls
that contain $\Re p$, $\Re q$, $\Im p$, $\Im q$, and with $\mathcal G^{\rm b}$ the
union of the balls that contain the two endpoints of $\Cell$; by (CP),
$\Lambda^{\rm b}\supset[p,q]$ and $\mathcal G^{\rm b}\supset\Cell$. The resulting box
$\Evans^{\rm b}$ is accepted as of type P if it is certified that
$\Im\Evans^{\rm b}>0$, or $\Im\Evans^{\rm b}<0$, or $\Re\Evans^{\rm b}>0$; otherwise it
is accepted as of type N if it is certified that $\Re\Evans^{\rm b}<0$. The first
round treats every segment; each later round treats the two halves, cut at the
midpoint, of every segment not accepted in the previous round, including those for
which the procedure failed. A segment is cut only if there are at most $K$
segments at that moment, and at most $J$ rounds are made; otherwise the
certification does not succeed. If every segment has been accepted, then at every
point $p$ where a segment of type N meets a segment of type P the procedure is run
with $\Lambda^{\rm b}$ the box of the balls that contain $\Re p$ and $\Im p$ and
with the same $\mathcal G^{\rm b}$, and the
indicator $\mathbf 1[\Im\Evans_3(p;G)<0]$ is taken to be $1$ if the resulting box
is certified to have negative imaginary part and $0$ if it is certified to have
positive imaginary part. The certification \textbf{succeeds} if every segment is
accepted and every such indicator is decided; its result is then the indicator sum
of Proposition~\ref{prop:winding}.

A \textbf{run} with at most $D$ halvings certifies a list of pairs
$(\Cell,\Rect')$: if the certification of $(\Cell,\Rect')$ does not succeed, $\Cell$
is halved at its midpoint and each half is certified with the same $\Rect'$, at
most $D$ times in succession; the certified pieces take the place of $\Cell$, and a
piece still not certified after $D$ halvings is reported as failed.

The run for \textup{(N1$'$)} uses the rectangle $\Rect$ and the $25$ intervals
\[
\Big[\frac{j}{20}d,\frac{j+1}{20}d\Big]\quad(0\le j\le14),\qquad
\Big[\frac{30+j}{40}d,\frac{31+j}{40}d\Big]\quad(0\le j\le9),
\]
whose union is $[0,\Gmax']$, with density $\nu=20$, at most $16$ rounds, at most
$6000$ segments and at most $4$ halvings. The run for \textup{(N2$'$)} uses $236$
intervals $\Cell_1,\dots,\Cell_{236}$, numbered from left to right, and rectangles
$\Rect_i=[\frac{211}{2500},x_i]\times[y_i,y_i']$, with density
$\nu=3/(x_i-\frac{211}{2500})$, at most $10$ rounds, at most $800$ segments and at
most $5$ halvings. These $236$ pairs, with every number an exact fraction, form the file
\texttt{sector3\_gap\_cells.json} of the reproduction package
(Section~\ref{app:reproduce}), and the file \texttt{sector3\_count\_cells.json}
lists the $25$ intervals of the first run, also numbered from left to right. In each of these lists the endpoints of
the intervals are rational, the first interval starts at $0$, the last ends at
$d=\Gmax'$, and consecutive intervals share an endpoint; every $\Rect_i$ has
rational vertices, satisfies $x_i>\frac{211}{2500}$ and $y_i<y_i'$, and lies in
$\Rect$. The re-check of Section~\ref{app:reproduce} confirms these facts, and that
every rectangle of the first run is $\Rect$, for the output files of the runs, in
exact rational arithmetic.

\begin{proposition}
\label{prop:runs}
\textup{(N1$'$)} and \textup{(N2$'$)} hold.
\end{proposition}
\begin{proof}
The proof is the computation just described, carried out under the trusted base of
Section~\ref{app:arith}. Each of the $25$ certifications of the first run and each
of the $236$ of the second succeeded at once, with result $1$, so no interval was
halved and the output files list exactly the intervals and rectangles described
above.

Consider a certification of $(\Cell,\Rect')$ that succeeds with result $1$. By
Lemma~I.\ref{I-lem:evans-basic}\textup{(1)} and its last assertion, $\Evans_3(\cdot;G)$ is continuous on
$\{\Re\lambda>0\}\supset\partial\Rect'$ for every $G\in\Cell$. By
Proposition~\ref{prop:enclosure}, the box of each accepted segment $[p,q]$
contains $\Evans_3(\lambda;G)$ for all $\lambda\in[p,q]$ and $G\in\Cell$, and each decided
indicator equals $\mathbf 1[\Im\Evans_3(p;G)<0]$ for every $G\in\Cell$; by (CP) for
the comparisons, a box accepted as of type P misses $(-\infty,0]$ and one accepted
as of type N misses $[0,\infty)$. The accepted segments are consecutive and cover
$\partial\Rect'$. By Proposition~\ref{prop:winding}, $\Evans_3(\cdot;G)\ne0$ on
$\partial\Rect'$ and $\wind(\Evans_3(\partial\Rect';G),0)=1$ for every $G\in\Cell$.
Applied to the first run, whose intervals cover $[0,\Gmax']$, this gives
\textup{(N1$'$)}; applied to the second, with the properties of the $\Cell_i$ and
$\Rect_i$ stated above, it gives \textup{(N2$'$)}.
\end{proof}

\subsection{The Closed-Form Inequality}

\begin{proposition}
\label{prop:N3}
\textup{(N3)} holds.
\end{proposition}
\begin{proof}
We first bound $\epsst$. Since $H(x)\le e^{2x}$, Equation~\eqref{eq:xayb} gives
$\epsst\le e^{2\max(x_a,y_b)}<e^{-49}<10^{-21}$, as $e^{49}>2.7^{49}>10^{21}$
(proof of Lemma~\ref{lem:apriori}).

We next bound $e^{a_0}$. From Definition~I.\ref{I-def:ring}, $a_0=-7304210039150669\cdot2^{-54}<-0.4054$. The partial
sum $\sum_{k\le5}0.4054^k/k!$ of the exponential series is a lower bound for
$e^{0.4054}$, and it exceeds $1.49989$, so $e^{a_0}<e^{-0.4054}<1/1.49989<0.66672$.
Squaring and multiplying, $0.66672^2<0.444516$, $0.66672^3<0.296368$,
$0.66672^6<0.087834$ and $0.66672^9<0.026032$.

We now bound $\Thl{\ell}$ for $\ell\in\{6,9\}$. Put $u=\pi\ell d$. Since $\pi<3.1416$ and $0<d<0.0200001$, we have $0<u<0.377$ for $\ell=6$ and
$0<u<0.5655$ for $\ell=9$; in particular $0<u<\pi$. For $u\ge0$,
$\sin u\ge u-u^3/6$, since the difference vanishes at $0$ and has derivative
$\cos u-1+u^2/2$, which is nonnegative (it vanishes at $0$ and its derivative $u-\sin u$ is
nonnegative). So $\Thl{\ell}=u/\sin u\le(1-u^2/6)^{-1}$ as
$u^2<6$. Hence $\Thl{6}\le(1-0.377^2/6)^{-1}<(1-0.023689)^{-1}<1.024264$
and $\Thl{9}<(1-0.0533)^{-1}<1.0564$. Therefore
\[
\begin{gathered}
e^{6a_0}\Thl{6}+\epsst<0.087834\cdot1.024264+10^{-21}<0.08997,\\
e^{9a_0}\Thl{9}+\epsst<0.026032\cdot1.0564+10^{-21}<0.0276,
\end{gathered}
\]
and both are less than $\frac9{100}$, which is \textup{(N3)}.
\end{proof}

\begin{proof}[Proof of Theorem~\ref{thm:N-holds}]
By Proposition~\ref{prop:runs}, \textup{(N1$'$)} and \textup{(N2$'$)} of
Definition~\ref{def:Nprime} hold; by Lemma~\ref{lem:Nprime-N}, \textup{(N1)} holds, and
then \textup{(N2)}; by Proposition~\ref{prop:N3}, \textup{(N3)} holds.
\end{proof}

\subsection{Reproduction}
\label{app:reproduce}

The computations of Section~\ref{app:runs} can be repeated with version~1.0.0 of
the package \texttt{euler-ring-certificate}, which is available under the Apache
License 2.0 in the public repository \url{https://github.com/mcoiculescu/NonUniqueEuler}
and contains the code, the two
certificate files of Section~\ref{app:runs} and the reference outputs. It requires
only python-flint~0.9.0, which bundles FLINT~3.6.0, under CPython~3.10 or later.
The program refuses to run with assertions disabled, and, unless explicitly
allowed, with another FLINT version. From the top directory of the package, the
command \mbox{\texttt{./reproduce.sh 8}} runs, with $8$ worker processes, among other
items: an environment check; the unit tests; the certification of the $25$
intervals of \textup{(N1$'$)} and of the $236$ intervals and rectangles of
\textup{(N2$'$)} (commands \texttt{ringcert sector3-count} and
\texttt{ringcert sector3-gap}); a re-check of the outputs in exact rational
arithmetic (\texttt{ringcert check}); and a comparison with the reference
outputs. A successful run reports all items as passed.

The outputs record, for every certified interval, its endpoints, its rectangle,
its result, and the numbers of segments and indicators. From the output files
alone, the re-check confirms that the intervals of each run cover $[0,\Gmax']$,
with consecutive intervals sharing an endpoint, that every result is $1$, that no
indicator is undecided and no interval failed, that every rectangle of the first
run is $\Rect$, and that every rectangle of the second run is nondegenerate, has
left edge on the line $\Re\lambda=\frac{211}{2500}$ and lies in $\Rect$; it also
checks that these output files are the ones recorded in the summary of the run.

The output files \texttt{sector3\_count.json} and \texttt{sector3\_gap.json}
contain no timing or environment data; on an Apple arm64 machine with
CPython~3.13.5 they are byte-identical to the reference outputs of the package.
The SHA-256 digests of the certificate files and of these outputs are the following.
\begin{center}
\scriptsize
\begin{tabular}{ll}
\texttt{sector3\_count\_cells.json} &
\texttt{518a0ca4f05fd24bae273aebea20862fad1c1beb46dbdee9f8b5fb76c8729dc1}\\
\texttt{sector3\_gap\_cells.json} &
\texttt{fb95a3f72eee1f017ddc5c26e7dfca52841759603e34b0131a440c8992086c91}\\
\texttt{sector3\_count.json} &
\texttt{2aa43b18d7289b7288f4bc660b1fdbec547458949ee675ef4b23f65cbcf1a1c9}\\
\texttt{sector3\_gap.json} &
\texttt{f06feec05cc6633b390bb9ee6191ce59c3bd5231a01ffa0f48263c8c9173a786}
\end{tabular}
\end{center}
The reference outputs were produced on that machine only. With another FLINT
version or on other hardware the numbers of segments may differ, and an interval
may be halved. The proof needs only that the re-check passes, that is, that the output files
have the properties listed above. If some intervals
were halved, the argument of Proposition~\ref{prop:runs} gives \textup{(N1$'$)},
and \textup{(N2$'$)} with the certified pieces, each paired with the rectangle of
its interval, in place of the $236$ intervals; the proof of
Lemma~\ref{lem:Nprime-N} does not use their number.

\section{Global Smooth Solutions}
\label{app:smooth}

In this appendix we prove Theorem~\ref{thm:smooth}. The argument is a Picard
iteration through linear transport: each iterate is transported by the
velocity field of the previous one, the iterates obey bounds that are uniform
in the iteration index on every bounded time interval, and they converge
together with all their spatial derivatives. We write the iteration in
$L^2$-Sobolev form, so that, besides elementary real analysis, it uses only
Section~I.\ref{I-sec:prelim}, Section~\ref{sec:II-prelim} and the facts on ordinary differential equations
collected in Lemma~\ref{lem:ode}. For
the classical theory of smooth solutions see \cite{Wolibner}, \cite{Holder} and
\cite{Kato1967}.

\subsection{Ordinary Differential Equations}

The following facts are used without proof.

\begin{lemma}[Theorems 2.2, 2.10, 2.13 and 2.17 in \cite{Teschl}]
\label{lem:ode}
Let $n\ge1$ and $k\ge1$ be integers, let $U\subset\R\times\R^n$ be open, and let
$F\in C^k(U;\R^n)$.
\begin{enumerate}[label=\textup{(\roman*)}]
\item For every $(s,y)\in U$ the problem $\dot x=F(t,x)$, $x(s)=y$, has a unique
  maximal solution, defined on an open interval containing $s$. We write
  $\flow_F(t,s,y)$ for its value at time $t$.
\item Around each point of $U$ there is an open set $I\times\mathcal U\subset U$,
  with $I$ an interval, such that $\flow_F$ and $\partial_t\flow_F$ are defined
  and of class $C^k$ on $I\times I\times\mathcal U$, and there $D_y\flow_F$
  solves the variational equation
  \[
  \partial_tD_y\flow_F(t,s,y)=D_xF\big(t,\flow_F(t,s,y)\big)\,D_y\flow_F(t,s,y).
  \]
\item If $U=\R\times\R^n$ and for every $T>0$ there are numbers $M(T)$, $L(T)$
  with $|F(t,x)|\le M(T)+L(T)|x|$ on $[-T,T]\times\R^n$, then every maximal
  solution is defined on all of $\R$.
\end{enumerate}
\end{lemma}

Item (i) is Theorem 2.2 of \cite{Teschl} (the Picard-Lindel\"of theorem; a
$C^1$ function is locally Lipschitz) together with Theorem 2.13 there, which
produces the maximal solution. Item (ii) is Theorem 2.10 of \cite{Teschl}, and
item (iii) is Theorem 2.17 there. We use these facts through the following
lemma.

\begin{lemma}
\label{lem:flow}
Let $u\in C^\infty(\R\times\R^2;\R^2)$ be bounded on $[-T,T]\times\R^2$ for every
$T>0$, and write $\flow_u$ for the solution map of $\dot x=u(t,x)$ from
Lemma~\ref{lem:ode}(i). Then:
\begin{enumerate}[label=\textup{(\alph*)}]
\item $\flow_u(t,s,y)$ is defined for all $(t,s,y)\in\R\times\R\times\R^2$, and
  $\flow_u\in C^\infty(\R\times\R\times\R^2)$;
\item $\flow_u(t,\tau,\flow_u(\tau,s,y))=\flow_u(t,s,y)$ for all $t,\tau,s,y$, and
  each map $\flow_u(t,s,\cdot)$ is a $C^\infty$ diffeomorphism of $\R^2$ with
  inverse $\flow_u(s,t,\cdot)$;
\item $|\flow_u(t,s,y)-y|\le|t-s|\sup_{[-T,T]\times\R^2}|u|$ whenever
  $s,t\in[-T,T]$;
\item $Y(t)=D_y\flow_u(t,s,y)$ solves $\partial_tY=\nabla u(t,\flow_u(t,s,y))\,Y$
  with $Y(s)=I$, and if $\divop u=0$, then $\det D_y\flow_u\equiv1$.
\end{enumerate}
\end{lemma}

\begin{proof}
Lemma~\ref{lem:ode}(iii), with $M(T)=\sup_{[-T,T]\times\R^2}|u|$ and $L(T)=0$,
shows that $\flow_u$ is defined on $\R\times\R\times\R^2$. Both sides of the
identity in (b) solve $\dot x=u(t,x)$ as functions of $t$ and agree at $t=\tau$,
so they agree by the uniqueness in Lemma~\ref{lem:ode}(i). With $t=s$ the
identity gives $\flow_u(s,\tau,\flow_u(\tau,s,y))=y$, so $\flow_u(\tau,s,\cdot)$ is
a bijection with inverse $\flow_u(s,\tau,\cdot)$. Part (c) follows from
$\flow_u(t,s,y)-y=\int_s^tu(\sigma,\flow_u(\sigma,s,y))\dd\sigma$.

We prove the smoothness in (a). Fix $k\ge1$ and $(t_*,s_*,y_*)$, let $J$ be the
closed interval with endpoints $s_*$ and $t_*$, and put
$\gamma(\tau)=\flow_u(\tau,s_*,y_*)$. For every $\tau\in J$,
Lemma~\ref{lem:ode}(ii) gives an open set
$I_\tau\times\mathcal U_\tau\ni(\tau,\gamma(\tau))$ on which $\flow_u$ is $C^k$ in
the sense stated there. The sets
$\{\sigma\in J\cap I_\tau:\gamma(\sigma)\in\mathcal U_\tau\}$, $\tau\in J$, are
relatively open in $J$ and cover $J$. By the Lebesgue number lemma there are
points $s_*=\tau_0,\tau_1,\dots,\tau_N=t_*$, monotone in the index, and points
$\tau'_i\in J$ such that $\tau_i,\tau_{i+1}\in I_{\tau'_i}$ and
$\gamma(\tau_i)\in\mathcal U_{\tau'_i}$ for $0\le i<N$. By (b), for $(t,s,y)$ near $(t_*,s_*,y_*)$,
\[
\flow_u(t,s,y)=\flow_u\Big(t,\tau_{N-1},\flow_u\big(\tau_{N-1},\tau_{N-2},\cdots
\flow_u(\tau_2,\tau_1,\flow_u(\tau_1,s,y))\cdots\big)\Big).
\]
The innermost map $(s,y)\mapsto\flow_u(\tau_1,s,y)$ is $C^k$ near
$(\tau_0,\gamma(\tau_0))$; each map $y\mapsto\flow_u(\tau_{i+1},\tau_i,y)$ is $C^k$
on $\mathcal U_{\tau'_i}$, which is a neighborhood of $\gamma(\tau_i)$, and it
maps $\gamma(\tau_i)$ to $\gamma(\tau_{i+1})$; and
$(t,x)\mapsto\flow_u(t,\tau_{N-1},x)$ is $C^k$ for $t\in I_{\tau'_{N-1}}\ni t_*$
and $x\in\mathcal U_{\tau'_{N-1}}$. By
continuity and the chain rule, $\flow_u$ is $C^k$ near $(t_*,s_*,y_*)$. Since $k$
is arbitrary, $\flow_u\in C^\infty$, and by (b) the maps $\flow_u(t,s,\cdot)$ are
$C^\infty$ diffeomorphisms.

We prove (d). By (a), $\flow_u$ is $C^\infty$. Differentiating
$\partial_t\flow_u(t,s,y)=u(t,\flow_u(t,s,y))$ in $y$ and exchanging the order
of the derivatives gives
\[
\partial_tD_y\flow_u(t,s,y)=\nabla u\big(t,\flow_u(t,s,y)\big)\,D_y\flow_u(t,s,y),
\]
and $Y(s)=I$ because $\flow_u(s,s,y)=y$. The matrix
$Y(t)$ is invertible by (b), and Jacobi's formula gives
$\partial_t\det Y=\det Y\operatorname{tr}(Y^{-1}\partial_tY)=\det
Y\operatorname{tr}\nabla u=(\divop u)\det Y$. If $\divop u=0$, then $\det Y$ is
constant in $t$ and equal to $\det Y(s)=1$.
\end{proof}

\subsection{Two Estimates on the Plane}

\begin{lemma}
\label{lem:grad-bound}
Let $g\in C^1_c(\R^2)$ be real with $g\not\equiv0$, and let $u=\KBS\ast g$. Then
for all $x\in\R^2$ and $i,l\in\{1,2\}$,
\[
|\partial_lu_i(x)|\le\|g\|_{L^\infty}\Big(2+\log_+\frac{\|\nabla g\|_{L^\infty}}{\|g\|_{L^\infty}}\Big)
+\frac{\|g\|_{L^1}}{2\pi},
\]
where $\|\nabla g\|_{L^\infty}=\sup_x|\nabla g(x)|$.
\end{lemma}

\begin{proof}
Since $\KBS$ is locally integrable and $g\in C^1_c$, differentiation under the
integral sign gives $\partial_lu_i(x)=\int(\KBS)_i(z)\,\partial_lg(x-z)\dd z$.
Fix $\delta\in(0,1]$. By Lemma~I.\ref{I-lem:BS-basic}(a), the part of the integral
over $\{|z|<\delta\}$ is at most $\delta\|\nabla g\|_{L^\infty}$. On
$\{|z|>\delta\}$ we write $\partial_lg(x-z)=-\partial_{z_l}[g(x-z)]$ and apply
the divergence theorem on $\{\delta<|z|<\rho'\}$, with $\rho'$ so large that
$g(x-z)=0$ for $|z|\ge\rho'$. The outward normal on $\{|z|=\delta\}$ is
$-z/|z|$, so
\begin{align*}
\int_{|z|>\delta}(\KBS)_i(z)\,\partial_lg(x-z)\dd z
&=\int_{|z|>\delta}\partial_l(\KBS)_i(z)\,g(x-z)\dd z\\
&\quad+\int_{|z|=\delta}(\KBS)_i(z)\,g(x-z)\frac{z_l}{|z|}\dd S .
\end{align*}
By Lemma~I.\ref{I-lem:BS-basic}(a), the boundary term is at most
$(2\pi\delta)^{-1}\|g\|_{L^\infty}2\pi\delta=\|g\|_{L^\infty}$, the volume integral
over $\{\delta<|z|<1\}$ is at most
$\|g\|_{L^\infty}\int_\delta^1r^{-1}\dd r=\|g\|_{L^\infty}\log\frac1\delta$, and the
volume integral over $\{|z|\ge1\}$ is at most $\|g\|_{L^1}/(2\pi)$. Hence
$|\partial_lu_i(x)|\le\delta\|\nabla g\|_{L^\infty}+\|g\|_{L^\infty}\big(1+\log\frac1\delta\big)+\|g\|_{L^1}/(2\pi)$.
The choice $\delta=\min\{1,\|g\|_{L^\infty}/\|\nabla g\|_{L^\infty}\}$ gives the
claim.
\end{proof}

\begin{lemma}
\label{lem:comm-R2}
Let $s\ge2$ be an integer, $g\in C_c^\infty(\R^2)$, and let
$u\in C^\infty(\R^2;\R^2)$ satisfy
$|\nabla u|_{\max}:=\max_{i,l}\sup_x|\partial_lu_i(x)|<\infty$ and
$\partial^\beta u\in L^2$ for $1\le|\beta|\le s$. Then for every multi-index
$\alpha$ with $|\alpha|=s$,
\[
\|[\partial^\alpha,u\cdot\nabla]g\|_{L^2}\le2^{s+1}\Big[\big(|\nabla u|_{\max}+\|\nabla u\|_{H^{s-1}}\big)\|g\|_{H^s}
+\|\nabla u\|_{H^{s-1}}\|\nabla g\|_{L^\infty}\Big].
\]
\end{lemma}

\begin{proof}
By the Leibniz formula,
$[\partial^\alpha,u\cdot\nabla]g=\sum_{l=1}^2\sum_{0<\beta\le\alpha}\binom\alpha\beta\partial^\beta u_l\,\partial_l\partial^{\alpha-\beta}g$,
and $\sum_l\sum_{\beta\le\alpha}\binom\alpha\beta=2^{s+1}$. It suffices to bound
each product by the bracket on the right. If $|\beta|=1$, the product is at
most $|\nabla u|_{\max}\|g\|_{H^s}$. If $\beta=\alpha$, it is at most
$\|\partial^\alpha u_l\|_{L^2}\|\partial_lg\|_{L^\infty}\le\|\nabla u\|_{H^{s-1}}\|\nabla g\|_{L^\infty}$,
since $\partial^\alpha u_l$ is a derivative of order $s-1$ of an entry of
$\nabla u$. If $2\le|\beta|\le s-1$, H\"older's inequality and (E1) of
Lemma~\ref{lem:embed} give the bound
$\|\partial^\beta u_l\|_{H^1}\|\partial_l\partial^{\alpha-\beta}g\|_{H^1}
\le\|\nabla u\|_{H^{|\beta|}}\|g\|_{H^{s-|\beta|+2}}\le\|\nabla u\|_{H^{s-1}}\|g\|_{H^s}$,
because $|\beta|\le s-1$ and $|\beta|\ge2$.
\end{proof}

\subsection{The Iteration}

Throughout the rest of this appendix, $\omega_0$ and $R_0$ are as in
Theorem~\ref{thm:smooth}, and we assume $\omega_0\not\equiv0$. We put
\begin{gather*}
A=\|\omega_0\|_{L^\infty},\qquad
N_0=\sup_x|\nabla\omega_0(x)|,\\
U_*=\Big(\frac{2A\|\omega_0\|_{L^1}}{\pi}\Big)^{1/2},\qquad
c_*=2+\frac{\|\omega_0\|_{L^1}}{2\pi A},
\end{gather*}
and, for $t\in\R$,
\begin{gather*}
\mathcal L(t)=\Big(\log_+\frac{N_0}{A}+c_*\Big)e^{2A|t|}-c_*,\qquad
\Psi(t)=Ae^{\mathcal L(t)},\\
V(t)=2A\big(\mathcal L(t)+c_*\big),\qquad R(t)=R_0+U_*|t| .
\end{gather*}
The functions $\mathcal L$, $\Psi$ and $V$ are even and nondecreasing in
$|t|$; moreover $\mathcal L\ge\mathcal L(0)=\log_+(N_0/A)\ge0$, and
$\mathcal L'=V$ on $(0,\infty)$. Since $\supp\omega_0$ is a compact subset of
$B_{R_0}$, there is
$R_0'<R_0$ with $\omega_0=0$ outside $\overline B_{R_0'}$.

\begin{definition}
\label{def:iterates}
Put $\omega^{[0]}(t,x)=\omega_0(x)$ for $(t,x)\in\R\times\R^2$. If
$\omega^{[j]}\in C^\infty(\R\times\R^2)$ is defined and
$v^{[j]}=\KBS\ast\omega^{[j]}$ satisfies the hypotheses of Lemma~\ref{lem:flow},
let $\flow^{[j]}=\flow_{v^{[j]}}$ and put
$\omega^{[j+1]}(t,x)=\omega_0\big(\flow^{[j]}(0,t,x)\big)$.
\end{definition}

\begin{lemma}
\label{lem:iterates}
The iterates $\omega^{[j]}$ are defined for every $j\ge0$, and:
\begin{enumerate}[label=\textup{(\roman*)}]
\item $\omega^{[j]}\in C^\infty(\R\times\R^2)$ is real, and $\omega^{[j]}(t,x)=0$
  for $|x|>R_0'+U_*|t|$;
\item $\|\omega^{[j]}(t)\|_{L^p}=\|\omega_0\|_{L^p}$ for every $t\in\R$ and every
  $p\in[1,\infty]$;
\item $v^{[j]}\in C^\infty(\R\times\R^2;\R^2)$, every derivative of $v^{[j]}$ is
  bounded on $[-T,T]\times\R^2$ for every $T>0$, $\divop v^{[j]}=0$, and
  $|v^{[j]}|\le U_*$;
\item $\flow^{[j]}$ is defined and $C^\infty$ on $\R\times\R\times\R^2$,
  $\det D_y\flow^{[j]}\equiv1$, $\omega^{[j+1]}(0)=\omega_0$, and
  $\partial_t\omega^{[j+1]}+v^{[j]}\cdot\nabla\omega^{[j+1]}=0$ on $\R\times\R^2$.
\end{enumerate}
\end{lemma}

\begin{proof}
We argue by induction on $j$. Items (i) and (ii) hold for $j=0$. Assume (i) and
(ii) for some $j$.

We prove (iii) for $j$. Fix $T>0$. Every derivative
$\partial_t^i\partial_x^\gamma\omega^{[j]}$ is continuous and vanishes outside
the compact set $[-T,T]\times\overline B_{R(T)}$ on $[-T,T]\times\R^2$, so it is
bounded there. When $(t,x)$ ranges over a bounded subset of
$(-T,T)\times\R^2$, the functions
$z\mapsto\KBS(z)\,\partial_t^i\partial_x^\gamma\omega^{[j]}(t,x-z)$ are dominated
by a constant multiple of the integrable function
$|\KBS(z)|\mathbf 1_{\{|z|\le\rho\}}$, for a suitable $\rho$. Differentiation under the integral sign
and dominated convergence give
$\partial_t^i\partial_x^\gamma v^{[j]}=\KBS\ast\partial_t^i\partial_x^\gamma\omega^{[j]}$,
continuous, and bounded on $[-T,T]\times\R^2$ by Lemma~I.\ref{I-lem:BS-basic}(b).
Lemma~I.\ref{I-lem:BS-basic}(d) gives $\divop v^{[j]}=0$, and
Lemma~I.\ref{I-lem:BS-basic}(b) with (ii) gives
\[
|v^{[j]}(t,x)|\le\big(2\|\omega^{[j]}(t)\|_{L^1}\|\omega^{[j]}(t)\|_{L^\infty}/\pi\big)^{1/2}=U_* .
\]

We prove (iv) for $j$. By (iii), Lemma~\ref{lem:flow} applies to $v^{[j]}$ and
gives the properties of $\flow^{[j]}$. Hence $\omega^{[j+1]}$ is defined, real
and $C^\infty$, and $\omega^{[j+1]}(0,x)=\omega_0(x)$. By Lemma~\ref{lem:flow}(b),
$\omega^{[j+1]}(t,\flow^{[j]}(t,0,a))=\omega_0(a)$ for all $t$ and $a$.
Differentiating in $t$ gives
$(\partial_t\omega^{[j+1]}+v^{[j]}\cdot\nabla\omega^{[j+1]})(t,\flow^{[j]}(t,0,a))=0$,
and $\flow^{[j]}(t,0,\cdot)$ is onto.

We prove (i) and (ii) for $j+1$. If $\omega^{[j+1]}(t,x)\ne0$, then
$a=\flow^{[j]}(0,t,x)$ satisfies $|a|\le R_0'$, and by Lemma~\ref{lem:flow}(c),
$|x-a|\le U_*|t|$; so $|x|\le R_0'+U_*|t|$. The substitution
$x=\flow^{[j]}(t,0,a)$, whose Jacobian determinant is $1$, gives
$\int|\omega^{[j+1]}(t,x)|^p\dd x=\int|\omega_0(a)|^p\dd a$ for $p<\infty$, and the
case $p=\infty$ follows because $\flow^{[j]}(0,t,\cdot)$ is a bijection.
\end{proof}

\begin{lemma}
\label{lem:iter-lip}
For every $j\ge0$, $t\in\R$ and $x\in\R^2$,
\[
|\nabla\omega^{[j]}(t,x)|\le\Psi(t),\qquad
\max_{i,l}|\partial_lv^{[j]}_i(t,x)|\le A\big(\mathcal L(t)+c_*\big),\qquad
\|\nabla v^{[j]}(t,x)\|_{\mathrm{op}}\le V(t).
\]
\end{lemma}

\begin{proof}
The third bound follows from the second, because the operator norm of a
$2\times2$ matrix is at most twice the largest modulus of its entries. We show,
by induction on $j$, that the first bound holds for all $t$ and that it implies
the second.

For $j=0$ we have
$|\nabla\omega_0|\le N_0\le\max\{A,N_0\}=Ae^{\mathcal L(0)}\le\Psi(t)$. Suppose
that $|\nabla\omega^{[j]}(\tau,x)|\le\Psi(\tau)$ for all $\tau$ and $x$. Since
$\mathcal L\ge0$, we have $\log_+(\Psi(\tau)/A)=\mathcal L(\tau)$. By
Lemma~\ref{lem:iterates}(ii), $\|\omega^{[j]}(\tau)\|_{L^\infty}=A$ and
$\|\omega^{[j]}(\tau)\|_{L^1}=\|\omega_0\|_{L^1}$. Thus Lemma~\ref{lem:grad-bound},
applied to $g=\omega^{[j]}(\tau)$, gives
\[
|\partial_lv^{[j]}_i(\tau,x)|\le A\big(2+\mathcal L(\tau)\big)+\frac{\|\omega_0\|_{L^1}}{2\pi}=A\big(\mathcal L(\tau)+c_*\big).
\]
Hence $\|\nabla v^{[j]}(\tau,x)\|_{\mathrm{op}}\le V(\tau)$.

Next, by the chain rule,
$\nabla\omega^{[j+1]}(t,x)=D_y\flow^{[j]}(0,t,x)^{\mathsf T}\,\nabla\omega_0\big(\flow^{[j]}(0,t,x)\big)$.
By Lemma~\ref{lem:flow}(d), $Y(\tau)=D_y\flow^{[j]}(\tau,t,x)$ solves
$\partial_\tau Y=\nabla v^{[j]}(\tau,\flow^{[j]}(\tau,t,x))Y$ with $Y(t)=I$, so
$\|Y(\tau)\|_{\mathrm{op}}\le1+\big|\int_t^\tau V(\sigma)\|Y(\sigma)\|_{\mathrm{op}}\dd\sigma\big|$,
and Gronwall's inequality gives
$\|Y(0)\|_{\mathrm{op}}\le\exp\big|\int_0^tV(\sigma)\dd\sigma\big|=e^{\mathcal L(t)-\mathcal L(0)}$,
because $\mathcal L'=V$ on $(0,\infty)$ and both functions are even. Hence
$|\nabla\omega^{[j+1]}(t,x)|\le N_0e^{\mathcal L(t)-\mathcal L(0)}\le Ae^{\mathcal L(t)}=\Psi(t)$,
since $N_0e^{-\mathcal L(0)}=\min\{N_0,A\}$.
\end{proof}

\begin{lemma}
\label{lem:iter-Hs}
For every integer $s\ge1$ and every $T>0$, and for all $j\ge0$ and $|t|\le T$,
$\|\omega^{[j]}(t)\|_{H^s}\le\mathcal M_s(T)$, where
\begin{gather*}
\mathcal M_1(T)^2=\|\omega_0\|_{L^2}^2+\pi R(T)^2\Psi(T)^2,\\
\mathcal M_s(T)^2=\big(\|\omega_0\|_{H^s}^2+\mathcal M_{s-1}(T)^2+1\big)e^{\mathcal B_s(T)T}+\mathcal M_{s-1}(T)^2\quad(s\ge2),\\
\mathcal B_s(T)=(s+1)2^{s+2}\big(V(T)+\mathcal M_{s-1}(T)+\mathcal M_{s-1}(T)\Psi(T)\big).
\end{gather*}
In particular $\mathcal M_s(T)$ depends only on $s$, $T$, $R_0$, $A$, $N_0$ and
the norms $\|\omega_0\|_{L^1}$ and $\|\omega_0\|_{H^{s'}}$, $0\le s'\le s$.
\end{lemma}

\begin{proof}
We argue by induction on $s$. For $s=1$, Lemma~\ref{lem:iterates}(i)--(ii) and
Lemma~\ref{lem:iter-lip} give $\|\omega^{[j]}(t)\|_{L^2}=\|\omega_0\|_{L^2}$ and
$\|\nabla\omega^{[j]}(t)\|_{L^2}^2\le|B_{R(T)}|\Psi(T)^2$.

Let $s\ge2$ and assume the bound at level $s-1$ for all iterates. For $j=0$
the bound holds because $\mathcal M_s(T)\ge\|\omega_0\|_{H^s}$. Let $j\ge0$ and
$|\alpha|=s$. Differentiating the equation of Lemma~\ref{lem:iterates}(iv) gives
\[
\partial_t\partial^\alpha\omega^{[j+1]}+v^{[j]}\cdot\nabla\partial^\alpha\omega^{[j+1]}
=-[\partial^\alpha,v^{[j]}\cdot\nabla]\omega^{[j+1]} .
\]
We multiply by $\partial^\alpha\omega^{[j+1]}$ and integrate; all functions are
smooth and real, and $\omega^{[j+1]}(t)$ vanishes outside $B_{R(T)}$ for
$|t|\le T$. Since $\divop v^{[j]}=0$, the transport term gives
$-\frac12\int(\divop v^{[j]})(\partial^\alpha\omega^{[j+1]})^2\dd x=0$. Put
$\mathcal H(t)=\sum_{|\alpha|=s}\|\partial^\alpha\omega^{[j+1]}(t)\|_{L^2}^2$ and
$Z(t)=\|\omega^{[j+1]}(t)\|_{H^s}$; there are $s+1$ multi-indices with
$|\alpha|=s$. Put $\mathcal M=\mathcal M_{s-1}(T)$. We apply Lemma~\ref{lem:comm-R2} with
$u=v^{[j]}(t)$ and $g=\omega^{[j+1]}(t)$, using $|\nabla v^{[j]}|_{\max}\le V(T)$
and $|\nabla\omega^{[j+1]}|\le\Psi(T)$ (Lemma~\ref{lem:iter-lip}), and
\[
\|\nabla v^{[j]}(t)\|_{H^{s-1}}=\|\omega^{[j]}(t)\|_{H^{s-1}}\le\mathcal M
\]
(Lemma~I.\ref{I-lem:BS-basic}(d) and the induction hypothesis). This gives, for
$|t|\le T$,
\begin{align*}
|\mathcal H'(t)|&\le(s+1)2^{s+2}\big[(V(T)+\mathcal M)Z+\mathcal M\Psi(T)\big]Z\\
&\le\mathcal B_s(T)(Z^2+1)\le\mathcal B_s(T)\big(\mathcal H+\mathcal M^2+1\big),
\end{align*}
because $Z\le Z^2+1$ and $Z^2\le\mathcal H+\|\omega^{[j+1]}(t)\|_{H^{s-1}}^2\le\mathcal H+\mathcal M^2$.
Since $\mathcal H(0)\le\|\omega_0\|_{H^s}^2$, Gronwall's inequality gives
$\mathcal H(t)\le(\|\omega_0\|_{H^s}^2+\mathcal M^2+1)e^{\mathcal B_s(T)|t|}$, and
$Z^2\le\mathcal H+\mathcal M^2\le\mathcal M_s(T)^2$.
\end{proof}

\subsection{Convergence}

\begin{lemma}
\label{lem:iter-cauchy}
Let $T>0$ and put $e_j(t)=\|\omega^{[j]}(t)-\omega^{[j-1]}(t)\|_{L^2}$ for $j\ge1$.
Then for every $j\ge1$ and every $t\in[-T,T]$,
\[
e_{j}(t)\le2\|\omega_0\|_{L^2}\frac{\big(2R(T)\Psi(T)|t|\big)^{j-1}}{(j-1)!} .
\]
\end{lemma}

\begin{proof}
Put $R=R(T)$ and $\Psi=\Psi(T)$. For $j=1$, $e_1(t)\le2\|\omega_0\|_{L^2}$ by
Lemma~\ref{lem:iterates}(ii). Let $j\ge1$. Subtracting the equations of
Lemma~\ref{lem:iterates}(iv) for $\omega^{[j+1]}$ and $\omega^{[j]}$ shows that
$\delta=\omega^{[j+1]}-\omega^{[j]}$ satisfies
\[
\partial_t\delta+v^{[j]}\cdot\nabla\delta=-\big(\KBS\ast(\omega^{[j]}-\omega^{[j-1]})\big)\cdot\nabla\omega^{[j]} .
\]
We multiply by $\delta$ and integrate; the transport term vanishes, because
$\divop v^{[j]}=0$ and $\delta(t)$ has compact support. For $|t|\le T$, $\nabla\omega^{[j]}(t)$ vanishes outside $B_R$ and is
bounded by $\Psi$, and Lemma~I.\ref{I-lem:BS-basic}(c), with $\rho=R$, gives
$\|\KBS\ast(\omega^{[j]}-\omega^{[j-1]})(t)\|_{L^2(B_R)}\le2Re_j(t)$. Hence
$|\frac{\dd}{\dd t}e_{j+1}^2|\le4R\Psi e_je_{j+1}$, and for all $\epsilon>0$,
$\big|\frac{\dd}{\dd t}(e_{j+1}^2+\epsilon)^{1/2}\big|\le2R\Psi e_j$. Since
$e_{j+1}(0)=0$, letting $\epsilon\to0$ gives
$e_{j+1}(t)\le2R\Psi\big|\int_0^te_j(\tau)\dd\tau\big|$ for $|t|\le T$, and the
claim follows by induction on $j$.
\end{proof}

\begin{proposition}
\label{prop:iter-limit}
There is a real $\omega\in C^\infty(\R\times\R^2)$ such that, with
$v=\KBS\ast\omega$, for every $T>0$ and every multi-index $\gamma$,
$\partial_x^\gamma\omega^{[j]}\to\partial_x^\gamma\omega$ and
$\partial_x^\gamma v^{[j]}\to\partial_x^\gamma v$ uniformly on $[-T,T]\times\R^2$
as $j\to\infty$. The function $\omega$ has properties \textup{(1)}--\textup{(3)}
of Theorem~\ref{thm:smooth}, and for all $t\in\R$, $x\in\R^2$ and
$p\in[1,\infty]$,
\begin{equation}
\label{eq:limit-bounds}
\begin{gathered}
\omega(t,x)=0\ \text{ for }|x|>R_0'+U_*|t|,\qquad
\|\omega(t)\|_{L^p}=\|\omega_0\|_{L^p},\qquad |\nabla\omega(t,x)|\le\Psi(t),\\
|v(t,x)|\le U_*,\qquad \max_{i,l}|\partial_lv_i(t,x)|\le A\big(\mathcal L(t)+c_*\big).
\end{gathered}
\end{equation}
\end{proposition}

\begin{proof}
We first prove the convergence. Fix $T>0$. By Lemma~\ref{lem:iter-cauchy} the
series $\sum_j\sup_{|t|\le T}e_j(t)$ is at most
$2\|\omega_0\|_{L^2}\exp(2R(T)\Psi(T)T)$. Thus $(\omega^{[j]})$ is a Cauchy sequence
for the norm $\sup_{|t|\le T}\|\cdot\|_{L^2}$. For an integer
$k\ge1$, (E3) of Lemma~\ref{lem:embed} and Lemma~\ref{lem:iter-Hs} give
\[
\sup_{|t|\le T}\|\omega^{[j]}(t)-\omega^{[i]}(t)\|_{H^k}^2\le3^k\cdot2\mathcal M_{2k}(T)\sup_{|t|\le T}\|\omega^{[j]}(t)-\omega^{[i]}(t)\|_{L^2},
\]
which tends to $0$ as $i,j\to\infty$. By (E2), for all $\gamma$ the functions
$\partial_x^\gamma\omega^{[j]}$, which are continuous on $\R\times\R^2$, form a
uniform Cauchy sequence on $[-T,T]\times\R^2$. Their uniform limits are
continuous; the limit of $\partial_x^\gamma\omega^{[j]}$ is $\partial_x^\gamma$ of
the limit $\omega$ of $\omega^{[j]}$, as one sees by passing to the limit in
$g(x+he_l)-g(x)=\int_0^h\partial_lg(x+\sigma e_l)\dd\sigma$, with
$g=\partial_x^\gamma\omega^{[j]}(t)$; and the limits for
different $T$ agree. The first line of Equation~\eqref{eq:limit-bounds} follows
from Lemma~\ref{lem:iterates}(i)--(ii) and Lemma~\ref{lem:iter-lip}, because the
convergence is uniform and the supports lie in a fixed ball for $|t|\le T$. In
particular $\omega(t)\in C_c^\infty(\R^2)$ for all $t$, and $\omega$ has
property (1) of Theorem~\ref{thm:smooth}, since $R_0'+U_*|t|<R(t)$.

We now treat the velocity. For a bounded measurable $g$ vanishing outside
$B_\rho$, Lemma~I.\ref{I-lem:BS-basic}(b) and $\|g\|_{L^1}\le\pi\rho^2\|g\|_{L^\infty}$
give $\|\KBS\ast g\|_{L^\infty}\le\sqrt2\,\rho\|g\|_{L^\infty}$. Applied to
$g=\partial_x^\gamma\omega^{[j]}(t)-\partial_x^\gamma\omega(t)$ with $\rho=R(T)$,
this shows that $\partial_x^\gamma v^{[j]}=\KBS\ast\partial_x^\gamma\omega^{[j]}$
converges to $\KBS\ast\partial_x^\gamma\omega$ uniformly on $[-T,T]\times\R^2$.
By Lemma~I.\ref{I-lem:BS-basic}(d) applied to $\omega(t)$, $\KBS\ast\partial_x^\gamma\omega(t)=\partial_x^\gamma v(t)$
and $\divop v=0$. The bounds on $v$ in Equation~\eqref{eq:limit-bounds} follow
from Lemma~\ref{lem:iterates}(iii) and Lemma~\ref{lem:iter-lip}.

We now derive the equation. For every $\gamma$, Lemma~\ref{lem:iterates}(iv)
gives
\[
\partial_x^\gamma\omega^{[j+1]}(t,x)=\partial^\gamma\omega_0(x)
-\int_0^t\partial_x^\gamma\big(v^{[j]}\cdot\nabla\omega^{[j+1]}\big)(\tau,x)\dd\tau .
\]
By the Leibniz formula the integrand is a finite sum of products of
$x$-derivatives of $v^{[j]}$ and of $\omega^{[j+1]}$, each converging uniformly
and bounded on $[-T,T]\times\R^2$. Passing to the limit,
\[
\partial_x^\gamma\omega(t,x)=\partial^\gamma\omega_0(x)-\int_0^t\partial_x^\gamma(v\cdot\nabla\omega)(\tau,x)\dd\tau,
\]
with an integrand that is continuous and bounded on $[-T,T]\times\R^2$. Hence
$\partial_t\partial_x^\gamma\omega$ exists, is continuous and bounded, and
equals $-\partial_x^\gamma(v\cdot\nabla\omega)$. With $\gamma=0$ this is
property (3).

Finally we prove that $\omega$ and $v$ are $C^\infty$ with bounded derivatives.
For $T>0$ and an integer $l\ge0$, let $\mathcal C^l(T)$ be the set of functions
$g$ on $(-T,T)\times\R^2$ such that $\partial_t^i\partial_x^\gamma g$ exists, is
continuous and is bounded, for all $i\le l$ and all multi-indices $\gamma$
(including $i=0$ and $\gamma=0$); here the
$x$-derivatives are taken first. We have shown $\omega\in\mathcal C^1(T)$. Suppose
$\omega\in\mathcal C^l(T)$. Every $\partial_t^i\partial_x^\gamma\omega$ with
$i\le l$ vanishes for $|x|\ge R(T)$, so differentiation under the integral sign,
as in the proof of Lemma~\ref{lem:iterates}(iii), gives
$\partial_t^i\partial_x^\gamma v=\KBS\ast\partial_t^i\partial_x^\gamma\omega$,
continuous and bounded by $\sqrt2R(T)\sup|\partial_t^i\partial_x^\gamma\omega|$;
so $v\in\mathcal C^l(T)$. By the Leibniz formula $v\cdot\nabla\omega\in\mathcal C^l(T)$,
and the identity $\partial_t\partial_x^\gamma\omega=-\partial_x^\gamma(v\cdot\nabla\omega)$
shows that $\partial_t^{i+1}\partial_x^\gamma\omega$ exists, is continuous and
is bounded for $i\le l$. Hence $\omega\in\mathcal C^{l+1}(T)$. By induction,
$\omega$ and $v$ lie in $\mathcal C^l(T)$ for all $l$ and $T$. By Schwarz's
theorem on mixed partial derivatives (if $\partial_ag$, $\partial_bg$ and
$\partial_b\partial_ag$ exist near a point and $\partial_b\partial_ag$ is
continuous at it, then $\partial_a\partial_bg$ exists at that point and equals
$\partial_b\partial_ag$ there), an induction on the order shows that every
partial derivative of $\omega$ or $v$, taken in any order, exists, is continuous,
and equals the corresponding $\partial_t^i\partial_x^\gamma$ derivative. Thus
$\omega,v\in C^\infty(\R\times\R^2)$, $v$ and every derivative of $v$ are bounded
on $[-T,T]\times\R^2$ for all $T$, and $\omega$ has property (2).
\end{proof}

\subsection{Proof of Theorem~\ref{thm:smooth}}

The proof of Theorem~\ref{thm:smooth} also gives the following explicit bounds.

\begin{lemma}
\label{lem:smooth-bounds}
In Theorem~\ref{thm:smooth}, assume $\omega_0\not\equiv0$ and put
$A=\|\omega_0\|_{L^\infty}$, $N_0=\sup_x|\nabla\omega_0(x)|$,
$U_*=\sqrt{2A\|\omega_0\|_{L^1}/\pi}$, $c_*=2+\|\omega_0\|_{L^1}/(2\pi A)$,
$\mathcal L(t)=(\log_+(N_0/A)+c_*)e^{2A|t|}-c_*$ and $\Psi(t)=Ae^{\mathcal L(t)}$.
Then for every $t\in\R$ and $x\in\R^2$,
\begin{gather*}
|\nabla\omega(t,x)|\le\Psi(t),\qquad|v(t,x)|\le U_*,\\
\max_{i,l}|\partial_lv_i(t,x)|\le A\big(\mathcal L(t)+c_*\big),\qquad
\|\nabla v(t,x)\|_{\mathrm{op}}\le2A\big(\mathcal L(t)+c_*\big).
\end{gather*}
\end{lemma}

\begin{proof}[Proof of Theorem~\ref{thm:smooth} and Lemma~\ref{lem:smooth-bounds}]
We first prove (5), for any $\omega$ with properties (1)--(3). Let $\tilde\omega$
be as in (5), and fix a compact interval $J\subset I$ containing $0$. Let
$\rho>0$ be such that $\omega(t)$ and $\tilde\omega(t)$ vanish outside $B_\rho$ for
$t\in J$; this is possible by (1) and the hypothesis on $\tilde\omega$. Put
$\delta=\omega-\tilde\omega\in C^1(J\times\R^2)$. Since
$\partial_t\tilde\omega=-v\cdot\nabla\tilde\omega+(v-\KBS\ast\tilde\omega)\cdot\nabla\tilde\omega$,
\[
\partial_t\delta+v\cdot\nabla\delta=-(\KBS\ast\delta)\cdot\nabla\tilde\omega
\qquad\text{on }J\times\R^2 .
\]
We multiply by $\delta$ and integrate in $x$. The transport term vanishes,
because $\divop v=0$ and $\delta(t)$ has compact support. The function
$\nabla\tilde\omega$ is continuous and vanishes outside $J\times B_\rho$, so it is
bounded on $J\times\R^2$ by some $L<\infty$, and
Lemma~I.\ref{I-lem:BS-basic}(c) gives $\|\KBS\ast\delta(t)\|_{L^2(B_\rho)}\le2\rho\|\delta(t)\|_{L^2}$.
Hence $|\frac{\dd}{\dd t}\|\delta(t)\|_{L^2}^2|\le4\rho L\|\delta(t)\|_{L^2}^2$ on
$J$. Since $\delta(0)=0$, Gronwall's inequality gives $\delta=0$ on $J$. Every
point of $I$ lies in such an interval $J$, which proves (5). In particular at
most one $\omega\in C^\infty(\R\times\R^2)$ has properties (1)--(3), since such
an $\omega$ satisfies the hypotheses on $\tilde\omega$ in (5) with $I=\R$.

If $\omega_0\equiv0$, the function $\omega\equiv0$ has properties (1)--(4), (6)
and (7). From now on $\omega_0\not\equiv0$, and $\omega$ is the function of
Proposition~\ref{prop:iter-limit}, which has properties (1)--(3). The bounds of
Lemma~\ref{lem:smooth-bounds} are contained in Equation~\eqref{eq:limit-bounds},
together with $\|M\|_{\mathrm{op}}\le2\max_{i,l}|M_{il}|$ for $2\times2$ matrices
$M$.

We prove (4). By (2), Lemma~\ref{lem:flow} applies to $v$; let
$\flow(t,y)=\flow_v(t,0,y)$ and $\flow_t=\flow(t,\cdot)$. By Lemma~\ref{lem:flow},
$\flow$ is defined for all $t$; by (a) and (b) there, $\flow$ and
$(t,x)\mapsto\flow_t^{-1}(x)=\flow_v(0,t,x)$ lie in $C^\infty(\R\times\R^2)$ and
each $\flow_t$ is a $C^\infty$ diffeomorphism; and by (d),
$\det D\flow_t\equiv1$ because $\divop v=0$. By (3),
$\frac{\dd}{\dd t}\omega(t,\flow(t,a))=(\partial_t\omega+v\cdot\nabla\omega)(t,\flow(t,a))=0$,
so $\omega(t)=\omega_0\circ\flow_t^{-1}$. For a continuous $F$ with $F(0)=0$,
the function $F(\omega(t))$ is continuous with compact support, and the change
of variables $x=\flow_t(a)$ gives $\int F(\omega(t))\dd x=\int F(\omega_0)\dd a$.

We prove (6). Let $\Orot_\theta\in SO(2)$. Since $(\Orot_\theta z)^\perp=\Orot_\theta z^\perp$,
we have $\KBS(\Orot_\theta z)=\Orot_\theta\KBS(z)$. Put
$\omega_\theta(t,x)=\omega(t,\Orot_\theta^{-1}x)$. The substitution
$y=\Orot_\theta y'$ gives
$\KBS\ast\omega_\theta(t,x)=\Orot_\theta v(t,\Orot_\theta^{-1}x)$, and
$\nabla\omega_\theta(t,x)=\Orot_\theta\nabla\omega(t,\Orot_\theta^{-1}x)$. Hence
$(\KBS\ast\omega_\theta)\cdot\nabla\omega_\theta=(v\cdot\nabla\omega)(t,\Orot_\theta^{-1}x)$,
and $\omega_\theta$ is a $C^\infty$ solution of the equation in (3) with datum
$\omega_0\circ\Orot_\theta^{-1}$, vanishing outside $B_{R(t)}$ at time $t$. If
$\omega_0\circ\Orot_\theta=\omega_0$, then $\omega_0\circ\Orot_\theta^{-1}=\omega_0$,
and (5) with $I=\R$ gives $\omega_\theta=\omega$, that is,
$\omega(t)\circ\Orot_\theta^{-1}=\omega(t)$ for all $t$, which is equivalent to
$\omega(t)\circ\Orot_\theta=\omega(t)$.

We prove (7). Let $T>0$, $\varphi\in\DT$, and let $\Ren\in C^1(\R)$ with
$\Ren'$ bounded and $\Ren(0)=0$. The function $\Ren(\omega)$ lies in
$C^1(\R\times\R^2)$, vanishes outside $B_{R(t)}$ at time $t$, and by (3)
satisfies $\partial_t\Ren(\omega)+v\cdot\nabla\Ren(\omega)=\Ren'(\omega)(\partial_t\omega+v\cdot\nabla\omega)=0$.
We multiply by $\varphi$ and integrate over $(0,T)\times\R^2$. Since
$\varphi(t,\cdot)=0$ for $t$ near $T$, an integration by parts in $t$ gives
$\int_0^T\!\int\partial_t\Ren(\omega)\,\varphi=-\int\Ren(\omega_0)\varphi(0,\cdot)-\int_0^T\!\int\Ren(\omega)\,\partial_t\varphi$,
and since $\divop v=0$, an integration by parts in $x$ gives
$\int(v\cdot\nabla\Ren(\omega))\varphi\dd x=-\int\Ren(\omega)\,v\cdot\nabla\varphi\dd x$.
Hence Equation~(I.\ref{I-eq:renorm}) holds, and with $\Ren(s)=s$ so does
Equation~(I.\ref{I-eq:weak}). Finally, $\|\omega(t)\|_{L^1}=\|\omega_0\|_{L^1}$ by
(4) with $F(s)=|s|$, and $\|\omega(t)\|_{L^\infty}=\|\omega_0\|_{L^\infty}$
because $\omega(t)=\omega_0\circ\flow_t^{-1}$. Moreover $\omega$ is continuous and
$\omega(t)$ vanishes outside $B_{R(T)}$ for $|t|\le T$, so $\omega$ is uniformly
continuous on $[0,T]\times\R^2$, and $t\mapsto\omega(t)$ is continuous from
$[0,T]$ into $(L^1\cap L^q)(\R^2)$ for all $q\in[1,\infty]$. Hence the
restriction of $\omega$ to $[0,T)$ lies in $L^\infty_t(L^1\cap L^\infty)_x$, and
in $X_q(T)$ for every $q\in(4/3,\infty)$. This proves (7) and completes the
proof.
\end{proof}

\section{Dependency Table}
\label{app:deps}

Table~\ref{tab:deps} lists, for each statement of this paper on the path to
Theorem~\ref{thm:main} whose proof cites another statement, the statements of
this paper and of Part~I that its proof uses and the parts of
Hypothesis~I.\ref{I-hyp:N} on which it depends,
directly or through its inputs. Statements whose proofs cite no other statement
have no row, and definitions are not listed as inputs. Theorem~\ref{thm:smooth}
and Lemma~\ref{lem:smooth-bounds} have a joint proof and share a row; consecutive
statements with the same inputs that depend on the same parts of
Hypothesis~I.\ref{I-hyp:N} also share a row. The rows are
ordered so that every input that has a row appears in an earlier row.
Statements of Part~I appear only as inputs, with their Part~I numbers; they
have no row, and a row is marked with the parts of Hypothesis~I.\ref{I-hyp:N}
that its Part~I inputs assume in their statements.

A statement that assumes some of \textup{(N1)}--\textup{(N3)} is marked with them,
and so is every statement whose inputs are marked. The objects of
Definition~I.\ref{I-def:eigendata} exist only under \textup{(N1)}--\textup{(N3)}, and
the seed of the approximate data of Definition~\ref{def:approx} is built from
them; so every statement about these objects, or about the approximate
solutions, depends on \textup{(N1)}--\textup{(N3)} through these definitions,
even when its proof cites none of the three. Theorem~\ref{thm:N-holds} proves
\textup{(N1)}--\textup{(N3)}, so it and Theorem~\ref{thm:main} depend on none of
them. A dash means that the statement depends on none of them.
Lemma~\ref{lem:ode} is quoted from \cite{Teschl} without proof, and only the proof of
Proposition~\ref{prop:runs} relies on machine computation.

\begingroup
\small
\captionsetup{type=table}
\caption{Dependencies of the statements on the path to Theorem~\ref{thm:main}, and the parts of Hypothesis~I.\ref{I-hyp:N} that they use.}
\label{tab:deps}
\setlength{\LTleft}{0pt}\setlength{\LTright}{\fill}\setlength{\LTpre}{0pt}
\begin{longtable}{p{0.25\textwidth}p{0.54\textwidth}p{0.12\textwidth}}
Statement & Inputs of the proof & Uses \\
\hline
\endfirsthead
\multicolumn{3}{l}{\emph{Table~\ref{tab:deps} (continued)}} \\
Statement & Inputs of the proof & Uses \\
\hline
\endhead
\raggedright Lem.~\ref{lem:flow}
  & \raggedright Lem.~\ref{lem:ode} & --- \\
\raggedright Lem.~\ref{lem:grad-bound}
  & \raggedright Lem.~I.\ref{I-lem:BS-basic} & --- \\
\raggedright Lem.~\ref{lem:comm-R2}
  & \raggedright Lem.~\ref{lem:embed} & --- \\
\raggedright Lem.~\ref{lem:iterates}
  & \raggedright Lem.~I.\ref{I-lem:BS-basic}, \ref{lem:flow} & --- \\
\raggedright Lem.~\ref{lem:iter-lip}
  & \raggedright Lem.~\ref{lem:flow}, \ref{lem:grad-bound}, \ref{lem:iterates} & --- \\
\raggedright Lem.~\ref{lem:iter-Hs}
  & \raggedright Lem.~I.\ref{I-lem:BS-basic}, \ref{lem:comm-R2}, \ref{lem:iterates}, \ref{lem:iter-lip} & --- \\
\raggedright Lem.~\ref{lem:iter-cauchy}
  & \raggedright Lem.~I.\ref{I-lem:BS-basic}, \ref{lem:iterates} & --- \\
\raggedright Prop.~\ref{prop:iter-limit}
  & \raggedright Lem.~I.\ref{I-lem:BS-basic}, \ref{lem:embed}, \ref{lem:iterates}, \ref{lem:iter-lip}, \ref{lem:iter-Hs}, \ref{lem:iter-cauchy} & --- \\
\raggedright Thm.~\ref{thm:smooth}; Lem.~\ref{lem:smooth-bounds}
  & \raggedright Prop.~\ref{prop:iter-limit}; Lem.~I.\ref{I-lem:BS-basic}, \ref{lem:flow} & --- \\
\raggedright Lem.~\ref{lem:Hk-calculus}
  & \raggedright Lem.~I.\ref{I-lem:BS-basic}, \ref{lem:embed} & --- \\
\raggedright Lem.~\ref{lem:duhamel}
  & \raggedright Lem.~I.\ref{I-lem:sector}, I.\ref{I-lem:semigroup-construct} & --- \\
\raggedright Lem.~\ref{lem:reduction}
  & \raggedright Thm.~\ref{thm:smooth}; Lem.~I.\ref{I-lem:ring-facts}, \ref{lem:flow} & --- \\
\raggedright Lem.~\ref{lem:bilinear}
  & \raggedright Lem.~I.\ref{I-lem:uniform}, \ref{lem:embed}, \ref{lem:Hk-calculus}, \ref{lem:reduction} & (N1)--(N3) \\
\raggedright Lem.~\ref{lem:energy}
  & \raggedright Prop.~I.\ref{I-prop:U}; Lem.~\ref{lem:embed}, \ref{lem:Hk-calculus}, \ref{lem:reduction} & (N1)--(N3) \\
\raggedright Lem.~\ref{lem:majorant}
  & \raggedright Prop.~I.\ref{I-prop:SU} & (N1)--(N3) \\
\raggedright Thm.~\ref{thm:quiet}
  & \raggedright Thm.~\ref{thm:smooth}; Prop.~I.\ref{I-prop:semigroup}; Lem.~\ref{lem:annuli}, \ref{lem:duhamel}, \ref{lem:reduction}, \ref{lem:bilinear}, \ref{lem:energy} & (N1)--(N3) \\
\raggedright Lem.~\ref{lem:running}
  & \raggedright Prop.~I.\ref{I-prop:semigroup} & (N1)--(N3) \\
\raggedright Lem.~\ref{lem:linear-part}
  & \raggedright Prop.~I.\ref{I-prop:semigroup}; Lem.~I.\ref{I-lem:uniform}, \ref{lem:running} & (N1)--(N3) \\
\raggedright Lem.~\ref{lem:window}
  & \raggedright Prop.~I.\ref{I-prop:semigroup}; Lem.~I.\ref{I-lem:uniform}, \ref{lem:annuli}, \ref{lem:duhamel}, \ref{lem:reduction}, \ref{lem:bilinear}, \ref{lem:energy}, \ref{lem:running}, \ref{lem:linear-part} & (N1)--(N3) \\
\raggedright Thm.~\ref{thm:active}
  & \raggedright Thm.~\ref{thm:smooth}; Prop.~I.\ref{I-prop:semigroup}; Lem.~I.\ref{I-lem:uniform}, \ref{lem:annuli}, \ref{lem:duhamel}, \ref{lem:reduction}, \ref{lem:bilinear}, \ref{lem:energy}, \ref{lem:running}, \ref{lem:linear-part}, \ref{lem:window} & (N1)--(N3) \\
\raggedright Cor.~\ref{cor:active-moment}
  & \raggedright Thm.~\ref{thm:active}; Lem.~I.\ref{I-lem:uniform} & (N1)--(N3) \\
\raggedright Lem.~\ref{lem:seed}
  & \raggedright Thm.~I.\ref{I-thm:evans-simple}; Lem.~I.\ref{I-lem:sector}, \ref{lem:flow} & (N1)--(N3) \\
\raggedright Prop.~\ref{prop:integrability}
  & \raggedright Lem.~I.\ref{I-lem:ring-facts} & --- \\
\raggedright Lem.~\ref{lem:adm}
  & \raggedright Thm.~\ref{thm:smooth}; Prop.~\ref{prop:integrability}; Lem.~I.\ref{I-lem:ring-facts}, \ref{lem:seed} & (N1)--(N3) \\
\raggedright Lem.~\ref{lem:regions-basic}
  & \raggedright Thm.~\ref{thm:smooth}; Lem.~I.\ref{I-lem:ring-facts}, \ref{lem:annuli}, \ref{lem:seed}, \ref{lem:adm} & (N1)--(N3) \\
\raggedright Lem.~\ref{lem:inner-circ}
  & \raggedright Thm.~\ref{thm:smooth}; Lem.~\ref{lem:regions-basic} & (N1)--(N3) \\
\raggedright Lem.~\ref{lem:speed}
  & \raggedright Lem.~I.\ref{I-lem:BS-basic}, I.\ref{I-lem:ring-facts}, \ref{lem:seed}, \ref{lem:inner-circ} & (N1)--(N3) \\
\raggedright Lem.~\ref{lem:clock}
  & \raggedright Prop.~I.\ref{I-prop:SU}; Lem.~\ref{lem:inner-circ} & (N1)--(N3) \\
\raggedright Lem.~\ref{lem:transfer}
  & \raggedright Lem.~\ref{lem:hollow}, \ref{lem:ext-multipole}, \ref{lem:regions-basic}, \ref{lem:speed} & (N1)--(N3) \\
\raggedright Lem.~\ref{lem:inner-field}
  & \raggedright Lem.~\ref{lem:ext-multipole}, \ref{lem:regions-basic}, \ref{lem:inner-circ}, \ref{lem:transfer} & (N1)--(N3) \\
\raggedright Prop.~\ref{prop:confine}
  & \raggedright Thm.~\ref{thm:smooth}; Lem.~\ref{lem:hollow}, \ref{lem:regions-basic}, \ref{lem:speed} & (N1)--(N3) \\
\raggedright Lem.~\ref{lem:field-bound}
  & \raggedright Lem.~\ref{lem:regions-basic}, \ref{lem:speed}, \ref{lem:transfer}, \ref{lem:inner-field} & (N1)--(N3) \\
\raggedright Lem.~\ref{lem:coupling}
  & \raggedright Thm.~\ref{thm:quiet}; Prop.~I.\ref{I-prop:semigroup}; Lem.~I.\ref{I-lem:BS-basic}, I.\ref{I-lem:semigroup-construct}, I.\ref{I-lem:uniform}, \ref{lem:reduction}, \ref{lem:regions-basic} & (N1)--(N3) \\
\raggedright Lem.~\ref{lem:detuning}
  & \raggedright Prop.~I.\ref{I-prop:U}; Lem.~\ref{lem:inner-circ} & (N1)--(N3) \\
\raggedright Lem.~\ref{lem:transfer-step}
  & \raggedright Lem.~\ref{lem:ext-multipole}, \ref{lem:regions-basic}, \ref{lem:inner-circ}, \ref{lem:transfer} & (N1)--(N3) \\
\raggedright Lem.~\ref{lem:O}
  & \raggedright Thm.~\ref{thm:quiet}; Lem.~\ref{lem:reduction}, \ref{lem:majorant}, \ref{lem:detuning}, \ref{lem:transfer-step} & (N1)--(N3) \\
\raggedright Prop.~\ref{prop:relay}
  & \raggedright Thm.~\ref{thm:quiet}, \ref{thm:active}; Lem.~\ref{lem:majorant}, \ref{lem:clock}, \ref{lem:transfer}, \ref{lem:detuning}, \ref{lem:transfer-step}, \ref{lem:O} & (N1)--(N3) \\
\raggedright Lem.~\ref{lem:Oprime}
  & \raggedright Thm.~\ref{thm:quiet}, \ref{thm:active}; Prop.~\ref{prop:relay}; Lem.~\ref{lem:detuning}, \ref{lem:transfer-step}, \ref{lem:O} & (N1)--(N3) \\
\raggedright Lem.~\ref{lem:drift}
  & \raggedright Thm.~\ref{thm:active}; Prop.~\ref{prop:confine}; Lem.~I.\ref{I-lem:uniform}, \ref{lem:running}, \ref{lem:inner-circ}, \ref{lem:coupling}, \ref{lem:detuning} & (N1)--(N3) \\
\raggedright Lem.~\ref{lem:handoff}
  & \raggedright Thm.~\ref{thm:quiet}; Prop.~I.\ref{I-prop:semigroup}; Lem.~I.\ref{I-lem:uniform}, \ref{lem:majorant}, \ref{lem:inner-circ}, \ref{lem:coupling} & (N1)--(N3) \\
\raggedright Lem.~\ref{lem:kicktrue}
  & \raggedright Thm.~\ref{thm:quiet}, \ref{thm:active}; Prop.~\ref{prop:confine}, \ref{prop:kick}; Lem.~I.\ref{I-lem:ring-facts}, I.\ref{I-lem:uniform}, \ref{lem:annuli}, \ref{lem:majorant}, \ref{lem:inner-circ}, \ref{lem:clock}, \ref{lem:coupling}, \ref{lem:detuning}, \ref{lem:Oprime}, \ref{lem:drift}; Cor.~\ref{cor:active-moment} & (N1)--(N3) \\
\raggedright Lem.~\ref{lem:time}
  & \raggedright Lem.~\ref{lem:detuning} & (N1)--(N3) \\
\raggedright Lem.~\ref{lem:quiet-time}
  & \raggedright Lem.~\ref{lem:time} & (N1)--(N3) \\
\raggedright Lem.~\ref{lem:theta-quiet}
  & \raggedright Lem.~\ref{lem:O}, \ref{lem:quiet-time} & (N1)--(N3) \\
\raggedright Lem.~\ref{lem:nonempty}
  & \raggedright Thm.~\ref{thm:quiet}, \ref{thm:active}; Prop.~I.\ref{I-prop:U}; Lem.~\ref{lem:inner-circ}, \ref{lem:Oprime} & (N1)--(N3) \\
\raggedright Lem.~\ref{lem:geometry}
  & \raggedright Prop.~\ref{prop:confine}; Lem.~I.\ref{I-lem:uniform}, \ref{lem:reduction}, \ref{lem:clock}, \ref{lem:field-bound}, \ref{lem:detuning} & (N1)--(N3) \\
\raggedright Prop.~\ref{prop:bootstrap}
  & \raggedright Thm.~\ref{thm:smooth}, \ref{thm:quiet}, \ref{thm:active}; Prop.~\ref{prop:confine}, \ref{prop:relay}; Lem.~I.\ref{I-lem:uniform}, \ref{lem:reduction}, \ref{lem:seed}, \ref{lem:regions-basic}, \ref{lem:speed}, \ref{lem:clock}, \ref{lem:inner-field}, \ref{lem:coupling}, \ref{lem:detuning}, \ref{lem:O}, \ref{lem:Oprime}, \ref{lem:drift}, \ref{lem:handoff}, \ref{lem:kicktrue}, \ref{lem:time}, \ref{lem:quiet-time}, \ref{lem:theta-quiet}, \ref{lem:nonempty}, \ref{lem:geometry}; Cor.~\ref{cor:active-moment} & (N1)--(N3) \\
\raggedright Prop.~\ref{prop:enclosure}; Lem.~\ref{lem:apriori}
  & \raggedright Lem.~\ref{lem:end-regions} & --- \\
\raggedright Lem.~\ref{lem:Nprime-N}
  & \raggedright Lem.~I.\ref{I-lem:evans-basic}, \ref{lem:apriori} & (N1) \\
\raggedright Prop.~\ref{prop:runs}
  & \raggedright Prop.~\ref{prop:enclosure}, \ref{prop:winding}; Lem.~I.\ref{I-lem:evans-basic}; computation under (TB1)--(TB3) & --- \\
\raggedright Prop.~\ref{prop:N3}
  & \raggedright Lem.~\ref{lem:apriori} & --- \\
\raggedright Prop.~\ref{prop:reduction}
  & \raggedright Thm.~\ref{thm:active}; Prop.~\ref{prop:integrability}, \ref{prop:bootstrap}; Lem.~I.\ref{I-lem:radial-stationary}, I.\ref{I-lem:uniform}, \ref{lem:adm}, \ref{lem:inner-circ}, \ref{lem:time}, \ref{lem:nonempty}; Cor.~I.\ref{I-cor:defect} & (N1)--(N3) \\
\raggedright Thm.~\ref{thm:N-holds}
  & \raggedright Prop.~\ref{prop:runs}, \ref{prop:N3}; Lem.~\ref{lem:Nprime-N} & --- \\
\raggedright Thm.~\ref{thm:main}
  & \raggedright Thm.~\ref{thm:N-holds}; Prop.~\ref{prop:reduction} & --- \\
\end{longtable}
\endgroup


\section*{Statements and Declarations}
The author has no competing interests to declare.

\begin{thebibliography}{99}

\bibitem{ABCDGJK}
Albritton, D., Bru\`e, E., Colombo, M., De Lellis, C., Giri, V., Janisch, M., Kwon, H.
\textit{Instability and Non-uniqueness for the 2D Euler Equations, after M. Vishik}.
Annals of Mathematics Studies. Princeton University Press. (2024).
arXiv:2112.04943. We use the numbering of the book.

\bibitem{PartI}
Coiculescu, M.P.
\textit{Non-Uniqueness of Weak Solutions to the Unforced Two-Dimensional Euler Equations with
Vorticity in $L^1\cap L^p$, I: A Compactness Criterion and the Linear Analysis}.
Companion paper.

\bibitem{DiPernaLions}
DiPerna, R.J., Lions, P.-L.
\textit{Ordinary Differential Equations, Transport Theory and Sobolev Spaces}.
Invent. Math. Vol. 98. No. 3. pp. 511-547. (1989).

\bibitem{Grenier}
Grenier, E.
\textit{On the Nonlinear Instability of Euler and Prandtl Equations}.
Comm. Pure Appl. Math. Vol. 53. No. 9. pp. 1067-1091. (2000).

\bibitem{FLINT}
Hart, W.B.
\textit{Fast Library for Number Theory: An Introduction}.
In: Fukuda, K., van der Hoeven, J., Joswig, M., Takayama, N. (eds.).
Mathematical Software -- ICMS 2010. Lecture Notes in Computer Science.
Springer. pp. 88-91. (2010).

\bibitem{Holder}
H\"older, E.
\textit{\"Uber die unbeschr\"ankte Fortsetzbarkeit einer stetigen ebenen Bewegung in
einer unbegrenzten inkompressiblen Fl\"ussigkeit}.
Math. Z. Vol. 37. No. 1. pp. 727-738. (1933).

\bibitem{Johansson2017}
Johansson, F.
\textit{Arb: Efficient Arbitrary-Precision Midpoint-Radius Interval Arithmetic}.
IEEE Trans. Comput. Vol. 66. No. 8. pp. 1281-1292. (2017).

\bibitem{pythonflint}
Johansson, F., Benjamin, O.
\textit{python-flint: Bindings for FLINT}. Version 0.9.0.
\url{https://github.com/flintlib/python-flint}. (2026).

\bibitem{Kato1967}
Kato, T.
\textit{On Classical Solutions of the Two-Dimensional Non-Stationary Euler Equation}.
Arch. Rational Mech. Anal. Vol. 25. No. 3. pp. 188-200. (1967).

\bibitem{Teschl}
Teschl, G.
\textit{Ordinary Differential Equations and Dynamical Systems}.
Graduate Studies in Mathematics. 140. American Mathematical Society. (2012).
We use the theorem numbering of the author's preliminary version, which the
author makes available with the permission of the publisher.

\bibitem{Vishik1}
Vishik, M.
\textit{Instability and Non-Uniqueness in the Cauchy Problem for the Euler Equations
of an Ideal Incompressible Fluid. Part I}.
arXiv:1805.09426. (2018).

\bibitem{Vishik2}
Vishik, M.
\textit{Instability and Non-Uniqueness in the Cauchy Problem for the Euler Equations
of an Ideal Incompressible Fluid. Part II}.
arXiv:1805.09440. (2018).

\bibitem{Wolibner}
Wolibner, W.
\textit{Un th\'eor\`eme sur l'existence du mouvement plan d'un fluide parfait,
homog\`ene, incompressible, pendant un temps infiniment long}.
Math. Z. Vol. 37. No. 1. pp. 698-726. (1933).

\end{thebibliography}


\end{document}
